\documentclass[11pt,a4paper]{ltjsarticle}

% 日本語版（著者の逐行確認用）．英語正本：02-v2-lg-fanout-boundary/main.tex（コミット d6753a9）．
% LuaLaTeX でコンパイルする：latexmk -lualatex main_JP.tex
\usepackage[T1]{fontenc}
\usepackage{lmodern}
\usepackage{microtype}
\usepackage{amsmath,amssymb,amsthm,mathtools}
\numberwithin{equation}{section}
\usepackage{enumitem}
\usepackage{array}
\usepackage{tabularx}
\usepackage{geometry}
\geometry{margin=1in}
\usepackage{cite}
\usepackage{hyperref}
\usepackage{url}
\usepackage{tikz}
\hypersetup{hidelinks}
\setlength{\emergencystretch}{3em}
\clubpenalty=5000
\widowpenalty=5000
\displaywidowpenalty=5000

\theoremstyle{plain}
\newtheorem{theorem}{定理}[section]
\newtheorem{lemma}[theorem]{補題}
\newtheorem{proposition}[theorem]{命題}
\newtheorem{corollary}[theorem]{系}
\theoremstyle{definition}
\newtheorem{definition}[theorem]{定義}
\theoremstyle{remark}
\newtheorem{remark}[theorem]{注意}
\newtheorem{example}[theorem]{例}
\newtheorem{openproblem}[theorem]{未解決問題}
\renewcommand{\proofname}{証明}
% ltjsarticle の \appendix は \thesection に「付録」を含めるので，英語版と同じ番号（A, A.6 など）になるよう戻す．
\makeatletter
\g@addto@macro\appendix{\gdef\thesection{\@Alph\c@section}}
\makeatother

\newcommand{\Wit}{\mathsf{W}}
\newcommand{\Cmcf}[2]{\mathcal{C}^{\mathrm{mcf}}_{#1,#2}}
\newcommand{\D}{\mathcal{D}}
\newcommand{\out}{\mathrm{out}}
\newcommand{\tuple}[1]{\vec{#1}}
\newcommand{\blank}{\square}
\newcommand{\eps}{\varepsilon}
\newcommand{\pausesym}{\bot}
\newcommand{\fan}{\operatorname{fan}}
\newcommand{\YSub}[1]{\mathsf{YSub}(#1)}
\newcommand{\YSL}[2]{\mathsf{YSL}(#1,#2)}
\newcommand{\evalmap}{\mathsf{ev}}

\newcommand{\cutpos}[1]{%
  \raisebox{-0.7ex}{%
    \tikz[baseline=(char.base)]{%
      \node[draw,circle,inner sep=0.4pt,font=\scriptsize]
        (char) {\ensuremath{#1}};%
    }%
  }%
}

\theoremstyle{plain}
\newtheorem{mainthm}{定理}
\renewcommand{\themainthm}{\Alph{mainthm}}
\title{固定有限半群型付けの下での\\
多重文脈自由言語の分布学習：\\
局所群と特徴データの長さ}
\author{Takayuki Kuriyama\\
独立研究者，東京，日本\\
\texttt{growup.kuriyama@gmail.com}}
\date{}

\begin{document}
\maketitle

\begin{abstract}
本稿は，前もって固定した型付け（有限半群の上への準同型）の値が一致するときにのみ構成素を比較するという条件の下で，
多重文脈自由言語の学習に要る正データを調べる．
極限同定は Yoshinaka の分布的枠組みから従う；本稿は特徴標本の総長を標的の文法の大きさと規則の厚さで評価し，
どの型付けでこの長さが多項式になるかを問う．
型付けが局所群，すなわちすべての冪等元 $e$ について $e\mathsf Se$ が群であるような有限半群 $\mathsf S$ の上への写像ならば，
固定したすべてのファンアウトとランクについて，文法の大きさの2乗と規則の厚さの1乗に比例する長さで足り，
これらのクラスは多項式時間・厚データ同定可能である．
証明は，交互に並ぶタプル成分を，非可換群の順序付き積を保ったまま同時に短縮する．
ファンアウト2の非文脈自由の標的は型付けの領域を分離する：一つは局所群型付けのクラスに入るが，
有界な窓と群への準同型を組み合わせたどの型付けのクラスにも入らず，もう一つはそのような型付けのクラスに非可換群についてのみ入る．
十分大きなアルファベット上では，どの型付けの下でも文法の大きさと規則の厚さの積に線形な長さが必要である．
ファンアウト1かつランク2以上では，局所群がちょうど境界になる：
擬多様体がそのすべての型付けで多項式のデータを与えるのは，それが局所群からなるとき，かつそのときに限る；
そうでなければ，クラス全体で特徴標本をもつすべての集合駆動型の学習者について，文法より指数的に長い標本を要する標的が存在する．
半束の上への型付けについては，これには三文字が必要である：二文字上では，そのような型付けはすべてのファンアウトで多項式のデータを与える．
ファンアウト1では，二元の半束の上への型付けが多項式のデータをもつのは，その単位元に写される文字がちょうど一つのとき，かつそのときに限る．
ファンアウト2での境界は語の文字の集合を記録する型付けに帰着し，三文字以上で未解決である；ファンアウト1の下界の族はそこへ持ち上がらないことが証明される．
\end{abstract}

\section{序論}
\label{sec:introduction}

正データ文法推論は，言語を正例のみから極限同定できるのはいつかを問う
\cite{Gold1967,Angluin1980}．
分布学習は，共通の文脈に現れる語，または語のタプルを，構成素の候補として扱う．
それは代入可能な
CFG \cite{ClarkEyraud2007,Yoshinaka2008} に始まり，多重文脈自由文法（MCFG）に拡張された．
MCFG のタプル値の非終端記号は，不連続な依存と交差依存を表す
\cite{SekiEtAl1991,Yoshinaka2009,Yoshinaka2011,Kallmeyer2010}．
Yoshinaka の有限シグネチャの枠組み \cite{Yoshinaka2015General} は，
多項式個の語からなる特徴標本による正データ同定を与える．
多項式個の語でもなお指数的に長くあり得るので，
本稿は有限標本 $K$ をその符号化総長
$\|K\|_+=\sum_{w\in K}(|w|+1)$ で測る．これは
\cite[Section~2.5.2]{EyraudHeinzYoshinaka2016} の厚データの観点に従う．
質的な同定は
Yoshinaka の枠組みから従う；本稿の主題は，標的の MCFG 表示に相対的な特徴標本の長さであり，
問いは次である：\emph{どの比較バイアスについてこの長さは多項式か？}

\paragraph{設定．}
データを見る前に，ファンアウトの上限 $f$ と全射な
有限半群準同型 $s:\Sigma^+\twoheadrightarrow\mathsf S$，すなわち
\emph{型付け}を固定する．再構成は，観測された空でないタプル成分を，
それらの $s$-値が一致するときにのみ比較し，標的クラス $\Cmcf{f}{s}$
は，この制限の下でタプル代入可能な $f$-MCFL からなる（定義~\ref{def:class-cmcf}）．
型付けは \cite{CosteTyping2004} の意味でのクラス単位のバイアスであり，
標的ごとに選ぶものではない；他のバイアスには，制御集合によって導出を制限するもの
\cite{Takada1988,Takada1995} や，因子の比較に用いる証拠を制限するもの
\cite{CosteGaretNicolas2012,CosteNicolas2019} がある．
上界は，観測された二つのタプルを型が一致するときにのみ連結する Yoshinaka の再構成作用素
$\mathcal B_{f,s}$ と，その変種 $\mathcal B^{\le r_{\max}}_{f,s}$（合成規則のランクを高々
$r_{\max}$ に制限したもの）について述べる（定義~\ref{def:reconstruction}）．
標的は既約な良い MCFG $G$（定義~\ref{def:grammar-properties}）で与え，その大きさ $\|G\|$
と規則の厚さ $\bar\theta_G$ で測る．後者は，与えられた規則を用いる最短の文の長さのうち最大のものである．
これは \cite{EyraudHeinzYoshinaka2016} と姉妹原稿
\cite{kuriyamaCFG} の非終端の厚さ
$t_G$（姉妹原稿では $\tau_G$ と書く）と多項式的に同値である．実際，
$t_G\le\theta_G\le2\|G\|\max\{1,t_G\}$ が成り立つ
（補題~\ref{lem:thickness-equivalence}）．

\paragraph{なぜ型付けか．}
型付けは，再構成が同定するクラスを広げる．すでに
ファンアウト2において，$\mathbb Z_2$ の型付けは，Yoshinaka の型なしのクラスの外にある非文脈自由な標的を許し，
その上で型なしの作用素は三語の標本から過汎化する
（命題~\ref{prop:typed-separation}）．
その代価はデータで払われる：特徴標本は，標的が実現するすべての型について，その型の語を含まなければならず，
ある型は長い語によってしか実現されないことがある
（命題~\ref{prop:rare-type}）．本稿の結果は，この代価がどの型付けについて多項式にとどまるかを述べる．

\paragraph{どの型付けが局所群か．}
有限半群 $\mathsf S$ が\emph{局所群}であるとは，$e\mathsf Se$ がすべてのべき等元 $e$ について
群であることをいう．
局所群は\emph{擬多様体} $\mathbf{LG}$
（部分半群，準同型像および有限直積について閉じたクラス）をなし，
有限半群が局所群であるのは，それが二元半束 $U_1=\{1,\zeta\}$ に同型な部分半群をもたないとき，
かつそのときに限る（定理~\ref{thm:lg-characterization}）．
窓（局所自明な半群，$\mathbf{LI}$），有限群（$\mathbf{Grp}$）およびそれらの積の除数
（$\mathbf{LI}\vee\mathbf{Grp}$）は局所群であり，包含は真である：
$\{\mathtt a,\mathtt b\}$ 上で，
$\mathtt{ab}$ の出現回数の偶奇を記録することは局所群型付けであるが，
接頭辞--接尾辞窓と群準同型の組み合わせのどれもそれを決定しない：
各窓の大きさと群準同型について，ある $m,n$ が存在して，
$\mathtt b^m\mathtt a^{n}\mathtt b^m$ と $\mathtt b^{2m}$ はその両方で一致するが，
その偶奇では異なる
（命題~\ref{prop:li-g-strict}）．DFA の遷移準同型は局所群である必要はない：
恒等写像として作用する文字と，少なくとも二つの状態の上の定値写像として作用する文字は，
$U_1$ の複製をなす．

\paragraph{主結果．}
定理~\ref{thm:A} と~\ref{thm:B} では，アルファベット $\Sigma$ とファンアウト
$f$ を固定する．

\begin{mainthm}[局所群]
\label{thm:A}
$r_{\max}\ge1$ とし，$s(\Sigma^+)\in\mathbf{LG}$ と仮定する．すべての
$L\in\Cmcf{f}{s}$ と，既約な良い表示 $G$ であって $L\cap\Sigma^+$ を表し，
ファンアウト $\le f$ かつランク $\le r_{\max}$ のものすべてについて，作用素
$\mathcal B^{\le r_{\max}}_{f,s}$ と $\mathcal B_{f,s}$ は，$L$ に対する共通の
特徴標本であって総長 $O(\|G\|^2\bar\theta_G)$ のものをもつ；
定数は $f$，$r_{\max}$，$s$ および $\Sigma$ のみに依存する
\textup(定理~\ref{thm:local-group-polydata}\textup)．
\end{mainthm}

定理~\ref{thm:A} の型付けの領域は，すでに非文脈自由の標的の上で異なる．
ファンアウト2の非文脈自由言語であって，ある局所群型付け $s$ について $\Cmcf{f}{s}$（$f\ge2$）に入るが，
窓の大きさと群をどう選んでもどの窓--群型付けのクラスにも入らないものがある
（命題~\ref{prop:lg-needed-target}）；また，$S_3$ 上の窓--群型付けのクラスには入るが，
どの窓--可換群型付けのクラスにも入らないものがある
（命題~\ref{prop:nonabelian-needed-target}）．したがって，ファンアウト2の非文脈自由の標的に限っても，
これらのクラスを覆う結果において，窓--群型付けは局所群型付けの代わりにならず，窓--可換群型付けは窓--群型付けの代わりにならない．
これらの分離は型付けの領域の範囲についてのものであり，特定の短縮の証明の必要性やその評価の最適性についてのものではない．

\begin{mainthm}[線形因子]
\label{thm:B}
\begin{enumerate}[label=\textup{(\roman*)},leftmargin=*]
\item $\Sigma$ が少なくとも $4f+1$ 文字をもつか，または $f=1$ かつ
$\Sigma$ が少なくとも2文字をもつとする．このとき，すべての型付け $s$ について，
$\Cmcf{f}{s}$ の有限の標的であって，ランク2の表示 $G$ が最小の既約な良い表示の大きさの
$7f+17$ 倍以内に収まるものが存在し，その特徴標本はすべて長さ $\Omega(\|G\|\bar\theta_G)$ をもつ．これは
$\mathcal B^{\le r}_{f,s}$ \textup($r\ge1$\textup) と，クラス全体で特徴標本をもつ
すべての集合駆動型作用素について成り立つ
\textup(命題~\ref{prop:linear-factor-all-fanout}\textup)．
\item 自明な型付けとランク
$\le r_{\max}$ の既約な良い表示について，作用素 $\mathcal B_{f,s_{\mathrm{triv}}}$ は
Yoshinaka の上界 $O(\|G\|\bar\theta_G)$ \cite{Yoshinaka2011} を満たす；
$r_{\max}\ge2$ かつ \textup{(i)} のアルファベットの条件が成り立つならば，この上界は
\textup{(i)} の標的によって定数倍を除いて達成される
\textup(系~\ref{cor:trivial-exact}\textup)．
\end{enumerate}
\end{mainthm}

\begin{mainthm}[ファンアウト1での境界]
\label{thm:C}
$f=1$ かつ $r_{\max}\ge2$ とする．
\begin{enumerate}[label=\textup{(\roman*)},leftmargin=*]
\item 有限半群の擬多様体 $\mathbf V$ が
$\mathbf{LG}$ に含まれるのは，それが\emph{一様多項式特徴データ長}をもつとき，
かつそのときに限る \textup(定義~\ref{def:uniform-polydata}\textup)．ここで一様多項式特徴データ長とは次のことである：
すべての空でないアルファベット $\Sigma$ と，すべての型付け $s$ であって $s(\Sigma^+)\in\mathbf V$ を満たすものについて，
多項式 $p$ が存在し（それは $s$ と $\Sigma$ には依存してよいが，
標的には依存しない），すべての $L\in\Cmcf{1}{s}$ が
$\mathcal B^{\le r_{\max}}_{1,s}$ に対する長さ高々 $p(\|G\|,\bar\theta_G)$ の特徴標本を，
既約な良い表示 $G$ であって $L\cap\Sigma^+$ を表しランク $\le r_{\max}$ のもののそれぞれについてもつ
\textup(定理~\ref{thm:fanout-one-data}\textup)．
\item $\mathbf{LG}$ の外では失敗は指数的である：三文字上の型付け $s_U$ で
$U_1$ の上へのものについて，標的 $L_n\in\Cmcf{1}{s_U}$ であって，大きさと規則の厚さが
$O(n)$ のランク2の既約な良い表示をもち，その
特徴標本がすべて長さ $3\cdot2^n-1\ge2^{(\|G_n\|+\bar\theta_{G_n})/17}$ 以上をもつものが存在する．これは
$\mathcal B^{\le r_{\max}}_{1,s_U}$ と，クラス全体で特徴標本をもつすべての集合駆動型作用素について成り立つ
\textup(定理~\ref{thm:fanout-one-data}\textup)．したがってデータの長さは，$\mathbf V$ へのすべての型付けについて $O(\|G\|^2\bar\theta_G)$ であるか，
そのいくつかについて $\|G\|+\bar\theta_G$ の指数であるかのどちらかであり，中間の増加は起きない．
\end{enumerate}
\end{mainthm}

したがって，$s(\Sigma^+)\in\mathbf{LG}$ ならば，クラス
$\Cmcf{f}{s}\cap\mathsf{MCFL}(f,r_{\max})$（ファンアウト
$\le f$ かつ規則ランク $\le r_{\max}$ の MCFG の言語）は，すべてのファンアウトで多項式時間・厚データ同定可能であり
\cite{Yoshinaka2008,EyraudHeinzYoshinaka2016}，
ファンアウト1では $\mathbf{LG}$ がちょうどこの性質の擬多様体の境界である
（系~\ref{cor:iptd}）．ファンアウト1の MCFG は文脈自由文法であるので，
定理~\ref{thm:C} は文脈自由言語についての主張である；それは
姉妹原稿の十分条件 $s(\Sigma^+)\in\mathbf{LI}\vee\mathbf{Grp}$
\cite{kuriyamaCFG} を特徴づけに変える．
本稿はこれを，定理~\ref{thm:A} を測る基準として含める：ファンアウト1では，局所群という十分条件は擬多様体の水準で弱められない．

\paragraph{二文字と孤立した文字．}
半束の上への型付けについては，定理~\ref{thm:C}(ii) の三文字を二文字に減らすことはできず，そのうち二つは型 $1$ でなければならない．
型がそれ自身の冪だけで実現され，それを含む他のすべての語が一つの型を共有する文字を\emph{孤立した}文字と呼ぶ（定義~\ref{def:isolated-letter}）．
ファンアウト1では，孤立した文字の型は，すべてのランクの reduced good な表示について，どのアルファベットの上でも，
$\|G\|$ と $\bar\theta_G$ の多項式の長さの代表元をもつ（定理~\ref{thm:two-letter-fanout-one}）．
したがって，$\{\mathtt a,\mathtt b\}$ の上では正の像が半束であるすべての型付けが，またどのアルファベットの上でも型 $1$ の文字がちょうど一つの $U_1$ の上へのすべての型付けが，
ファンアウト1で絶対定数による長さ $O(r_{\max}(\|G\|+2)^{23}\bar\theta_G^{\,8})$ の特徴標本をもつ（系~\ref{cor:two-letter-fanout-one}）；
二文字の上では同じことがすべてのファンアウト $f\ge2$ で成り立つ（系~\ref{cor:two-letter-semilattice}）．
型 $1$ の二つの文字はすでに指数的なデータを強いる（系~\ref{cor:typing-level-obstruction}）ので，
これはファンアウト1での $U_1$ の上への型付けを分類する（系~\ref{cor:u1-classification}）．
これらのデータは局所群の形ではない：対数的な大きさの文法について，それは $\bar\theta_G$ の2乗になりうる（命題~\ref{prop:binary-u1-quadratic}）．
この評価は代表元の長さ $\tau$ についてのものであり，型付きの厚さについてではない．型付きの厚さは，そのような型付けの下で $\Sigma^+$ についてすでに指数的である（命題~\ref{prop:typed-thickness-exponential}）．

\paragraph{残る問い．}
定理~\ref{thm:A} は，定理~\ref{thm:C}(i) の \emph{if} の方向を
すべてのファンアウトで与える．$\mathbf{LG}$ がファンアウト $f\ge2$ でもなお境界であるかどうか
は未解決である（未解決問題~\ref{op:fanout-two}）．
定理~\ref{thm:C}(ii) の族はそれを決定できない：すべての固定した型付けの下で，
その成員のうち有限個を除くすべてはファンアウト2でクラスを離れる
（定理~\ref{thm:doubling-no-lift}）．
問いは内容型付け，すなわち語の文字の集合を記録する型付けに帰着する（命題~\ref{prop:sl-reduction}）；
二文字の上ではそれらはすべてのファンアウトで多項式のデータをもつ（定理~\ref{thm:two-letter-content}）ので，
問いは三文字以上で未解決である．

\paragraph{方針．}
文法を型で細分すると，二元の型付けについてさえ厚さが指数的に増え得る
（命題~\ref{prop:typed-thickness-exponential}）．したがって本稿は
特徴データを元の表示を用いて評価する．所定の型をもつ短いタプルは，
長い部分導出を短いものに置き換えて得る．置き換えた個々の群の値は変わってよいが，
それらを合わせた効果は変わらないようにする．ファンアウト1では，置き換える部分木は一つの連続した
ブロックを占め，接頭辞積に対する鳩の巣論法で足りる．ファンアウト
$f\ge2$ では，部分木は最大 $f$ 個のブロックを異なる成分に送り，それらは
他の部分木のブロックと交互に並ぶ；非可換群では，
同時置換の効果はこの順序に依存する．例えば，二つの部分木が $(x_1,x_2)$ と $(y_1,y_2)$ を書き，
根の成分が $\cdots x_1\cdots y_1\cdots$ と $\cdots y_2\cdots x_2\cdots$ なら，
二つの部分木は二つの積の中で逆の順序に現れ，一方の積を保つ置換の選び方が他方の積を変え得る；
部分木が多ければ，あらゆる交差の型が生じる．鍵となる段階は，
すべての成分の順序付き積を同時に保つ置換の同時選択である
（命題~\ref{prop:block-replacement-interface}）．
局所群では，値はブロックの境界をまたぐ因子にも依存する；
ブロックを互いに遠く離して選び，変化しない境界の文字で埋めることにより，
同じ段階に帰着する（補題~\ref{lem:lg-buffered-owners}）．下界は，
その部分集合がすべてクラスにとどまる有限族
（命題~\ref{prop:linear-factor-all-fanout}）と，
長い語によってしか実現されない型（命題~\ref{prop:rare-type}）を用いる．

\paragraph{先行研究との関係．}
Yoshinaka は $(k,\ell)$-代入可能な CFG の特徴データを
厚さのパラメータで評価し \cite{Yoshinaka2008}，多次元
代入可能な MCFG を正データから学習し \cite{Yoshinaka2009,Yoshinaka2011}，
有限の証拠の個数をもつ有限シグネチャの学習原理を与えた
\cite{Yoshinaka2015General}；質的な結果と，ここで用いる
再構成作用素は彼のものである．有限半群の
構造理論と局所群の表現は古典的である
\cite{Straubing1985,Tilson1987,AlmeidaEtAl2009}．本稿の寄与は，
符号化長の解析と上記の境界の結果である．他の
弱文脈依存の形式体系の正データからの学習は
\cite{OatesEtAl2006} にあり，質問や双対文脈集合を用いた MCFG の学習は
\cite{Yoshinaka2010,ClarkYoshinaka2014,KanazawaYoshinaka2023} にある；
\cite{DuncanEtAl2026Substitution} の MCFL の
代入補題は所属の条件であり，学習の基準ではない．

\paragraph{姉妹 CFG 原稿との関係．}
姉妹原稿 \cite{kuriyamaCFG} はファンアウト1を扱う．要するに，それは
$\mathbf{LI}\vee\mathbf{Grp}$ の型付けについての多項式データを
十分条件として，また $U_1$ の上への一つの型付けについての指数的データを証明する；ここでは
ファンアウト1の境界をちょうど（$\mathbf{LG}$）決定し，
$\mathbf{LG}$ に対する上界と線形因子の下界はすべてのファンアウトで成り立つ．
表~\ref{tab:companion-overlap} はその結果を本稿の結果から区別する（第2列の番号は
姉妹原稿のものを指す）．姉妹原稿のどの結果も前提としては用いない：
その括弧付き倍増族
\cite[Proposition~6.4]{kuriyamaCFG} の代わりに，三文字の変種を用い，それは
付録~\ref{app:bracketed-family} で証明する．また，その表現力についての結果（カウンタと
Dyck の障害，合同文法との比較）は用いない．
姉妹原稿に対応物をもたない結果は表に挙げない：分離する標的
（命題~\ref{prop:lg-needed-target} と~\ref{prop:nonabelian-needed-target}），
$\|G\|\bar\theta_G$ の必要性（命題~\ref{prop:linear-factor-all-fanout}），
孤立文字と二文字（定理~\ref{thm:two-letter-fanout-one}，系~\ref{cor:two-letter-fanout-one}
と~\ref{cor:two-letter-semilattice}，命題~\ref{prop:binary-u1-quadratic}），
およびファンアウト2での冪の移送と内容型付けへの帰着
（定理~\ref{thm:doubling-no-lift}，命題~\ref{prop:sl-reduction}）．

\begin{table}[htbp]
\caption{姉妹 CFG 原稿 \cite{kuriyamaCFG}
（ファンアウト1）の結果と本稿の結果．}
\label{tab:companion-overlap}
\centering
\footnotesize
\setlength{\tabcolsep}{3pt}
\renewcommand{\arraystretch}{1.15}
\begin{tabularx}{\linewidth}{@{}>{\raggedright\arraybackslash}p{0.27\linewidth}>{\raggedright\arraybackslash}p{0.25\linewidth}>{\raggedright\arraybackslash}X@{}}
\hline
主題 & 姉妹原稿 \cite{kuriyamaCFG} & 本稿\\
\hline
型付き再構成，Gold 学習器，型付きの厚さによる証拠の上界
 & 定理~4.2，定理~5.2--5.3，系~5.4；補題~6.2，系~6.3
 & タプル版（第~\ref{sec:learner}~節，
   定理~\ref{thm:data-from-typed-thickness}）；質的な結果は
   Yoshinaka に帰する；新しくない\\
窓および窓--群の上界
 & 命題~3.5，補題~7.1，定理~7.2，補題~7.4，系~7.5
 & $f=1$ では新しくない；$f\ge2$ では新しい：交互に並ぶ
   成分にわたる所有者ごとの置換（命題~\ref{prop:block-replacement-interface}）\\
局所群型付け
 & 扱わない（$\mathbf{LI}\vee\mathbf{Grp}$ のみ；補題~7.6）
 & データ長について新しい，すべての $f$ で（定理~\ref{thm:lg-characterization}
   と~\ref{thm:local-group-polydata}）；$\mathbf{LG}$ の $U_1$ による特徴づけと
   半群の代数は古典的である \cite{AlmeidaEtAl2009}\\
$U_1$ の上への型付けの下での倍増族
 & 命題~6.4，系~6.6，付録~C（五文字，括弧）
 & 同じ上界・下界をもつ三文字の変種，ここで証明する
   （命題~\ref{prop:three-letter-family}）\\
稀な型による下界
 & 命題~6.4(iii)--(iv)，一つの語について
 & すべての実現型とすべての $f$（命題~\ref{prop:rare-type}）\\
ファンアウト1での擬多様体の境界
 & $\mathbf{LI}\vee\mathbf{Grp}$ に対する十分性（系~7.5）；必要性は主張しない（注意~7.7）
 & 新しい（定理~\ref{thm:fanout-one-data}）\\
$U_1$ の上への個々の型付け
 & 一つの型付け（命題~6.4）
 & 新しい：ファンアウト1で，型 $1$ の文字がちょうど一つのとき，かつそのときに限り多項式のデータ
   （系~\ref{cor:u1-classification}）；文字が互いに素な型 $1$ の二つの語があれば指数的なデータ
   （系~\ref{cor:typing-level-obstruction}）\\
型付きのクラスと型なしのクラス
 & $f=1$ での分離（命題~6.4 と~8.4）
 & $f\ge2$ で，Yoshinaka のクラスの外にある非文脈自由な成員
   （命題~\ref{prop:typed-separation}）\\
\hline
\end{tabularx}
\end{table}

\paragraph{主張の範囲．}
すべての上界・下界は，述べた作用素に対する特徴標本の符号化総長に関するものであり，
述べた表示に相対的である；それらは文法の固有の大きさも Gold の収束の段階も評価しない．
型付けは標的を見る前に固定する；それを細分するとクラスは大きくなるが，
再構成作用素が変わるので，本稿はそれを推論も最適化もしない．
第~\ref{sec:fanout-one}~節の終わりの表~\ref{tab:scope} は，各主張について，
固定するパラメータ，作用素および付帯条件を挙げ，
第~\ref{sec:conclusion}~節は指数を論じる．

\paragraph{構成．}
第~\ref{sec:prelim}~節は設定を固定する．
第~\ref{sec:learner}~節は再構成作用素を思い出し，型付きのクラスを型なしのクラスから分離する．
第~\ref{sec:char-data}~節は特徴データを型付きの厚さと代表元の長さに帰着させる．
第~\ref{sec:upper}~節は定理~\ref{thm:A} を証明して分離する標的を与え，第~\ref{sec:linear}~節は定理~\ref{thm:B} を証明し，
第~\ref{sec:fanout-one}~節は定理~\ref{thm:C} を証明して，$U_1$ の上への型付けを分類して，ファンアウト1での孤立した文字と二文字を扱う．
第~\ref{sec:fanout-two}~節はファンアウト2を扱う：持ち上がらないこと，内容型付けへの帰着，二文字．
第~\ref{sec:conclusion}~節は未解決問題と指数を論じる．付録は技術的な証明を含む．

\section{準備}
\label{sec:prelim}

全体を通じて，$\Sigma$ は有限アルファベット，$\eps$ は空語，
$\Sigma^+:=\Sigma^*\setminus\{\eps\}$ とする．$\Sigma=\emptyset$ ならば，言語は
$\emptyset$ と $\{\eps\}$ のみであり，直接扱う．
有限の $K\subseteq\Sigma^*$ について $\|K\|_+:=\sum_{w\in K}(|w|+1)$ とおき，
タプル $\tuple u=(u_1,\ldots,u_d)$ について $\|\tuple u\|_1:=\sum_i|u_i|$ とおく．以下の
すべての多項式の上界・下界は，それが個数に関するものか
$\|\cdot\|_+$ に関するものかを明記する．次の表は複数の節にわたって用いる記号を定める；それ以外のもの
（滑動因子ラベル付け，補助的な準同型，詰め物の長さ）はすべて，
それが現れるところで定義する．
\begin{center}
\footnotesize
\renewcommand{\arraystretch}{1.12}
\begin{tabularx}{\linewidth}{@{}lX@{}}
\hline
記号 & 意味\\
\hline
$G$, $G_{\mathrm{in}}$ & （既約な良い）MCFG；任意の標準の入力 MCFG\\
$B,C,X,Y$ & 非終端記号；$X,Y$ は指定した箇所では型付き非終端記号も表す\\
$\|G\|$, $\theta_G$, $t_G$ & 符号化の大きさ（定義~\ref{def:standard-mcfg}），規則の厚さ，非終端記号の厚さ\\
$s,\ \widehat s,\ \mathsf S$ & 正の型付け，その新単位元による完備化，その正の像\\
$\preceq$ & 核による細分（定義~\ref{def:semigroup-typing}），$\Sigma^*$ 上のモノイド準同型の間でも用いる\\
$\mathbb G$, $\mathbb H$ & 有限群（$\mathbb G$ は型付けの群成分，$\mathbb H$ は Rees 行列の群または補助的な群）\\
$g$ & 有限群へのモノイド準同型 $\Sigma^*\to\mathbb G$\\
$\mathsf T$, $\mathsf N$ & 有限半群；補助的な準同型 $\lambda$ の行き先の有限モノイド\\
$\lambda$, $\lambda_{k,\nu}$ & $\widehat s$ を細分する補助的な有限モノイド準同型（第~\ref{sec:char-data}~節）；滑動因子準同型（定義~\ref{def:local-group-typing}）\\
$T$ & 導出木\\
$n$ & 文法の族の添字\\
$t_{\mathrm{typ}}$, $\tau$, $\mu_m$ & 最短の型付き産出，代表長，型 $m$ の最短の語（第~\ref{sec:char-data}~節）；これらはいずれも $G$ の厚さではない\\
$\mathsf S_L$ & 実現された正の型の集合 $s(L\cap\Sigma^+)$\\
$\mathcal B_{f,s}$, $\mathcal B^{\le r}_{f,s}$ & 上限のない再構成作用素とランクに上限のある再構成作用素（定義~\ref{def:reconstruction}）\\
\hline
\end{tabularx}
\end{center}

\subsection{正データ学習}

\begin{definition}[テキスト，TxtEx 同定，保守性]
休止記号 $\pausesym\notin\Sigma$ を固定する．
$L\subseteq\Sigma^*$ に対する\emph{テキスト}とは，$L\cup\{\pausesym\}$ 上の無限列であって，
$L$ のすべての元が現れるものである；有限の接頭辞の\emph{内容}とは，それに現れる
$\Sigma^*$ の元の集合である．\emph{学習器}とは，有限の接頭辞から文法への計算可能な写像
である．それが $L$ を \emph{TxtEx 同定する}~\cite{Gold1967}
とは，$L$ に対するすべてのテキスト上で，その出力がある段階以降，一つの固定した文法
$H$ で $L(H)=L$ なるものであることをいう；それがクラスを TxtEx 同定するとは，各成員についてそうであることをいう．
それが\emph{保守的}であるとは，その仮説を，現在の仮説が生成しないデータ
$w\in\Sigma^*$ を受けたときにのみ変えることをいう．仮説
$H_0$ で $L(H_0)=\emptyset$ なるものを固定する．
\end{definition}

\begin{definition}[集合駆動型作用素と特徴標本]
\label{def:characteristic-sample}
\emph{集合駆動型再構成作用素} $\mathcal B$ は，各有限の
$K\subseteq\Sigma^*$ を，文法 $\mathcal B(K)$ で，集合
$K$ のみに依存するものに写す；それが\emph{標本無矛盾}であるとは，$K\subseteq L(\mathcal B(K))$ がすべての
$K$ について成り立つことをいう．有限の $\mathcal F\subseteq L$ が\emph{
$\mathcal B$ と $L$ に対する特徴標本}であるとは，$L(\mathcal B(K))=L$ が，すべての有限の $K$ で
$\mathcal F\subseteq K\subseteq L$ なるものについて成り立つことをいう（\cite{deLaHiguera1997} を参照）．
\end{definition}

この定義は生成される言語のみを制約する：$\mathcal B(K)$ は，
$K$ が大きくなるとき，構文的に安定する必要はない．

\subsection{固定有限半群型付け}
\label{subsec:semigroup-typing}

\begin{definition}[型付け，新単位元による完備化，細分]
\label{def-monoid-morphism}
\label{def:semigroup-typing}
\label{def:h-refinement}
\emph{型付け}とは，全射な半群準同型
$s:\Sigma^+\twoheadrightarrow\mathsf S$ で有限半群の上へのものであり，
$\mathsf S$ の乗積表と文字の像によって与えられ，したがって
$s(w)$ は時間 $O(|w|)$ で計算可能である；$\mathsf S$ をその\emph{正の
像}という．（より大きな有限半群への準同型は，像への
余制限で置き換える．）$\mathsf S^\sharp:=\mathsf
S\uplus\{\mathbf1_\sharp\}$ を，\emph{新しい}単位元を添加して得られるものとする．
これは $\mathsf S$ がすでに単位元をもつ場合でもそうする．また
$\widehat{s}:\Sigma^*\twoheadrightarrow\mathsf S^\sharp$ を，モノイド準同型であって
$\widehat{s}(\eps)=\mathbf1_\sharp$ かつ $\widehat{s}|_{\Sigma^+}=s$ なるものとする．
$\tuple w\in(\Sigma^*)^d$ について
$\widehat{s}^{(d)}(\tuple w):=(\widehat{s}(w_1),\ldots,\widehat{s}(w_d))$ とおく；
$(\Sigma^+)^d$ 上ではこれは $s^{(d)}(\tuple w)$ である．型付け $s,s'$ について，$s\preceq
s'$ と書くのは $\ker s'\subseteq\ker s$ が $\Sigma^+$ 上で成り立つときとし，$(s,s')$ を型付け
$w\mapsto(s(w),s'(w))$ で $\mathsf S\times\mathsf T$ におけるその像の上へのものとする．ここで $\mathsf T:=s'(\Sigma^+)$ である；したがって
$s,s'\preceq(s,s')$．一元の型付けを $s_{\mathrm{triv}}$ と書く．
\end{definition}

新単位元による完備化は，文法の雛形（空の終端記号の隙間を含む）を評価し，
第~\ref{sec:refinement}~節の細分においてタプルの産出に型を付ける；その新しい単位元は正の語の型ではない．
$\mathbf{LG}$ のような代数的条件は $\mathsf S$ に関するものであり，
$\mathsf S^\sharp$ に関するものではない；実際，$\mathsf S^\sharp$ は決して局所群ではない．なぜなら，
$\{\mathbf1_\sharp,e\}\cong U_1$ が任意のべき等元 $e$（$\mathsf S$ の極小イデアルの）について成り立つからである．有限モノイド準同型 $q:\Sigma^*\to\mathsf N$ は
正の型付け $q|_{\Sigma^+}$ を定める：
補題~\ref{lem:semigroup-transfer}(i) により，それはその型付けの新単位元による完備化と同じ標的クラスと
集合駆動型再構成の規則を与える．
補題~\ref{lem:semigroup-transfer}(ii) はさらに，\emph{良い}文法についてそれらの
刈り込んだ産出型細分を同一視する．この最後の
同一視は，空の成分の産出をもつ任意の標準の文法については主張しない．

\begin{definition}[接頭辞--接尾辞型付け]
\label{def:prefix-suffix-typing}
$k,\ell\ge0$ について $m_0:=\max\{1,k+\ell\}$ とおき，$|w|\ge m_0$ について
$\operatorname{pref}_k(w)$，$\operatorname{suf}_\ell(w)$ を，示された長さの接頭辞と
接尾辞とする（長さ0の窓は $\eps$ である）．次のように定める．
\[
 \widehat{s}_{k,\ell}(w)=
 \begin{cases}
   w, & |w|<m_0,\\
   \langle\operatorname{pref}_k(w),\operatorname{suf}_\ell(w)\rangle,
      & |w|\ge m_0,
 \end{cases}
\]
ここで短い語と対には互いに素な標識を付ける．また
$M_{k,\ell}:=\widehat{s}_{k,\ell}(\Sigma^*)$ とおく．短い $x,y$ と対
$\langle u,v\rangle,\langle u',v'\rangle$ について次のように定める．
\[
\begin{aligned}
 x\cdot y&:=\widehat{s}_{k,\ell}(xy), &
 x\cdot\langle u,v\rangle&:=\langle\operatorname{pref}_k(xu),v\rangle,\\
 \langle u,v\rangle\cdot y&:=\langle u,\operatorname{suf}_\ell(vy)\rangle, &
 \langle u,v\rangle\cdot\langle u',v'\rangle&:=\langle u,v'\rangle.
\end{aligned}
\]
$s_{k,\ell}:=\widehat{s}_{k,\ell}|_{\Sigma^+}$ とおく；型付け $s$ であって，
$s\preceq s_{k,\ell}$ がある $k,\ell\ge0$ について成り立つものを\emph{窓型付け}という．
\end{definition}

$x,y\in\Sigma^*$ の長さについての場合分けにより
$\widehat{s}_{k,\ell}(xy)=\widehat{s}_{k,\ell}(x)\cdot\widehat{s}_{k,\ell}(y)$ を得る．
$\widehat{s}_{k,\ell}$ は $M_{k,\ell}$ の上への写像であるので，積は結合的であり，
$\widehat{s}_{k,\ell}(\eps)$ は単位元である．したがって $\widehat{s}_{k,\ell}$ は
有限モノイド準同型である．短い語の標識は $\eps$ を
$\widehat{s}_{k,\ell}(\Sigma^+)$ から分離するので，$\widehat{s}_{k,\ell}$ は
完備化 $\widehat{s_{k,\ell}}$ と同型である．

\subsection{多重文脈自由文法}

\begin{definition}[標準の MCFG 表示と大きさ]
\label{def:standard-mcfg}
MCFG \cite{SekiEtAl1991,Kallmeyer2010} とは
$G=(V,\Sigma,P,S,\fan)$ であって，$V$ は有限，$S\in V$，$\fan(S)=1$ かつ
$\fan(A)\ge1$ なるものである．ランク $r\ge0$ の規則とは
$A\to\rho(B_1,\ldots,B_r)$ であって，雛形
$\rho=(\alpha_1,\ldots,\alpha_{d_0})$，$d_0=\fan(A)$，$d_j=\fan(B_j)$ をもち，ここで
各 $\alpha_i$ は
$\Gamma_\rho:=\Sigma\uplus\Xi_1\uplus\cdots\uplus\Xi_r$ 上の語であり，
$\Xi_j:=\{\xi_j^1,\ldots,\xi_j^{d_j}\}$ であり，すべての変数は
$\alpha_1\cdots\alpha_{d_0}$ に高々一回現れる（\emph{線形性}）．$\tuple u_j\in(\Sigma^*)^{d_j}$ について，
$\rho(\tuple u_1,\ldots,\tuple u_r)$ は $u_{j,\kappa}$ を
$\xi_j^\kappa$ に代入する；$r=0$ のとき規則を\emph{定数規則}という．タプル
言語 $L_A(G)\subseteq(\Sigma^*)^{\fan(A)}$ は，規則の下で閉じた最小の族をなし，
$L(G):=\{w:(w)\in L_S(G)\}$ とする．
$A\Rightarrow_G^*\tuple u$ を\emph{終端導出}（$A$ を根とする有限の導出
木で，下から上へ評価するもの）について書く；$S$ からのものを\emph{成功する}導出という．
$G$ が $f$-MCFG であるとは，$\fan(A)\le f$ がすべての $A$ について成り立つことをいい，規則ランクの上限
$r_{\max}$ をもつとは，すべてのランクが高々 $r_{\max}$ であることをいう．このような表示で，
それ以上の制限のないものを\emph{標準}という．その\emph{大きさ}は
\[
 \|G\|:=|V|+\sum_{(A\to\rho(B_1,\ldots,B_r))\in P}
 \Bigl(1+r+\textstyle\sum_{i}|\alpha_i|\Bigr),
\]
であり，ここで $|\alpha_i|$ は終端記号と変数の出現を数える．すべての
上界・下界はこの記号単位の大きさを用いる；二進の
非終端記号名によるビット単位の符号化は，それを高々因子 $O(\log\|G\|)$ だけ変える．
\end{definition}

$\mathsf{MCFL}(f,r)$ は，ファンアウト $\le f$ かつ
規則ランク $\le r$ の MCFG の言語のクラスである；$f\text{-}\mathrm{MCFL}=\bigcup_{r<\infty}\mathsf{MCFL}(f,r)$
であり，言語の最小の規則ランクは，そのファンアウト $\le f$ の
すべての表示にわたってとる．等しいファンアウト $d$ の非終端記号について，$A\to B$ は
$A\to(\xi_1^1,\ldots,\xi_1^d)(B)$ の略記であり，$A\to\tuple c$ は定数規則の略記である．規則の
変数は常に $\xi_j^\kappa$ である；例における終端記号は
\texttt{typewriter} 書体で書く．ファンアウトとランクは異なる資源である．固定したファンアウトが少なくとも
2のとき，規則ランクの階層は一般には潰れないが，隣り合う
水準が一致することはある；特に
$\mathsf{MCFL}(2,2)=\mathsf{MCFL}(2,3)$
\cite{RambowSatta1999,GomezRodriguezEtAl2009}．これに対し，ファンアウト1の
CFG は二分化できる．したがって以下の再構成は，
二分標準形も既知のランクの上限も仮定しない．

\begin{definition}[文法の性質]
\label{def:grammar-properties}
規則が\emph{非削除}であるとは，すべての変数がその
雛形にちょうど一回現れることをいう；非削除の規則が\emph{非置換}であるとは，各 $j$ について
変数 $\xi_j^1,\ldots,\xi_j^{d_j}$ がこの順に
$\alpha_1\cdots\alpha_{d_0}$ に現れることをいい，\emph{非併合}であるとは，さらにどの $\alpha_i$ も
因子 $\xi_j^\kappa\xi_j^{\kappa+1}$ を含まないことをいう．
\cite[Section~2.3]{Yoshinaka2011} に従い，規則が \emph{$\eps$-free}（そこでは「$\lambda$-free」）
であるとは，どの $\alpha_i$ も空でないことをいう．規則，雛形または文法が\emph{良い}とは，それが
$\eps$-free，非削除，非置換かつ非併合であることをいう．非終端記号 $A$
が\emph{到達可能}であるとは，鎖 $S=A_0,\ldots,A_m=A$ であって，各
$A_{i+1}$ が $A_i$ に対するある規則の子であるものが存在することをいい，\emph{生産的}であるとは
$L_A(G)\ne\emptyset$ であることをいう；$G$ が\emph{既約}であるとは，すべての非終端記号がその両方であることをいう．
\end{definition}

$\eps$-free な文法では，導出されるすべてのタプルのすべての成分は空でない
（導出に関する帰納法）．$\eps$-free 性は雛形の成分に関するものであり，
$\eps\notin L(G)$ とは別のものである．

\subsection{タプル文脈と代入可能性}

\emph{アリティ $d$ の文文脈}とは，語
$E=v_0\blank_1v_1\cdots v_{d-1}\blank_dv_d$ であって，各穴が一回現れ，
$v_0,v_d\in\Sigma^*$ かつ $v_1,\ldots,v_{d-1}\in\Sigma^+$ なるものである；$\mathcal E_d$ は
これらの\emph{許容}文脈の集合である \cite[Section~3.1]{Yoshinaka2011}．
$E[\tuple w]$ は $\blank_i$ を $w_i$ で埋めたものであり，$|E|_{\mathrm{fix}}:=\sum_{i=0}^d|v_i|$ とし，
$E$ が $\tuple w$ を $L$ において\emph{受理する}とは $E[\tuple w]\in L$ であることをいう．
\emph{$d$-タプル分布}（$\tuple x\in(\Sigma^*)^d$ の）は
$\D_L^{(d)}(\tuple x):=\{E\in\mathcal E_d:E[\tuple x]\in L\}$ である；ファンアウト1では
それは両側の分布 $\{(u,v):uxv\in L\}$ である．

\begin{definition}[タプル出現]
\label{def:tuple-occurrence}
$w=a_1\cdots a_n$ について，切れ目の位置に番号を付けて
$\cutpos{0}a_1\cutpos{1}\cdots a_n\cutpos{n}$ とし，
$w[i:i']:=a_{i+1}\cdots a_{i'}$ と書く．\emph{アリティ $d$ のタプル出現}（$w$ における）とは
列 $0\le i_1<i_1'<i_2<i_2'<\cdots<i_d<i_d'\le n$ である；それは
$\tuple x=(w[i_1:i_1'],\ldots,w[i_d:i_d'])\in(\Sigma^+)^d$ と文脈
$E\in\mathcal E_d$（これらの区間を
$\blank_1,\ldots,\blank_d$ で置き換えて得られるもの）とを定め，$E[\tuple x]=w$ となる．有限の $K$ について，
$\mathrm{Obs}_d(K):=\{\tuple x\in(\Sigma^+)^d:E[\tuple x]\in K\text{ for some
}E\in\mathcal E_d\}$ は，このような出現から生じるタプルの集合であり，
固定したファンアウトの上限 $f\ge1$ について
$\mathrm{Obs}(K):=\bigcup_{1\le d\le f}\mathrm{Obs}_d(K)$ とおく．
\end{definition}

\begin{definition}[$(f,s)$-タプル代入可能性と標的クラス]
\label{def:tuple-substitutable}
\label{def:class-cmcf}
$f\ge1$ とし，$s$ を型付けとする．言語 $L\subseteq\Sigma^*$ が
\emph{$(f,s)$-タプル代入可能}であるとは，すべての $1\le d\le f$ とすべての
$\tuple x,\tuple y\in(\Sigma^+)^d$ について次が成り立つことをいう．
\[
 s^{(d)}(\tuple x)=s^{(d)}(\tuple y)\ \text{ and }\
 \D_L^{(d)}(\tuple x)\cap\D_L^{(d)}(\tuple y)\ne\emptyset
 \ \Longrightarrow\
 \D_L^{(d)}(\tuple x)=\D_L^{(d)}(\tuple y).
\]
$\Cmcf{f}{s}$ は $(f,s)$-タプル代入可能な $f$-MCFL のクラスである．
$\tuple x,\tuple y\in(\Sigma^+)^d$ について $\tuple x\equiv_{L,s}^d\tuple y$（または
$\tuple x\equiv_{L,s}\tuple y$）と書くのは，$s^{(d)}(\tuple x)=s^{(d)}(\tuple y)$ かつ
$\D_L^{(d)}(\tuple x)=\D_L^{(d)}(\tuple y)$ であるときとする；これは同値関係である．
\end{definition}

空の成分は決して比較しない．これは
\cite[Section~3.1 and footnote~3]{Yoshinaka2011} に従う（CFG については \cite{Clark2013} を参照）．
また $\eps\in L$ は開始記号で扱う；そうしなければすでに
$\mathtt a^+$ が成り立たなくなる．なぜなら $\mathtt a$ と $\eps$ は文脈を共有するが，
その分布は共有しないからである．バイアスは $s$ と $f$ の両方からなる：
\begin{enumerate}[label=\textup{(\alph*)},leftmargin=*]
\item $s\preceq s'$ ならば，等しい $s'$-型は等しい $s$-型であるので，
$\Cmcf{f}{s}\subseteq\Cmcf{f}{s'}$；特に
$\Cmcf{f}{s_{\mathrm{triv}}}\subseteq\Cmcf{f}{s}$．これはクラスを比較している
だけである：第~\ref{sec:learner}~節の作用素 $\mathcal B_{f,s}$ と $\mathcal B_{f,s'}$ は
一般には異なる（それらは
$\ker s=\ker s'$ のとき一致する）．
\item $\Cmcf{f+1}{s}\cap f\text{-}\mathrm{MCFL}\subseteq\Cmcf{f}{s}$ であるが，
クラスは $f$ について単調でない：自明な型付けについて
$\{w\mathtt{\#}w^R:w\in\{\mathtt a,\mathtt b\}^*\}$ は
$1$-次元代入可能であるが $2$-次元代入可能でない
\cite[Section~3.1, Example~5]{Yoshinaka2011}．
\item すべての $\mathsf{Acc}\subseteq\mathsf S^\sharp$ について，
$\widehat{s}^{-1}(\mathsf{Acc})\in\Cmcf{f}{s}$：それは正則であり，
成分ごとに等しい型は，すべての文脈で等しい全体の型を与える．
\end{enumerate}

\begin{proposition}[共通部分と可換群によるふるい分け]
\label{prop:abelian-coset-filter}
$f$ と $s$ を固定する．
\begin{enumerate}[label=(\roman*),leftmargin=*]
\item $L_1,L_2$ が $(f,s)$-タプル代入可能ならば，$L_1\cap L_2$ もそうである．
\item $L$ が $(f,s)$-タプル代入可能であり，$\mathsf A$ が可換群，
$g:\Sigma^*\to\mathsf A$ がモノイド準同型，$\mathsf H\le\mathsf A$ かつ
$\gamma\in\mathsf A$ ならば，$L\cap g^{-1}(\gamma+\mathsf H)$ は
$(f,s)$-タプル代入可能である．
\item $L\in\Cmcf{f}{s}$，$\mathsf{Acc}\subseteq\mathsf S^\sharp$ であり，かつ (ii) において
$\mathsf A$ が有限ならば，
$L\cap\widehat{s}^{-1}(\mathsf{Acc})\cap g^{-1}(\gamma+\mathsf H)\in\Cmcf{f}{s}$．
\end{enumerate}
無限の $\mathsf A$ については，(ii) の共通部分は
$\Cmcf{f}{s}$ に属する，それが $f$-MCFL であるときにはいつでも．
\end{proposition}

\begin{proof}
(i) 等しい型の $\tuple x,\tuple y$ を $L_1\cap L_2$ において受理する文脈は，
各 $L_i$ において両方を受理するので，それらの分布は各 $L_i$ において一致し，したがって
共通部分において一致する．(ii) ある文脈 $E$ が
$\tuple x,\tuple y$ を $g^{-1}(\gamma+\mathsf H)$ において受理するならば，
$\delta:=\sum_i(g(y_i)-g(x_i))\in\mathsf H$ であり，可換性により
$g(E'[\tuple y])-g(E'[\tuple x])=\delta$ がすべての
$E'\in\mathcal E_d$ について成り立つので，剰余類への所属は両方向に移る．したがって
$g^{-1}(\gamma+\mathsf H)$ は $(f,s_{\mathrm{triv}})$-，したがって
$(f,s)$-タプル代入可能である；(i) を適用する．(iii) 言語
$\widehat{s}^{-1}(\mathsf{Acc})$ は上の項目 (c) により $(f,s)$-タプル代入可能であるので，
(i) と (ii) により共通部分もそうである；それは $f$-MCFL と
正則言語の共通部分であるので，$f$-MCFL である \cite{SekiEtAl1991}．
\end{proof}

(ii) と (iii) の分解は，クラスの成員を構成するための道具である；
それらは再構成には与えない．

\section{再構成と質的学習}
\label{sec:learner}

この節は，後で解析する再構成作用素を固定し，
その質的な保証を記録する．そのいずれも新しいとは主張しない．
固定した $f$，ランクの上限および型付けについて，本稿の作用素は
\cite{Yoshinaka2015General} における有限シグネチャ上の代入可能な言語に対する分布的再構成の型に
従う．本稿は
\cite[Example~3]{Yoshinaka2015General} の MCFG シグネチャを，タプルのソートを $s$-型で細分して用いる
（\cite[Section~8]{Yoshinaka2015General} を参照）．関連する質的な
結果は次のとおりである．
\begin{center}
\footnotesize
\begin{tabular}{@{}lll@{}}
\hline
本稿 & 関連する先行結果 & 本稿での扱い\\
\hline
命題~\ref{prop:reconstruction-soundness}（健全性） & \cite[Lemma~2]{Yoshinaka2015General} & ---\\
命題~\ref{prop:completeness}（完全性） & \cite[Lemma~3]{Yoshinaka2015General} & 同じ型の任意の代表\\
定理~\ref{thm:exact-reconstruction}（正確な再構成） & \cite[Lemmas~2--3]{Yoshinaka2015General} & 命題~\ref{prop:completeness} と同様\\
系~\ref{cor:gold-explanatory}（同定） & \cite[Theorem~1]{Yoshinaka2015General} & クラス全体で非有界なランク\\
\hline
\end{tabular}
\end{center}
この表は類似物を記録するのであって，文字どおりの同一性を記録するのではない：この
節のすべての主張は，具体的な文法について
付録~\ref{app:qualitative-construction-details}
と~\ref{app:typed-witness-aux-proofs} で直接証明し，シグネチャの
枠組みへの埋め込みは行わない．$\Cmcf{f}{s}$ の全体を同定することは，その成員が非有界な
最小ランクをもつので，標本に依存する上限 $\|K\|_+$ を用い，これは
\cite{Yoshinaka2015General} の有界アリティの仮定の外にある．固定した有限
シグネチャについて，その論文は特徴
標本の語の\emph{個数}を文法について多項式に評価する；それはそれらの語の総
\emph{長さ}を評価しない．それが本稿の主題である．

\subsection{集合駆動型再構成}

雛形と文脈の間のすべての相互作用は一つの観察に基づく．

\begin{lemma}[誘導される子の文脈]
\label{lem:filling-identity}
$\rho=(\alpha_1,\ldots,\alpha_{d_0})$ をランク $r$ の良い雛形で，子の
アリティが $d_1,\ldots,d_r$ のものとし，$E\in\mathcal E_{d_0}$ とし，$1\le j\le r$ とし，
$\tuple u_i\in(\Sigma^+)^{d_i}$ を $i\ne j$ について固定する．$E_j$ を，
$E[(\alpha_1,\ldots,\alpha_{d_0})]$ において $\xi_j^\kappa$ を
$\blank_\kappa$ で，他のすべての変数 $\xi_i^{\kappa'}$ を $u_{i,\kappa'}$ で置き換えて得られる語とする．このとき
$E_j\in\mathcal E_{d_j}$ であり，
$E_j[\tuple x]=E[\rho(\tuple u_1,\ldots,\tuple u_{j-1},\tuple x,
\tuple u_{j+1},\ldots,\tuple u_r)]$ がすべての
$\tuple x\in(\Sigma^*)^{d_j}$ について成り立つ．
\end{lemma}


\begin{lemma}[合成は同値性を保つ]
\label{lem:witnessed-composition}
$L\subseteq\Sigma^*$ とし，$\rho$ をランク $r$ の良い雛形とし，
$\tuple x_j,\tuple u_j\in(\Sigma^+)^{d_j}$ が
$\tuple u_j\equiv_{L,s}\tuple x_j$ をすべての $j$ について満たすとする．このとき
$\rho(\tuple u_1,\ldots,\tuple u_r)\equiv_{L,s}\rho(\tuple x_1,\ldots,\tuple x_r)$．
\end{lemma}

\begin{definition}[合成証拠]
\label{def:composition-witness}
$K$ を有限とし，$r\ge1$ とする．\emph{ランク $r$ の合成証拠}（$K$ における）
とは三つ組 $(\tuple z,\rho,(\tuple u_1,\ldots,\tuple u_r))$ であって，$\rho$
がランク $r$ の良い雛形で，出力アリティ $d_0\le f$ と子のアリティ
$d_1,\ldots,d_r\le f$ をもつものであり，$\tuple z\in\mathrm{Obs}_{d_0}(K)$，
$\tuple u_j\in(\Sigma^+)^{d_j}$，かつ
$\tuple z=\rho(\tuple u_1,\ldots,\tuple u_r)$ なるものである．
\end{definition}

同値な言い方をすれば，合成証拠とは，
観測された $\tuple z$ の成分を，終端記号の隙間と，変数でラベル付けした空でない区間とに
切り分けたもので，良い雛形の条件に従うものである．$E[\tuple z]\in K$ ならば，
補題~\ref{lem:filling-identity} は $E_j[\tuple u_j]=E[\tuple z]\in K$ を
$E_j$ が許容であるという形で与えるので，すべての $\tuple u_j\in\mathrm{Obs}(K)$ である．子は
$\tuple z$ の互いに素で空でない部分を占めるので，
\begin{equation}
\label{eq:sample-rank}
 r\le\textstyle\sum_j\|\tuple u_j\|_1\le\|\tuple z\|_1\le\max_{w\in K}|w|
 \le\|K\|_+ .
\end{equation}

\begin{definition}[上限付きと上限なしの再構成作用素]
\label{def:reconstruction}
$f$ と型付け $s$ を固定する．$r_{\max}\ge1$ について
$\mathcal B_{f,s}^{\le r_{\max}}(K)=\widehat G_{\le r_{\max}}(K)$ を，
開始記号 $\widehat S$，非終端記号 $[\tuple x]$（アリティ
$d$，各 $\tuple x\in\mathrm{Obs}_d(K)$，$1\le d\le f$ について一つ）および次の規則をもつ文法とする．
\begin{enumerate}[leftmargin=*]
\item[\textnormal{(Start)}] $\widehat S\to[(w)]$（$w\in K\cap\Sigma^+$ について），および
      $\widehat S\to\eps$（$\eps\in K$ のとき）；
\item[\textnormal{(Const)}] $[\tuple x]\to\tuple x$（すべての
      $\tuple x\in\mathrm{Obs}(K)$ について）；
\item[\textnormal{(Comp)}] $[\tuple z]\to\rho([\tuple u_1],\ldots,[\tuple u_r])$
      （ランク $1\le r\le r_{\max}$ のすべての合成証拠について）；
\item[\textnormal{(Link)}] $[\tuple x]\to[\tuple y]$（
      $\tuple x,\tuple y\in\mathrm{Obs}_d(K)$ であって
      $s^{(d)}(\tuple x)=s^{(d)}(\tuple y)$ かつ $E[\tuple x],E[\tuple y]\in K$
      がある $E\in\mathcal E_d$ について成り立つものについて）．
\end{enumerate}
$K\ne\emptyset$ について $\mathcal B_{f,s}(K):=\widehat G(K):=
\widehat G_{\le\|K\|_+}(K)$ とおく；$\mathcal B_{f,s}(\emptyset)$ は規則をもたない．
固定した符号化の順序により，両方の作用素は決定的かつ集合駆動型となる．
\end{definition}

\eqref{eq:sample-rank} により，$\mathcal B_{f,s}(K)$ は，$K$ におけるすべての
合成証拠について Comp 規則を含む．両方の作用素は標本無矛盾である．それは
$\widehat S\to[(w)]\to(w)$ と $\widehat S\to\eps$ による．開始以外のすべての雛形
は $\eps$-free であるので，$\widehat S\Rightarrow^*\eps$ は規則
$\widehat S\to\eps$ によってのみ起こる．自明な型付けと $\eps\notin K$ について，Comp/Const，Link
および Start は
\cite[Section~4.1]{Yoshinaka2011} の Types~I--III の規則に対応し，$\mathcal B_{f,s}$ は
\cite[Corollary~2]{Yoshinaka2011} の非有界ランクの学習器
$\mathcal A(p,*)$ の作用素に対応する．

\subsection{産出型細分と証拠}
\label{sec:refinement}

\begin{definition}[産出型細分]
\label{def:output-refinement}
$G=(V,\Sigma,P,S,\fan)$ を有限の非削除 MCFG とし，$s$ を型付けとする．
雛形 $\rho$ と $\mathbf t_j\in(\mathsf S^\sharp)^{d_j}$ について，
$\out_\rho(\mathbf t_1,\ldots,\mathbf t_r)\in(\mathsf S^\sharp)^{d_0}$ を，その
第 $i$ 成分が $\alpha_i$ の像であるものとする．ここで像はモノイド準同型
$\Gamma_\rho^*\to\mathsf S^\sharp$ で $a\mapsto\widehat{s}(a)$ と
$\xi_j^\kappa\mapsto(\mathbf t_j)_\kappa$ を拡張するものによる．\emph{完全な産出型
細分} $G^{\widehat{s}}$ は，新しい開始記号 $\widetilde S$，
非終端記号 $A_{\mathbf t}$（$A\in V$ と
$\mathbf t\in(\mathsf S^\sharp)^{\fan(A)}$ について），開始規則
$\widetilde S\to S_{(m)}$（$m\in\mathsf S^\sharp$ について）をもち，また，すべての規則
$A\to\rho(B_1,\ldots,B_r)$（$r=0$ を含む）とすべての
$\mathbf t_j\in(\mathsf S^\sharp)^{\fan(B_j)}$ について規則
$A_{\out_\rho(\mathbf t_1,\ldots,\mathbf t_r)}\to
\rho(B_{1,\mathbf t_1},\ldots,B_{r,\mathbf t_r})$ をもつ．\emph{刈り込んだ
細分} $\widetilde G_0$ は，$\widetilde S$ からのある成功する導出に現れる非終端記号と規則を
ちょうど残す．$G^{\widehat{s}}$ と
明示しない限り，\emph{型付き非終端記号}と\emph{型付き規則}とは，$\widetilde G_0$ の開始以外の
非終端記号と規則のことであり，
$L_X:=\{\tuple u:X\Rightarrow_{\widetilde G_0}^*\tuple u\}$ とする．
\end{definition}

型付き規則はその元の規則の雛形をもつので，$\widetilde G_0$ は
$G$ が良いときにはいつでも良い；各 $L_X$ は空でなく，そのタプルは
$G$ が $\eps$-free ならば空でない成分をもつ．

\begin{proposition}[産出型の不変性]
\label{prop:output-invariants}
\begin{enumerate}[label=(\roman*),leftmargin=*]
\item $A_{\mathbf t}\Rightarrow^*\tuple u$ が $G^{\widehat{s}}$ において成り立つならば，
$\widehat{s}^{(\fan(A))}(\tuple u)=\mathbf t$ であり，添字を消すと
$G$-導出（$\tuple u$ の）が得られる．
\item すべての終端導出 $A\Rightarrow_G^*\tuple u$ は
$A_{\widehat{s}^{(\fan(A))}(\tuple u)}\Rightarrow_{G^{\widehat{s}}}^*\tuple u$ に一意に持ち上がる；
$w$ の成功する導出は $\widetilde S\to S_{(\widehat{s}(w))}$ を通じて一意に持ち上がる．
\item $L(\widetilde G_0)=L(G^{\widehat{s}})=L(G)$．
\item すべての型付き非終端記号 $X$ について，すべての終端導出
$X\Rightarrow^*_{G^{\widehat s}}\tuple u$ は $\widetilde G_0$ の規則のみを用いる；したがって
$L_X=\{\tuple u:X\Rightarrow^*_{G^{\widehat s}}\tuple u\}$．
\end{enumerate}
\end{proposition}

\begin{lemma}[型付き非終端記号の標準文脈]
\label{lem:identity-canonical-contexts}
$G$ を良いとする．すべての型付き非終端記号 $X$（アリティ $d$ の）は，文脈
$\chi(X)\in\mathcal E_d$ で $\chi(X)[\tuple u]\in L(G)$ がすべての
$\tuple u\in L_X$ について成り立つものをもつ；$\widetilde S\to X$ が $\widetilde G_0$ にあるならば，
$\chi(X)=\blank_1$ ととってよい．
\end{lemma}

\begin{lemma}[型付き非終端記号は分布的に一様である]
\label{lem:typed-homogeneity}
$G$ をファンアウト $\le f$ の良い文法で $L:=L(G)\in\Cmcf{f}{s}$ なるものとする．すべての型付き非終端記号 $X$（アリティ $d$ の）
について，$L_X$ は一つの $\equiv_{L,s}^d$-類に含まれる．その類を
$\operatorname{Cl}_{L,s}(X)$ と書く．
\end{lemma}

意味論的な類 $\operatorname{Cl}_{L,s}(X)$ は，定義~\ref{def:reconstruction} の構文的な
再構成の非終端記号 $[\tuple x]$ とは別のものである．

これら三つの主張の証明は
付録~\ref{app:typed-witness-aux-proofs} にある．
補題~\ref{lem:typed-homogeneity} の一様性は，合同的な
MCFG~\cite{ClarkYoshinaka2016} のそれより弱い．後者の文脈は空の区切りをもち得る．

\begin{definition}[標準タプル，標準文脈と証拠集合]
\label{def:target-canonical}
\label{def:witness-set}
$G$ を良いとする．すべての型付き非終端記号 $X$ について，\emph{標準タプル}
$\omega(X)\in L_X$ と\emph{標準文脈} $\chi(X)$ を
補題~\ref{lem:identity-canonical-contexts} のとおりに固定する．ただし，
$\chi(X)=\blank_1$ とする，$\widetilde S\to X$ が $\widetilde G_0$ にあるときにはいつでも．$\Wit(\widetilde G_0)$
を次の語からなるものとする．
\[
 \chi(X)[\omega(X)]\ \ (X\text{ typed}),\qquad
 \chi(X)[\rho(\omega(Y_1),\ldots,\omega(Y_q))]\ \
 (X\to\rho(Y_1,\ldots,Y_q)\text{ typed},\ q\ge0).
\]
\end{definition}

$\omega(X)$ と $\rho(\omega(Y_1),\ldots,\omega(Y_q))$ は $L_X$ に属するので，
$\Wit(\widetilde G_0)$ は $L(G)$ の有限部分集合であり，
$L(G)\ne\emptyset$ ならば空でなく，$|\Wit(\widetilde G_0)|$ は高々
型付き非終端記号と規則の個数である．標準の対象は
標的の性質であって，作用素の入力ではない：質的にはどの選び方でもよく，
第~\ref{sec:char-data}~節では短いものを選ぶ．

\subsection{質的な保証}
\label{sec:completeness}
\label{sec:poly-time}

\begin{proposition}[すべての仮説の健全性]
\label{prop:reconstruction-soundness}
$L$ を $(f,s)$-タプル代入可能とし，$K\subseteq L$ を有限とする．すべての
$r_{\max}\ge1$ について，$[\tuple x]\Rightarrow^*\tuple u$ が
$\widehat G_{\le r_{\max}}(K)$ において成り立つならば，$\tuple u\equiv_{L,s}\tuple x$ である．
したがって $L(\mathcal B_{f,s}^{\le r_{\max}}(K))\subseteq L$ かつ
$L(\mathcal B_{f,s}(K))\subseteq L$．
\end{proposition}

\begin{proposition}[証拠の割り当てからの完全性]
\label{prop:completeness}
$G$ をファンアウト $\le f$ の有限の良い MCFG とし，$r_G$ を $\widetilde G_0$ の型付き規則のランクの最大値
とし（なければ $r_G:=0$），
各型付き非終端記号 $X=A_{\mathbf t}$ に，タプル
$\iota(X)\in(\Sigma^+)^{\fan(A)}$ で $\widehat s^{(\fan(A))}(\iota(X))=\mathbf t$ なるもの
と，文脈 $\chi(X)\in\mathcal E_{\fan(A)}$ とを割り当てる．ただし $\chi(X)=\blank_1$ とする，
$\widetilde S\to X$ が $\widetilde G_0$ にあるならば．$K$ が
$\chi(X)[\iota(X)]$ を各 $X$ について含み，
$\chi(X)[\rho(\iota(Y_1),\ldots,\iota(Y_q))]$ を各
型付き規則 $X\to\rho(Y_1,\ldots,Y_q)$，$q\ge0$ について含むとする．このとき
\[
 L(G)\subseteq L\bigl(\mathcal B_{f,s}^{\le r_{\max}}(K)\bigr)
 \quad(r_{\max}\ge\max\{1,r_G\}),\qquad
 L(G)\subseteq L\bigl(\mathcal B_{f,s}(K)\bigr).
\]
\end{proposition}

代入可能性も $\iota(X)\in L_X$ も
命題~\ref{prop:completeness} では仮定しない；$\iota=\omega$ が
$\Wit(\widetilde G_0)\subseteq K$ の場合である．

正確な再構成には標的の良い表示が必要である；次の
正規化は大きさ，ファンアウトおよびランクを制御するが，厚さは制御しない．
$G_{\mathrm{good}}$ の規則の
厚さは $G_{\mathrm{in}}$ のそれを超え得るので，以下のすべての
データの評価は与えられた既約な良い表示に相対的に述べ，
正規化を通じて移すことはしない．

\begin{proposition}[ランクを保つ良い正規化]
\label{prop:good-normalization}
$G_{\mathrm{in}}$ を有限の標準 $f$-MCFG で，規則ランクが高々
$r_{\max}$ のものとする．有限の良い $f$-MCFG $G_{\mathrm{good}}$ で，規則ランクが
高々 $r_{\max}$ であり，$L(G_{\mathrm{good}})=L(G_{\mathrm{in}})\setminus\{\eps\}$
であり，かつ，固定した $f,r_{\max}$ について
$\|G_{\mathrm{good}}\|=O_{f,r_{\max}}(\|G_{\mathrm{in}}\|)$ なるものが存在する．
$L(G_{\mathrm{in}})\cap\Sigma^+\ne\emptyset$ ならば，$G_{\mathrm{good}}$ は既約である．
\end{proposition}

\begin{theorem}[固定した型付けに対する正確な再構成]
\label{thm:exact-reconstruction}
$f\ge1$ と型付け $s$ を固定し，$L\in\Cmcf{f}{s}$ とする．
$L\cap\Sigma^+=\emptyset$ ならば $\mathcal F_L:=L$ かつ $r_G:=0$ とおく．そうでなければ $G$
を既約な良いファンアウト $\le f$ の表示（$L\cap\Sigma^+$ の）とし
（命題~\ref{prop:good-normalization}），$r_G$ を
命題~\ref{prop:completeness} のとおりとし，
$\mathcal F_L:=\Wit(\widetilde G_0)\cup(\{\eps\}\cap L)$ とおく．このとき
$\mathcal F_L\subseteq L$ は有限であり，すべての有限の $K$ で
$\mathcal F_L\subseteq K\subseteq L$ なるものは次を満たす．
\[
 L(\mathcal B_{f,s}(K))=L,\qquad
 L(\mathcal B_{f,s}^{\le r_{\max}}(K))=L\quad
 \text{for every }r_{\max}\ge\max\{1,r_G\}.
\]
$L\in\mathsf{MCFL}(f,r)$ ならば，$G$ をランク $\le r$ のものに選べるので，
$r_{\max}=\max\{1,r\}$ で足りる（$r=0$ のときも）．
\end{theorem}

\begin{theorem}[規則ランクの上限の下での多項式時間構成]
\label{thm:poly}
$f,r_{\max}\ge1$ と $s$ を固定する．$\widehat G_{\le r_{\max}}(K)$ の符号化は
時間 $(1+\|K\|_+)^{O(fr_{\max})}$ で計算可能である；そのファンアウトは高々 $f$
であり，そのランクは高々 $r_{\max}$ である．
\end{theorem}

\begin{definition}[保守的逐次学習器]
\label{def:conservative-learner}
学習器 $\mathcal A_{f,s}$（それぞれ $\mathcal A_{f,s}^{\le r_{\max}}$）
は $H_0$ から始め，$\pausesym$ と，仮説が生成するすべてのデータの上ではその仮説を保ち，
それ以外の場合はそれを $\mathcal B_{f,s}(K)$
（それぞれ $\mathcal B_{f,s}^{\le r_{\max}}(K)$）で置き換える．ここで $K$ はそれまでに読んだ
内容である．これは \cite[Algorithm~1]{Yoshinaka2011} の方策である．
\end{definition}

\begin{corollary}[Gold 同定]
\label{cor:gold-explanatory}
\label{cor:capped-gold-poly}
$f\ge1$ と $s$ を固定する．
\begin{enumerate}[label=\textup{(\roman*)}]
\item $\mathcal A_{f,s}$ は $\Cmcf{f}{s}$ を TxtEx 同定する．
\item 固定した $r_{\max}$ について，$\mathcal A_{f,s}^{\le r_{\max}}$ は
$\Cmcf{f}{s}\cap\mathsf{MCFL}(f,r_{\max})$ を TxtEx 同定し，各データを，
$\|K\|_+$ について多項式の時間で処理する．ここで $K$ はそれまでに読んだ内容である．
\end{enumerate}
\end{corollary}

証明は付録~\ref{app:qualitative-construction-details} にある
（命題~\ref{prop:good-normalization} については付録~\ref{app:normalization-details}）．
これらの主張のいずれも，特徴データの長さも収束の
段階も評価しない．

\subsection{型付きのクラスと型なしのクラス}
\label{subsec:typed-separation}

正則な道の言語は，すべてのクラスの単純な成員を与える．

\begin{lemma}[道の言語]\label{lem:app-lf-path}
$R$ を，頂点 $p$ から頂点 $p'$ への道のラベルの集合とする．ここで道は
有限の有向グラフにおけるもので，そのグラフではすべての文字がちょうど一つの辺のラベルである．
$\tuple x,\tuple y\in(\Sigma^+)^d$ と $E\in\mathcal E_d$ が
$E[\tuple x],E[\tuple y]\in R$ を満たすならば，$F[\tuple x]\in R\iff F[\tuple y]\in R$ が
すべての $F\in\mathcal E_d$ について成り立つ．
\end{lemma}

\begin{proof}
$\Sigma^+$ の語が $R$ に属するのは，各文字の辺の終点が
次の文字の辺の始点であり，最初の辺が $p$ から始まり，最後の辺が
$p'$ で終わるとき，かつそのときに限る．第 $h$ の穴（$E$ の）の直前には
$v_{h-1}$ の最後の文字があるか，または頂点 $p$ がある（$h=1$ かつ $v_0=\eps$ のとき）．直後についても
同様である；内部の区切りは空でない．したがって $E$ は各穴の両端の頂点を
定め，$x_h$ と $y_h$ は同じ二つの頂点の間の道のラベルである．
したがって，すべての $x_h$ を $y_h$ で $F[\tuple x]$ において置き換えることは
三つの条件をすべて保ち，逆も同様である．
\end{proof}

クラスは，すでにファンアウト2において，また
文脈自由言語の外において，型付けとともに真に大きくなる；そのとき型なしの作用素は過汎化する．

\begin{proposition}[型なしのクラスを超える型付きのクラス]
\label{prop:typed-separation}
$\Sigma=\{\mathtt a,\mathtt b,\mathtt c,\mathtt d,\mathtt\#_1,\mathtt\#_2,\mathtt\#_3\}$ とし，
$s:\Sigma^+\twoheadrightarrow\mathbb Z_2$ を，文字 $\mathtt a$ と $\mathtt c$ を
$2$ を法として数えるものとし，次のようにおく．
\[
 L_{\mathrm{sep}}:=\{\mathtt a^m\mathtt\#_1\mathtt b^n\mathtt\#_2\mathtt c^{m'}\mathtt\#_3\mathtt d^n:
 m,n\ge0,\ m'\in\{m,m-1\},\ m'\ge0\}.
\]
\begin{enumerate}[label=\textup{(\roman*)},leftmargin=*]
\item $L_{\mathrm{sep}}$ は規則ランク2の $2$-MCFL であり，文脈自由でない．
\item $L_{\mathrm{sep}}$ は $(f,s)$-タプル代入可能である，すべての $f\ge1$ について；したがって
$L_{\mathrm{sep}}\in\Cmcf{f}{s}$ がすべての $f\ge2$ について成り立つ．
\item $L_{\mathrm{sep}}$ は $(f,s_{\mathrm{triv}})$-タプル代入可能でない，どの $f\ge1$ についても：
語 $\mathtt a$ と $\mathtt a\mathtt a$ は文脈
$\blank_1\mathtt\#_1\mathtt\#_2\mathtt c\mathtt\#_3$ を共有するが，文脈
$\blank_1\mathtt\#_1\mathtt\#_2\mathtt\#_3$ は共有しない．それに応じて，標本
$K=\{\mathtt a\mathtt\#_1\mathtt\#_2\mathtt c\mathtt\#_3,\
\mathtt a\mathtt a\mathtt\#_1\mathtt\#_2\mathtt c\mathtt\#_3,\
\mathtt a\mathtt\#_1\mathtt\#_2\mathtt\#_3\}\subseteq L_{\mathrm{sep}}$ の上で，型なしの作用素
$\mathcal B_{f,s_{\mathrm{triv}}}$ は
$\mathtt a\mathtt a\mathtt\#_1\mathtt\#_2\mathtt\#_3\notin L_{\mathrm{sep}}$ を導出するが，一方
$\mathcal B_{f,s}$ は $\mathtt a$ を $\mathtt a\mathtt a$ と連結しない．
\end{enumerate}
\end{proposition}

\begin{proof}
$\psi:\Sigma^*\to\mathbb Z^2$ を，$\psi(\mathtt a)=(1,0)$，
$\psi(\mathtt c)=(-1,0)$，$\psi(\mathtt b)=(0,1)$，$\psi(\mathtt d)=(0,-1)$ かつ
$\psi(\mathtt\#_i)=(0,0)$ なる準同型とし，
$R:=\mathtt a^*\mathtt\#_1\mathtt b^*\mathtt\#_2\mathtt c^*\mathtt\#_3\mathtt d^*$ とする．このとき
$L_{\mathrm{sep}}=R\cap\psi^{-1}(P)$，ただし $P:=\{(0,0),(1,0)\}$ であり，$R$ は道の言語
（$0$ から $3$ への）であって，グラフは $\{0,1,2,3\}$ 上のもので，ループ $\mathtt a,\mathtt b,\mathtt c,\mathtt d$
を $0,1,2,3$ にもち，辺 $\mathtt\#_h:h-1\to h$ をもつ；すべての文字はちょうど一つの辺のラベルである．

(i) 規則 $E\to(\eps,\eps)\mid(\mathtt a\xi_1^1,\mathtt c\xi_1^2)(E)$，
$D\to(\eps,\eps)\mid(\mathtt b\xi_1^1,\mathtt d\xi_1^2)(D)$ および
\[
 S\to(\xi_1^1\mathtt\#_1\xi_2^1\mathtt\#_2\xi_1^2\mathtt\#_3\xi_2^2)(E,D)\mid
 (\mathtt a\xi_1^1\mathtt\#_1\xi_2^1\mathtt\#_2\xi_1^2\mathtt\#_3\xi_2^2)(E,D)
\]
は，ファンアウト2かつ規則ランク2の標準 MCFG であって $L_{\mathrm{sep}}$ を生成するものをなす．なぜなら $E$
は対 $(\mathtt a^k,\mathtt c^k)$ を導出し，$D$ は対 $(\mathtt b^n,\mathtt d^n)$ を導出するからである．
$L_{\mathrm{sep}}$ の語のうち，$\mathtt a$ と $\mathtt c$ が合わせて偶数回
現れるものは，ちょうど $m'=m$ なるものである；これは正則
言語との共通部分であり，標識を消すと結果は
$\{\mathtt a^m\mathtt b^n\mathtt c^m\mathtt d^n:m,n\ge0\}$ の上に写り，これは文脈自由でない．文脈自由
言語は両方の操作について閉じている．

(ii) $1\le d\le f$ とし，$\tuple x,\tuple y\in(\Sigma^+)^d$ が成分ごとに等しい
$s$-型をもつとし，$E\in\mathcal E_d$ が両方を受理するとする．
$(\delta_1,\delta_2):=\sum_i\bigl(\psi(x_i)-\psi(y_i)\bigr)=\psi(E[\tuple x])-\psi(E[\tuple y])
\in P-P$ とおく．したがって $\delta_2=0$ かつ $\delta_1\in\{-1,0,1\}$．$\psi$ の第1座標は
すべての語の上で $s$ と $2$ を法として合同であるので，$\delta_1$ は偶数である；したがって
$\psi(F[\tuple x])=\psi(F[\tuple y])$ がすべての $F\in\mathcal E_d$ について成り立つ．
補題~\ref{lem:app-lf-path} により，$F[\tuple x]\in R$ であるのは $F[\tuple y]\in R$ であるとき，かつそのときに限る．したがって $F$ が
$\tuple x$ を受理するのは，それが $\tuple y$ を受理するとき，かつそのときに限る．(i) と合わせて，$L_{\mathrm{sep}}\in\Cmcf{f}{s}$ が $f\ge2$ について成り立つ．

(iii) 二つの文脈は
$\mathtt a\mathtt\#_1\mathtt\#_2\mathtt c\mathtt\#_3,\
\mathtt a\mathtt a\mathtt\#_1\mathtt\#_2\mathtt c\mathtt\#_3\in L_{\mathrm{sep}}$（$m'=m$ と
$m'=m-1$），$\mathtt a\mathtt\#_1\mathtt\#_2\mathtt\#_3\in L_{\mathrm{sep}}$ および
$\mathtt a\mathtt a\mathtt\#_1\mathtt\#_2\mathtt\#_3\notin L_{\mathrm{sep}}$ を与える．$K$ の上で，最初の
文脈は，自明な型付けについて Link 規則 $[(\mathtt a)]\to[(\mathtt a\mathtt a)]$ を与え，
$((\mathtt a\mathtt\#_1\mathtt\#_2\mathtt\#_3),(\xi_1^1\mathtt\#_1\mathtt\#_2\mathtt\#_3),
((\mathtt a)))$ はランク1の合成証拠であるので，
$\widehat S\to[(\mathtt a\mathtt\#_1\mathtt\#_2\mathtt\#_3)]\to[(\mathtt a)]\mathtt\#_1\mathtt\#_2
\mathtt\#_3\to[(\mathtt a\mathtt a)]\mathtt\#_1\mathtt\#_2\mathtt\#_3\to
\mathtt a\mathtt a\mathtt\#_1\mathtt\#_2\mathtt\#_3$ となる．$s$ の下では $\mathtt a$ と
$\mathtt a\mathtt a$ の型は異なるので，そのような Link 規則は存在しない．
\end{proof}

$s$ の正の像は群 $\mathbb Z_2$ であり，これは局所群であるので，
定理~\ref{thm:A} は $L_{\mathrm{sep}}$ に長さ
$O(\|G\|^2\bar\theta_G)$ の特徴データを，その既約な良い表示のそれぞれについて，すべてのファンアウト
$f\ge2$ で与える．一方 $L_{\mathrm{sep}}$ は，すべてのファンアウトで Yoshinaka の型なしの学習器
\cite{Yoshinaka2011} のクラスの外にある．

\section{特徴データの評価}
\label{sec:char-data}

この節を通じて，$f$，$r_{\max}\ge1$，$\Sigma$ および型付け
$s:\Sigma^+\twoheadrightarrow\mathsf S$ を固定する；
$G=(V,\Sigma,P,S,\fan)$ はファンアウト $\le f$ かつ規則
ランク $\le r_{\max}$ の既約な良い MCFG であり，$\widetilde G_0$ はその刈り込んだ $\widehat s$-細分である．
ランクの上限の下では，定理~\ref{thm:exact-reconstruction} がすでに
$O(\|G\|)$ 個の特徴的な語を与える；課題はそれらの総長を
\emph{$G$ を用いて}評価することである．すべての評価は一つの図式に従う．
\begin{enumerate}[label=\textup{(\arabic*)},leftmargin=*]
\item \emph{証拠の長さ}（補題~\ref{lem:witness-length-template}）．
補助的な有限モノイド準同型 $\lambda$（$s$ を細分する）について，すべての
非終端記号のすべてのタプルを，同じ非終端記号の，同じ
$\lambda$-ベクトルをもつ総長高々 $\beta$ のタプルと交換できるならば，$\Wit(\widetilde G_0)$ は
特徴的であり，総長 $O(\|G\|(\bar\theta_G+\beta))$ をもつ．
\item \emph{極小反例}（第~\ref{subsec:group-invariants}~節）．
根の $\lambda$-ベクトルを指定した最短の導出において，重い部分木の大きな反鎖は，
そのベクトルを保つ同時置換を許すであろう；
したがって反鎖は小さく，補題~\ref{lem:heavy-node-count} はこれを
$\beta=O(\|G\|\bar\theta_G)$ または $O(\|G\|^2\bar\theta_G)$ に変える．
\item \emph{場合．}$\lambda$ と置換の段階だけが変わる：窓
型付けは標識付きの境界の文字を保つ；窓--群型付けはさらに
順序付き群積を保つ
（命題~\ref{prop:block-replacement-interface}）；局所群型付けは
滑動因子の座標を通じてこれらに帰着する
（第~\ref{subsec:algebraic-typed-data}~節）．
\end{enumerate}
完全な産出型細分は主たる厚さの尺度としては用いない．
それは $G$ より指数的に厚くなり得るからである
（命題~\ref{prop:typed-thickness-exponential}）．この
節の証明は付録~\ref{app:typed-witness-aux-proofs} にある．

\subsection{厚さの尺度と証拠の長さ}
\label{subsec:thickness-contexts}

規則 $R$（$G$ の）について次のようにおく．
\[
 \theta_G(R):=\min\{\,|w|:w\in L(G)\text{ has a successful derivation using }R\,\},
\]
これは導出された\emph{文}の長さであり，導出の長さではない．
$\theta_G:=\max_{R\in P}\theta_G(R)$，$\bar\theta_G:=\max\{1,\theta_G\}$ および
$t_G:=\max_{A\in V}\min\{\|\tuple w\|_1:\tuple w\in L_A(G)\}$ とおく（規則の厚さと
非終端記号の厚さ；$\theta_G$ は
\cite[Section~4.3]{Yoshinaka2011} の厚さである）．


標的 $L$ が $\eps$ を含むならば，以下のすべての主張は
$L\cap\Sigma^+$ に適用し，データ $\eps$ を加える．これは
定理~\ref{thm:exact-reconstruction} と同様である；これは符号化長に1を加える．

\begin{lemma}[厚さの補題]
\label{lem:thickness-equivalence}
\label{lem:thickness-local-bounds}
\label{lem:untyped-contexts}
$G$ を既約かつ良いとする．
\begin{enumerate}[label=\textup{(\roman*)},leftmargin=*]
\item すべての規則は，高々 $\theta_G$ 個の終端記号を明示的に書く．
\item $t_G\le\theta_G$：すべての $A\in V$ は，ある $\tuple u\in L_A(G)$ で
$\|\tuple u\|_1\le\theta_G$ なるものをもつ．
\item $\theta_G\le 2\|G\|\max\{1,t_G\}$．
\item すべての $A\in V$（ファンアウト $d$ の）は文脈 $E_A\in\mathcal E_d$ で，
$|E_A|_{\mathrm{fix}}\le\theta_G$ であり，ある成功する導出から
$A$-部分木の一つの出現を削除して得られるものをもつ；$E_S=\blank_1$ であり，
$E_A[\tuple u]\in L(G)$ がすべての $\tuple u\in L_A(G)$ について成り立つ．すべての型付き非終端記号
$X=A_{\mathbf t}$ について，選び方 $\chi(X):=E_A$ は
定義~\ref{def:target-canonical} を満たし，$X$ について\emph{導出的}である：それは
$\widetilde G_0$ のある成功する導出から，
$X$-部分木の一つの出現を削除して得られる．
\end{enumerate}
\end{lemma}


したがって $\theta_G$ と $t_G$ は既約な良い
表示の上で多項式的に同値である；Yoshinaka の効率の基準はデータを標的の大きさと
規則の厚さで測る \cite[Definition~1]{Yoshinaka2011}；
\cite{EyraudHeinzYoshinaka2016} も参照．

有限モノイド準同型 $\lambda:\Sigma^*\to\mathsf N$ について，$s\preceq\lambda$ と書くのは
$\ker(\lambda|_{\Sigma^+})\subseteq\ker s$ であるときとする；$\widehat s$ は $\eps$ を分離するので，
これは，$\lambda$ も
$\eps$ を $\Sigma^+$ から分離するときにはいつでも，定義~\ref{def:semigroup-typing} の関係である．
$\beta\ge0$ について次を考える．
\begin{itemize}[leftmargin=4em]
\item[$(\mathrm R_\beta^\lambda)$] すべての $B\in V$ と $\tuple u\in L_B(G)$ について，
$\tuple u'\in L_B(G)$ で
$\lambda^{(\fan(B))}(\tuple u')=\lambda^{(\fan(B))}(\tuple u)$ かつ
$\|\tuple u'\|_1\le\beta$ なるものが存在する．
\end{itemize}
次のようにおく．
\[
 C_s(G):=|V||\mathsf S^\sharp|^f+|P||\mathsf S^\sharp|^{fr_{\max}}
 \;\le\;|\mathsf S^\sharp|^{fr_{\max}}\,\|G\|,
\]
ここで $|V|+|P|\le\|G\|$ を用いた；これは型付き非終端記号の個数と型付き
規則の個数を合わせて上から抑える．型付き規則は，その元の規則と，
高々 $r_{\max}$ 個の子の型ベクトルとによって定まるからである．

\begin{lemma}[証拠の長さ]
\label{lem:witness-length-template}
\label{lem:type-preserving-replacement}
$\lambda:\Sigma^*\to\mathsf N$ を有限モノイド準同型で $s\preceq\lambda$ なるものとする．
\begin{enumerate}[label=\textup{(\roman*)},leftmargin=*]
\item \emph{（型を保つ交換．）}$X=B_{\mathbf t}$ が型付き
非終端記号であり，$\tuple u\in L_X$，$\tuple u'\in L_B(G)$ かつ
$\lambda^{(\fan(B))}(\tuple u')=\lambda^{(\fan(B))}(\tuple u)$ ならば，
$\tuple u'\in L_X$．したがって $(\mathrm R_\beta^\lambda)$ の下で，標準
タプルであって $\|\omega(X)\|_1\le\beta$ をすべての $X$ について満たすものが存在する．
\item \emph{（長さ．）}$\chi(X):=E_B$（$X=B_{\mathbf t}$ について）を
補題~\ref{lem:thickness-equivalence}(iv) のとおりとし，$\iota$ はすべての型付き $X$ に，
そのアリティの空でない語のタプル $\iota(X)$ で $\|\iota(X)\|_1\le b$ なるものを割り当てるとする．
$\Wit_\iota$ を語 $\chi(X)[\iota(X)]$ と，
$\chi(X)[\rho(\iota(Y_1),\ldots,\iota(Y_q))]$（型付き規則
$X\to\rho(Y_1,\ldots,Y_q)$，$q\ge0$ について）からなるものとする；したがって $\Wit_\omega=\Wit(\widetilde G_0)$．
$\Wit_\iota$ のすべての
語は長さ高々 $2\bar\theta_G+r_{\max}b$ をもち，
\[
 |\Wit_\iota|\le C_s(G),\qquad
 \|\Wit_\iota\|_+\le C_s(G)\,\bigl(1+2\bar\theta_G+r_{\max}b\bigr).
\]
\item \emph{（細分の厚さ．）}(ii) において，各規則の語
$\chi(X)[\rho(\omega(Y_1),\ldots,\omega(Y_q))]$ は，
$\widetilde G_0$ の成功する導出で型付き規則 $X\to\rho(Y_1,\ldots,Y_q)$ を用いるものの産出である；
したがって $\theta_{\widetilde G_0}\le2\bar\theta_G+r_{\max}b$．
\item \emph{（特徴標本．）}$L(G)=L\cap\Sigma^+$ がある
$L\in\Cmcf{f}{s}$ について成り立つとし，\textup{(ii)} において $\iota(X)\in\operatorname{Cl}_{L,s}(X)$ がすべての
$X$ について成り立つとする（補題~\ref{lem:typed-homogeneity}），例えば $\iota=\omega$．このとき
$\Wit_\iota\subseteq L$ であり，$\Wit_\iota\cup(\{\eps\}\cap L)$ は
$\mathcal B^{\le r_{\max}}_{f,s}$ と $\mathcal B_{f,s}$ に対する特徴標本である．
\end{enumerate}
\end{lemma}

補題~\ref{lem:witness-length-template} の項目 (i)--(iii) は
代入可能性を用いない．


\subsection{型付きの厚さと特徴データ}
\label{subsec:typed-thickness}

$\widetilde G_0$ は $G$ の刈り込んだ完全な産出型細分であった
（定義~\ref{def:output-refinement}）．次のように定める．
\[
 t_{\mathrm{typ}}(G,s):=\max_{X}\min_{\tuple u\in L_X}\|\tuple u\|_1,
\]
ここで $X$ は開始以外の型付き非終端記号を動く．また
$N_{\mathrm{typ}}(G,s)$ でそれらの個数と，開始以外の
型付き規則の個数との和を表す．特に $N_{\mathrm{typ}}\le C_s(G)\le
|\mathsf S^\sharp|^{fr_{\max}}\|G\|$．


\begin{theorem}[型付きの厚さからの特徴データ]
\label{thm:data-from-typed-thickness}
$L\in\Cmcf{f}{s}$ とし，$G$ を既約な良いファンアウト $\le f$，
ランク $\le r_{\max}$ の文法で $L(G)=L\cap\Sigma^+$ なるものとする．
最短の標準タプルと，
補題~\ref{lem:thickness-equivalence}(iv) の導出的文脈 $E_B$ をとる．このとき
$K=\Wit(\widetilde G_0)\cup(\{\eps\}\cap L)$ は
$\mathcal B^{\le r_{\max}}_{f,s}$ と $\mathcal B_{f,s}$ の両方に対する特徴標本であり，
\[
 |\Wit(\widetilde G_0)|\le N_{\mathrm{typ}},\qquad
 \|\Wit(\widetilde G_0)\|_+\le
 N_{\mathrm{typ}}\bigl(1+2\bar\theta_G+r_{\max}t_{\mathrm{typ}}\bigr).
\]
標本全体 $K$ については，長さの上界に追加の項
$\mathbf 1_{\{\eps\in L\}}$ が加わる．
\end{theorem}

数値的な長さの評価は任意の型付けについて成り立つ；代入可能性は
特徴標本であるという主張にのみ用いる．


自明な型付けについては，$\eps$-free な $G$ のすべての非終端記号は一つの
実現された型ベクトルをもつので，$t_{\mathrm{typ}}=t_G$ かつ
$N_{\mathrm{typ}}=|V|+|P|\le\|G\|$．この定理は
$O(\|G\|\bar\theta_G)$ の型なしの尺度を回復する．ただし Yoshinaka の元の
標本はより鋭い評価 $|K_G|\le|P|$ と
$\|K_G\|\le|P|\theta_G$ をもつ \cite[Lemma~7 and Theorem~1]{Yoshinaka2011}．
自明でない型付けは，代わりに $t_{\mathrm{typ}}$ を指数的にし得る．

\paragraph{稀な型からの表示に依存しない下界．}
$L\subseteq\Sigma^*$ と $m\in\mathsf S$ について次のようにおく．
\[
 \mu_m(L):=\min\{|w|:w\in L\cap\Sigma^+,\ s(w)=m\},
 \qquad (\min\emptyset:=\infty).
\]
実現型の空でない集合 $\mathsf S_L:=s(L\cap\Sigma^+)$ について，
$\mu^\star(L,s):=\max_{m\in \mathsf S_L}\mu_m(L)$ とおく；
$\mu^\star(L,s)=0$ とおく，$\mathsf S_L=\emptyset$ のときは．型付き開始非終端記号
$S_{(m)}$ が各 $m\in \mathsf S_L$ について存在し，
$L_{S_{(m)}}=\{(w):w\in L\cap\Sigma^+,\ s(w)=m\}$
（命題~\ref{prop:output-invariants}）．したがって，
$L\cap\Sigma^+\ne\emptyset$ であるときにはいつでも，
\begin{equation}\label{eq:root-thickness}
 \mu^\star(L,s)\le t_{\mathrm{typ}}(G,s)
 \qquad\text{for every reduced good presentation }G.
\end{equation}

\begin{lemma}[産出型全体の削除]
\label{lem:type-deletion}
$L\in\Cmcf{f}{s}$ かつ $m\in\mathsf S$ ならば，
$L^{\setminus m}:=L\setminus s^{-1}(m)$ も
$\Cmcf{f}{s}$ に属する．さらに，$L\setminus\{\eps\}\in\Cmcf{f}{s}$．
\end{lemma}

\begin{proof}
命題~\ref{prop:abelian-coset-filter}(iii) を，自明な群と
$\mathsf{Acc}=\mathsf S^\sharp\setminus\{m\}$，それぞれ $\mathsf{Acc}=\mathsf S$ をとって適用する．
\end{proof}

本稿のすべての下界は一つの原理に基づく．

\begin{lemma}[判別部分標的]
\label{lem:distinguishing}
$L,L'\in\Cmcf{f}{s}$ で $L'\subsetneq L$ とする．
$L$ のすべての特徴標本で $\mathcal B^{\le r}_{f,s}$ $(r\ge1)$ または $\mathcal B_{f,s}$ に対するもの，および，
$L$ のすべての特徴標本で，集合駆動型作用素であって $L'$ に対しても特徴標本をもつものに対するものは，
$L\setminus L'$ の語を含む．
\end{lemma}

\begin{proof}
$F\subseteq L'$ を有限とする．標準の作用素については，
標的 $L'$ に適用した命題~\ref{prop:reconstruction-soundness} は
$L(\mathcal B(F))\subseteq L'\ne L$ を与えるので，$F$ は $L$ に対して特徴的でない．
もう一方の作用素については，$F'$ を $L'$ に対して特徴的なものとする；このとき
$F\cup F'\subseteq L'\subseteq L$ であるので，もし $F$ が $L$ に対して特徴的であれば，
$F\cup F'$ の上の出力は $L$ と $L'$ の両方となってしまう．
\end{proof}

\begin{proposition}[稀な型は正データを強いる]
\label{prop:rare-type}
$L\in\Cmcf{f}{s}$ とし，$m\in \mathsf S_L$ とする．
$L$ の特徴標本であって
$\mathcal B^{\le r}_{f,s}$（$r\ge1$）または $\mathcal B_{f,s}$ に対するものはすべて，
型 $m$ の語を含む．同じことは，任意の集合駆動型
作用素で，$L$ と
$L^{\setminus m}$ の両方に対して特徴標本をもつものについて成り立ち，特に，クラス
$\Cmcf{f}{s}$ 全体に対して特徴標本をもつ作用素について成り立つ．
$\eps\in L$ ならば，どちらの標準の作用素に対する標本も
$\eps$ を含まなければならない；同じことは，
$L\setminus\{\eps\}$ に対しても特徴標本を許す集合駆動型作用素について成り立つ．
したがって，標準の作用素と，クラス全体の
集合駆動型作用素について，すべての特徴標本 $F$ は次を満たす．
\[
 \|F\|_+\ge
 \sum_{m\in \mathsf S_L}\bigl(\mu_m(L)+1\bigr)
 +\mathbf 1_{\{\eps\in L\}}.
\]
\end{proposition}


\begin{proposition}[完全な産出型細分の指数的な膨張]
\label{prop:typed-thickness-exponential}
$\Sigma=\{\mathtt a,\mathtt c\}$ とし，$s_{\mathtt c}$ を型付けで
$U_1=\{1,\zeta\}$（$\zeta^2=\zeta$，第~\ref{sec:fanout-one}~節）の上へのものであって，$s_{\mathtt c}(\mathtt a)=1$ かつ
$s_{\mathtt c}(\mathtt c)=\zeta$ なるものとする．$n\ge1$ について，ファンアウト1，ランク2の文法
\[
 S\to SS\mid\mathtt a\mid\mathtt c\mid A_0\mid\cdots\mid A_n,\qquad
 A_0\to\mathtt a,\qquad
 A_i\to A_{i-1}A_{i-1}\mid\mathtt c\quad(1\le i\le n)
\]
は既約かつ良く，$\Sigma^+$ を生成し，$\|G_n\|=O(n)$ と
$\theta_{G_n}\le2$ をもつが，$t_{\mathrm{typ}}(G_n,s_{\mathtt c})\ge2^n$ かつ
$\theta_{\widetilde G_{n,0}}\ge2^n$ である．
\end{proposition}


この膨張は証拠の構成に関わるのであって，データに関わるのではない．なぜなら，
型付き非終端記号に用いるタプルはその産出である必要はないからである．型付き
非終端記号 $X$ について，補題~\ref{lem:typed-homogeneity} は $L_X$ を一つの類
$\operatorname{Cl}_{L,s}(X)$ に置く；\emph{代表長}は次である．
\[
 \tau_X:=\min\{\|\tuple v\|_1:\tuple v\in\operatorname{Cl}_{L,s}(X)\},\qquad
 \tau(G,s):=\max_X\tau_X\quad(\max\emptyset:=0).
\]

\begin{corollary}[同値な代表からの特徴データ]
\label{cor:distribution-representatives}
定理~\ref{thm:data-from-typed-thickness} の仮定の下で，
\[
 \mu^\star(L,s)\le\tau(G,s)\le t_{\mathrm{typ}}(G,s),
\]
であり，$L$ は特徴標本 $K$ を，
$\mathcal B^{\le r_{\max}}_{f,s}$ と $\mathcal B_{f,s}$ の両方に対して，次を満たすものとしてもつ．
\[
 \|K\|_+\le N_{\mathrm{typ}}\bigl(1+2\bar\theta_G+r_{\max}\tau(G,s)\bigr)
 +\mathbf 1_{\{\eps\in L\}}.
\]
\end{corollary}


命題~\ref{prop:typed-thickness-exponential} では，$L=\Sigma^+$ であり，すべての
許容文脈がすべてのタプルを受理するので，二つのタプルが同値であるのはちょうど
それらの型が一致するときである；一文字が両方の型を代表し，
$\tau(G_n,s_{\mathtt c})=1$ である．それは $t_{\mathrm{typ}}(G_n,s_{\mathtt c})\ge2^n$ であるにもかかわらずである．
すると系~\ref{cor:distribution-representatives} は長さ $O(\|G_n\|)$ の特徴
データを与える．この系は $\tau$ そのものを評価しない；第~\ref{sec:upper}~節の
証明は $t_{\mathrm{typ}}$ を評価し，したがって $G$ の上で働く．
これに対して定理~\ref{thm:two-letter-fanout-one} は，ファンアウト1で，孤立した文字の型について $\tau$ を直接評価し，
したがって $t_{\mathrm{typ}}$ が指数的になりうる二文字の半束型付けと，型 $1$ の文字が一つの $U_1$ の上への型付けについて評価する（系~\ref{cor:two-letter-fanout-one}）．

\section{上界：すべてのファンアウトにおける局所群}
\label{sec:upper}

この節は定理~\ref{thm:A} を証明する．多重文脈自由文法に
固有の困難は，部分木が複数のブロックを
異なる根成分に与え，それらが他の部分木のブロックと交互に並ぶことである；
非可換群では，複数の部分木を一度に置き換える効果は
この順序に依存する．

\subsection{極小反例による短い代表}
\label{subsec:group-invariants}

$\lambda:\Sigma^*\to\mathsf N$ を有限モノイド準同型とする．ベクトル
$\mathbf t\in\mathsf N^d$ が $B$（ファンアウト $d$ の）で\emph{実現可能}であるとは，
$\mathbf t=\lambda^{(d)}(\tuple u)$ がある $\tuple u\in L_B(G)$ について成り立つことをいう．
\emph{$\lambda$-極小}導出（$(B,\mathbf t)$ に対する）とは，$B$ からの終端導出で，
根の $\lambda$-ベクトルが $\mathbf t$ であり，まず産出の総長を，次に
節点の個数を最小にするものである．閾値 $b$ について，節点が
\emph{重い}とは，その産出の総長が $b$ より大きいことをいう．

$n_\lambda(G)$ を対 $(C,\lambda^{(\fan(C))}(\tuple u))$ の個数とする．ここで $C\in V$ かつ
$\tuple u\in L_C(G)$ であり，これらを\emph{実現ラベル}という；したがって $n_\lambda(G)\le|V||\mathsf N|^f$．

\begin{lemma}[重い節点の個数と短縮原理]
\label{lem:heavy-node-count}
\begin{enumerate}[label=\textup{(\roman*)},leftmargin=*]
\item $b\ge\bar\theta_G$ かつ $c\ge0$ とし，$T$ を $\lambda$-極小
（$(B,\mathbf t)$ に対して）とする．$T$ の重い節点のすべての反鎖が高々 $c$ 個の
元をもつならば，
\[
 \|\text{root tuple of }T\|_1\le(r_{\max}+1)\,b\,\max\{1,c\,n_\lambda(G)\}.
\]
\item \emph{（短縮原理．）}$c,c'\ge0$ かつ $b=\bar\theta_G$ とする．次を仮定する：
すべての $\lambda$-極小な $T$ と，すべての反鎖 $Q$（$T$ の重い節点の）について，ある
$Q_0\subseteq Q$ で $|Q\setminus Q_0|\le c'$ なるものが存在して，$|Q_0|>c$ ならば，
ある空でない $J\subseteq Q_0$ の節点における部分木を同時に，
同じ非終端記号からの導出であって総長高々
$\bar\theta_G$ のタプルの導出で置き換えることが，根の $\lambda$-ベクトルを変えずにできる．このとき，
$\lambda$-極小な導出のすべての重い反鎖は高々 $c+c'$ 個の元をもち，
$(\mathrm R_\beta^\lambda)$ が
$\beta=(r_{\max}+1)\,\bar\theta_G\max\{1,(c+c')\,n_\lambda(G)\}$ について成り立つ．
\end{enumerate}
\end{lemma}

\begin{proof}
(i) 各節点に，その非終端記号とその産出の $\lambda$-ベクトルとをラベルとして付ける；
実現ラベルのみが現れる．ある祖先と真の子孫が同じ
ラベルをもつならば，祖先の部分木を子孫の部分木で置き換えることは，根の
$\lambda$-ベクトルを保ち（雛形は $\lambda$ を準同型的に評価する），根の産出を
長くせず（線形性により子孫の産出は祖先の産出の一部である），
節点を取り除くので，極小性に反する．したがってすべての
根から葉への道は高々 $n_\lambda(G)$ 個の節点をもつ．根が軽ければ，その産出は
長さ高々 $b$ である．そうでなければ，重い節点は祖先について閉じた部分木をなし，
その葉は反鎖をなすので，重い節点は高々 $c\,n_\lambda(G)$ 個である．
根の産出のすべての終端記号は，重い節点が書くもの（各々高々
$\bar\theta_G\le b$ 個，補題~\ref{lem:thickness-equivalence}(i)）であるか，または
重い節点の軽い子の中にある（高々 $r_{\max}$ 個の子で，長さは高々 $b$）．

(ii) $|Q_0|>c$ ならば，置換は $B$ からの導出（$J$ の節点は
比較不能である）で，根の $\lambda$-ベクトルが $\mathbf t$ であり，根の産出が真に短いものを与える．
なぜなら，置き換えたすべての産出は総長が $\bar\theta_G$ より大きく，
線形性と非削除性により根の産出にちょうど一回現れるからである．これは
極小性に反するので，$|Q|\le c+c'$．$\tuple u\in L_B(G)$ については，(i) を $\lambda$-極小な
導出（$(B,\lambda^{(\fan(B))}(\tuple u))$ に対する）に適用する．
\end{proof}

\paragraph{群の定数．}
有限可換群 ${\mathbb H}$ について，$\mathsf d({\mathbb H})$ は
${\mathbb H}$ 上の零和自由な列（空でない部分列で和が
単位元となるものをもたない列）の長さの最大値である；したがって $\mathsf d({\mathbb H})+1$ は Davenport 定数である；
$\mathsf d(\{1\})=0$．有限群 ${\mathbb G}$ と $f\ge1$ について次のようにおく．
\[
 N_{{\mathbb G},f}=\max\{3,|{\mathbb G}|^f\},\quad
 \mathcal P_f=f^{f+1}f!,\quad
 M_{{\mathbb G},f}=\mathcal P_f\,N_{{\mathbb G},f}^{\,2^{f^2}},
\]
\[
 c_{{\mathbb G},f}:=
 \begin{cases}
 \mathsf d({\mathbb G}^f),&{\mathbb G}\text{ abelian},\\
 |{\mathbb G}|-1,&{\mathbb G}\text{ nonabelian},\ f=1,\\
 M_{{\mathbb G},f}-1,&\text{otherwise}.
 \end{cases}
\]
$\mathcal P_f$ は一つの所有者のブロックの型の個数を上から抑え，$N_{{\mathbb G},f}\ge3$ は
付録~\ref{app:window-group} の順序の分類で必要となる．
二重指数 $2^{f^2}$ は，補題~\ref{lem:order-homogenize} の Erd\H{o}s--Szekeres 型の
抽出を，ブロックの位置の順序対（最大 $f^2$ 個）のそれぞれについて一回ずつ適用することから生じる；
それが必要であるかどうかはわからない．
可換群については順序の情報は必要なく，
$c_{{\mathbb G},f}=\mathsf d({\mathbb G}^f)$ は $f$ について線形である
（注意~\ref{rem:davenport-size}）．ファンアウト1では，すべての所有者は一つの行の上に一つのブロック
をもち，均質化は必要なく，接頭辞積の鳩の巣論法
だけで $|{\mathbb G}|-1$ が得られる（補題~\ref{lem:block-replacement}(c)）；
$\mathsf d({\mathbb G})\le|{\mathbb G}|-1$ であるので，可換の場合の値がこれより大きくなることはない．

\paragraph{同時ブロック置換．}
幅 $f$ の，${\mathbb G}$ 上の\emph{ブロック置換族}
（定義~\ref{def:block-family}）は $d\le f$ 個の行からなり，各行は
${\mathbb G}$ における固定された定数と\emph{ブロック}の順序付き積である；各ブロックは
\emph{所有者}を $[m]:=\{1,\ldots,m\}$ にもち，元の値と置換の値をもち，各所有者は
高々 $f$ 個のブロックを所有し，それらは他の所有者のブロックと
任意に交互に並び得る．$\operatorname{val}_\varrho(J)$ は行 $\varrho$ の値であって，ちょうど
$J\subseteq[m]$ の所有者のブロックがその置換の値をとるときのものである．

\begin{proposition}[同時ブロック置換]
\label{prop:block-replacement-interface}
幅 $f$ の，有限群 ${\mathbb G}$ 上のブロック置換族で，
$m>c_{{\mathbb G},f}$ 個の所有者をもつものにおいて，空でない $J\subseteq[m]$ で
$\operatorname{val}_\varrho(J)=\operatorname{val}_\varrho(\varnothing)$ がすべての行 $\varrho$ について成り立つものが存在する．
\end{proposition}

\begin{proof}
これは補題~\ref{lem:block-replacement} である．可換な ${\mathbb G}$ については，
${\mathbb G}^d\le {\mathbb G}^f$ における所有者の差のベクトルの，空でない零和部分列を選ぶ．
$f=1$ では，一つの行の接頭辞積の状態に対する鳩の巣論法が
所有者の区間を選ぶ．任意の ${\mathbb G}$ と $f$ については，どの因子も交換しない：
共通のブロックの型をもつ所有者（型は高々 $\mathcal P_f$ 通り）の間で，すべての対ごとの順序
比較を均質化し（補題~\ref{lem:order-homogenize}），これによりすべての行は
連続した順位のかたまりに分解される（補題~\ref{lem:group-block-order}）；すると
接頭辞積の境界の状態に対する鳩の巣論法が，すべての
順序付きの行の積を保つ所有者の区間を与える（補題~\ref{lem:group-interval}）．
\end{proof}

個々のブロックや所有者の値は置換によって保たれる必要はなく，
所有者の寄与は交換できない：二つの行が
$\mathfrak b_{1,1}\mathfrak b_{2,1}$ と $\mathfrak b_{2,2}\mathfrak b_{1,2}$ であり，値 $(12)$ を
所有者~1 のブロックに，$(23)$ を所有者~2 のブロックにもつ（$S_3$ において）ならば，行の値は
$(12)(23)\ne(23)(12)$ である．
例~\ref{ex:nonabelian-reverse-owner} は，
所有者の向きが逆の行を保つ置換を示す．

\begin{lemma}[窓--群不変量に対する短い代表]
\label{lem:marked-heavy-window-group}
$k,\ell\ge0$，有限群 ${\mathbb G}$（自明でもよい）および準同型
$g:\Sigma^*\to {\mathbb G}$ を固定する．$\lambda:=(\widehat s_{k,\ell},g):\Sigma^*\to
\mathsf N:=M_{k,\ell}\times {\mathbb G}$ とおく．このとき $(\mathrm R_\beta^\lambda)$ が次とともに成り立つ．
\[
 \beta=b_{k,\ell,{\mathbb G}}:=(r_{\max}+1)\,\bar\theta_G\,\bigl(c_{{\mathbb G},f}+f(k+\ell+1)\bigr)\,|V||\mathsf N|^f
 \;\le\;(r_{\max}+1)\bigl(c_{{\mathbb G},f}+f(k+\ell+1)\bigr)|\mathsf N|^f\,\|G\|\,\bar\theta_G .
\]
\end{lemma}

\begin{proof}
補題~\ref{lem:heavy-node-count}(ii) の仮定を，
$c=c_{{\mathbb G},f}$ と $c'=f(k+\ell+1)$ について確かめる；$c+c'\ge1$ かつ $n_\lambda(G)\le|V||\mathsf N|^f$ であるので，それは
評価を与え，二番目の不等式は $|V|\le\|G\|$ を用いる．$T$ を $\lambda$-極小で，
重さの閾値が $b=\bar\theta_G$ のものとする．
長さが少なくとも $m_0=\max\{1,k+\ell\}$ の各根成分において，最初の $k$ 個
と最後の $\ell$ 個の位置に標識を付け（$k=\ell=0$ ならば任意の一つの位置），
より短い各成分においてはすべての位置に標識を付ける：標識は高々 $c'$ 個である．$Q$ を
重い節点の反鎖とする．比較不能な部分木は互いに素な産出をもつので，高々 $c'$ 個の節点
（$Q$ の）が標識を含む；$Q_0$ を残りのものとする．

$|Q_0|>c_{{\mathbb G},f}$ とする．$v\in Q_0$（非終端記号 $C_v$ をもつ）について，
$L_{C_v}(G)$ のタプルで総長高々 $\bar\theta_G$ のものを固定する
（補題~\ref{lem:thickness-equivalence}(ii)）．線形性と非削除性により，各
$v\in Q_0$ は高々 $f$ 個の互いに素で連続したブロックを根
成分に与え，各根成分の $g$-値は，固定された部分と
これらのブロックの $g$-値の順序付き積である．これは
${\mathbb G}$ 上のブロック置換族で所有者を $Q_0$ とするものであるので，
命題~\ref{prop:block-replacement-interface} は空でない
$J\subseteq Q_0$ で，その同時置換がすべての根の $g$-値を保つものを与える．
置き換えたブロックは標識を含まない：短い根成分には触れず，
長い成分では，それらは標識付きの接頭辞と接尾辞の真に間にあり，
それらの接頭辞と接尾辞はその文字と両端からの相対的な位置を保ち，一方，成分は
その $k+\ell$ 個の標識付きの文字（$k=\ell=0$ ならば少なくとも一つ）を保ち，長いままである．したがって
根の $\widehat s_{k,\ell}$-ベクトルも保たれ，それとともに根の
$\lambda$-ベクトルも保たれる．
\end{proof}

\subsection{窓型付けと窓--群型付け：明示的な定数}
\label{subsec:boundary-polydata}

窓--群型付けは局所群型付けである
（命題~\ref{prop:li-g-strict}(i)）ので，定理~\ref{thm:A} はそれらを扱う；この
小節は，置換の議論の基本の場合を明示的な定数とともに与える．
型付け $s$ が\emph{窓--群型付け}であるとは，$\widehat s\preceq(\widehat s_{k,\ell},g)$ が，ある $k,\ell\ge0$，
有限群 ${\mathbb G}$ および準同型 $g:\Sigma^*\to{\mathbb G}$ について成り立つことをいう；自明な
${\mathbb G}$ については，これは $s$ が窓型付けであることを意味する．

\begin{proposition}[窓--群型付けに対する多項式の特徴データ]
\label{thm:window-group-polydata}
${\mathbb G}$ を有限群（自明でもよい），$g:\Sigma^*\to {\mathbb G}$ を準同型とし，
$\widehat s\preceq(\widehat s_{k,\ell},g)$ とする；$L\in\Cmcf{f}{s}$ かつ
$L(G)=L\cap\Sigma^+$ とする．このとき $\Wit(\widetilde G_0)\cup(\{\eps\}\cap L)$ は
$\mathcal B^{\le r_{\max}}_{f,s}$ と
$\mathcal B_{f,s}$ の両方に対する特徴標本であり，$|\Wit(\widetilde G_0)|\le C_s(G)=O(\|G\|)$ かつ
\[
 \begin{gathered}
 \|\Wit(\widetilde G_0)\|_+\le K_{k,\ell,{\mathbb G}}\,\|G\|^2\bar\theta_G,\\
 K_{k,\ell,{\mathbb G}}:=|\mathsf S^\sharp|^{fr_{\max}}
 \bigl(3+r_{\max}(r_{\max}+1)(c_{{\mathbb G},f}+f(k+\ell+1))(|M_{k,\ell}||{\mathbb G}|)^f\bigr).
 \end{gathered}
\]
\end{proposition}

\begin{proof}
補題~\ref{lem:marked-heavy-window-group} と
補題~\ref{lem:witness-length-template}(i),(ii),(iv) を
$\lambda=(\widehat s_{k,\ell},g)$ と $b=b_{k,\ell,{\mathbb G}}$ について用いると，標本と
$\|\Wit\|_+\le C_s(G)(1+2\bar\theta_G+r_{\max}b_{k,\ell,{\mathbb G}})$ が得られる；ここに
$C_s(G)\le|\mathsf S^\sharp|^{fr_{\max}}\|G\|$，$1+2\bar\theta_G\le3\|G\|\bar\theta_G$ および
$|V|\le\|G\|$ を代入する．
\end{proof}

群は非可換でもよく，そのブロックは成分にまたがって交互に並んでもよい．
自明な ${\mathbb G}$，すなわち $s\preceq s_{k,\ell}$ については，$c_{{\mathbb G},f}=0$ である．

\begin{remark}[可換群]
\label{rem:window-abelian}
${\mathbb G}$ が可換ならば，補題~\ref{lem:block-replacement}(a) は
$c_{{\mathbb G},f}=\mathsf d({\mathbb G}^f)$ を与え，
$\mathsf d({\mathbb G}^f)=O(\exp({\mathbb G})\,f\log|{\mathbb G}|)$ である
（注意~\ref{rem:davenport-size}；$\exp({\mathbb G})$ は
元の位数の最小公倍数である）．したがって $K_{k,\ell,{\mathbb G}}$ は $f$ について多項式である，因子
$|\mathsf S^\sharp|^{fr_{\max}}$ と $(|M_{k,\ell}||{\mathbb G}|)^f$ を除いて．これに対し，非可換な
${\mathbb G}$ と $f\ge2$ については，$c_{{\mathbb G},f}=M_{{\mathbb G},f}-1$ は $f$ について二重指数的に増大する．
\end{remark}

ファンアウト1では，固定した窓に対する多項式時間・厚データ学習は
Yoshinaka~\cite[Theorem~1]{Yoshinaka2008} による；上の評価は
既約な良い表示に相対的なものであり，正規化の後でその結果を改良するとは主張しない．
新しいのは，窓と順序付きの群の値が，最大 $f$ 個の交互に並ぶ成分にわたって
\emph{同時に}保たれることである．

\subsection{局所群型付け}
\label{subsec:algebraic-typed-data}

$\mathbf{LI}$ は，有限半群 $\mathsf T$ で
$e\mathsf Te=\{e\}$ をすべてのべき等元 $e\in \mathsf T$ について満たすものの擬多様体であり，$\mathbf{Grp}$ は有限
群のそれであり，$\mathbf{LI}\vee\mathbf{Grp}$ は
積 $\mathsf T\times {\mathbb G}$（$\mathsf T\in\mathbf{LI}$ かつ ${\mathbb G}$ は有限群）の因子からなる \cite{Costa2002LI,Costa2002LIG}．
補題~\ref{lem:mcfg-li-group-window}（付録~\ref{app:local-group}）により，
$\mathsf S\in\mathbf{LI}\vee\mathbf{Grp}$ をもつすべての型付けは窓--群型付けであり，$k$，${\mathbb G}$
および $g$ は $s$ のみに依存するので，命題~\ref{thm:window-group-polydata} がそれに適用される．

有限半群が\emph{局所群}である，$\mathsf T\in\mathbf{LG}$，とは，$e\mathsf Te$ が
すべてのべき等元 $e\in\mathsf T$ について群であることをいう．型付けが\emph{局所群型付け}であるとは，
その正の像 $\mathsf S$ が局所群であることをいう；この条件を
$\mathsf S^\sharp$ の上で調べることは決してない．
命題~\ref{prop:li-g-strict} により，$\mathbf{LI}\vee\mathbf{Grp}\subsetneq\mathbf{LG}$．

\begin{definition}[滑動因子準同型]
\label{def:local-group-typing}
$k\ge1$，有限群 ${\mathbb G}$ および写像 $\nu:\Sigma^{k+1}\to {\mathbb G}$（
\emph{滑動因子ラベル付け}）を固定する．$|w|\ge2k$ について，
$\operatorname{fac}_k(w)=(F_1,\ldots,F_{|w|-k})$ を長さ $k+1$ の因子（$w$ の）を
その開始位置の順に並べたものとし，
$\gamma_{k,\nu}(w):=\nu(F_1)\nu(F_2)\cdots\nu(F_{|w|-k})$（
${\mathbb G}$ における順序付き積）とおく；$\gamma_{k,\nu}(w):=1_{\mathbb G}$ とおく，$|w|<2k$ については．\emph{滑動因子
準同型}は $\lambda_{k,\nu}(w):=\bigl(\widehat s_{k,k}(w),\gamma_{k,\nu}(w)\bigr)$ である；
補題~\ref{lem:local-group-morphism} により，それは有限モノイド
$M_{k,\nu}$ の上への準同型である．型付け $s$ が\emph{局所群決定的}であるとは，
$\widehat s\preceq\lambda_{k,\nu}$ が，ある $k$，${\mathbb G}$ および $\nu$ について成り立つことをいう．
\end{definition}

滑動因子準同型は，各文字単独によってではなく，
滑動因子によって定まる群の寄与を符号化できる．例えば，
$\lambda_{1,\nu}$ で $\nu(\mathtt{ab})=1\in\mathbb Z_2$ かつ
他の対の上では零であるものは，$\mathtt{ab}$ の出現の偶奇を記録する．
$\{\mathtt a,\mathtt b\}$ の上でさえ，これは，
固定した窓と有限群準同型の組み合わせのどれによっても定まらない
（命題~\ref{prop:li-g-strict}(ii)）．以下の局所群の
特徴づけは，新しい学習器も標的に固有の型付けも
必要とせずに，この追加の能力を表す．


\begin{theorem}[局所群型付けに対する多項式の特徴データ]
\label{thm:local-group-polydata}
$\mathsf S\in\mathbf{LG}$ とする．すべての既約な良いファンアウト $\le f$，
ランク $\le r_{\max}$ の文法 $G$ は
$t_{\mathrm{typ}}(G,s)=O_{f,r_{\max},s,\Sigma}(\|G\|\bar\theta_G)$ を満たす．さらに
$L\in\Cmcf{f}{s}$ かつ $L(G)=L\cap\Sigma^+$ ならば，
$\Wit(\widetilde G_0)\cup(\{\eps\}\cap L)$ は
$\mathcal B^{\le r_{\max}}_{f,s}$ と $\mathcal B_{f,s}$ の両方に対する特徴標本であり，
\[
 |\Wit(\widetilde G_0)|=O(\|G\|),\qquad
 \|\Wit(\widetilde G_0)\|_+=O_{f,r_{\max},s,\Sigma}(\|G\|^2\bar\theta_G).
\]
\end{theorem}

\begin{proof}
定理~\ref{thm:lg-characterization} により，$\widehat s\preceq\lambda:=\lambda_{k,\nu}$
が，$k=|\mathsf S|$ と，ある ${\mathbb G},\nu$（$s$ のみに依存する）について成り立つ．
補題~\ref{lem:local-group-shortest} は $(\mathrm R_\beta^\lambda)$ を
$\beta=\beta_\lambda=O(\|G\|\bar\theta_G)$ について与える：重い部分木は，その非終端記号の
任意の短い産出で置き換えられる．ただしその前に，所有者を間引いて，
そのブロックが互いに少なくとも $2k$ 文字離れるようにし，各ブロックを
両側で $k$ 個の変化しない文字だけ広げておく
（補題~\ref{lem:lg-buffered-owners}）；広げたブロックはその窓の
型を保つので，順序付きの滑動因子の積は
補題~\ref{lem:local-group-blocks} と
命題~\ref{prop:block-replacement-interface} によって制御される．この段階は
代入可能性を用いない．補題~\ref{lem:witness-length-template}(i) は
$t_{\mathrm{typ}}\le\beta_\lambda$ を与え，
定理~\ref{thm:data-from-typed-thickness} は
$\|\Wit\|_+\le N_{\mathrm{typ}}(1+2\bar\theta_G+r_{\max}\beta_\lambda)
=O(\|G\|^2\bar\theta_G)$ を与える．
\end{proof}

Rees の座標と滑動因子の座標は古典的である；それらと
所有者ごとのブロックの短縮との組み合わせが定量的な段階である．したがって局所群型付けと
窓--群型付けは $\|G\|$ について同じ二次の依存性をもつ；
命題~\ref{thm:window-group-polydata} は後者に対する明示的な定数を与える．

\subsection{非文脈自由の標的の上での型付けの領域の分離}
\label{subsec:features-needed}

定理~\ref{thm:A} は窓を超える二つの要素をもつ：滑り因子の情報（窓--群型付けではなく局所群）と，
交互に並ぶ成分にわたる順序付きの（非可換でもよい）群の積である．
命題~\ref{prop:li-g-strict} は対応する半群を分離する．次の標的はそれらを言語の上で分離する：
いずれもファンアウト2の非文脈自由言語で，ある型付けのクラスには入り，
弱い種類のどの型付け（それぞれ，どの窓--群型付け，どの窓--可換群型付け）のクラスにも入らない．したがって，その種類のどの型付けについても，弱い方の定理はそれを扱えない．
どちらも $\varsigma w\mathtt\$x_\times$ の形をもつ：旗 $\varsigma=\mathtt r$ は $w$ を制約せず，
$\varsigma=\mathtt l$ は型付けに依る条件を $w$ に課し，
$x_\times\in C_\times:=\{\mathtt{1}^m\mathtt{2}^n\mathtt{3}^m\mathtt{4}^n:m,n\ge1\}$
が標的を非文脈自由にする．$\varphi_\times$ を，
$\mathtt 1,\mathtt 3,\mathtt 2,\mathtt 4$ を $(1,0),(-1,0),(0,1),(0,-1)\in\mathbb Z^2$ に，
他の文字を $0$ に写すものとし，$\#_{\mathtt{ab}}(w)$ を $w$ における因子 $\mathtt{ab}$ の出現数とする．

\begin{lemma}[旗による分離]
\label{lem:flag-separation}
$s'$ を型付け，$u,v\in\Sigma^+$ で $s'(u)=s'(v)$，$z\in\Sigma^*$ とする．$L$ が
$\mathtt ruz$，$\mathtt rvz$ と $\mathtt luz$ を含み $\mathtt lvz$ を含まないなら，
すべての $f\ge1$ について $L\notin\Cmcf{f}{s'}$．
\end{lemma}

\begin{proof}
$\mathtt r\blank_1z\in\mathcal E_1$ は $u$ と $v$ を受理し，$\mathtt l\blank_1z$ は
$u$ を受理するが $v$ を受理しない；したがって $\D^{(1)}_L(u)\ne\D^{(1)}_L(v)$ であり，
すべての $(f,s')$ の条件はアリティ1を含む．
\end{proof}

\paragraph{標的．}
$\Sigma_{\mathrm A}=\{\mathtt l,\mathtt r,\mathtt a,\mathtt b,\mathtt c,\mathtt\$,
\mathtt 1,\mathtt 2,\mathtt 3,\mathtt 4\}$，
\[
 L_{\mathrm A}:=\{\varsigma w\mathtt\$x_\times:\varsigma\in\{\mathtt l,\mathtt r\},\
 w\in\{\mathtt a,\mathtt b,\mathtt c\}^+,\ x_\times\in C_\times,\
 \varsigma=\mathtt r\text{ または }\#_{\mathtt{ab}}(w)\text{ は偶数}\}
\]
とし，$s_{\mathrm A}:=\lambda_{1,\nu}|_{\Sigma_{\mathrm A}^+}$（正の像へ余制限），
$\nu:\Sigma_{\mathrm A}^2\to\mathbb Z_2$ は $\nu(\mathtt{ab})=1$，それ以外で $\nu=0$
とする（定義~\ref{def:local-group-typing}）；これは最初と最後の文字と，
$|w|\ge2$ なら $\#_{\mathtt{ab}}(w)$ の偶奇を記録する．
$\Sigma_{\mathrm B}=\Sigma_{\mathrm A}\setminus\{\mathtt c\}$，$\alpha=(12)$，
$\beta=(23)\in S_3$，$g_{\mathrm B}:\Sigma_{\mathrm B}^*\to S_3$ を
$\mathtt a\mapsto\alpha$，$\mathtt b\mapsto\beta$，他の文字 $\mapsto1$ とし，
\[
 L_{\mathrm B}:=\{\varsigma w\mathtt\$x_\times:\varsigma\in\{\mathtt l,\mathtt r\},\
 w\in\{\mathtt a,\mathtt b\}^+,\ x_\times\in C_\times,\
 \varsigma=\mathtt r\text{ または }g_{\mathrm B}(w)=1\}，
\]
$s_{\mathrm B}:=(\widehat s_{1,1},g_{\mathrm B})|_{\Sigma_{\mathrm B}^+}$（像へ余制限）とする．
$\mathbf{Ab}$ は有限可換群の擬多様体，
$\mathbf{LI}\vee\mathbf{Ab}$ は $\mathsf T\in\mathbf{LI}$ と有限可換群 ${\mathbb G}$ による
$\mathsf T\times {\mathbb G}$ の因子からなる．

\begin{proposition}[局所群を必要とする標的]
\label{prop:lg-needed-target}
\begin{enumerate}[label=\textup{(\roman*)},leftmargin=*]
\item $L_{\mathrm A}\in\mathsf{MCFL}(2,2)$ であり，$L_{\mathrm A}$ は文脈自由でない．
\item すべての $f\ge2$ について $L_{\mathrm A}\in\Cmcf{f}{s_{\mathrm A}}$ であり，
$s_{\mathrm A}(\Sigma_{\mathrm A}^+)\in\mathbf{LG}$．
\item ある $k,\ell$ と有限群への準同型 $g'$ について
$\widehat{s'}\preceq(\widehat s_{k,\ell},g')$ なら（特に
$s'(\Sigma_{\mathrm A}^+)\in\mathbf{LI}\vee\mathbf{Grp}$ なら），
すべての $f\ge1$ について $L_{\mathrm A}\notin\Cmcf{f}{s'}$．
\end{enumerate}
したがって $s_{\mathrm A}(\Sigma_{\mathrm A}^+)\in\mathbf{LG}\setminus(\mathbf{LI}\vee\mathbf{Grp})$ であり，
$L_{\mathrm A}$ は定理~\ref{thm:local-group-polydata} で扱えるが，どの型付けについても
命題~\ref{thm:window-group-polydata} では（補題~\ref{lem:mcfg-li-group-window} を併せても）扱えない．
\end{proposition}

\begin{proposition}[非可換群を必要とする標的]
\label{prop:nonabelian-needed-target}
\begin{enumerate}[label=\textup{(\roman*)},leftmargin=*]
\item $L_{\mathrm B}\in\mathsf{MCFL}(2,2)$ であり，$L_{\mathrm B}$ は文脈自由でない．
\item すべての $f\ge2$ について $L_{\mathrm B}\in\Cmcf{f}{s_{\mathrm B}}$ であり，$s_{\mathrm B}$ は
群 $S_3$ による窓--群型付けである．
\item ある $k,\ell$ と有限\emph{可換}群への準同型 $g'$ について
$\widehat{s'}\preceq(\widehat s_{k,\ell},g')$ なら（特に
$s'(\Sigma_{\mathrm B}^+)\in\mathbf{LI}\vee\mathbf{Ab}$ なら），
すべての $f\ge1$ について $L_{\mathrm B}\notin\Cmcf{f}{s'}$．
\end{enumerate}
したがって $L_{\mathrm B}$ は非可換の定数 $c_{S_3,f}=M_{S_3,f}-1$ をもつ
命題~\ref{thm:window-group-polydata} で扱えるが，どの型付けについても，その自明群や
可換群の場合（注意~\ref{rem:window-abelian}）では扱えない．
\end{proposition}

明示的な分離する語を含む証明は付録~\ref{app:separating-targets} にある．どちらの命題もクラスへの所属についてのものであって，データの長さについてのものではない：
定理~\ref{thm:A} の型付けの領域が非文脈自由の標的の上で異なることを示すのであって，
その指数が鋭いことを示すのではない．各標的は正則な旗付きの前半と固定した非文脈自由の後半 $C_\times$ の連結であり，
分離は前半から来る．
定理~\ref{thm:A} の多重文脈自由に固有の内容は，任意の表示に対するブロック置換の議論にある．

\section{下界：すべてのファンアウトにおける線形因子}
\label{sec:linear}

この節は定理~\ref{thm:B} を証明する．

定理~\ref{thm:data-from-typed-thickness} の上界は $\|G\|$ とともに増大し，
命題~\ref{prop:rare-type} の稀な型による下界は増大しない．それにもかかわらず，積 $\|G\|\bar\theta_G$ は，
すべてのファンアウトで，少なくとも
$4f+1$ 文字（ファンアウト1では2文字）のアルファベット上のすべての型付けの下で，かつ，最小のものの定数倍以内の表示について
必要である．

\begin{proposition}[線形因子はすべてのファンアウトで必要である]
\label{prop:linear-factor-all-fanout}
$f\ge1$ とし，$\Sigma$ が $4f+1$ 個の相異なる文字
$\mathtt a_1,\ldots,\mathtt a_{2f},\mathtt d_1,\ldots,\mathtt d_{2f},\mathtt z$ を含むとし，
$s$ を任意の型付けとする．整数 $q\ge n\ge1$ について $W:=2fq+2^n+2f$ とおき，
\[
 w_i:=\mathtt a_1^i\mathtt d_1\mathtt a_2^i\mathtt d_2\cdots\mathtt a_{2f}^i\mathtt d_{2f}
 \mathtt z^{2f(q-i)+2^n}\quad(0\le i\le q),\qquad
 L^{(f)}_{q,n}:=\{w_0,\ldots,w_q\}.
\]
\begin{enumerate}[label=\textup{(\roman*)},leftmargin=*]
\item $L^{(f)}_{q,n}$ のすべての部分集合は
$\Cmcf{f}{s_{\mathrm{triv}}}\subseteq\Cmcf{f}{s}$ に属する．
\item $\mathcal B_{f,s}$ について，$\mathcal B^{\le r}_{f,s}$ \textup($r\ge1$\textup) について，
および，
$\Cmcf{f}{s}$ 全体で特徴標本をもつすべての集合駆動型作用素について，$L^{(f)}_{q,n}$ の唯一の特徴標本は
$L^{(f)}_{q,n}$ 自身である．
\item $L^{(f)}_{q,n}$ は既約な良い表示 $G^{(f)}_{q,n}$ で，ファンアウト
$f$ かつ規則ランク2であって，
$\|G^{(f)}_{q,n}\|=(7f+11)q+6n+3f+10\le(7f+17)(q+1)$ かつ
$\bar\theta_{G^{(f)}_{q,n}}=W$ なるものをもつ．既約な良いファンアウト $\le f$ のすべての表示
$G$（$L^{(f)}_{q,n}$ の）は $\bar\theta_G=W$ かつ $\|G\|\ge q+1$ をもつ．
\item $f=1$ については，\textup{(i)}--\textup{(iii)} は，二つの
文字 $\mathtt a,\mathtt d$ を含む任意のアルファベット上で，語
$w_i=\mathtt a^i\mathtt d\mathtt a^i\mathtt d\mathtt a^{2(q-i)+2^n}$ について成り立つ．この語は
$\mathtt a_1,\mathtt a_2,\mathtt z$ を $\mathtt a$ と，
$\mathtt d_1,\mathtt d_2$ を $\mathtt d$ と同一視して得られ，文法も同じように
同一視する．
\end{enumerate}
したがって，\textup{(ii)} のすべての特徴標本 $F$ は次を満たす．
\[
 \|F\|_+=(q+1)(W+1)\ \ge\ \tfrac1{7f+17}\,\|G^{(f)}_{q,n}\|\,
 \bar\theta_{G^{(f)}_{q,n}}.
\]
\end{proposition}

各語 $w_i$ は一つの標識の位置と二つの標識の距離によって定まり，これが (i) を与える（付録~\ref{app:linear-factor-fanout}）．$q=n$ とすると $\|G^{(f)}_{n,n}\|=\Theta(n)$ かつ
$\bar\theta_{G^{(f)}_{n,n}}=2^n+\Theta(n)$ である；$n$ を固定すると，
$\bar\theta_{G^{(f)}_{q,n}}=\Theta(\|G^{(f)}_{q,n}\|)$ である．(iii) により，これらの表示は
最小の既約な良い表示の定数倍以内にあるので，この下界は
$G$ に詰め物をして得たものではない；自明な型付けでは，この命題は
Yoshinaka の上界 $\|K_G\|\le|P|\theta_G$ が定数倍を除いて達成されることを示す．
$G^{(f)}_{q,n}$ の非終端記号のすべての産出は一つの総長をもつので，
この族については $t_{\mathrm{typ}}=\bar\theta_G=W$ であり，この族は，
命題~\ref{thm:window-group-polydata} と
定理~\ref{thm:local-group-polydata} におけるさらなる因子 $\|G\|$ については何も述べない．証明は
付録~\ref{app:linear-factor-fanout} にある．

\begin{corollary}[自明な型付け]
\label{cor:trivial-exact}
$f$ と $r_{\max}\ge1$ を固定する．自明な型付け，すべての
$L\in\Cmcf{f}{s_{\mathrm{triv}}}$，およびすべての既約な良いランク $\le r_{\max}$ の
表示 $G$（$L\cap\Sigma^+$ の）について，作用素 $\mathcal B_{f,s_{\mathrm{triv}}}$ は
総長 $O(\|G\|\bar\theta_G)$ の特徴標本をもつ．これは Yoshinaka の
標本 \cite[Lemma~7 and Theorem~1]{Yoshinaka2011} と同様である．さらに，
$r_{\max}\ge2$ であり，かつ $\Sigma$ が少なくとも $4f+1$ 文字をもつとするか，または
$r_{\max}\ge2$，$f=1$ であり，かつ $\Sigma$ が少なくとも2文字をもつとする．このとき，この上界は
定数倍を除いて達成される：
命題~\ref{prop:linear-factor-all-fanout} の
標的 $L^{(f)}_{q,n}$ と表示 $G^{(f)}_{q,n}$ について，すべての特徴標本は
長さ少なくとも $\|G^{(f)}_{q,n}\|\bar\theta_{G^{(f)}_{q,n}}/(7f+17)$ をもつ．
\end{corollary}

\begin{proof}
すべての型付き非終端記号は型なしの非終端記号であるので，
補題~\ref{lem:thickness-equivalence} により $t_{\mathrm{typ}}(G,s_{\mathrm{triv}})\le\bar\theta_G$ であり，
定理~\ref{thm:data-from-typed-thickness} は $O(\|G\|\bar\theta_G)$ を与える．
命題~\ref{prop:linear-factor-all-fanout} は二番目の主張を与える；
表示 $G^{(f)}_{q,n}$ は規則ランク2をもち，それゆえそこでは
$r_{\max}\ge2$ が必要である．
\end{proof}

特徴標本は，標的の有限個の語を加えても特徴的なままであるので，
\emph{最短の}特徴標本の長さだけが意味をもつ；
この系は，それをすべての標的について上から抑え，その上界が一つの族によって達成されることを示す．

\section{ファンアウト1での境界}
\label{sec:fanout-one}

この節は定理~\ref{thm:C} を証明する．

\paragraph{三文字の倍増族．}
$\mathtt a,\mathtt b,\mathtt c$ を相異なる文字とし，$U_1=\{1,\zeta\}$ を，
単位元 $1$ と $\zeta^2=\zeta$ をもつものとし，
$s_U:\{\mathtt a,\mathtt b,\mathtt c\}^+\twoheadrightarrow U_1$ を，$\mathtt a,\mathtt b$
を $1$ に，$\mathtt c$ を $\zeta$ に送るものとする；したがって $s_U(w)=1$ となるのは，ちょうど
$\mathtt c$ を含まない空でない語である．$n\ge1$ について $G_n$ をファンアウト1の文法
\[
 S\to E_n,\qquad E_0\to\mathtt a,\qquad
 E_i\to\mathtt bE_{i-1}E_{i-1}\mathtt b\mid\mathtt cE_{i-1}\mathtt c\quad(1\le i\le n),
\]
とし，$L^{(i)}:=L_{E_i}(G_n)$ とおく．したがって $L^{(0)}=\{\mathtt a\}$ かつ
$L^{(i)}=\mathtt bL^{(i-1)}L^{(i-1)}\mathtt b\cup\mathtt cL^{(i-1)}\mathtt c$ である．また
$L_n:=L^{(n)}=L(G_n)$，$W_0:=\mathtt a$ および $W_i:=\mathtt bW_{i-1}W_{i-1}\mathtt b$ とおく．
これは
姉妹 CFG 原稿 \cite[Proposition~6.4]{kuriyamaCFG} の括弧付き倍増族の三文字の変種であり，そこでは
開き括弧と閉じ括弧が異なる文字である；その性質はすべて
ここで証明する．

\begin{proposition}[三文字の倍増族]
\label{prop:bracketed-family}
\label{prop:three-letter-family}
すべての $n\ge1$ について：
\begin{enumerate}[label=\textup{(\roman*)},leftmargin=*]
\item $G_n$ は既約かつ良く，ファンアウト1かつ規則ランク2であり，
$\|G_n\|=13n+7$ かつ $\bar\theta_{G_n}=4n$ をもつ．
\item $L_n$ は Clark--Eyraud の意味で代入可能である：すべての $x,y\in\Sigma^+$ について，
$\D^{(1)}_{L_n}(x)\cap\D^{(1)}_{L_n}(y)\ne\emptyset$ ならば
$\D^{(1)}_{L_n}(x)=\D^{(1)}_{L_n}(y)$ である．したがって $L_n\in\Cmcf{1}{s}$ が，すべての
型付け $s$（すべてのアルファベット $\Sigma\supseteq\{\mathtt a,\mathtt b,\mathtt c\}$ 上の）について成り立つ；
特に $L_n\in\Cmcf{1}{s_U}$．
\item $W_n$ は $L_n$ の語で $\mathtt c$ を含まない唯一のものであり，したがって
$s_U$-型 $1$ の唯一のものであり，$|W_n|=3\cdot2^n-2$．
\item すべての特徴標本 $F$（$L_n$ の）で，
$\mathcal B^{\le r}_{1,s_U}$ $(r\ge1)$ または $\mathcal B_{1,s_U}$ に対するもの，あるいは，
$\Cmcf{1}{s_U}$ 全体で特徴標本をもつ任意の集合駆動型作用素に対するものは，
$W_n$ を含む；したがって $\|F\|_+\ge3\cdot2^n-1$．
\end{enumerate}
\end{proposition}

証明は付録~\ref{app:bracketed-family} にある：(iii) が成り立つのは，
$\mathtt c$ を含まない導出が完全二分木であるからであり，(iv) は
(iii) と命題~\ref{prop:rare-type} から従う．


\begin{definition}[一様多項式特徴データ長]
\label{def:uniform-polydata}
$f,r_{\max}\ge1$ とする．有限半群の擬多様体 $\mathbf V$ が
\emph{ファンアウト $f$ かつ
ランク $r_{\max}$ で一様多項式特徴データ長}をもつとは，すべての空でない有限アルファベット $\Sigma$ と，すべての
全射な型付け $s:\Sigma^+\twoheadrightarrow\mathsf S$ で
$\mathsf S\in\mathbf V$ なるものについて，ある多項式 $p=p_{f,r_{\max},s,\Sigma}$ が存在して，
すべての $L\in\Cmcf{f}{s}$ と，すべての既約な良い
ファンアウト $\le f$，ランク $\le r_{\max}$ の表示 $G$（$L\cap\Sigma^+$ の）について，
作用素 $\mathcal B^{\le r_{\max}}_{f,s}$ が特徴標本
$F$ を $L$ に対して $\|F\|_+\le p(\|G\|,\bar\theta_G)$ なるものとしてもつことをいう．
\end{definition}

量化子は書かれたとおりの順序である：$p$ は $f$，
$r_{\max}$，$\Sigma$ および型付け $s$ に依存してよいが，標的 $L$ にも
表示 $G$ にも依存しない．

\begin{theorem}[ファンアウト1：データ長に対する擬多様体の境界]
\label{thm:fanout-one-data}
$r_{\max}\ge2$ とし，$\mathbf V$ を有限
半群の擬多様体とする．
\begin{enumerate}[label=\textup{(\roman*)},leftmargin=*]
\item $\mathbf V$ がファンアウト1かつランク $r_{\max}$ で一様多項式特徴データ長をもつのは，
$\mathbf V\subseteq\mathbf{LG}$ であるとき，かつそのときに限る．
\item $\mathbf V\subseteq\mathbf{LG}$ ならば，
定義~\ref{def:uniform-polydata} において
$p(\|G\|,\bar\theta_G)=O_{r_{\max},s,\Sigma}(\|G\|^2\bar\theta_G)$ ととってよい．
\item $\mathbf V\not\subseteq\mathbf{LG}$ ならば，
命題~\ref{prop:bracketed-family} の型付け $s_U$ は正の像
$U_1\in\mathbf V$ をもち，標的 $L_n\in\Cmcf{1}{s_U}$ は既約な良い
ランク2の表示 $G_n$ で $\|G_n\|=13n+7$ かつ
$\bar\theta_{G_n}=4n$ なるものをもつが，一方 $L_n$ のすべての特徴標本は
符号化総長少なくとも $3\cdot2^n-1\ge2^{(\|G_n\|+\bar\theta_{G_n})/17}$ をもつ．この下界は，
$\mathcal B^{\le r_{\max}}_{1,s_U}$ について，$\mathcal B_{1,s_U}$ について，および，
$\Cmcf{1}{s_U}$ 全体で特徴標本をもつすべての集合駆動型作用素について成り立つ．
\end{enumerate}
\end{theorem}

\begin{proof}
(ii)，したがって (i) の \emph{if} の方向は，
$f=1$ とした定理~\ref{thm:local-group-polydata} である．なぜならそのとき，すべての
$\mathsf S\in\mathbf V$ は $\mathbf{LG}$ に属するからである．(iii) については，
$\mathsf S\in\mathbf V\setminus\mathbf{LG}$ をとる．
定理~\ref{thm:lg-characterization}(ii) により，$\mathsf S$ は
$U_1$ と同型な部分半群をもち，擬多様体は部分半群について閉じているので
$U_1\in\mathbf V$ である．型付け $s_U$ は $U_1$ の上へ写す；
残りの主張は命題~\ref{prop:bracketed-family}(i) と~(iv) であり，
$\|G_n\|=O(n)$ と $\theta_{G_n}\le4n$ を用いる．
$3\cdot2^n-1$ は $\|G_n\|$ と
$\bar\theta_{G_n}$ のどんな多項式でも抑えられないので，(i) の \emph{only if} の方向が従う．
\end{proof}

\begin{remark}[定理~\ref{thm:fanout-one-data} の範囲]
\label{rem:fanout-one-data}
(a) 下界はランク2を用いる：倍増規則
$E_i\to\mathtt bE_{i-1}E_{i-1}\mathtt b$ は二分である．
線形 CFG で表される標的については，姉妹原稿は，すべての固定した型付けについて，
文法の大きさのみについて多項式の特徴データを与える
\cite[Theorem~8.3]{kuriyamaCFG}．
(b) この定理は擬多様体に関するものである．どの個々の型付けが
多項式のデータをもつかは，$U_1$ の上への型付け（系~\ref{cor:u1-classification}）を除き，ここでは特徴づけない．
(c) 上界は固定した作用素
$\mathcal B^{\le r_{\max}}_{1,s}$ に対するものであるが，下界は，クラス全体で特徴標本をもつ
すべての集合駆動型作用素について成り立つ．
\end{remark}

\paragraph{多項式時間・厚データ．}
Eyraud，Heinz および Yoshinaka \cite[Definition~2.12]{EyraudHeinzYoshinaka2016} は，
\cite{Yoshinaka2008} に従い，クラスが多項式時間・
厚データ同定可能（IPTtD）であるとは，あるアルゴリズムが，与えられた標本の大きさについて多項式の時間で，
それと無矛盾な仮説を返し，かつ，標的のすべての表現 $R$
が，大きさが $R$ の大きさとその厚さについて多項式の特徴標本をもつことをいう，と定める．
補題~\ref{lem:thickness-equivalence}(ii),(iii) により，彼らの
厚さ（非終端記号の最短の産出のうち最長のもの，$t_G$）と本稿の規則の
厚さは，$\|G\|$ を与えたとき多項式的に同値である．標本の大きさ
$\sum_{w\in K}|w|$ と $\|K\|_+$ は高々語の個数だけ異なる．彼らの
定義はアルゴリズムに，固定した表現のクラスにおける仮説を返すことを求める；
$\mathcal B^{\le r_{\max}}_{f,s}$ の生の出力は
規則 $\widehat S\to\eps$ を含み得るので，まずそれを取り除く．

\begin{lemma}[既約な良い仮説]
\label{lem:reduced-good-hypothesis}
$f$，$r_{\max}\ge1$ および $s$ を固定し，$K\subseteq\Sigma^*$ を有限で
$K\cap\Sigma^+\ne\emptyset$ なるものとし，$G^+_K$ を，
$\mathcal B^{\le r_{\max}}_{f,s}(K)$ から規則 $\widehat S\to\eps$ を除いたものとする．このとき
$G^+_K$ はファンアウト $\le f$ かつランク $\le r_{\max}$ の既約な良い MCFG であり，
$L(G^+_K)=L(\mathcal B^{\le r_{\max}}_{f,s}(K))\setminus\{\eps\}$ であり，$G^+_K$ は
$K$ から時間 $(1+\|K\|_+)^{O(fr_{\max})}$ で計算可能である．
\end{lemma}

\begin{proof}
Start と Link の規則は恒等の雛形をもち，Const の規則は
空でない語のタプルを書き，Comp の規則は良い雛形をもつ
（定義~\ref{def:composition-witness}）；したがって $G^+_K$ のすべての規則は良く，
どの非終端記号も空の成分をもつタプルを導出しない．したがって $\eps$ は
取り除いた規則によってのみ導出され，これが言語を与える．すべての $[\tuple x]$ は
その Const 規則により生産的であり，$\widehat S$ は $K$ が
空でない語をもつので生産的である．到達可能性については，$\tuple x\in\mathrm{Obs}_d(K)$ とし，例えば
$w=E[\tuple x]\in K$ で $E=v_0\blank_1v_1\cdots\blank_dv_d\in\mathcal E_d$ とする．このとき $w\ne\eps$ であり，
$(w)\in\mathrm{Obs}_1(K)$ が文脈 $\blank_1$ によって成り立ち，雛形
$\rho=(v_0\xi_1^1v_1\cdots\xi_1^dv_d)$ は良い：それは $\eps$-free であり，各
変数を一回ずつ順に用い，$v_1,\ldots,v_{d-1}$ が
空でないので非併合である．したがって $((w),\rho,(\tuple x))$ はランク1の合成証拠であり，
$\widehat S\to[(w)]\to\rho([\tuple x])$ は $[\tuple x]$ に到達する．ファンアウト，ランクおよび
実行時間は定理~\ref{thm:poly} のものである．
\end{proof}

\begin{corollary}[多項式時間・厚データ同定]
\label{cor:iptd}
$f\ge1$，$r_{\max}\ge1$ および型付け $s$ を固定する．$\mathcal R$ を
対 $(G,z)$ からなるものとする．ここで対は，ビット $z\in\{0,1\}$ と，既約な良いファンアウト $\le f$，
ランク $\le r_{\max}$ の MCFG $G$ または規則をもたない固定した文法 $G_\emptyset$ のいずれかとからなる；対は
$L(G)\cup\{\eps:z=1\}$ を表す．
$\Cmcf{f}{s}\cap\mathsf{MCFL}(f,r_{\max})$ の成員は
$(G,[\eps\in L])$ で表される．ここで $G$ は既約な良いファンアウト $\le f$，ランク $\le r_{\max}$ の
表示（$L\cap\Sigma^+$ の）であるか，または $G=G_\emptyset$ である，$L\cap\Sigma^+=\emptyset$ ならば．
\begin{enumerate}[label=\textup{(\roman*)},leftmargin=*]
\item $s(\Sigma^+)\in\mathbf{LG}$ ならば，
補題~\ref{lem:reduced-good-hypothesis} の $(G^+_K,[\eps\in K])$ を返し，または
$(G_\emptyset,[\eps\in K])$ を返す（$K\cap\Sigma^+=\emptyset$ ならば）アルゴリズムは，
$\Cmcf{f}{s}\cap\mathsf{MCFL}(f,r_{\max})$ を，クラス $\mathcal R$ について，多項式時間・厚データ同定する．
\item $f=1$ かつ $r_{\max}\ge2$ ならば，どのアルゴリズムも，その実行
時間がどうであれ，IPTtD のデータ条件を
$\Cmcf{1}{s_U}\cap\mathsf{MCFL}(1,r_{\max})$ について，クラス $\mathcal R$ に関して満たさない．
\item したがって，$r_{\max}\ge2$ と擬多様体 $\mathbf V$ について，
次は同値である：$\mathbf V\subseteq\mathbf{LG}$；すべての型付け $s$
で $s(\Sigma^+)\in\mathbf V$ なるものについて，クラス
$\Cmcf{1}{s}\cap\mathsf{MCFL}(1,r_{\max})$ はクラス $\mathcal R$ について IPTtD である．
\end{enumerate}
\end{corollary}

\begin{proof}
(i) 補題~\ref{lem:reduced-good-hypothesis} により，仮説は
$\mathcal R$ に属し，$\|K\|_+$ について多項式の時間で計算され，
$L(\mathcal B^{\le r_{\max}}_{f,s}(K))$ を表し，これは $K$ を含む（各 $w\in K\cap\Sigma^+$
はその Start と Const の規則によって導出される）．
定理~\ref{thm:local-group-polydata} は標本 $F$ で，長さ
$O(\|G\|^2\bar\theta_G)$ をもち，$L(\mathcal B^{\le r_{\max}}_{f,s}(K))=L$ が
$F\subseteq K\subseteq L$ であるときにはいつでも成り立つものを与えるので，そのとき仮説は $L$ を表す；また
$\bar\theta_G\le2\|G\|\max\{1,t_G\}$．
(ii) \cite[Definition~2.12]{EyraudHeinzYoshinaka2016} では，アルゴリズムは
標本を読み，標本は文字列の有限集合であるので，それは集合駆動型作用素である．そのような
アルゴリズムはクラスのすべての成員に対して特徴標本をもつ．なぜなら，すべての成員 $L$ で $L\cap\Sigma^+\ne\emptyset$ なるものは，既約な良い
表示でランク $\le r_{\max}$ のものをもつからである（命題~\ref{prop:good-normalization}）．
命題~\ref{prop:bracketed-family} の標的 $L_n$ はクラスに属し，
$L_n\setminus\{W_n\}$ も属する：それは
補題~\ref{lem:type-deletion} により $\Cmcf{1}{s_U}$ に属し，有限である．
命題~\ref{prop:rare-type} により，$L_n$ のすべての特徴標本は
$W_n$ を含み，その長さは $3\cdot2^n-2$ であるが，一方 $\|G_n\|=O(n)$ かつ $t_{G_n}\le\theta_{G_n}\le4n$．
(iii) は，$f=1$ とした (i) と (ii) から従う．なぜなら，$\mathbf{LG}$ に
含まれない擬多様体は $U_1$ を含むからである
（定理~\ref{thm:fanout-one-data}）．
\end{proof}

仮説はすべての標本の上で $\mathcal R$ に属する；
特徴標本を含まない標本の上では，その言語はクラスに属する必要はない．これは
\cite{Yoshinaka2011} の学習器と同様である．項目~(i) はすべてのファンアウトで成り立つ．
項目~(ii) はファンアウト1でのみ証明されている；ファンアウト2でのその類似は未解決問題~\ref{op:fanout-two} の一部であり，
族 $L_n$ はそれを証明できない（定理~\ref{thm:doubling-no-lift}）．

\paragraph{$U_1$ の上への個々の型付け．}
定理~\ref{thm:fanout-one-data} は擬多様体についての主張である．
ファンアウト1の個々の型付けについては，どの語が $U_1$ の写し $\{e,z\}$ を実現するかが問題になる：
文字が互いに素な型 $e$ の語は指数的なデータを強いる．
一方，型がそれ自身の冪だけで実現される孤立した文字（定義~\ref{def:isolated-letter}）は多項式のデータを与える（下の定理~\ref{thm:two-letter-fanout-one}）．
$U_1$ の上への型付けについては，この二つの場合で尽くされる（系~\ref{cor:u1-classification}）．

\begin{corollary}[ファンアウト1での型付けの水準の障害]
\label{cor:typing-level-obstruction}
$r_{\max}\ge2$ とし，$s:\Sigma^+\twoheadrightarrow\mathsf S$ を型付けとする．
相異なる冪等元 $e,z\in\mathsf S$ が $ez=ze=z$ をみたし，
語 $u,v,w\in\Sigma^+$ が $s(u)=s(v)=e$，$s(w)=z$ をみたし，\emph{文字が互いに素}である
（$uvw$ の中で同じ文字が二度現れない）とする．
$H_n$ を，三つの新しい文字の上の命題~\ref{prop:three-letter-family} の言語 $L_n$ の
$\mathtt a\mapsto u$，$\mathtt b\mapsto v$，$\mathtt c\mapsto w$ による像とし，
$G''_n$ を $G_n$ の像，すなわち規則 $S\to E_n$，$E_0\to u$ と
$E_i\to vE_{i-1}E_{i-1}v\mid wE_{i-1}w$ をもつ文法とする．
このとき $H_n\in\Cmcf{1}{s}$ であり，$G''_n$ は $H_n$ のランク2の reduced good な表示で
\[
 \|G''_n\|=(2|v|+2|w|+9)n+|u|+6,\qquad
 \bar\theta_{G''_n}\le2|u|+2|v|+4n|w|
\]
をみたし，$\mathcal B^{\le r_{\max}}_{1,s}$，$\mathcal B_{1,s}$，あるいは $\Cmcf{1}{s}$ 全体で特徴標本をもつ任意の集合駆動型の作用素についての
$H_n$ のすべての特徴標本の総符号長は $(|u|+2|v|)2^n-2|v|$ より大きい．
したがってこの一つの型付けについて，定義~\ref{def:uniform-polydata} のようにファンアウト1，ランク $r_{\max}$ での
特徴データの長さを抑える多項式は存在しない．
\end{corollary}

\begin{proof}
$L_n$ は有限で Clark--Eyraud 代入可能である（命題~\ref{prop:three-letter-family}）ので，
補題~\ref{lem:app-u1-code} から $H_n$ も Clark--Eyraud 代入可能である；
それは有限なので $H_n\in\Cmcf{1}{s}$ である．
文法は good であり，$E_i\Rightarrow^*w^iuw^i$ なので reduced である．
大きさは定義~\ref{def:standard-mcfg} から従う：開始規則の大きさは $3$，基底の規則は $1+|u|$，
各段の二つの規則は $2|v|+5$ と $2|w|+3$ である．
他の選択をすべて規則 $E_j\to wE_{j-1}w$ で補うと，各規則は長さ高々 $|u|+2n|w|$（開始規則と単項の規則），
あるいは $2|u|+2|v|+(2n+2i-4)|w|$（段 $i$ の二項の規則）の成功した使用をもつ．
$\{e,z\}$ では $e$ は単位元で $z$ は吸収元なので，$H_n$ の語が型 $e$ をもつのは，それが
$\mathtt c$ を含まない $L_n$ の語，すなわち $W_n$ の像であるとき，かつそのときに限る．
像 $h(W_n)$ の長さは $|h(W_i)|=2|h(W_{i-1})|+2|v|$ から $(|u|+2|v|)2^n-2|v|$ であり，
命題~\ref{prop:rare-type} はそれをそのようなすべての標本に入れることを強いる．
\end{proof}

$u,v,w$ が文字ならば，これは型 $e$ の二つの文字の場合である；
特に，型 $1$ の文字を少なくとも二つもつ $U_1$ の上へのすべての型付けに適用できる．全射性から型 $\zeta$ の文字があるからである．
型 $e$ の文字がないときにも適用できる：
単位元 $1$，$g^2=1$ と零元 $z$ をもつ $\mathsf S=\{1,g,z\}$ と，
型付け $\mathtt a,\mathtt b,\mathtt c,\mathtt d\mapsto g$，$\mathtt t\mapsto z$ について，
$u=\mathtt{ab}$，$v=\mathtt{cd}$，$w=\mathtt t$，$e=1$ ととればよい．
文字が互いに素であることなしに，型 $e$ の可換でない二つの語が常に十分かどうかは一般には未解決である；
二つの語が可換であるのはそれらが一つの語の冪であるとき，かつそのときに限るので，そのような語がないことは，
型 $e$ のすべての語が一つの原始的な語の冪であることを意味する．
$U_1$ の上への型付けではこの問いは生じない：型 $1$ の文字が二つあればこの系が適用でき，
一つ（例えば $\mathtt a$）ならば型 $1$ の語は $\mathtt a$ の冪であり，データは多項式である（系~\ref{cor:u1-classification}）．

$\Sigma=\{\mathtt a,\mathtt b\}$ とし，$s_{U,2}:\Sigma^+\twoheadrightarrow U_1$ を $\mathtt a$ を
$1$ に，$\mathtt b$ を $\zeta$ に写すもの（二文字の型付け）とする．
型 $1$ の文字が一つしかないこの型付けには系~\ref{cor:typing-level-obstruction} は適用できない；
それでも，そのデータは対数的な大きさの表示について $\bar\theta_G$ の2乗になりうる．

\begin{proposition}[ファンアウト1での二文字の型付けの2次のデータ]
\label{prop:binary-u1-quadratic}
$m,n\ge1$ について，$L_{m,n}\subseteq\{\mathtt a,\mathtt b\}^+$ を
$\{\mathtt a^{2^m},\mathtt b\}$ の語 $2^n$ 個の連接全体の集合とする．
\begin{enumerate}[label=\textup{(\roman*)},leftmargin=*]
\item $L_{m,n}\in\Cmcf{1}{s_{U,2}}$．
\item $L_{m,n}$ は規則のランクが2のファンアウト1の reduced good な表示 $G_{m,n}$ で
$\|G_{m,n}\|=6m+6n+9$，$\bar\theta_{G_{m,n}}=2^m+2^n-1$ をみたすものをもつ．
\item $\mathcal B_{1,s_{U,2}}$，$\mathcal B^{\le r}_{1,s_{U,2}}$（$r\ge1$），
および $\Cmcf{1}{s_{U,2}}$ 全体で特徴標本をもつすべての集合駆動型の作用素について，
$L_{m,n}$ のすべての特徴標本 $F$ は $\|F\|_+>2^{m+n}$ をみたす．
\end{enumerate}
したがって $m=n$ のとき $\|G_{n,n}\|=12n+9$，$\bar\theta_{G_{n,n}}=2^{n+1}-1$，
$\|F\|_+>\bar\theta_{G_{n,n}}^{\,2}/4$ である；特に，この型付けについてファンアウト1では，
定理~\ref{thm:local-group-polydata} と対照的に，$p$ を多項式とする $p(\|G\|)\,\bar\theta_G$ の形の評価は存在しない．
\end{proposition}

証明は付録~\ref{app:u1-typings} にある．純な語 $\mathtt a^{2^{m+n}}$ は塊の長さと塊の個数の積であり，
一方 $\bar\theta_G$ はそれらの和程度である．長い塊を使う最短の文は，他のすべての位置を $\mathtt b$ で埋めるからである．
$\mathtt a$ と $\mathtt b$ 以外の文字を含む語は $L_{m,n}$ の中で空の分布をもつので，
この命題は，$\mathtt a$ が型 $1$ の唯一の文字で $\mathtt b$ が型 $\zeta$ をもつ $U_1$ の上へのすべての型付けについて，そのまま成り立つ．
それでもファンアウト1でのこれらの型付けのデータは多項式であり（下の系~\ref{cor:two-letter-fanout-one}），
$s_{U,2}$ についてはすべてのファンアウト $f\ge2$ でもそうである（系~\ref{cor:two-letter-semilattice}）．

\paragraph{孤立した文字と二文字．}
定理~\ref{thm:fanout-one-data}(iii) の族は三文字を使い，そのうち二つは型 $1$ である．
型がそれ自身の冪だけで実現される文字は，どのアルファベットの上でも，すでに多項式のデータを与える．

\begin{definition}[孤立した文字]
\label{def:isolated-letter}
文字 $\mathtt a\in\Sigma$ が型付け $s$ について\emph{孤立している}とは，
$s^{-1}(s(\mathtt a))=\mathtt a^+$ であり，$\mathtt a$ を含み $\mathtt a^+$ に入らないすべての語が一つの共通の型をもつことである．
\end{definition}

例えば $\Sigma=\{\mathtt a,\mathtt b\}$ とし，$c$ を\emph{内容型付け}，すなわち $c(w)$ が $w$ の文字の集合であるような，
和集合についての三元の半束 $\{\{\mathtt a\},\{\mathtt b\},\Sigma\}$ の上への型付けとする；
$\mathtt a$ を含み $\mathtt a^+$ の外にある語の内容は $\Sigma$ なので，両方の文字が $c$ について孤立している．
$s:\Sigma^+\twoheadrightarrow U_1$ で $\mathtt a$ が型 $1$ の唯一の文字ならば，$s^{-1}(1)=\mathtt a^+$ であり，
他のすべての語は型 $\zeta$ をもつので，どのアルファベットの上でも $\mathtt a$ は孤立している．
三文字以上の上では，$\mathtt{ab}$ と $\mathtt{ac}$ の内容が異なるので，内容型付けについて孤立した文字はない．

\begin{theorem}[孤立した文字：ファンアウト1での短い純な代表元]
\label{thm:two-letter-fanout-one}
$\mathtt a$ を型付け $s:\Sigma^+\twoheadrightarrow\mathsf S$ について孤立した文字とし，
$L\in\Cmcf{1}{s}$，$G$ を $L\cap\Sigma^+$ のファンアウト1で任意のランクの reduced good な表示とする．
$G$ の，あるいはその二分化 $G'$（付録~\ref{app:two-letter-fanout-one}）の，$\mathtt a^+$ の生成語をもつすべての非終端記号 $A$ について，
型付き非終端記号 $A_{s(\mathtt a)}$ は
\[
 \tau_{A_{s(\mathtt a)}}\le C(\|G\|+2)^{22}\bar\theta_G^{\,8}
\]
をみたす．ここで $C<10^{43}$ は $\Sigma$，$s$，ランクによらない絶対定数である．
特に，$L$ が $\mathtt a^+$ の語を含めば，その最短のものの長さは高々 $C(\|G\|+2)^{22}\bar\theta_G^{\,8}$ である．
\end{theorem}

この定理は孤立した文字の型だけを評価する．他の型は，次の系の二つの場合について，補題~\ref{lem:app-h-other} で評価される．
また，$\mathtt a$ を含み $\mathtt a^+$ の外にあるすべての語が一つの型をもつという仮定は本質的である：
三文字以上の内容型付けではこれが成り立たず，データは指数的になりうる（下を見よ）．

\begin{corollary}[孤立した文字からの多項式のデータ]
\label{cor:two-letter-fanout-one}
$r_{\max}\ge2$ とし，$s$ を次のいずれかとする．
\begin{enumerate}[label=\textup{(\roman*)},leftmargin=*]
\item 正の像が半束である $\Sigma=\{\mathtt a,\mathtt b\}$ 上の型付け（例えば $c$ や $s_{U,2}$）．
\item 任意の有限アルファベットの上の型付け $s:\Sigma^+\twoheadrightarrow U_1$ で，型 $1$ の文字がちょうど一つのもの．
\end{enumerate}
$L\cap\Sigma^+$ の任意のランクのファンアウト1の reduced good な表示 $G$ をもつすべての $L\in\Cmcf{1}{s}$ は，
$\mathcal B^{\le r_{\max}}_{1,s}$ と $\mathcal B_{1,s}$ に共通の特徴標本 $K$ で
\[
 \|K\|_+\le C'r_{\max}(\|G\|+2)^{23}\bar\theta_G^{\,8}+\mathbf 1_{\{\eps\in L\}}
\]
をみたすものをもつ．ここで $C'<10^{45}$ は絶対定数である．
\end{corollary}

\begin{corollary}[ファンアウト1での $U_1$ の上への型付け]
\label{cor:u1-classification}
$r_{\max}\ge2$ とし，$s:\Sigma^+\twoheadrightarrow U_1$ を型付けとする．
定義~\ref{def:uniform-polydata} が各型付けに求める評価が，ファンアウト1，ランク $r_{\max}$ で $s$ について成り立つのは，
型 $1$ の文字がちょうど一つのとき，かつそのときに限る．より詳しく：
\begin{enumerate}[label=\textup{(\roman*)},leftmargin=*]
\item 型 $1$ の文字がちょうど一つならば，データはすべてのランクの表示について系~\ref{cor:two-letter-fanout-one} のとおりに評価され，
命題~\ref{prop:binary-u1-quadratic} により，$p$ を多項式とするどの $p(\|G\|)\bar\theta_G$ でも評価されない；
\item 型 $1$ の文字が少なくとも二つあれば，系~\ref{cor:typing-level-obstruction} の標的は，
そのランク2の reduced good な表示の大きさと厚さについて指数的なデータを要する．
\end{enumerate}
\end{corollary}

\begin{proof}
$s$ は全射で $U_1$ は文字の型で生成されるので，型 $1$ の文字と型 $\zeta$ の文字がある．
(i) は，系~\ref{cor:two-letter-fanout-one}(ii) と，命題~\ref{prop:binary-u1-quadratic} およびその後の注意を，型 $\zeta$ の文字を $\mathtt b$ として使ったものである．
(ii) について，$\mathtt a\ne\mathtt b$ が型 $1$，$\mathtt c$ が型 $\zeta$ をもつとする；
冪等元 $e:=1$ と $z:=\zeta$ は相異なり $ez=ze=z$ をみたし，$u:=\mathtt a$，$v:=\mathtt b$，$w:=\mathtt c$ は
$s(u)=s(v)=e$，$s(w)=z$ をみたす文字が互いに素な語なので，系~\ref{cor:typing-level-obstruction} が適用できる．
\end{proof}

定理~\ref{thm:two-letter-fanout-one} と系~\ref{cor:two-letter-fanout-one} の証明は付録~\ref{app:two-letter-fanout-one} にある．
考え方は次のとおりである．非終端記号の最短の生成語 $\mathtt a^n$ の，節点の数が最小の導出 $T$ をとる．
一つのラベルの距離 $d$ の二つの出現から，純な塊の消去によって，$\mathtt a^d$ と可換で $z_0z_0\sim\mathtt a^{c_0}z_0$ をみたす，
$\mathtt a^+$ の外の基底 $z_0$ が得られる．
$T$ の重い道の上のラベルの $\mathtt a^+$ の外の短い生成語は，両側をそれぞれ高々 $d$ 文字で詰めると $\mathtt a^ez_0$ に潰れる；
孤立性が使われるのはちょうどここであり，それは $\mathtt a$ を含み $\mathtt a^+$ の外にあるすべての語を比較可能にする．
$\mathtt a^fz_0$ が与えられた分布の類に入る指数 $f$ の集合は，ある $p_0$ を法とする剰余類である；
$p_0$ が存在すれば $\mathtt a\sim\mathtt a^{1+p_0}$ となり，純な語は短い同値な代表元をもつ．
そうでなければ所属は符号の価値で決まる．各符号はそれまでの符号の和を超えるように選ばれ，これは符号一つあたり高々 $3$ 倍の費用で済む．
これらの符号の置換を，倍増族の場合と同様に，文法の大きさに対して数え上げる．

系~\ref{cor:two-letter-fanout-one} は，すべてのアルファベットについて量化する定理~\ref{thm:fanout-one-data}(i) に矛盾しない：
半束は $\mathbf{LG}$ に含まれない擬多様体をなし，失敗は三文字の上で起きる．
二文字上の半束の上へのすべての型付けの正の像は $c$ の像の商なので，系の (i) はそれらすべてを覆う．
三文字以上では，内容型付け $c_\Sigma$ について孤立が成り立たないことは必然である：
命題~\ref{prop:three-letter-family}(ii) から $\Sigma\supseteq\{\mathtt a,\mathtt b,\mathtt c\}$ について $L_n\in\Cmcf{1}{c_\Sigma}$ であり，
$W_n$ はその内容 $\{\mathtt a,\mathtt b\}$ の唯一の語なので，命題~\ref{prop:rare-type} は長さ少なくとも $3\cdot2^n-1$ のデータを強いる．
証明の中で失敗が起きるのは符号を潰すところである：一つの非終端記号の内容 $\{\mathtt a,\mathtt c\}$ の短い生成語と
内容 $\{\mathtt a,\mathtt b\}$ の生成語は同値である必要がない．
したがって，ファンアウト1では，$U_1$ の上への型付けが多項式のデータをもつのは型 $1$ の文字がちょうど一つのとき，かつそのときに限る（系~\ref{cor:u1-classification}）；
そのときデータは多項式であるが，$p(\|G\|)\,\bar\theta_G$ の形ではない（命題~\ref{prop:binary-u1-quadratic}）．

表~\ref{tab:scope} は，各主な主張について，固定するパラメータ，
作用素および付帯条件をまとめる．

\begin{table}[htbp]
\caption{主な主張の範囲．「標準」は
$\mathcal B^{\le r_{\max}}_{f,s}$ と $\mathcal B_{f,s}$ の両方（一つの共通の標本で）を意味する；
「クラス全体」は，クラス全体で特徴標本をもつ
すべての集合駆動型作用素を意味する．別に述べない限り，$G$ は $L\cap\Sigma^+$ の既約な良い表示であって，
ファンアウト $\le f$ かつランク $\le r_{\max}$ のものである．}
\label{tab:scope}
\centering
\footnotesize
\setlength{\tabcolsep}{3pt}
\renewcommand{\arraystretch}{1.15}
\begin{tabularx}{\linewidth}{@{}>{\raggedright\arraybackslash}p{0.2\linewidth}>{\raggedright\arraybackslash}p{0.22\linewidth}>{\raggedright\arraybackslash}p{0.17\linewidth}>{\raggedright\arraybackslash}X@{}}
\hline
主張 & 固定するもの & 作用素 & 評価と条件\\
\hline
定理~\ref{thm:A}（定理~\ref{thm:local-group-polydata}）
 & $f$，$r_{\max}$，$\Sigma$，$s$ で $s(\Sigma^+)\in\mathbf{LG}$ のもの
 & 標準
 & 上界 $O(\|G\|^2\bar\theta_G)$，すべての $L\in\Cmcf{f}{s}$ とすべての $G$ について\\
定理~\ref{thm:B}(i)（命題~\ref{prop:linear-factor-all-fanout}）
 & $f$，任意の $s$
 & 標準（$r\ge1$），クラス全体
 & 有限の標的に対する下界 $\Omega(\|G\|\bar\theta_G)$；$|\Sigma|\ge4f+1$（$|\Sigma|\ge2$，ただし $f=1$ の場合）を要する；
   ランク2の表示\\
定理~\ref{thm:B}(ii)（系~\ref{cor:trivial-exact}）
 & $f$，$r_{\max}\ge1$，$s_{\mathrm{triv}}$
 & $\mathcal B_{f,s_{\mathrm{triv}}}$
 & 上界 $O(\|G\|\bar\theta_G)$（Yoshinaka \cite{Yoshinaka2011}）；$r_{\max}\ge2$ かつ $|\Sigma|\ge4f+1$（$\ge2$，ただし $f=1$ の場合）のとき達成される\\
定理~\ref{thm:C}(i)（定理~\ref{thm:fanout-one-data}）
 & $f=1$，$r_{\max}\ge2$，一つの擬多様体
 & $\mathcal B^{\le r_{\max}}_{1,s}$（定義~\ref{def:uniform-polydata}）
 & 一様多項式データをもつのは $\mathbf V\subseteq\mathbf{LG}$ のとき，かつそのときに限る\\
定理~\ref{thm:C}(ii)
  & $f=1$，三文字上の $s_U$
 & 標準，クラス全体
 & 下界 $3\cdot2^n-1$，大きさと厚さが $O(n)$ の表示に対して\\
系~\ref{cor:typing-level-obstruction}
 & $f=1$，$r_{\max}\ge2$，文字が互いに素な $u,v,w$ をもつ型付け
 & 標準，クラス全体
 & 下界 $(|u|+2|v|)2^n-2|v|$，大きさと厚さが $O(n)$ の表示に対して\\
命題~\ref{prop:binary-u1-quadratic}
 & $f=1$，$s_{U,2}$，あるいは型 $1$ の文字が一つの $U_1$ の上への任意の $s$
 & 標準（$r\ge1$），クラス全体
 & 下界 $\bar\theta_G^{\,2}/4$，$\|G\|=O(\log\bar\theta_G)$\\
系~\ref{cor:two-letter-semilattice}（定理~\ref{thm:two-letter-content}）
 & $f\ge2$，$r_{\max}\ge2$，$\Sigma=\{\mathtt a,\mathtt b\}$，半束の上への $s$
 & 標準
 & 上界 $O_{f,r_{\max}}(\|G\|^3\bar\theta_G^{\,4}\log\bar\theta_G)$；Parikh 飽和でなければ $O_{f,r_{\max}}(\|G\|^2\bar\theta_G)$\\
命題~\ref{prop:sl-reduction}
 & $f$，$r_{\max}$
 & $\mathcal B^{\le r_{\max}}_{f,s}$
 & 境界が $f$ で動くのは，すべての $c_\Sigma$ が $f$ で一様な多項式のデータをもつとき，かつそのときに限る\\
系~\ref{cor:two-letter-fanout-one}
 & $f=1$，$r_{\max}\ge2$；$\Sigma=\{\mathtt a,\mathtt b\}$ で半束の上への $s$，あるいは任意の $\Sigma$ で型 $1$ の文字が一つの $U_1$ の上への $s$
 & 標準
 & 上界 $O(r_{\max}(\|G\|+2)^{23}\bar\theta_G^{\,8})$（絶対定数），すべての $L\in\Cmcf{1}{s}$ と任意のランクのすべての $G$ について\\
系~\ref{cor:u1-classification}
 & $f=1$，$r_{\max}\ge2$，$U_1$ の上への $s$
 & 標準，クラス全体
 & 型 $1$ の文字がちょうど一つのとき，かつそのときに限り多項式のデータ\\
系~\ref{cor:iptd}
 & $f$，$r_{\max}$，$s$
 & 集合駆動型
 & $\mathbf{LG}$ について IPTtD，すべての $f$ で；$f=1$ ではちょうど $\mathbf{V}\subseteq\mathbf{LG}$ について\\
定理~\ref{thm:doubling-no-lift}
 & 任意の $f\ge2$，任意の $s$
 & なし（所属）
  & 定理~\ref{thm:C}(ii) の族とその文字が互いに素な像は $\Cmcf{f}{s}$ を離れ，その強制された語はそこで決してその型の唯一の語ではない；$f\ge2$ でのデータの評価はない\\
命題~\ref{prop:lg-needed-target}，\ref{prop:nonabelian-needed-target}
 & 字母と標的 $L_{\mathrm A}$，$L_{\mathrm B}$
 & なし（所属）
 & 非文脈自由，ファンアウト2，ランク2；$f\ge2$ で $\Cmcf{f}{s_{\mathrm A}}$，$\Cmcf{f}{s_{\mathrm B}}$ に入り，すべての $f\ge1$ でそれぞれどの窓--群型付け，どの窓--可換群型付けのクラスにも入らない；データの評価はない\\
未解決
 & 
 & 
 & $\mathbf{LG}$ が $f\ge2$ で境界であるかどうか（未解決問題~\ref{op:fanout-two}），すなわち $|\Sigma|\ge3$ のすべての $c_\Sigma$ が $f$ で一様な多項式のデータをもつかどうか\\
\hline
\end{tabularx}
\end{table}

\section{ファンアウト2：持ち上がらないことと二文字}
\label{sec:fanout-two}
\label{subsec:power-transfer}

この節は，定理~\ref{thm:C}(ii) の指数的な族と，
その文字が互いに素な準同型による像が，ファンアウト
2には持ち越されないことを示し（定理~\ref{thm:doubling-no-lift}），
未解決問題~\ref{op:fanout-two} を内容型付けに帰着させ（命題~\ref{prop:sl-reduction}），
すべてのファンアウト $f\ge2$ で二文字のアルファベットを決着させる（定理~\ref{thm:two-letter-content}）．

\begin{lemma}[二つのタプル成分にわたる冪の移送]
\label{lem:power-transfer}
$f\ge2$，$L\in\Cmcf{f}{s}$，$u\in\Sigma^+$ とし，
$i_0,p\ge1$ を $s(u)$ の最終的な指数と周期とする．すなわち
$s(u)^{t+p}=s(u)^t$ が $t\ge i_0$ について成り立つとする．
$i,j\ge i_0$，$\delta\in p\mathbb Z_{>0}$ であり，かつ $L$ のある語が
因子 $u^{i+\delta+j+1}$ を含むとする．このとき，すべての
$w_{\mathrm L},w_{\mathrm R}\in\Sigma^*$ と $w_{\mathrm M}\in\Sigma^+$ について，
\[
 w_{\mathrm L}u^iw_{\mathrm M}u^{\delta+j}w_{\mathrm R}\in L\quad\Longleftrightarrow\quad
 w_{\mathrm L}u^{i+\delta}w_{\mathrm M}u^jw_{\mathrm R}\in L.
\]
\end{lemma}
\begin{proof}
与えられた語を $yu^{i+\delta+j+1}z$ と書く．
タプル $\tuple x=(u^i,u^{\delta+j})$ と
$\tuple x'=(u^{i+\delta},u^j)$ は，最終的な周期性により，成分ごとに同じ $s$-型をもつ．
それらはまた，共通の許容な
受理文脈 $E=y\blank_1u\blank_2z$ をもつ．なぜなら
$E[\tuple x]=E[\tuple x']=yu^{i+\delta+j+1}z\in L$ だからである．
したがって $(f,s)$-条件はそれらの二成分の分布を等しくする．
それらを許容文脈 $w_{\mathrm L}\blank_1w_{\mathrm M}\blank_2w_{\mathrm R}$ で試す；
条件 $w_{\mathrm M}\ne\eps$ は本質的である．
\end{proof}

\begin{theorem}[倍増族は持ち上がらない]
\label{thm:doubling-no-lift}
$s:\Sigma^+\twoheadrightarrow\mathsf S$ を任意の固定した有限半群型付けとし，
$u_{\mathtt a},u_{\mathtt b},u_{\mathtt c}\in\Sigma^+$ を文字が互いに素なものとし，$\varphi$ を
準同型 $\mathtt a\mapsto u_{\mathtt a}$，$\mathtt b\mapsto u_{\mathtt b}$，
$\mathtt c\mapsto u_{\mathtt c}$ とし，$H_n:=\varphi(L_n)$ とする．ここで
$L_n$ は命題~\ref{prop:three-letter-family} の族である\textup(したがって $H_n=L_n$ である，
$\varphi$ が恒等写像のときは\textup)．$n\ge2|\mathsf S|+2$ かつ $f\ge2$ とする．
すべての $L\in\Cmcf{f}{s}$ で $\varphi(W_n)$ を含むものは，同じ
型，同じ Parikh ベクトルをもち，$u_{\mathtt c}$ の文字を含まない別の語を含む．特に
$H_n\notin\Cmcf{f}{s}$ であり，$\Cmcf{f}{s}$ のどの成員においても $\varphi(W_n)$ は
その型の唯一の語ではない．
\end{theorem}
\begin{proof}
準同型 $\varphi$ は単射であり，$H_n$ の語で
$u_{\mathtt c}$ の文字を含まないものは，$L_n$ の語で $\mathtt c$ を含まないものの像であるので，
$\varphi(W_n)$ はそのような唯一の語である．最終的な指数 $i_0\ge1$ と周期
$p\ge1$ を $s(u_{\mathtt b})$ について，$i_0+p\le|\mathsf S|+1$ となるように選ぶ；このとき $n\ge2i_0+p+1$．
$i=i_0$，$\delta=p$ および $j=n-i-\delta-1\ge i_0$ とおく．語
$\varphi(W_n)=u_{\mathtt b}\,\varphi(W_{n-1})\,\varphi(W_{n-1})\,u_{\mathtt b}$ は
$u_{\mathtt b}^nu_{\mathtt a}$ で始まるので，$u_{\mathtt b}^{i+\delta+j+1}$ を含む．
$u_{\mathtt b}^i$ を，最初のものの後の $i$ 個の $u_{\mathtt b}$ の複製とし，
$u_{\mathtt b}^{\delta+j}$ を，$\delta+j\le n-1$ 個の複製であって，最初の
$u_{\mathtt a}$（$\varphi(W_{n-1})$ の二番目の複製の）の直前にあるものとする．そしてそれの前には少なくとも
$n-1$ 個の $u_{\mathtt b}$ の複製がある．このとき
\[
 \varphi(W_n)=w_{\mathrm L}u_{\mathtt b}^i w_{\mathrm M}u_{\mathtt b}^{\delta+j}w_{\mathrm R}
\]
であり，ここで $w_{\mathrm L}=u_{\mathtt b}$ であり，$w_{\mathrm M}\ne\eps$ は最初の
$u_{\mathtt a}$ を含む．$L\in\Cmcf{f}{s}$ が $\varphi(W_n)$ を含むとする．
補題~\ref{lem:power-transfer} により，
$w':=w_{\mathrm L}u_{\mathtt b}^{i+\delta}w_{\mathrm M}u_{\mathtt b}^jw_{\mathrm R}\in L$．
二つの埋めたタプルは成分ごとに等しい型をもつので，$s(w')=s(\varphi(W_n))$；$w'$ は
$\varphi(W_n)$ の Parikh ベクトルをもち，$u_{\mathtt c}$ の文字を含まないが，
$u_{\mathtt b}^{n+\delta}u_{\mathtt a}$ で始まり，$u_{\mathtt b}^nu_{\mathtt a}$ で始まるのではない．したがって
$w'\ne\varphi(W_n)$ である．なぜなら $u_{\mathtt a}$ と $u_{\mathtt b}$ は共通の文字をもたないからである．
$L=H_n$ については，これは上で示した $\varphi(W_n)$ の一意性に反する．
\end{proof}

\begin{remark}
冪の移送は語の総長と Parikh ベクトルを保つので，
稀な型の代表が短いことは示さない；それは
命題~\ref{prop:three-letter-family} の族とその文字が互いに素な像を，
語を加えた後でも，排除するだけである．
\end{remark}

\paragraph{半束への帰着．}
$\mathbf{SL}$ を有限の半束（可換で冪等な半群）の擬多様体とする．
空でない有限のアルファベット $\Sigma$ について，\emph{内容型付け} $c_\Sigma$ は
$w\in\Sigma^+$ を $w$ に現れる文字の集合に写す．これは和集合についての $\Sigma$ の空でない部分集合の上への準同型である．

\begin{lemma}[粗い型付けは特徴標本を受け継ぐ]
\label{lem:typing-monotone}
$s\preceq t$ を型付け，$f,r\ge1$，$L\in\Cmcf{f}{s}$ とする．すべての有限の $K\subseteq L$ について
\[
 L\bigl(\mathcal B^{\le r}_{f,t}(K)\bigr)\subseteq
 L\bigl(\mathcal B^{\le r}_{f,s}(K)\bigr)\subseteq L
\]
であり，$\mathcal B_{f,t}$ と $\mathcal B_{f,s}$ についても同様である．
したがって $\mathcal B^{\le r}_{f,t}$（あるいは $\mathcal B_{f,t}$）についての $L$ のすべての特徴標本は
$\mathcal B^{\le r}_{f,s}$（あるいは $\mathcal B_{f,s}$）についての特徴標本である．
\end{lemma}

\begin{proof}
二つの文法は同じ非終端記号と同じ Start，Const，Comp 規則をもつ（定義~\ref{def:reconstruction}）．
$t$-文法の Link 規則は，共通の文脈をもち $t$-型が等しい観測されたタプルを結ぶ；
$\ker t\subseteq\ker s$ なのでそれらの $s$-型も等しく，それは $s$-文法の Link 規則である．
これが最初の包含を与える；二番目は命題~\ref{prop:reconstruction-soundness} である．
$F$ が $t$-作用素について特徴的で $F\subseteq K\subseteq L$ ならば，左辺は $L$ に等しく，したがって中辺もそうである．
\end{proof}

\begin{proposition}[ファンアウトの境界の半束への帰着]
\label{prop:sl-reduction}
$f\ge1$，$r_{\max}\ge1$ を固定し，一様な多項式の特徴データの長さをファンアウト $f$，ランク $r_{\max}$ で読む
（定義~\ref{def:uniform-polydata}）．
\begin{enumerate}[label=\textup{(\roman*)},leftmargin=*]
\item $\mathbf V\not\subseteq\mathbf{LG}$ であるすべての擬多様体は $\mathbf{SL}$ を含む．
したがって，一様な多項式の特徴データの長さをもつ擬多様体がちょうど $\mathbf{LG}$ に含まれるものであるのは，
$\mathbf{SL}$ がそれをもたないとき，かつそのときに限る．
\item $\mathbf{SL}$ が一様な多項式の特徴データの長さをもつのは，すべての空でない有限のアルファベット $\Sigma$ について
内容型付け $c_\Sigma$ が定義~\ref{def:uniform-polydata} の条件をみたすとき，かつそのときに限る．
\end{enumerate}
\end{proposition}

\begin{proof}
(i) 有限半群 $\mathsf S\notin\mathbf{LG}$ はアルファベット $\mathsf S$ の上のその評価型付けの正の像なので，
$U_1$ と同型な部分半群をもつ（定理~\ref{thm:lg-characterization}）；したがって $U_1\in\mathbf V$ である．
すべての有限の半束 $\mathsf H$ は $x\mapsto(\chi_z(x))_{z\in\mathsf H}$ によって $U_1^{\mathsf H}$ に埋め込まれる．
ここで $zx=z$ ならば $\chi_z(x)=1$，そうでなければ $\chi_z(x)=\zeta$ である
（$zxy=z$ であるのは $zx=z$ かつ $zy=z$ のとき，かつそのときに限るので準同型である）．したがって $\mathbf{SL}\subseteq\mathbf V$ である．
定義~\ref{def:uniform-polydata} の性質は，量化される型付けが少なくなるので，擬多様体からその部分擬多様体へ受け継がれる．
$\mathbf{SL}$ がそれをもたなければ，$\mathbf V\not\subseteq\mathbf{LG}$ であるすべての擬多様体もそれをもたない．
一方 $\mathbf V\subseteq\mathbf{LG}$ であるすべての擬多様体は定理~\ref{thm:local-group-polydata} からそれをもつ．
逆に，$U_1\notin\mathbf{LG}$ なので $\mathbf{SL}\not\subseteq\mathbf{LG}$ である．

(ii) $c_\Sigma$ の像は半束であり，一方の向きを与える．
他方の向きについて，$\mathsf S\in\mathbf{SL}$ とする $s:\Sigma^+\twoheadrightarrow\mathsf S$ をとる．
$\mathsf S$ は可換で冪等なので，$s(w)$ は $w$ に現れる文字 $a$ についての $s(a)$ の積であり，
$s\preceq c_\Sigma$，$\Cmcf{f}{s}\subseteq\Cmcf{f}{c_\Sigma}$ である（定義~\ref{def:class-cmcf} の後の項目~(a)）．
$L\in\Cmcf{f}{s}$ と表示 $G$ について，長さ高々 $p_{c_\Sigma}(\|G\|,\bar\theta_G)$ の
$\mathcal B^{\le r_{\max}}_{f,c_\Sigma}$ についての特徴標本は，補題~\ref{lem:typing-monotone} から
$\mathcal B^{\le r_{\max}}_{f,s}$ についての特徴標本である；$p_{s}:=p_{c_\Sigma}$ ととればよい．
\end{proof}

定理~\ref{thm:fanout-one-data} により，$r_{\max}\ge2$ のときファンアウト1で $\mathbf{SL}$ はこの性質をもたない．
ただし二文字の上では，$\mathbf{SL}$ のすべての型付けはそこで多項式のデータをもつ（系~\ref{cor:two-letter-fanout-one}）．
ファンアウト2ではこの問いは未解決であり（未解決問題~\ref{op:fanout-two}），
次の定理はすべてのファンアウト $f\ge2$ で二文字のアルファベットを決着させる．
稀な型の下界は内容型付けに移る：$s\preceq c_\Sigma$ について，$s$-型 $m$ に含まれるすべての内容の型は，
最短の語の長さが少なくとも $\mu_m(L)$（$s$ について計算したもの）である．

\paragraph{すべてのファンアウト $f\ge2$ での二文字．}
語の $\Sigma^+$ への所属がその Parikh ベクトルのみによるとき，$L$ は \emph{Parikh 飽和}であるという．

\begin{theorem}[すべてのファンアウト $f\ge2$ での二文字の内容型付け]
\label{thm:two-letter-content}
$\Sigma=\{\mathtt a,\mathtt b\}$，$c:=c_\Sigma$，$f\ge2$，$r_{\max}\ge2$ とし，
$D_f:=(6f+1)3^f$，
$C_{f,r_{\max}}:=(r_{\max}+1)\bigl(4r_{\max}(3f+1)+D_f\bigr)+3$，
$H_m:=\sum_{k=1}^m1/k\le1+\ln m$（$H_0:=0$）とおく．
すべての $L\in\Cmcf{f}{c}$ と，$L\cap\Sigma^+$ のファンアウト $\le f$，ランク $\le r_{\max}$ のすべての reduced good な表示 $G$ について，
$\tau(G,c)\le C_{f,r_{\max}}|V|^2\bar\theta_G^{\,4}H_{\bar\theta_G}$ であり，作用素
$\mathcal B^{\le r_{\max}}_{f,c}$ と $\mathcal B_{f,c}$ は共通の特徴標本 $K$ で
\[
 \|K\|_+\le4^{fr_{\max}}\bigl(3+r_{\max}C_{f,r_{\max}}\bigr)\|G\|^3\bar\theta_G^{\,4}H_{\bar\theta_G}
 +\mathbf 1_{\{\eps\in L\}}
\]
をみたすものをもつ．
$L$ が Parikh 飽和でなければ $t_{\mathrm{typ}}(G,c)\le D_f(r_{\max}+1)|V|\bar\theta_G$ であり，
$\|K\|_+\le4^{fr_{\max}}\bigl(3+r_{\max}D_f(r_{\max}+1)\bigr)\|G\|^2\bar\theta_G
+\mathbf 1_{\{\eps\in L\}}$ ととれる．
\end{theorem}

\begin{corollary}[二文字の上の半束型付け]
\label{cor:two-letter-semilattice}
$f\ge2$，$r_{\max}\ge2$ とし，$s:\{\mathtt a,\mathtt b\}^+\twoheadrightarrow\mathsf S$ を
$\mathsf S\in\mathbf{SL}$ である型付け，例えば $s_{U,2}$（$\mathtt a\mapsto1$，$\mathtt b\mapsto\zeta$）とする．
このとき，定義~\ref{def:uniform-polydata} が各型付けに求める評価が，この固定した型付けについて，
$\mathcal B^{\le r_{\max}}_{f,s}$ と $\mathcal B_{f,s}$ に対し，定理~\ref{thm:two-letter-content} の評価で，ファンアウト $f$，ランク $r_{\max}$ で成り立つ．
\end{corollary}

\begin{proof}
命題~\ref{prop:sl-reduction}(ii) の証明と同様に $s\preceq c$，$\Cmcf{f}{s}\subseteq\Cmcf{f}{c}$ である；
定理~\ref{thm:two-letter-content} と補題~\ref{lem:typing-monotone} を適用する．
\end{proof}

証明は付録~\ref{app:two-letter-content} にある．それは，すべてのファンアウト $f\ge2$ で成り立つアリティ1と2の条件を，
一つの文字の連を通して使う：ある文字 $x$ について $x^4$ が $L$ の語に現れれば，
所属は $x$-骨格（$x$ の各極大の連を一つの $x$ で置き換えたもの）と文字 $x$ の個数のみによる．
また，連 $xx$ の隣の他方の文字の二つの連は，Parikh 飽和を強いるか，禁じられたより長い連に合流する．
飽和していなければ，これは最短の型付き導出の重い節点の反鎖を抑える；
飽和していれば，与えられた Parikh ベクトルの語の節点数最小の導出を直接解析する．
この定理はファンアウトを分離しない：系~\ref{cor:two-letter-fanout-one} により，二文字の半束型付けはファンアウト1でも多項式のデータをもつので，
二文字の上ではそれらはすべてのファンアウトで多項式のデータをもつ．
三文字以上では，内容 $\{\mathtt a,\mathtt b\}$ の長い成分が長い連を含むとは限らず，問いは未解決である．

\section{結論と未解決問題}
\label{sec:conclusion}

定理~\ref{thm:data-from-typed-thickness} と系~\ref{cor:distribution-representatives} は，
特徴データを最短の型付き産出と代表長に帰着させる．局所群型付けはすべてのファンアウトで
$O(\|G\|^2\bar\theta_G)$ のデータをもち（定理~\ref{thm:A}），積 $\|G\|\bar\theta_G$ は必要であり（定理~\ref{thm:B}），
ファンアウト1では $\mathbf{LG}$ が擬多様体に対する正確な境界で，その外では失敗は指数的である
（定理~\ref{thm:C}，系~\ref{cor:iptd}）．半束の上への型付けではこの失敗には三文字が要り，
ファンアウト1で $U_1$ の上への型付けが多項式のデータをもつのは型 $1$ の文字がちょうど一つのときである
（系~\ref{cor:two-letter-fanout-one}，\ref{cor:two-letter-semilattice} と~\ref{cor:u1-classification}）．
ファンアウト2では，ファンアウト1の指数的な族は持ち上がらない（定理~\ref{thm:doubling-no-lift}）．

\begin{openproblem}[ファンアウト2での境界]
\label{op:fanout-two}
$f\ge2$ と $r_{\max}\ge2$ を固定する．$\mathbf{LG}$ はファンアウト $f$ における擬多様体の境界であるか．
すなわち，$\mathbf{LG}$ に含まれない有限半群のすべての擬多様体は，
ファンアウト $f$ かつランク $r_{\max}$ で一様多項式特徴データ長をもたないか
（定義~\ref{def:uniform-polydata}）？
\end{openproblem}

$\mathbf{LG}$ に含まれる擬多様体は，すべてのファンアウトで一様多項式
特徴データ長をもつ
（定理~\ref{thm:local-group-polydata}）．
$\mathbf{LG}$ に含まれないすべての擬多様体は $U_1$ を含む（定理~\ref{thm:lg-characterization}）ので，
$U_1$ の上への一つの型付けがファンアウト $f$ で超多項式のデータを必要とするならば，そのとき直ちに
答えは肯定的である；命題~\ref{prop:rare-type} により，標的の族で，
ある型の最短の語が $\|G\|$ と
$\bar\theta_G$ について超多項式的に長いものがあれば十分であり，定理~\ref{thm:doubling-no-lift} は
倍増族を排除する．否定的な答えは，境界が
ファンアウトに依存することを示すであろう．
命題~\ref{prop:sl-reduction} により，この問題は内容型付け $c_\Sigma$ がファンアウト $f$ で一様な多項式のデータをもつかどうかを問う；
$|\Sigma|=2$ ではすべてのファンアウトでそれをもつ（系~\ref{cor:two-letter-fanout-one} と定理~\ref{thm:two-letter-content}）ので，
問題は少なくとも三文字のアルファベットに関わる．

稀な型 $1$ には型 $\zeta$ の文字が多く要る．型付け
$s:\Sigma^+\twoheadrightarrow U_1=\{1,\zeta\}$ について $\Sigma_1:=\{a\in\Sigma:s(a)=1\}$ とおくと
$s^{-1}(1)=\Sigma_1^+$ である．$\#_\zeta(w)$ で $w$ の中の型 $\zeta$ の文字の個数を表し，
$\kappa_\zeta(L):=\sup\{\#_\zeta(w):w\in L\}$ とおく．

\begin{proposition}[少ない $\zeta$ の文字からの短い型 $1$ の語]
\label{prop:zeta-capacity}
$s:\Sigma^+\twoheadrightarrow U_1$ を型付けとし，$G$ を，任意のファンアウトとランク $\le r_{\max}$ の
既約で $\eps$-free かつ非削除の MCFG で $\kappa_\zeta(L(G))<\infty$ をみたすものとし，この数を $\kappa_\zeta$ と書く．このとき
\[
 \mu_1(L(G))\le\max\bigl\{\bar\theta_G,\ \kappa_\zeta\,\|G\|\bigl(\|G\|+r_{\max}\bar\theta_G\bigr)\bigr\}.
\]
\end{proposition}

\begin{proof}
$\mu:=\mu_1(L(G))$ が有限で $\bar\theta_G$ を超えると仮定する．非終端記号 $A$ について，
$\lambda(A)$ を $\Sigma_1$ 上の $L_A(G)$ のタプルの総長の最小値（なければ $\infty$）とし，
$\operatorname{cap}(A):=\sup\{\#_\zeta(\tuple y):\tuple y\in L_A(G)\}$ とおく．長さ $\mu$ の
$\Sigma_1^+$ の語の導出木のうち，節点数が最小のもの $T$ をとる．$G$ は非削除なので，各節点 $\eta$
（ラベル $A_\eta$）は語の中に文字が現れる $\Sigma_1$ 上のタプル $\tuple x_\eta$ を導出し，
$\eta$ の部分木を $\Sigma_1$ 上のより短い $L_{A_\eta}(G)$ のタプルの導出で置き換えると
型 $1$ のより短い語が得られてしまう．したがって $|\tuple x_\eta|=\lambda(A_\eta)$ であり，
$\lambda$ は節点からその子へ増えない．ある道の上で一つのラベルが二度現れれば，上の部分木を
下の部分木で置き換えると，長さ高々 $\mu$ で節点の少ない型 $1$ の語が得られる；
したがって各道の節点は高々 $|V|$ 個である．

$G$ は既約なので，各 $A$ は $L(G)$ の中に外側の文脈をもち，$\operatorname{cap}(A)\le\kappa_\zeta$ である；
また $G$ は非削除なので，すべての規則 $A\to\rho(B_1,\ldots,B_m)$ について
$\operatorname{cap}(A)\ge\sum_j\operatorname{cap}(B_j)$ である．$\lambda(A)>\bar\theta_G$ のとき $A$ を
重いと呼ぶ．左辺が $A$ の規則は長さ高々 $\bar\theta_G$ の語の成功する導出で使われ，そこで $A$ から
導出されるタプルの総長は $\lambda(A)$ より小さいので，それは $\Sigma_1$ 上にない；したがって重いラベルの容量は少なくとも $1$ である．
$T$ の重い節点は根を含む部分木をなし，優加法性を用いた下からの帰納法により，重い節点の下にある，重い子をもたない
重い節点の個数は，そのラベルの容量以下である．したがって $T$ はそのような節点を高々 $\kappa_\zeta$ 個，
重い節点を高々 $\kappa_\zeta|V|$ 個もつ．語のすべての終端記号は，重い節点の規則の中（節点あたり高々
$\|G\|$ 個）にあるか，重い節点の高々 $r_{\max}$ 個の重くない子の下（子あたり高々 $\bar\theta_G$ 個）にある．
$|V|\le\|G\|$ とあわせて評価を得る．
\end{proof}

代入可能性もファンアウトも用いていない．したがって，命題~\ref{prop:rare-type} と $U_1$ の上への
型付けの型 $1$ を通じて未解決問題~\ref{op:fanout-two} に肯定的に答える族は，$\kappa_\zeta(L_n)$ が
$\|G_n\|$ と $\bar\theta_{G_n}$ について超多項式でなければならない．この議論は
命題~\ref{prop:sl-reduction} の内容型付け $c_\Sigma$ には適用されない：内容 $\{\mathtt a,\mathtt b\}$ を
避ける短い導出は単に $\mathtt b$ を省けばよく，それは容量を消費しない．

関連する量，すなわち保護付きタプル代入を保つ有限モノイド観測の大きさについては，
それが擬多様体の上で安定するアリティは $1$，$2$ または $\infty$ である \cite{KuriyamaSCL2026}．

\paragraph{指数の鋭さ．}
$O(\|G\|^2\bar\theta_G)$ において，窓--群型付けと局所群型付けのいずれについても，
一つの因子 $\|G\|$ は証拠の語を数え，もう一つは $\lambda$-極小な導出における根から葉への
道を節点のラベルの個数で抑える
（補題~\ref{lem:heavy-node-count}）．道の因子は証拠の
構成にとっては本物である：文法 $S\to SS\mid\mathtt a\mid\mathtt b\mid A_0\mid\cdots\mid A_n$，
$A_0\to\mathtt b$，$A_i\to\mathtt a\mid A_{i-1}\mathtt b$（$1\le i\le n$）は
$\{\mathtt a,\mathtt b\}^+$ を $\|G_n\|=10n+16$ かつ $\bar\theta_{G_n}=2$ で生成するが，
$\mathtt b^{n+1}$ は $A_n$ の産出で $s_{1,1}$-型
$\langle\mathtt b,\mathtt b\rangle$ をもつ唯一のものであるので，$t_{\mathrm{typ}}(G_n,s_{1,1})\ge n+1$．しかし，この
標的については $\tau(G_n,s_{1,1})\le2$ である
（系~\ref{cor:distribution-representatives}）ので，この例は
証拠の方法を制限するのであって，データを制限するのではない．自明な型付けについては，上界
$O(\|G\|\bar\theta_G)$ は標識族によって達成される
（系~\ref{cor:trivial-exact}）；
$\mathsf S\in\mathbf{LG}$ をもつすべての型付けについて，残る隔たりは一つの因子
$\|G\|$ であり，これは既知の例と上界との間の隔たりであって，証明された
比ではない．どちらの下界の議論もそれを閉じない：固定した型付けについて，
稀な型の議論は高々 $|\mathsf S|(t_{\mathrm{typ}}+1)$ を与え，標識
族は $O(\|G\|)$ 個の語で長さ $\bar\theta_G=t_{\mathrm{typ}}$ のものからなる．
非可換
群については，定数 $M_{\mathbb G,f}$ は $f$ について二重指数的である．これは
Ramsey 型の均質化（補題~\ref{lem:order-homogenize}）による；$|\mathbb G|$ 個の所有者は
$f=1$ で足り，$\mathrm{poly}(|\mathbb G|^f)$ 個の所有者が $f\ge2$ について足りるかどうかは，
可換群の場合のように，未解決である．ファンアウト1の孤立した文字については，系~\ref{cor:two-letter-fanout-one} の指数 $(23,8)$ は
付録~\ref{app:two-letter-fanout-one} の数え上げの議論から来るものであり，命題~\ref{prop:binary-u1-quadratic} の下界 $\bar\theta_G^{\,2}/4$ からは遠い．

\section*{謝辞}
著者は，Makoto Kanazawa 氏と Ryo Yoshinaka 氏に，有益な議論と
励ましに感謝する．

\section*{利益相反の申告}
著者は，本論文で報告した研究に影響を与えたと見なされ得る，
既知の競合する金銭的利害も
個人的関係もないことを申告する．

\section*{資金}
本研究は，公的，
商業的，非営利のいずれの部門の資金提供機関からも特定の助成を受けていない．

\section*{データの利用可能性}
本理論論文に関連する経験的データセットはない．すべての証明は
完全な形で与えている．

\appendix
\section{有限群における同時ブロック置換}
\label{app:group-blocks}

この付録は命題~\ref{prop:block-replacement-interface} を証明する．これは
補題~\ref{lem:marked-heavy-window-group} と~\ref{lem:local-group-shortest} で用いる．
可換な場合は零和の議論である．一般の場合は決して
因子を交換しない；それは三つの部分からなる：均質な点の
配置に対する順序の分類（補題~\ref{lem:group-block-order}），
境界の状態による区間の選択（補題~\ref{lem:group-interval}），および
均質な添字の Ramsey 型の抽出（補題~\ref{lem:two-sequence} と~\ref{lem:order-homogenize}）．

\paragraph{一般の場合の道案内．}
補題~\ref{lem:block-replacement}(b) の証明は，この付録の終わりにあり，
以下の段階をこの順に適用する；各段階は，それが捨てない所有者を
変えずに保ち，捨てた所有者のブロックを定数に吸収する
（それらはすべての行で元の値を保つ）．
\begin{center}
\footnotesize
\setlength{\tabcolsep}{3pt}
\renewcommand{\arraystretch}{1.15}
\begin{tabularx}{\linewidth}{@{}>{\raggedright\arraybackslash}p{0.2\linewidth}>{\raggedright\arraybackslash}X>{\raggedright\arraybackslash}p{0.3\linewidth}@{}}
\hline
段階 & 入力 $\to$ 出力 & 得られる不変性\\
\hline
1. 共通の型
 & $m\ge M_{{\mathbb G},f}$ 個の所有者 $\to$ $m/\mathcal P_f$ 個の所有者
 & ブロックの個数，行の個数および所有者内の順序が同じ\\
2. 補題~\ref{lem:two-sequence}，\ref{lem:order-homogenize}
 & $\to N_{{\mathbb G},f}=\max\{3,|{\mathbb G}|^f\}$ 個の所有者
 & 各比較「ブロック $t$（$i$ の）がブロック $u$（$j$ の）の左にある」は $i<j$ にわたって一定\\
3. 補題~\ref{lem:group-block-order}
 & すべての行 $\to$ 合計で高々 $f$ 個の類の領域，各々は $N_{{\mathbb G},f}$ 個のかたまりの連なり
 & 所有者ごとに一つのかたまり，増加または減少；領域の内部に他の点はない\\
4. 補題~\ref{lem:group-interval}
 & 領域 $\to$ 所有者の区間 $J=(a,b]$
 & 状態 $(\eta^{\mathcal Q}_p)_{\mathcal Q}$，$\eta^{\mathcal Q}_p=P_pQ_p^{-1}$ が $a$ と $b$ で一致し，したがってすべての順序付きの行の積が保たれる\\
\hline
\end{tabularx}
\end{center}
$f=2$ と ${\mathbb G}=S_3$ について，例~\ref{ex:nonabelian-reverse-owner} は
段階~3 の出力（行~1 上の一つの増加領域，行~2 上の一つの減少領域）と，
段階~4 における $\eta$ の使い方を示す：置き換えた所有者は
自身の寄与を保たず，二つの順序付きの行の積だけが保たれる．

\begin{definition}[ブロック置換族]
\label{def:block-family}
${\mathbb G}$ を有限群とし，$f\ge1$ とする．
幅 $f$ の，${\mathbb G}$ 上の，添字集合 $[m]$ をもつ\emph{ブロック置換族}は，$d\le f$ 本の\emph{行}からなる；
行 $\varrho$ の上には，ブロックの左から右への列 $\mathfrak b_{\varrho,1},\ldots,\mathfrak b_{\varrho,q_\varrho}$
と定数 $a_{\varrho,0},\ldots,a_{\varrho,q_\varrho}\in {\mathbb G}$ がある；すべてのブロック $\mathfrak b$ について，
\emph{所有者} $o(\mathfrak b)\in[m]$ と二つの値 $x_{\mathfrak b},x'_{\mathfrak b}\in {\mathbb G}$ がある；すべての
添字は高々 $f$ 個のブロックを所有する．$J\subseteq[m]$ について次のようにおく．
\[
 \operatorname{val}_\varrho(J):=a_{\varrho,0}\,y_{\varrho,1}\,a_{\varrho,1}\cdots y_{\varrho,q_\varrho}\,a_{\varrho,q_\varrho},
 \qquad
 y_{\varrho,j}:=
 \begin{cases}
 x'_{\mathfrak b_{\varrho,j}},& o(\mathfrak b_{\varrho,j})\in J,\\
 x_{\mathfrak b_{\varrho,j}},& o(\mathfrak b_{\varrho,j})\notin J.
 \end{cases}
\]
\end{definition}

\begin{lemma}[ブロック置換]
\label{lem:block-replacement}
幅 $f$ の，${\mathbb G}$ 上のブロック置換族において，空でない
$J\subseteq[m]$ で $\operatorname{val}_\varrho(J)=\operatorname{val}_\varrho(\emptyset)$ がすべての
行 $\varrho$ について成り立つものが存在する．ただし，次のいずれかが成り立つとする：
\textup{(a)} ${\mathbb G}$ が可換であり $m>\mathsf d({\mathbb G}^f)$，または
\textup{(b)} ${\mathbb G}$ が任意であり $m\ge M_{{\mathbb G},f}$，または
\textup{(c)} ${\mathbb G}$ が任意であり，$f=1$ かつ $m\ge|{\mathbb G}|$．ここで $\mathsf d$ と
$M_{{\mathbb G},f}$ は第~\ref{subsec:group-invariants}~節のとおりである．
\end{lemma}

\begin{proof}[\textup{(a)} の証明]
${\mathbb G}$ を加法的に書く．可換性により
$\operatorname{val}_\varrho(J)-\operatorname{val}_\varrho(\emptyset)=\sum_{i\in J}\delta_{i,\varrho}$ であり，ここで
$\delta_{i,\varrho}$ は $x'_{\mathfrak b}-x_{\mathfrak b}$ の和であって，ブロック $\mathfrak b$（行 $\varrho$ の，
$i$ が所有するもの）にわたるものである．列 $(\delta_i)_{i\in[m]}$（$\delta_i:=(\delta_{i,\varrho})_\varrho\in {\mathbb G}^d$）
は長さ $m>\mathsf d({\mathbb G}^f)\ge\mathsf d({\mathbb G}^d)$ をもつ（${\mathbb G}^d$ は ${\mathbb G}^f$ に埋め込まれる）ので，
空でない零和部分列 $(\delta_i)_{i\in J}$ をもつ．
\end{proof}

\begin{remark}[零和定数の大きさ]
\label{rem:davenport-size}
接頭辞和の鳩の巣論法は $\mathsf d({\mathbb G}^f)\le|{\mathbb G}|^f-1$ を与えるが，はるかに多くのことが成り立つ．
${\mathbb G}\cong C_{q_1}\oplus\cdots\oplus C_{q_\ell}$ と書き，$q=\exp({\mathbb G})\ge2$ とする．このとき
\[
 f\sum_{i=1}^{\ell}(q_i-1)\ \le\ \mathsf d({\mathbb G}^f)\ \le\ (q-1)+q\log\frac{|{\mathbb G}|^f}{q},
\]
ここで下界は基底からなる列によって達成され，上界は
van Emde Boas と Kruyswijk による \cite[Proposition~5.1.8 and
Theorem~5.5.5]{GeroldingerHalterKoch2006}；$p$-群については下界は
等号である \cite[Theorem~5.5.9]{GeroldingerHalterKoch2006}．したがって
$\mathsf d({\mathbb G}^f)=O(\exp({\mathbb G})\,f\log|{\mathbb G}|)$；${\mathbb G}=C_2$ については，反鎖の上界は
$2^f-1$ から $f$ に下がる．長さ $\mathsf d({\mathbb G}^d)$ の零和自由な列は存在するので，
この議論はこれ以上よくはならない；それらを実現する文法があるとは主張しない．
\end{remark}

\subsection{任意の有限群：順序の均質化}
\label{app:window-group}

\begin{lemma}[均質なブロック順序の補題]
\label{lem:group-block-order}
$\mathcal T$ を型の有限集合とし，$n\ge3$ とし，$(t,i)$，$t\in\mathcal T$，$i\in[n]$ を，直線上の相異なる
点とする．\textup{(i)} $(t,i)$ と $(u,i)$ の順序が
$i$ に依存せず，かつ \textup{(ii)} 各順序対 $(t,u)$（
$t=u$ を含む）について，$(t,i)$ と $(u,j)$ の順序がすべての $i<j$ について同じであるとする．このとき $\mathcal T$ は
類に分割され，次が成り立つ：各類の点は，他の点を含まない
連続した領域を占める；各類は向き
（$i$ について増加または減少）をもつ；そしてその領域は $n$ 個の連続した
\emph{かたまり}からなり，かたまりはその向きで添字ごとに一つずつあり，各かたまりは，その類の
すべての型の点をちょうど一つずつ，添字に依存しない順序で含む．
\end{lemma}

\begin{proof}
$t=u$ とした (ii) により，各型の点は $i$ について単調である．相異なる
$t,u$ について，$X$ は $(t,i)<(u,j)$ を，$Y$ は $(u,i)<(t,j)$ を，$i<j$ について意味するとする；どちらも
(ii) により一定である．

\emph{両方が増加．}$X\wedge Y$ は，両方の型の添字 $i$ のすべての点を，
添字 $j>i$ の点より前に置く（\emph{順位によって交互に並ぶ}）．$X\wedge\neg Y$ の下では，
異なる添字のすべての比較は $t$ を $u$ の左に置き，等しい添字についても同様である．
なぜなら $(u,i)<(t,i)$ は $(t,1)<(t,2)<(u,1)<(t,1)$ を与えるからである；したがって $t$ と $u$ は
\emph{分離される}．$\neg X\wedge Y$ は対称であり，$\neg X\wedge\neg Y$ は
$(t,1)<(t,2)<(u,1)<(u,2)<(t,1)$ を与える．

\emph{$t$ が増加，$u$ が減少．}$X\wedge Y$ は
$(t,2)<(u,3)<(u,1)<(t,2)$ を与え，$\neg X\wedge\neg Y$ は
$(t,3)<(u,2)<(t,1)<(t,3)$ を与える．$X\wedge\neg Y$ の下では，異なる
添字の比較は $t$ を $u$ の左に置き，$(u,i)<(t,i)$ は
$(t,1)<(u,2)<(u,1)<(t,1)$ を与える；したがってそれらは分離され，$\neg X\wedge Y$ は
対称である．添字を逆にすれば，二つの減少する型の場合が扱える．

したがって二つの型は分離されるか，または同じ向きをもち交互に並ぶ．
交互に並ぶことは推移的である：$t,u$ と $u,v$ が交互に並び（例えばすべて増加とする），
しかし $t,v$ が分離されるならば，$(v,1)<(u,2)<(t,3)<(v,1)$ となる（$t$ が $v$ の左にあるとき），
そうでなければ $(t,1)<(u,2)<(v,3)<(t,1)$ となる．したがって交互に並ぶことは同値
関係である．類の外の型は類のすべての型から分離され，
それらすべての同じ側にある．そうでなければ交互に並ぶ二つの型が
分離されることになるからである；したがって類の領域は他の点を含まない．類の内部では，
交互に並ぶことがかたまりを与え，(i) がかたまりの内部の順序を与える．
\end{proof}

\begin{lemma}[有限群における同時区間短縮]
\label{lem:group-interval}
${\mathbb G}$ を有限群とし，幅 $f$ の，
${\mathbb G}$ 上のブロック置換族で，$n\ge\max\{3,|{\mathbb G}|^f\}$ 個の添字をもつものを考える．ここで，すべての添字は同じ
個数のブロックを，同じ行の上に，同じ添字内の順序でもち，かつ，
共通の行の上の二つのブロックの位置 $t,u$ のすべてについて，
「ブロック $t$（添字 $i$ の）がブロック $u$（添字 $j$ の）の左にある」の真理値がすべての
$i<j$ について同じであるとする．このとき，添字の空でない区間 $J$ で
$\operatorname{val}_\varrho(J)=\operatorname{val}_\varrho(\emptyset)$ をすべての行 $\varrho$ について満たすものがある．
\end{lemma}

\begin{proof}
補題~\ref{lem:group-block-order} を各行に適用する；類の
領域は合計で $c\le f$ 個ある．領域 $\mathcal Q$ を固定し，そのかたまりに $1,\ldots,n$ と
左から右へ番号を付ける．かたまりの間の定数をその左のかたまりに割り当て，
領域の両端の固定された定数をその最初または最後の
かたまりに吸収する．各かたまりは，\emph{同じ}
固定された定数を同じ順序でもつ元の積と置換の積をもつ．$P_p$ をかたまり $1,\ldots,p$ の元の値とし，$Q_p$ を
それらがすべて置き換えられたときの値とし（$P_0=Q_0=1_{\mathbb G}$），
$\eta^{\mathcal Q}_p:=P_pQ_p^{-1}$ とおく．ちょうどかたまり $p+1,\ldots,p'$ を置き換えると，
順序付きの値 $P_p\,(Q_p^{-1}Q_{p'})\,(P_{p'}^{-1}P_n)$ が得られ，これが $P_n$ に等しいのは，
$P_pQ_p^{-1}\cdot Q_{p'}P_{p'}^{-1}=1_{\mathbb G}$ のとき，すなわち
$\eta^{\mathcal Q}_p=\eta^{\mathcal Q}_{p'}$ のとき，かつそのときに限る；どの因子も交換しない．
$\eta^{\mathcal Q}_p=P_pQ_p^{-1}$ における順序は，置き換えたかたまりが
元の接頭辞と元の接尾辞の間にある位置によって強制される；
元 $Q_p^{-1}P_p$ は非可換群ではこれを検出する必要はない．
$\mathcal Q$ が増加ならば，$(a,b]$ の所有者はちょうどかたまり $(a,b]$ を占める；それが
減少ならば，かたまり $(n-b,n-a]$ を占める．$j=0,\ldots,n$ について，状態を
${\mathbb G}^c$ のベクトルで，その $\mathcal Q$-成分が増加の領域については $\eta^{\mathcal Q}_j$，
減少の領域については $\eta^{\mathcal Q}_{n-j}$ であるものとする．$n+1>|{\mathbb G}|^c$ であるので，二つの状態が
$a<b$ で一致する．このとき所有者 $a+1,\ldots,b$ を置き換えると，すべての領域の
値が保たれ，領域の外では何も変わらない．
\end{proof}

\begin{example}[逆の向きと非可換な値]
\label{ex:nonabelian-reverse-owner}
$S_3$ において $\sigma_1=(12)$，$\sigma_2=(23)$ とする．三つの所有者が，一つのブロック $B_i$ を
行~1 上に，一つのブロック $W_i$ を行~2 上にもつ（したがって族の幅は2）：
\[
 \text{row 1: }B_1\,\sigma_1\,B_2\,\sigma_1\,B_3,\qquad
 \text{row 2: }W_3\,\sigma_2\,W_2\,\sigma_2\,W_1,
\]
したがって行~1 は一つの増加の類であり，行~2 は一つの減少の類である．元は
$B_i=W_i=1$ であり，両方の行は値 $1$ をもつ．所有者~$1$ には
自明な置換 $B'_1=B_1$，$W'_1=W_1$ を与え，ちょうど所有者
$J=\{2,3\}$ を $B'_2=\sigma_2$，$W'_2=\sigma_2\sigma_1\sigma_2$，$B'_3=\sigma_1\sigma_2\sigma_1$，
$W'_3=\sigma_1$ で置き換える．すると行~1 は
$\sigma_1\cdot\sigma_2\sigma_1\cdot\sigma_1\sigma_2\sigma_1=1$ となり，行~2 は
$\sigma_1\sigma_2\cdot\sigma_2\sigma_1\sigma_2\cdot\sigma_2=1$ となる．これらは順に計算したものである．証明の
記法では，$(a,b]=(1,3]$ であり，行~1 上で $\eta_1=\eta_3=1$ であり，一方
行~2 上で $\eta_{3-b}=\eta_0=1=\eta_2=\eta_{3-a}$ である：一つの所有者の区間が
二つの異なる境界添字の向きに適合する．ただし，置き換えたどちらの所有者も
自身の寄与を保たない．この例は仕組みを示すためのものであり，所有者が三つなので
命題~\ref{prop:block-replacement-interface} の閾値には遠く及ばない；良い所有者の区間は直接の計算で見つけている．
\end{example}

\begin{lemma}[二つの単調列]
\label{lem:two-sequence}
$m\ge3$，$m'\ge(m-1)^2+1$ とし，$\mathrm{lt}:\{(i,j):1\le i<j\le m'\}\to\{0,1\}$
を，各引数について固定した向きで単調であるとする：各固定した $i$ について
$j$ について単調であり，各固定した $j$ について $i$ について単調であり，これら二つの
向きのそれぞれは固定した引数のすべての値について同じである．このとき $\mathrm{lt}$ は，
対 $i<j$（ある $I\subseteq[m']$ で $|I|=m$ なるものの）の上で一定である．
\end{lemma}

\begin{proof}
必要ならば $\mathrm{lt}$ を $1-\mathrm{lt}$ で置き換えて，行は $j$ について非減少であるとする．
大きさが少なくとも $m$ の添字集合を見つければ十分である．一定性は
部分集合に引き継がれるからである．

列が $i$ について非増加ならば，$\mathrm{nx}(i):=\min\{j>i:\mathrm{lt}(i,j)=1\}$ とする（
なければ $\infty$）；$\mathrm{nx}$ は非減少である．鎖 $i_1=1$，$i_{r+1}=\mathrm{nx}(i_r)$ は
$\mathrm{lt}(i_a,i_b)=1$ を $a<b$ について満たす；それが $m$ 項をもてば終わりである．そうでなければ，その高々
$m-1$ 個の項は $[m']$ を高々 $m-1$ 個の区間 $[i_r,i_{r+1})$，$[i_L,m']$ に切り分け，その一つ
は少なくとも $m'/(m-1)>m-1$ 個，したがって少なくとも $m$ 個の添字をもち，その内部では $j<\mathrm{nx}(i_r)\le \mathrm{nx}(i)$ が
$\mathrm{lt}(i,j)=0$ を与える．

列が非減少ならば，$i^\ast$ を
$\mathrm{lt}(i^\ast,i^\ast+1)=1$ なる最小の添字とする（なければ $i^\ast=m'$）．このとき $\mathrm{lt}(i,j)=1$ が $i^\ast\le i<j$ について成り立ち，
$\mathrm{lt}(i,j)\le \mathrm{lt}(j-1,j)=0$ が $i<j\le i^\ast$ について成り立ち，$[1,i^\ast]$，
$[i^\ast,m']$ の一方は少なくとも $(m'+1)/2\ge m$ 個の添字をもつ．なぜなら $(m-1)^2+1\ge2m-1$ が
$m\ge3$ について成り立つからである．
\end{proof}

\begin{lemma}[多項式的な順序の均質化]
\label{lem:order-homogenize}
$n_0\ge3$，$m_0:=n_0$ および $m_{q+1}:=(m_q-1)^2+1$ とし，したがって $m_q\le n_0^{2^q}$ である．$\mathcal T$ を
型の有限集合で，各型が一つの行に割り当てられたものとし，$i\in[n]$ と $t\in \mathcal T$ について
$p_t(i)$ を $t$ の行の上の点とし，各行の上の点は相異なるとする．
$\mathcal O$ を，共通の行の上の型の順序対 $(t,u)$（$t=u$ を許す）の集合とする，
$n\ge m_{|\mathcal O|}$ ならば，$I\subseteq[n]$ で $|I|=n_0$ なるものが存在して，
すべての $(t,u)\in\mathcal O$ について，$p_t(i)<p_u(j)$ の真理値はすべての $i<j$（
$I$ における）について同じである．
\end{lemma}

\begin{proof}
対を一つずつ扱い，対 $(t,t)$ を先に扱う；各段階は，
大きさが少なくとも $(m-1)^2+1$ の添字集合から大きさが少なくとも $m$ の部分集合に移り，それ以前の
均質性は部分集合に引き継がれる．$(t,t)$ については，Erd\H{o}s--Szekeres が長さ $m$ の単調な
部分列を与える．$t\ne u$（両方ともすでに単調）については，現在の
添字に $1,\ldots,m'$ と番号を付け直し，$\mathrm{lt}(i,j):=[p_t(i)<p_u(j)]$ を $i<j$ についておく．固定した $i$ について，
$\mathrm{lt}$ は $j$ について $p_u$ の向きに単調である；固定した $j$ について，$i$ について
$p_t$ と逆の向きに単調である．$m\ge m_0=n_0\ge3$ であるので，
補題~\ref{lem:two-sequence} が適用できる．
\end{proof}

\begin{proof}[補題~\ref{lem:block-replacement}\textup{(c)} の証明]
ブロックを所有しない添字は単独で選べるので，すべての添字がブロックを一つ所有するとする；
$f=1$ ではすべてのブロックは一つの行の上にある．添字を
そのブロックの位置で番号付け直し，かたまりごとに一つのブロックをもつ一つの増加の領域をなすようにする．
補題~\ref{lem:group-interval} の証明と同様に，ちょうど添字
$a+1,\ldots,b$ を置き換えることが行の値を保つのは $\eta_a=\eta_b$ のとき，かつそのときに限り，$m+1>|{\mathbb G}|$
個の値 $\eta_0,\ldots,\eta_m\in{\mathbb G}$ はすべて相異なることはあり得ない．
\end{proof}

\begin{proof}[補題~\ref{lem:block-replacement}\textup{(b)} の証明]
ブロックを所有しない添字は単独で選べるので，すべての添字がブロックを所有するとする．
添字の\emph{型}は，そのブロックの個数，各ブロックの行，
および共有する行の上でのそれらの順序を記録する；型は高々 $f\cdot f^f\cdot f!=\mathcal P_f$
通りであるので，少なくとも $N_{{\mathbb G},f}^{2^{f^2}}$ 個の添字が一つの型を共有する．
他のすべての添字のブロックを定数に吸収する：これらの添字は決して置き換えられないので，その
ブロックはすべての行で元の値を保つ．共通の型のブロックは
高々 $f$ 個の型と，共通の行の上の高々 $f^2$ 個の順序対を与えるので，
$n_0=N_{{\mathbb G},f}$ とした補題~\ref{lem:order-homogenize} は，$N_{{\mathbb G},f}$
個の添字で，順序の比較が均質なものを与える．この抽出で捨てた添字のブロックも
同じように定数に吸収する．共通の
型は一つの添字のブロックの順序を固定し，これは
補題~\ref{lem:group-block-order} の仮定~(i) であり，補題~\ref{lem:order-homogenize} は
仮定~(ii) を与える；$i>j$ についての比較は，
逆の対についての比較の否定である．点は相異なるからである．残りの
$N_{{\mathbb G},f}$ 個の添字は，このとき
補題~\ref{lem:group-interval} の仮定を満たし，その区間が求める~$J$ である．
\end{proof}

固定した $f$ について，$M_{{\mathbb G},f}=\mathcal P_f\max\{3,|{\mathbb G}|^f\}^{2^{f^2}}$ は $|{\mathbb G}|$ について多項式である；その
次数 $f\cdot2^{f^2}$ は $f^2$ について指数的であるので，$M_{{\mathbb G},f}$ そのものは $f$ について二重
指数的である．

\section{局所群の座標と MCFG のブロック短縮}
\label{app:local-group}

$\mathbf{LG}$ は常に正の像 $\mathsf S=\widehat s(\Sigma^+)$ の上で調べる．
古典的な恒等式 $\mathbf{LG}=\mathbf{Grp}*\mathbf D$ はそれを
滑動窓と遅延の理論の中に位置づけ \cite{Straubing1985,Tilson1987}，
$U_1=\{1,\zeta\}$（$\zeta^2=\zeta$）の排除は \cite[Proposition~2.3]{AlmeidaEtAl2009} である．
$x^\omega$ で有限半群の元 $x$ のべき等な冪を表す．
同値な言い方をすれば，$\mathbf{LG}$ は一つの擬恒等式
$(x^\omega yx^\omega)^\omega=x^\omega$ によって定義される：べき等元 $e=x^\omega$ について，有限モノイド
$e\mathsf Se=\{eye:y\in\mathsf S\}$ が群であるのはちょうど $e$ がその唯一のべき等元であるときであり，
すべてのべき等元は $x^\omega$ の形をもつ．
以下の座標はこの代数の自己完結した実装であって，
新しい特徴づけではない；MCFG の学習に固有なのは，群の
値を固定された区切りと置換可能なブロックとに分解することである
（補題~\ref{lem:local-group-blocks}）．全体を通じて，$k\ge1$，${\mathbb G}$ は有限群，
$\lambda_{k,\nu}$ は定義~\ref{def:local-group-typing} のとおりとし，語が
\emph{短い}とは $|w|<2k$ であること（$\widehat s_{k,k}$ によって正確に記録される）をいい，そうでないとき\emph{長い}
という．また $x,y\in\Sigma^k$ について
$\Delta_\nu(x,y):=\prod_{i=1}^{k}\nu\bigl((xy)[i..i+k]\bigr)$ とおく．これは，
$k$ 個の因子（長さ $k+1$ の，$xy$ の因子で $x$ の中で始まるもの）にわたる順序付き積である．

\emph{ポンピング．}有限半群 $\mathsf T$ について，少なくとも
$|\mathsf T|$ 個の元のすべての積は $\mathsf T^1E(\mathsf T)\mathsf T^1$ に属する
\cite[Lemma~2.2]{AlmeidaEtAl2009}，\cite{Pin1986}：その最初の
$|\mathsf T|$ 個の接頭辞積のうち二つが一致する，$t_i=t_j$（$i<j$）であって，そのとき
$t_i=t_ix^\omega$ が積 $x$（因子 $i+1,\ldots,j$ の）について成り立つか；またはこれらの接頭辞
積が $\mathsf T$ を尽くし，そのうちの一つがべき等であるかのいずれかである．以下の三つの使用はすべて
（補題~\ref{lem:lg-rees-window-coordinate} と~\ref{lem:mcfg-li-group-window}，
定理~\ref{thm:lg-characterization}）この事実に基づく．

\begin{lemma}[滑動因子写像は準同型である]
\label{lem:local-group-morphism}
$\lambda_{k,\nu}(uv)$ は $\lambda_{k,\nu}(u)$ と $\lambda_{k,\nu}(v)$ の関数である；
したがって $\lambda_{k,\nu}$ は有限モノイド $M_{k,\nu}$ で，高々
$|\Sigma|^{<2k}+|\Sigma|^{2k}|{\mathbb G}|$ 個の元をもつものの上への準同型である．長い $u,v$ について，
\[
 \lambda_{k,\nu}(uv)
 =\bigl(\langle\operatorname{pref}_ku,\operatorname{suf}_kv\rangle,\;
   \gamma_{k,\nu}(u)\,\Delta_\nu(\operatorname{suf}_ku,\operatorname{pref}_kv)\,
   \gamma_{k,\nu}(v)\bigr).
\]
\end{lemma}

\begin{proof}
窓の成分は準同型である．長い $u,v$ について，長さ $k+1$ の
因子（$uv$ の）は，順に，$u$ のもの，$k$ 個の因子で，
最後の $k$ 個の位置（$u$ の）で始まるもの（$\operatorname{suf}_k(u)\operatorname{pref}_k(v)$ の内部にある），および
$v$ のものである．$u$ が短く $v$ が長いならば，$u$ の中で始まる因子は
$u\operatorname{pref}_k(v)$ の内部にあり，他のものは $v$ のものである；$u$ が長く
$v$ が短い場合は対称である；そして両方が短いならば，$uv$ は正確に定まる．
\end{proof}

\begin{lemma}[ブロック分解]
\label{lem:local-group-blocks}
$w=x_0\mathfrak b_1x_1\cdots\mathfrak b_qx_q$ と $w'=x_0\mathfrak b'_1x_1\cdots\mathfrak b'_qx_q$ が
$\widehat s_{k,k}(\mathfrak b'_j)=\widehat s_{k,k}(\mathfrak b_j)$ をすべての $j$ について満たすとする．このとき
$\widehat s_{k,k}(w')=\widehat s_{k,k}(w)$ であり，$w$ が長いならば，
$z_0,\ldots,z_q\in {\mathbb G}$ で，$x_j$ と
$\mathfrak b_j$ の窓の型のみに依存し，次を満たすものが存在する．
\[
 \gamma_{k,\nu}(w)=z_0\,\gamma_{k,\nu}(\mathfrak b_1)\,z_1\cdots\gamma_{k,\nu}(\mathfrak b_q)\,z_q,
 \qquad
 \gamma_{k,\nu}(w')=z_0\,\gamma_{k,\nu}(\mathfrak b'_1)\,z_1\cdots\gamma_{k,\nu}(\mathfrak b'_q)\,z_q .
\]
\end{lemma}

\begin{proof}
窓についての主張は $\widehat s_{k,k}$ が準同型であることから成り立つ．短いブロックは
その置換と等しく，固定された語に吸収されるので，すべてのブロックは長いとする．
長さ $k+1$ の因子は，あるブロックの内部にあるか，または，それが触れるすべてのブロックと
長さ高々 $k$ の接頭辞または接尾辞で交わる；二番目の種類の因子は，
固定された語とブロックの最初と最後の $k$ 文字によって綴られるので，それらは
$w$ と $w'$ で同じであり，同じ順序で現れる．開始位置の順に並べると，
$\mathfrak b_j$ の内部の因子は一つの連なりをなし，その積は $\gamma_{k,\nu}(\mathfrak b_j)$ である；
$z_j$ を，連続する連なりの間の因子の順序付き積とする．
\end{proof}

\begin{lemma}[滑動因子の座標]
\label{lem:lg-rees-window-coordinate}
$s$ を型付けで，その正の像 $\mathsf S$ のすべてのべき等元がその極小
イデアル $\mathsf K=\mathcal M({\mathbb H};I,\Lambda;\mathsf P)$ に属するものとする．これは群
${\mathbb H}$ 上の Rees 行列半群で，サンドイッチ行列 $\mathsf P=(p_{ji})_{j\in\Lambda,i\in I}$ をもち，その元
$(i,\gamma,j)\in I\times{\mathbb H}\times\Lambda$ は
$(i,\gamma,j)(i',\gamma',j')=(i,\gamma p_{ji'}\gamma',j')$ によって掛け合わされる．$k\ge1$ が
$\mathsf S^k\subseteq\mathsf K$ を満たすとする（ポンピングにより，
$\mathsf S^{|\mathsf S|}\subseteq\mathsf S^1E(\mathsf S)\mathsf S^1\subseteq\mathsf K$ であるので，$k=|\mathsf S|$ でよい；
そのような最小の $k$ は最小の窓を与える）．
$\nu:\Sigma^{k+1}\to {\mathbb H}$ で，すべての $w$ で $|w|\ge2k$ なるものが次を満たすものが存在する．
\[
 s(w)=\bigl(\operatorname{row}(s(\operatorname{pref}_kw)),\;
 \operatorname{grp}(s(\operatorname{pref}_kw))\,\gamma_{k,\nu}(w),\;
 \operatorname{col}(s(\operatorname{suf}_kw))\bigr),
 \tag{\ensuremath{\star}}
\]
ここで $\operatorname{row},\operatorname{grp},\operatorname{col}$ は
$\mathsf K$ の元の三つの座標である．特に $s(w)$ は，その最初と最後の
$k$ 文字と，${\mathbb H}$ における順序付き積とによって定まる．
\end{lemma}

\begin{proof}
\emph{(1) 座標．}$i_0\in I$ を固定する；$j\in\Lambda$ について，元
$r_j:=(i_0,p_{ji_0}^{-1},j)$ は，
$\mathsf K$ の元で $\Lambda$-座標 $j$ をもつものに対する右単位元である．$y\in \mathsf S$ について $r_jy=(i',\gamma',j')$ と書き，
$\operatorname{act}(j,y):=p_{ji'}\gamma'$ とおく；このとき
$(i,\gamma,j)y=(i,\gamma\operatorname{act}(j,y),j')$ であるので，右からの乗法は
$I$-座標を固定する．双対的に，左からの乗法は $\Lambda$-座標（$\mathsf K$ の元の）を固定する：
任意の $j''\in\Lambda$ について，べき等元
$e':=(i,p_{j''i}^{-1},j'')\in\mathsf K$ は $e'(i,\gamma,j)=(i,\gamma,j)$ を満たすので，
$y(i,\gamma,j)=(ye')(i,\gamma,j)$ であり $ye'\in\mathsf K$ である．そして $\mathsf K$ の二つの
元の Rees 積は，その右の因子の $\Lambda$-座標をもつ．

\emph{(2) 群の座標．}$w=a_1\cdots a_m$ を文字 $a_t$ からなるものとし，$m\ge2k$，
$s(a_1\cdots a_k)=(i,\gamma,j_k)$ とする．$t\ge k$ について，$\Lambda$-座標 $j_t$
（$s(a_1\cdots a_t)$ の）は $s(a_{t-k+1}\cdots a_t)$ のそれに等しい．これは $\mathsf S^k\subseteq\mathsf K$
と (1) による．
$\nu(F):=\operatorname{act}\bigl(\operatorname{col}(s(F[1..k])),s(F[k+1])\bigr)$ とおく．
右作用の恒等式を $t=k,\ldots,m-1$ について順に適用すると
$s(w)=(i,\,\gamma\prod_{t=k}^{m-1}\nu(a_{t-k+1}\cdots a_{t+1}),\,j_m)$ が得られ，その
積は $\gamma_{k,\nu}(w)$ である．
\end{proof}

\begin{theorem}[滑動因子の座標による局所群の判定条件]
\label{thm:lg-characterization}
型付け $s$ で正の像 $\mathsf S$ をもつものについて，次は同値である：
\textup{(i)} $\mathsf S\in\mathbf{LG}$；\textup{(ii)} $\mathsf S$ は
$U_1$ と同型な部分半群をもたない；\textup{(iii)} $\mathsf S$ のすべてのべき等元は $\mathsf S$ の極小イデアルに属する；
\textup{(iv)} $s$ は局所群決定的であり，$k=|\mathsf S|$，または任意の $k$ で $\mathsf S^k=\mathsf K$ なるものをとってよい．
\end{theorem}

\begin{proof}
\cite[Proposition~2.3, items (1), (3) and (4)]{AlmeidaEtAl2009} により，(i) と (ii)
はそれぞれ，$\mathsf S^m=\mathsf K$ がある $m$ について成り立つことと同値である．$\mathsf S^m=\mathsf K$ ならば，
$e=e^m\in\mathsf K$ が各べき等元 $e$ について成り立つ；逆に，$E(\mathsf S)\subseteq\mathsf K$ ならば，
ポンピングにより $\mathsf S^k\subseteq\mathsf S^1E(\mathsf S)\mathsf S^1\subseteq\mathsf K$ が $k=|\mathsf S|$ について成り立ち，
$\mathsf S^k\supseteq\mathsf K$ である，$\mathsf S^k$ はイデアルであるので．したがって (i)--(iii) は同値である．

(iii)$\Rightarrow$(iv)：補題~\ref{lem:lg-rees-window-coordinate} は，
$k=|\mathsf S|$（または任意の $k$ で $\mathsf S^k=\mathsf K$ なるもの）をとって，$s(w)$ を $|w|\ge2k$ について
$\lambda_{k,\nu}(w)$ から定め，$\widehat s_{k,k}$ はより短い語と $\eps$ を正確に記録する；
したがって $\widehat s\preceq\lambda_{k,\nu}$．

(iv)$\Rightarrow$(i)：$\mathsf S$ は $\mathsf T:=\lambda_{k,\nu}(\Sigma^+)$ の準同型像であり，
$\mathbf{LG}$ は準同型像について閉じているので，$\mathsf T\in\mathbf{LG}$ であることを示せば十分である．
短い元はべき等でない（二乗すると，正確に記録された語が変わるか，それが
長くなる）．$e=(\langle p,q\rangle,\gamma_e)$ をべき等とする．
補題~\ref{lem:local-group-morphism} により，$e\mathsf Te$ のすべての元は窓
$\langle p,q\rangle$ をもち，この窓をもつ長い元は
$(\langle p,q\rangle,\gamma)(\langle p,q\rangle,\gamma')=(\langle p,q\rangle,\gamma\Delta\gamma')$，
$\Delta:=\Delta_\nu(q,p)$ によって掛け合わされる．この演算は $\{\langle p,q\rangle\}\times {\mathbb G}$ を，
群で ${\mathbb G}$ と同型なもの（$\gamma\mapsto\gamma\Delta$ を通じて）にし，$e\mathsf Te$ はその有限で空でない
部分半群であるので，部分群である．
\end{proof}

\begin{lemma}[局所自明--群型付けは窓--群細分を許す]
\label{lem:mcfg-li-group-window}
$\Sigma\ne\emptyset$ とする．$\mathsf S\in\mathbf{LI}\vee\mathbf{Grp}$ ならば，
$k\ge1$，有限群 ${\mathbb G}$ および準同型 $g:\Sigma^*\to {\mathbb G}$ で
$\widehat s\preceq(\widehat s_{k,k},g)$ なるものが存在する；$\mathsf S\in\mathbf{LI}$ ならば，${\mathbb G}$ は
自明にとれる．
\end{lemma}

\begin{proof}
有限の $\mathsf S'\in\mathbf{LI}$，有限群 ${\mathbb G}$，部分半群
$\mathsf T\le \mathsf S'\times {\mathbb G}$ および全射準同型 $\vartheta:\mathsf T\to\mathsf S$ を選ぶ（
$\mathsf S\in\mathbf{LI}$ については $\mathsf S'=\mathsf S$，${\mathbb G}=\{1\}$）．すべての文字について，
$\mathsf T$ における逆像（その $s$-像の）を一つ選ぶことにより，準同型 $s':\Sigma^+\to\mathsf S'$ と $g$ で
$\vartheta(s'(w),g(w))=s(w)$ なるものが定まる；$g(\eps)=1_{\mathbb G}$ とおく．部分半群
$\mathsf S_0:=s'(\Sigma^+)\le \mathsf S'$ は局所自明である；$k:=|\mathsf S_0|$ とする．
べき等元 $e,e'$（$\mathsf S_0$ の）と $y\in \mathsf S_0^1$ について，
$eye'=ey(e'ee')=(e(ye')e)e'=ee'$ が成り立つ．
$|w|\ge2k$ について $w=pxz$ と書く．ただし $|p|=|z|=k$ である．ポンピングにより，
$s'(p)=y_1ey_2$ かつ $s'(z)=y_3e'y_4$ であって，$y_1,\ldots,y_4\in \mathsf S_0^1$ であり，$e,e'$ は
べき等である．したがって $s'(w)=y_1e\,(y_2\,s'(x)\,y_3)\,e'y_4=y_1ee'y_4$（$s'(\eps):=1$ とする）であり，これは
$p=\operatorname{pref}_k(w)$ と $z=\operatorname{suf}_k(w)$ のみに依存する．
より短い語は $\widehat s_{k,k}$ によって正確に記録され，これは
$\eps$ も分離する．したがって等しい $(\widehat s_{k,k},g)$-値は等しい $(s',g)$-値と
等しい $s$-値を導く．
\end{proof}

\begin{proposition}[窓--群型付けは局所群型付けであり，その包含は真である]
\label{prop:li-g-strict}
\begin{enumerate}[label=\textup{(\roman*)},leftmargin=*]
\item $\widehat s\preceq(\widehat s_{k,\ell},g)$ が有限
群への準同型 $g$ について成り立つならば，$\widehat s\preceq\lambda_{\hat k,\nu}$ が $\hat k=\max\{k,\ell,1\}$ と
$\nu(F):=g(\text{last letter of }F)$ について成り立つ．さらに $\mathbf{LI}\vee\mathbf{Grp}\subseteq\mathbf{LG}$．
\item すでに $\Sigma=\{\mathtt a,\mathtt b\}$ と，
型付け $s_{\mathtt{ab}}:=\lambda_{1,\nu}|_{\Sigma^+}$ で
$\nu(\mathtt{ab})=1\in\mathbb Z_2$ かつそれ以外では $\nu=0$ なるものについて，
$s_{\mathtt{ab}}(\Sigma^+)\in\mathbf{LG}\setminus(\mathbf{LI}\vee\mathbf{Grp})$．
\end{enumerate}
\end{proposition}

\begin{proof}
(i) $|w|\ge2\hat k$ について $g(w)=g(\operatorname{pref}_{\hat k}w)\,\gamma_{\hat k,\nu}(w)$ であり，
$\operatorname{pref}_{\hat k},\operatorname{suf}_{\hat k}$ は
$\widehat s_{k,\ell}(w)$ を定める；より短い語は正確に記録される．包含が成り立つのは，
$\mathbf{LG}$ が $\mathbf{LI}$（$e\mathsf Te=\{e\}$）と
$\mathbf{Grp}$ を含む擬多様体であるからである．
(ii) $\mathbf{LG}$ への所属は
定理~\ref{thm:lg-characterization} から従う．$k,\ell\ge0$ とし，
$g':\Sigma^*\to {\mathbb G}'$ を任意の有限群への写像とする．
$m\ge k+\ell+1$ と $n\ge1$ で $g'(\mathtt a)$ の位数で割り切れるものを選び，
$u=\mathtt b^m\mathtt a^{n}\mathtt b^m$ と
$v=\mathtt b^{2m}$ とおく．それらは等しい長さ $k$ の接頭辞と
長さ $\ell$ の接尾辞をもち，$s_{k,\ell}$ について長く，
$g'(u)=g'(\mathtt b)^{2m}=g'(v)$ を満たすが，
$\#_{\mathtt{ab}}(u)=1$ かつ $\#_{\mathtt{ab}}(v)=0$ である．
したがって $s_{\mathtt{ab}}$ は窓--群型付けではなく，
補題~\ref{lem:mcfg-li-group-window} は
$\mathbf{LI}\vee\mathbf{Grp}$ を排除する．
\end{proof}

\begin{lemma}[間引きと詰め物をした所有者]
\label{lem:lg-buffered-owners}
$k\ge1$ を固定する．高々 $f$ 個の
語（\emph{行}）の中の，対ごとに互いに素で空でない区間を考える．各区間は $m$ 個の所有者のいずれか一つに属し，各
所有者は高々 $f$ 個の区間をもつとし，どの区間もその行の
最初または最後の $k$ 文字と交わらないとする．このとき，集合 $Q_1$ で少なくとも
$m/(4kf+1)$ 個の所有者からなり，次の性質をもつものが存在する．$Q_1$ の所有者の区間の任意の空でない
置換語について，対ごとに
互いに素な\emph{広げたブロック}で，所有者ごとに高々 $f$ 個のものが存在して，
\textup{(a)} 各広げたブロックは，
$Q_1$ のちょうど一つの所有者の区間を含み，$Q_1$ の所有者のすべての区間はちょうど一つの広げた
ブロックに含まれる；\textup{(b)} 広げたブロックの内部でその
所有者の区間を置き換えると，長さ少なくとも $2k+1$ の語で，ブロックと同じ最初と最後の
$k$ 文字をもつものが得られ，ブロック自身も長さ少なくとも $2k+1$ をもつ；かつ
\textup{(c)} すべての $J\subseteq Q_1$ について，$J$ の所有者の区間を置き換えることは，
それらの広げたブロックを \textup{(b)} のように置き換えることに等しい．
\end{lemma}

\begin{proof}
二つの相異なる所有者を，一方のある区間と
他方のある区間が同じ行にあって，その間の文字が $2k$ 個未満であるとき結ぶ．
区間 $[l,r)$（切れ目の位置で）の右の隣の区間は，
$2k$ 個の切れ目 $r,\ldots,r+2k-1$ のいずれかで始まり，互いに素であることと
空でないことから，切れ目ごとに高々一つの区間しかない；左についても同様である．
したがって各所有者は高々 $4kf$ 個の隣をもち，貪欲に所有者を選んで
その閉近傍を削除すると，独立集合 $Q_1$ で，少なくとも
$m/(4kf+1)$ 個の所有者からなるものが得られる．

$Q_1$ の所有者のすべての区間を両側で $k$ 文字だけ広げる；
境界の仮定により，広げたものはその行の内部にとどまる．
$Q_1$ の相異なる所有者の区間は少なくとも $2k$ 文字離れているので，それらを広げたものは
互いに素な内部をもつ；同じ所有者の重なり合う広げたものを併合する．
これは (a) を与え，所有者ごとのブロックの個数を増やさない．
$[l,r)$ を併合した広げたものの和とし，$l_0$ をその中の区間の左端の最小値とする．
したがって $l=l_0-k$．$Q_1$ の所有者のどの区間も
$[l,l_0)$ と交わらない：同じ所有者のものであれば，より小さな左端をもって併合されていたはずであり，
$Q_1$ の別の所有者のものであれば，少なくとも $2k$ 文字だけ $l_0$ より前で終わるはずである．
最後の $k$ 文字についても対称である．したがって所有者の
区間だけを置き換えることは最初と最後の $k$ 文字を保ち，置換は
これら $2k$ 文字の他に少なくとも一つの空でない語を含み，これが
(b) を与える．ブロックの内部にある $Q_1$ の外の所有者の文字は決して置き換えられないので，
(c) がすべての $J\subseteq Q_1$ について成り立つ．
\end{proof}

\begin{lemma}[所定の局所群型をもつ短いタプル]
\label{lem:local-group-shortest}
$G$ をファンアウト $\le f$ かつランク $\le r_{\max}$ の既約な良い文法とし，
$\lambda=\lambda_{k,\nu}$ を群 ${\mathbb G}$ をもつものとする．
$C^{\mathrm{buf}}_{k,{\mathbb G},f}:=2kf+(4kf+1)c_{{\mathbb G},f}$ と
$\bar n_\lambda:=|V||M_{k,\nu}|^f$ とおく．すべての $B\in V$ と，すべての実現可能な
$\mathbf t\in M_{k,\nu}^{\fan(B)}$ について，$\tuple u\in L_B(G)$ で
$\lambda^{(\fan(B))}(\tuple u)=\mathbf t$ かつ次を満たすものが存在する．
\[
 \|\tuple u\|_1\le\beta_\lambda:=(r_{\max}+1)\,\bar\theta_G\,
 C^{\mathrm{buf}}_{k,{\mathbb G},f}\,\bar n_\lambda .
\]
\end{lemma}

\begin{proof}
補題~\ref{lem:heavy-node-count}(ii) の仮定を，
$c=(4kf+1)c_{{\mathbb G},f}$ と $c'=2kf$ について確かめる．したがって $c+c'=C^{\mathrm{buf}}_{k,{\mathbb G},f}\ge1$；
$n_\lambda(G)\le\bar n_\lambda$ とともに，それは評価を与える．$T$ を $\lambda$-極小で，重さの
閾値が $b=\bar\theta_G$ のものとする．長さが少なくとも $2k$ の各根成分において，最初
と最後の $k$ 文字に標識を付け，より短い各成分においてはすべての文字に標識を付ける：
標識は高々 $2kf$ 個である．$Q$ を重い節点の反鎖とする．線形性と
非削除性により，それらの成分は根
成分の互いに素で空でない区間を，節点ごとに高々 $f$ 個占め，高々 $2kf$ 個の節点が標識を含む；
$Q_0$ を残りのものとする．したがってその区間は，
すべての根成分の最初と最後の $k$ 文字を避け，長さが少なくとも $2k$ の成分の中にある．

$|Q_0|>(4kf+1)c_{{\mathbb G},f}$ とする．補題~\ref{lem:lg-buffered-owners} は集合
$Q_1\subseteq Q_0$ で $c_{{\mathbb G},f}$ 個より多い所有者からなるものを与える．$v\in Q_1$（非終端記号
$C_v$ をもつ）について，$L_{C_v}(G)$ のタプルで総長高々 $\bar\theta_G$ のものを固定する
（補題~\ref{lem:thickness-equivalence}(ii)）；その成分は空でなく，それには
窓や群の条件を課さない．
補題~\ref{lem:lg-buffered-owners}(b) により，すべての広げたブロックとその置換は
長く，同じ最初と最後の $k$ 文字をもち，したがって同じ
$\widehat s_{k,k}$-型をもつ．したがって補題~\ref{lem:local-group-blocks} を
広げたブロックについて適用すると，すべての長い根成分の $\gamma_{k,\nu}$ を，
ブロックの値と，どの所有者を置き換えるかに依存しない定数との
順序付き積として書ける．これは幅 $f$ の，${\mathbb G}$ 上のブロック置換族で，
$c_{{\mathbb G},f}$ 個より多い所有者をもつものであり，
命題~\ref{prop:block-replacement-interface} は空でない $J\subseteq Q_1$ で，
すべての根の $\gamma_{k,\nu}$-値を保つものを与える．

$J$ の節点における部分木を，選んだタプルの導出で置き換える．
これは $B$ からの導出である．$J$ は反鎖だからである．そして
補題~\ref{lem:lg-buffered-owners}(c) により，その根成分は
広げたブロックの置換によるものである．短い根成分は置き換えた文字を含まず，
変わらない；長い根成分は標識付きの最初と最後の $k$ 文字を保ち，
長いままであるので，その $\widehat s_{k,k}$-型は変わらず，その
$\gamma_{k,\nu}$-値は $J$ の選び方によって保たれる．したがって根の
$\lambda$-ベクトルは変わらない．
\end{proof}

代入可能性は補題~\ref{lem:local-group-shortest} では用いない．

\section{分離する標的の証明}
\label{app:separating-targets}

この付録は命題~\ref{prop:lg-needed-target} と~\ref{prop:nonabelian-needed-target} を証明する．
言語 $L\subseteq\Sigma^*$ が\emph{狭義2-局所的}であるとは，
$\Sigma\cup\{\vdash,\dashv\}$ 上の対の集合 $P$ があって，$w\in L$ であることと
$\vdash w\dashv$ の隣り合う記号の対がすべて $P$ に入ることが同値であることをいう．

\begin{lemma}[狭義2-局所的な言語は境界で型付けされる]
\label{lem:sl2-boundary-control}
狭義2-局所的な $L$ はすべての $f$ について $(f,s_{1,1})$-タプル代入可能である；
特に $L\in\Cmcf{f}{s_{1,1}}$ であり，$\widehat s_{1,1}\preceq\widehat s$ なら $L\in\Cmcf{f}{s}$．
\end{lemma}

\begin{proof}
$\tuple x,\tuple y\in(\Sigma^+)^d$ が等しい $s_{1,1}$-型をもつとする．すなわち
対応する成分は同じ文字であるか，最初と最後の文字を共有する．$E_0[\tuple x],E_0[\tuple y]\in L$ とし，
$E[\tuple x]\in L$ と仮定する．$E$ の内部の区切りは空でないので，
$\vdash E[\tuple y]\dashv$ の隣り合う対は，ある $y_i$ の内部にあって $E_0[\tuple y]\in L$ に現れるか，
ある $y_i$ の最初または最後の文字を $\vdash E\dashv$ の記号につなぐか，
$\vdash E\dashv$ の中にあるかであり，後の二つは $\vdash E[\tuple x]\dashv$ に現れる．
よって $E[\tuple y]\in L$；対称性と $\widehat s_{1,1}\preceq\widehat s$ で証明が終わる（正規言語は MCFL である）．
\end{proof}

\begin{proof}[命題~\ref{prop:lg-needed-target} と~\ref{prop:nonabelian-needed-target} の証明]
$\diamond\in\{\mathrm A,\mathrm B\}$ と書き，$\Sigma_\diamond^\circ$ を $\{\mathtt a,\mathtt b,\mathtt c\}$
または $\{\mathtt a,\mathtt b\}$，$P_\diamond(w)$ を「$\#_{\mathtt{ab}}(w)$ は偶数」
または「$g_{\mathrm B}(w)=1$」とする．

(i) $W$ を正規言語
$\{\varsigma w:\varsigma\in\{\mathtt l,\mathtt r\},w\in(\Sigma_\diamond^\circ)^+,
\varsigma=\mathtt r\text{ または }P_\diamond(w)\}$ の右線形（ファンアウト1，ランク1，良い）文法とする；
それは最後の文字と偶奇，あるいは $S_3$ での値を追跡する．これに
\[
 C\to(\xi_1^1\mathtt\$\mathtt 1,\mathtt 3)(W)\mid(\xi_1^1\mathtt 1,\xi_1^2\mathtt 3)(C),\quad
 D\to(\mathtt 2,\mathtt 4)\mid(\xi_1^1\mathtt 2,\xi_1^2\mathtt 4)(D),\quad
 S\to(\xi_1^1\xi_2^1\xi_1^2\xi_2^2)(C,D)
\]
を加える．すると $L_C=\{(\varsigma w\mathtt\$\mathtt 1^m,\mathtt 3^m)\}$，$L_D=\{(\mathtt 2^n,\mathtt 4^n)\}$ であり，
$S$ は $L_\diamond$ を生成する；開始の型板で隣り合う変数は異なる子に属するので文法は良い．
正規言語 $\mathtt{ra\$}\mathtt 1^+\mathtt 2^+\mathtt 3^+\mathtt 4^+$ との交わりを取り
$\mathtt{ra\$}$ を消すと $C_\times$ が得られ，これは文脈自由でない；
文脈自由言語は両方の操作で閉じているので，$L_\diamond$ は文脈自由でない．

(ii) $L_\diamond=Q_\diamond\cap\Pi_\diamond\cap\varphi_\times^{-1}(0)$，ここで
$Q_\diamond=\{\mathtt l,\mathtt r\}(\Sigma_\diamond^\circ)^+\mathtt\$\mathtt 1^+\mathtt 2^+\mathtt 3^+\mathtt 4^+$，
$\Pi_\diamond$ は $\mathtt r$ で始まるか，$\mathtt l$ で始まり語全体が $P_\diamond$ を満たす空でない語の集合である．
$Q_\diamond$ の上では語全体の条件 $P_\diamond$ は $w$ のそれである：因子 $\mathtt{ab}$ は
$\varsigma$，$\mathtt\$$ や数字に触れず，$g_{\mathrm B}$ は他のすべての文字で自明である．
$Q_\diamond$ は狭義2-局所的であり（その対は $\varsigma\in\{\mathtt l,\mathtt r\}$，$x,y\in\Sigma_\diamond^\circ$ について
$\vdash\varsigma$，$\varsigma x$，$xy$，$x\mathtt\$$，$\mathtt{\$1}$，
$\mathtt{11},\mathtt{12},\mathtt{22},\mathtt{23},\mathtt{33},\mathtt{34},\mathtt{44}$，$\mathtt 4\dashv$），
$\widehat s_{1,1}\preceq\widehat s_\diamond$ なので補題~\ref{lem:sl2-boundary-control} が使える．
$\Pi_\diamond$ は $s_\diamond$-類の和である：$s_\diamond$ は最初の文字と，長さ2以上の語については
$\#_{\mathtt{ab}}$ の偶奇，あるいは $g_{\mathrm B}$ の値を記録する（一文字の語 $\mathtt l,\mathtt r$ は $\Pi_\diamond$ に入る）；
したがって定義~\ref{def:class-cmcf} の後の項目 (c) が使える．
命題~\ref{prop:abelian-coset-filter} が $\varphi_\times^{-1}(0)$ と交わりを扱う；(i) と合わせて
$f\ge2$ について $L_\diamond\in\Cmcf{f}{s_\diamond}$．$s_{\mathrm A}$ は局所群で決まるので，
定理~\ref{thm:lg-characterization} によりその正の像は $\mathbf{LG}$ に入る．

(iii)（$\mathrm A$ の場合）．$\hat k=\max\{k,\ell,1\}$，$q\ge2$ を $g'$ の群の指数の倍数である偶数，
$u=\mathtt a^{\hat k}(\mathtt{abc})^q\mathtt a^{\hat k}$，
$v=\mathtt a^{\hat k}(\mathtt{bac})^q\mathtt a^{\hat k}$ とする．どちらも $\widehat s_{k,\ell}$ について長く，
接頭辞 $\mathtt a^k$ と接尾辞 $\mathtt a^\ell$ をもち，$g'(u)=g'(\mathtt a)^{2\hat k}=g'(v)$；
したがって $s'(u)=s'(v)$．しかし $\#_{\mathtt{ab}}(u)=q$ は偶数で，$\#_{\mathtt{ab}}(v)=1$
（$\mathtt a^{\hat k}\mathtt b$ の所だけ）．$z=\mathtt{\$1234}$ として補題~\ref{lem:flag-separation} が使える．
正の像が $\mathbf{LI}\vee\mathbf{Grp}$ に入る型付けは，補題~\ref{lem:mcfg-li-group-window} により仮定を満たす．

(iii)（$\mathrm B$ の場合）．$\hat k\ge\max\{k,\ell,1\}$ を偶数とし，
$u=\mathtt a^{\hat k}(\mathtt{ab})^3\mathtt a^{\hat k}$，
$v=\mathtt a^{\hat k}(\mathtt{ab})^2\mathtt{ba}\,\mathtt a^{\hat k}$ とする．これらは同じ接頭辞 $\mathtt a^k$，
接尾辞 $\mathtt a^\ell$ と Parikh ベクトルをもつので，可換な $g'$ について $s'(u)=s'(v)$．
$\alpha^2=\beta^2=1$，$\alpha\beta$ の位数は $3$，$\beta\alpha=(\alpha\beta)^{-1}$ なので，
$g_{\mathrm B}(u)=(\alpha\beta)^3=1$，$g_{\mathrm B}(v)=(\alpha\beta)^2(\alpha\beta)^{-1}=\alpha\beta\ne1$．
$z=\mathtt{\$1234}$ として補題~\ref{lem:flag-separation} が使える．
$\mathbf{LI}\vee\mathbf{Ab}$ については，補題~\ref{lem:mcfg-li-group-window} の証明を可換群の因子で行えば可換な $g'$ が得られる．

\emph{帰結．} $s_{\mathrm A}(\Sigma_{\mathrm A}^+)$ が $\mathbf{LI}\vee\mathbf{Grp}$ に入れば，
$s'=s_{\mathrm A}$ とした (iii) は (ii) に矛盾する．
命題~\ref{thm:window-group-polydata} とその自明群・可換群の場合は，除外された種類の型付けを仮定する．
\end{proof}

\section{質的な再構成の補助的な検証}
\label{app:qualitative-construction-details}

\begin{proof}[補題~\ref{lem:filling-identity} の証明]
恒等式は代入が可換であることから成り立つ．線形性と
非削除性により，各 $\blank_\kappa$ は $E_j$ にちょうど一回現れ，
非置換性により順序 $\kappa=1,\ldots,d_j$ で現れる．なぜなら $E$ は
$\alpha_1,\ldots,\alpha_{d_0}$ を左から右へ置くからである．$\xi_j^\kappa$ と
$\xi_j^{\kappa+1}$ を考える．両方が一つの $\alpha_i$ にあるならば，非併合性により，それらの間に終端記号
または変数 $\xi_i^{\kappa'}$（$i\ne j$）があり，それは
空でない語になる．そうでなければ，それらは $\alpha_i,\alpha_{i'}$（$i<i'$）にあり，
それらの間の部分は区切り $v_i\in\Sigma^+$（$E$ の）を含む．したがって
$E_j$ のすべての内部の区切りは空でない．
\end{proof}

\begin{proof}[補題~\ref{lem:witnessed-composition} の証明]
$\tuple z_j:=\rho(\tuple u_1,\ldots,\tuple u_j,\tuple x_{j+1},\ldots,\tuple x_r)$ とおく．
$E\in\mathcal E_{d_0}$ について，補題~\ref{lem:filling-identity} は許容な $E_j$ で
$E[\tuple z_{j-1}]=E_j[\tuple x_j]$ かつ $E[\tuple z_j]=E_j[\tuple u_j]$ なるものを与える；
$\tuple u_j\equiv_{L,s}\tuple x_j$ であるので，$E$ が $\tuple z_{j-1}$ を受理するのは，それが
$\tuple z_j$ を受理するとき，かつそのときに限る．したがって $\D_L(\tuple z_0)=\D_L(\tuple z_r)$．型が一致するのは，
$\widehat s^{(d_0)}(\rho(\cdots))=\out_\rho$（子の型の）だからである．
\end{proof}

\begin{proof}[命題~\ref{prop:reconstruction-soundness} の証明]
$[\tuple x]$ からの導出の高さに関する帰納法による．Const は $\tuple x$ を導出する．
Link $[\tuple x]\to[\tuple y]$ は，同じ型の $\tuple x,\tuple y$ で共通の文脈を
$K\subseteq L$ にもつものに対するものであるので，代入可能性により $\tuple x\equiv_{L,s}\tuple y$；これを
$[\tuple y]$ についての仮定と組み合わせる．Comp
$[\tuple z]\to\rho([\tuple u_1],\ldots,[\tuple u_r])$ で $\tuple z=\rho(\tuple u_1,
\ldots,\tuple u_r)$ なるものについては，子は $\tuple u_j$ と同値なタプルを導出し，
補題~\ref{lem:witnessed-composition} はこれを $\tuple z$ に移す．
$\widehat S\to[(w)]$ で $w\in K\subseteq L$ なるものについては，導出された $(v)$ は
$(v)\equiv_{L,s}(w)$ を満たし，$\blank_1$ は $(w)$ を受理するので，$v\in L$；$\widehat S\to\eps$ が現れるのは
$\eps\in K\subseteq L$ のときに限る．上限なしの文法は上限 $\|K\|_+$ の場合である．
\end{proof}

\begin{proof}[命題~\ref{prop:completeness} の証明]
これは具体的な文法についての \cite[Lemma~3]{Yoshinaka2015General} の議論である．
各 $\iota(X)$ と，各
$\tuple z_R:=\rho(\iota(Y_1),\ldots,\iota(Y_q))$（型付き規則
$R:X\to\rho(Y_1,\ldots,Y_q)$ について）は，$\chi(X)$ を通じて観測される．それらは同じ型をもち
（子の型の $\out_\rho$ は $X$ の型である），共通の文脈
$\chi(X)$ をもつので，Link は $[\iota(X)]\to[\tuple z_R]$ を与える．$q=0$ ならば，Const は
$[\tuple z_R]\to\tuple z_R$ を与える；そうでなければ $(\tuple z_R,\rho,(\iota(Y_j))_j)$ は
ランク $q\le r_G$（かつ \eqref{eq:sample-rank} により $q\le\|K\|_+$）の合成証拠であるので，
Comp は $[\tuple z_R]\to\rho([\iota(Y_1)],\ldots,[\iota(Y_q)])$ を与える．帰納法により，
$X\Rightarrow^*_{\widetilde G_0}\tuple u$ は $[\iota(X)]\Rightarrow^*\tuple u$ を導く．
$w\in L(G)=L(\widetilde G_0)$（命題~\ref{prop:output-invariants}(iii)）について，
成功する導出は開始の連結 $\widetilde S\to X$ で始まる；このとき
$\chi(X)=\blank_1$ であるので，$\iota(X)=(w')$ で $w'\in K$ であり，Start は
$\widehat S\to[\iota(X)]\Rightarrow^*(w)$ を与える．
\end{proof}

\begin{proof}[定理~\ref{thm:exact-reconstruction} の証明]
$L\cap\Sigma^+=\emptyset$ ならば，$K=L\subseteq\{\eps\}$ であり，
定義~\ref{def:reconstruction} が主張を与える．そうでなければ，$\mathcal F_L\subseteq L$
は有限であり，$\iota=\omega$ とした命題~\ref{prop:completeness}（標準タプルは
命題~\ref{prop:output-invariants}(i) により正しい型をもつ）は
$L\cap\Sigma^+\subseteq L(\mathcal B(K))$ を与え，命題~\ref{prop:reconstruction-soundness}
は $L(\mathcal B(K))\subseteq L$ を与え，$\eps\in L(\mathcal B(K))$ であるのは $\eps\in K$ のとき，かつそのときに限り，それは
$\eps\in L$ のとき，かつそのときに限る．ランクを保つ正規化が最後の文を与える．
\end{proof}

\begin{proof}[定理~\ref{thm:poly} と系~\ref{cor:gold-explanatory} の証明]
\cite[Lemma~6]{Yoshinaka2011} と同様である：出現は標本の語と高々
$2f$ 個の切れ目の位置であり，ランク $\le r_{\max}$ の合成証拠は高々
$2fr_{\max}$ 個をさらに加えるので，$(1+\|K\|_+)^{O(fr_{\max})}$ 個の候補の規則があり，各々は
多項式時間で検査される．同定は，
\cite[Corollary~1 and Theorem~1]{Yoshinaka2011} の保守的学習器の議論である：
定理~\ref{thm:exact-reconstruction} の特徴
標本が現れると，すべての新しい
仮説は正しく，決して変えられない以前の仮説は，
命題~\ref{prop:reconstruction-soundness} により健全であり，また，
それ以降のすべてのデータがそれによって生成されるので完全である．所属は，固定した $f$ と $r_{\max}$ について，多項式
時間で決定可能である \cite[Proposition~4]{Yoshinaka2011}．これらの
いずれも収束の段階を評価しない．
\end{proof}

\section{型付き証拠と特徴データの証明}
\label{app:typed-witness-aux-proofs}

\begin{proof}[命題~\ref{prop:output-invariants} の証明]
すべての規則の適用において，準同型の性質は次を与える．
\[
 \widehat s^{(\fan(A))}(\rho(\tuple u_1,\ldots,\tuple u_r))=
 \out_\rho\bigl(\widehat s^{(d_1)}(\tuple u_1),\ldots,\widehat s^{(d_r)}(\tuple u_r)\bigr).
\]  高さに関する帰納法は，(i) と，(ii) における持ち上げの存在と一意性を与える．
各型のラベルはその産出によって強制されるからである．ラベルを消すと
$L(G^{\widehat s})\subseteq L(G)$ が得られ，(ii) は逆を与え，刈り込みは
すべての成功する導出のすべての規則を保つ．(iv) については，終端導出を
$X$ の出現（$\widetilde G_0$ の成功する導出における）に挿入する：型付き規則は型付きの
記号のみに依存するので，結果は $G^{\widehat s}$ の成功する導出であり，刈り込みは
そのすべての規則を残す．（$G^{\widehat s}$ の規則で，その子が終端導出をもたないもの，例えば
$A_{(0)}\to\mathtt aB_{(0)}$ で $B$ が別の型の語のみを導出する場合のものは，残されない；
それはどの終端導出にも現れない．）
\end{proof}

定義~\ref{def:class-cmcf}，\ref{def:reconstruction}
と~\ref{def:output-refinement} は，任意の有限モノイド準同型
$q:\Sigma^*\to\mathsf N$ を $\widehat s$ の代わりに用いても文字どおりに適用できる（タプルの型は $q^{(d)}$ とする）；
そのクラスを $\Cmcf{f}{q}$ と書く．クラスと作用素は空でない
成分の型のみを用いるので，$q|_{\Sigma^+}$ のみに依存する．

\begin{lemma}[モノイドによる実装]
\label{lem:semigroup-transfer}
$q:\Sigma^*\to \mathsf N$ を有限モノイド準同型とし，$\mathsf S_0:=q(\Sigma^+)$ かつ
$s_0:=q|_{\Sigma^+}$ とする．
\begin{enumerate}[label=\textup{(\roman*)}]
\item $(\Sigma^+)^d$ 上で，$q^{(d)}$-，$s_0^{(d)}$- および
$\widehat{s_0}^{(d)}$-値の等しさは一致する．したがって $\Cmcf{f}{q}=\Cmcf{f}{s_0}$ であり，
$q$ で作った作用素と $\widehat{s_0}$ で作った作用素は，すべての
標本の上で同じ文法を返す．
\item すべての良い MCFG $G$ について，$G$ の刈り込んだ細分で $q$ の下のものと
$\widehat{s_0}$ の下のものは一致する（刈り込まないものは一致する必要はない）；したがって同じ標準
タプル，文脈および証拠の語を選んでよい．
\end{enumerate}
\end{lemma}

\begin{proof}
(i) 三つの写像は $\Sigma^+$ 上で一致する．(ii) $\eps$-free 性と
命題~\ref{prop:output-invariants}(i) により，どちらの準同型の下でも，成功する
導出のすべての型付き非終端記号は $A_{\mathbf t}$ で $\mathbf t\in \mathsf S_0^{\fan(A)}$ なるものであり，
用いられるすべての開始の連結は $\widetilde S\to S_{(m)}$ で $m\in \mathsf S_0$ なるものである．子が
$\mathsf S_0$ で添字付けられているとき，空でない各雛形成分は両方の準同型によって $\mathsf S_0$ において同一に評価されるので，
両方の細分は $\mathsf S_0$ で添字付けられた記号の間で同じ規則をもち，
それらの記号がすべての成功する導出を担う．
\end{proof}

\begin{proof}[補題~\ref{lem:identity-canonical-contexts} の証明]
$X$ を含む成功する型付き導出をとり，根から
$X$ の出現への道を，$\blank_1$ から始めてたどる；各規則で，
（空でない）兄弟の産出を固定して補題~\ref{lem:filling-identity} を適用する．これは
$\chi(X)\in\mathcal E_d$ を与え，$X$-部分木を任意の $X$-導出で置き換えると
$\chi(X)[\tuple u]$ の成功する導出が得られる．$\widetilde S\to X$ が規則ならば，
長さ0の道をとる．
\end{proof}

\begin{proof}[補題~\ref{lem:typed-homogeneity} の証明]
$\tuple u,\tuple v\in L_X$ について，命題~\ref{prop:output-invariants}(i) は等しい
型を与え，$\chi(X)$ は両方を受理し，アリティは高々 $f$ である；代入可能性は
$\D_L(\tuple u)=\D_L(\tuple v)$ を与える．
\end{proof}

\begin{proof}[補題~\ref{lem:thickness-equivalence} の証明]
規則 $R$ で左辺が $A$ のものについて，成功する導出 $D_R$ で
長さ $\theta_G(R)$ の，$R$ を用いるものを固定する；$G$ は既約であるのでそれは存在する．
非削除性により，$R$ が書くすべての終端記号と，$A$ のその出現で生成される
タプルのすべての成分は，文の中にちょうど一回残る．これは
(i) を与え，すべての生産的な $A$ は規則をもつので (ii) を与える．

(iv) $A=S$ については $\blank_1$ をとる．そうでなければ，$A$-部分木（$D_R$ の）を削除し，
その成分の場所に順序付きの穴を置く；根への道に沿って
補題~\ref{lem:filling-identity} を繰り返し適用すると（兄弟は
空でない産出をもつ），$E_A\in\mathcal E_d$ が示され，その固定された文字は
$D_R$ の文に属する．任意の $A$-導出を挿入すると成功する導出が得られるので，
$E_A$ は $L_A(G)\supseteq L_X$ を受理する
（命題~\ref{prop:output-invariants}(i)）．開始の
連結をもつ型付き非終端記号は $S_{\mathbf t}$ の形をもち，$E_S=\blank_1$．導出性については，
任意の $\widetilde G_0$-導出 $D_X$（$X$ からの）のラベルを消したものを
穴に挿入する；得られた成功する導出の一意な持ち上げ
（命題~\ref{prop:output-invariants}(ii)）は，すべての節点に，その
非終端記号とその産出の $\widehat s$-ベクトルとをラベルとして付け，それは $D_X$ と同様であるので，持ち上げは
穴に $D_X$ を含み，その $X$-部分木を削除すると $E_A$ が残る．

(iii) $S$ から $A$ への最短の依存の道をたどる；それは
非終端記号を繰り返さないので，その規則は相異なり，$R$ とも異なる．道の上のすべての
兄弟と $R$ のすべての子を最短の終端タプル（長さ
$\le t_G$）で完成させ，$R$ を適用する．用いた相異なる規則は高々 $\|G\|$ 個の
終端記号を書き，高々 $\|G\|$ 個の子の出現をもつので，文の長さは
高々 $\|G\|+\|G\|t_G\le2\|G\|\max\{1,t_G\}$ である．
\end{proof}

\begin{proof}[補題~\ref{lem:witness-length-template} の証明]
(i) $G$ は $\eps$-free であるのですべての成分は空でなく，したがって $s\preceq\lambda$
と命題~\ref{prop:output-invariants}(i) は
$\widehat s^{(\fan(B))}(\tuple u')=\widehat s^{(\fan(B))}(\tuple u)=\mathbf t$ を与える．
命題~\ref{prop:output-invariants}(ii) により，$G$-導出（$\tuple u'$ の）は
$G^{\widehat s}$-導出（$X$ からの）に持ち上がるので，
命題~\ref{prop:output-invariants}(iv) により $\tuple u'\in L_X$．二番目の主張については，
$(\mathrm R_\beta^\lambda)$ をある $\tuple u\in L_X\ne\emptyset$ に適用する．

(ii) $\Wit_\iota$ の語は $\chi(X)[\iota(X)]$ または規則の語である．文脈は
高々 $\bar\theta_G$ 個の固定された文字をもち，型付き規則はその元の規則の雛形をもち，
したがって高々 $\bar\theta_G$ 個の文字を書き
（補題~\ref{lem:thickness-equivalence}(i)），高々 $r_{\max}$ 個の
子のタプルがある．型付き非終端記号ごと，および型付き規則ごとに一つの語がある．

(iii) 型付き規則を，$\widetilde G_0$-導出（
$\omega(Y_j)\in L_{Y_j}$ の）に適用し，ラベルを消し，その結果を，
成功する $G$-導出で $E_B$ を得たものに挿入する．
補題~\ref{lem:thickness-equivalence}(iv) の証明と同様に，一意な持ち上げは型付き
導出を穴に含み，$G^{\widehat s}$ の成功する導出のすべての規則は，
$\widetilde G_0$ の定義により残される．開始の連結 $\widetilde S\to S_{\mathbf t}$ は，
$\omega(S_{\mathbf t})$ の成功する導出によって用いられ，その長さは高々 $b$ である．

(iv) $\chi(X)$ は $L_X$ のすべてのタプルを受理し，$\iota(X)$ はそれらと同値であるので，
$\chi(X)[\iota(X)]\in L$．型付き規則 $X\to\rho(Y_1,\ldots,Y_q)$ と
$\tuple u_j\in L_{Y_j}$ について，補題~\ref{lem:witnessed-composition} は
$\rho(\iota(Y_1),\ldots,\iota(Y_q))\equiv_{L,s}\rho(\tuple u_1,\ldots,\tuple u_q)\in L_X$ を与えるので，
規則の語も $L$ に属する．同値なタプルは等しい型をもつので，
$\widehat s^{(\fan(B))}(\iota(X))=\mathbf t$ が $X=B_{\mathbf t}$ について成り立つ．したがって
命題~\ref{prop:completeness}（$r_G\le r_{\max}$，かつ
$\chi(S_{\mathbf t})=E_S=\blank_1$）と命題~\ref{prop:reconstruction-soundness}
は $L(\mathcal B(K))\cap\Sigma^+=L\cap\Sigma^+$ を，すべての有限の $K$ で
$\Wit_\iota\subseteq K\subseteq L$ なるものについて，両方の作用素について与える；Start 規則が $\eps$ を扱う．
\end{proof}

\begin{proof}[定理~\ref{thm:data-from-typed-thickness} の証明]
補題~\ref{lem:witness-length-template}(ii),(iv) を，$\iota=\omega$ と
$b=t_{\mathrm{typ}}$ について適用し，上界 $C_s(G)$ を用いる代わりに，開始以外の型付き非終端記号と規則を数える．
\end{proof}

\begin{proof}[命題~\ref{prop:rare-type} の証明]
補題~\ref{lem:distinguishing} を $L'=L^{\setminus m}$ に適用する．これは
補題~\ref{lem:type-deletion} により $\Cmcf{f}{s}$ に属する．また，$\eps\in L$ ならば
$L'=L\setminus\{\eps\}$ に適用する；標準の作用素は，定理~\ref{thm:exact-reconstruction} により，その
両方に対して特徴標本をもつ．相異なる正の型は
相異なる標本の語を要求し，型 $m$ の語の費用は少なくとも $\mu_m(L)+1$ である．
\end{proof}

\begin{proof}[命題~\ref{prop:typed-thickness-exponential} の証明]
すべての雛形は空でなく，各変数を順に一回ずつ用いるので，$G_n$ は
良い；それは既約であり，$\Sigma^+$ を $S\to SS\mid\mathtt a\mid\mathtt c$ を通じて生成する．
各規則は長さ高々 $2$ の成功する使用をもつ：すべての $A_i$ を
$\mathtt c$（$i\ge1$）または $\mathtt a$（$i=0$）で完成させる．$A_n$ の産出が型 $1$ をもつのは，それが
$\mathtt c$ を含まないとき，かつそのときに限るので，$(A_n)_{(1)}$ の唯一の産出は $\mathtt a^{2^n}$ であり，
$t_{\mathrm{typ}}(G_n,s_{\mathtt c})\ge2^n$；型付き規則
$S_{(1)}\to(A_n)_{(1)}$ は（$S\to A_n$ を通じて）残され，それを用いるすべての成功する
導出は長さ少なくとも $2^n$ をもつ．
\end{proof}

\begin{proof}[系~\ref{cor:distribution-representatives} の証明]
$L_X\subseteq\operatorname{Cl}_{L,s}(X)$ は $\tau\le t_{\mathrm{typ}}$ を与える．
補題~\ref{lem:witness-length-template}(ii),(iv) を，$\iota(X)$ を最短のタプル（
$\operatorname{Cl}_{L,s}(X)$ の）とし $b=\tau$ として適用し，
定理~\ref{thm:data-from-typed-thickness} と同様に数える；$K:=\Wit_\iota\cup(\{\eps\}\cap L)$．
最後に，$m\in \mathsf S_L$ について，型付き開始非終端記号 $X=S_{(m)}$ は
$\chi(X)=\blank_1$ をもつので，$\blank_1\in\D_L(\iota(X))$；したがって $\iota(X)=(v)$ で
$v\in L$ かつ $s(v)=m$ であり，$\tau_X\ge\mu_m(L)$．
\end{proof}

\section{良い正規化の詳細}
\label{app:normalization-details}

以下の四つの変換は命題~\ref{prop:good-normalization} を与える；各々は，
大きさの上界を読み取れるタプル言語の恒等式を満たす．
非削除の変換と非併合の変換の正しさは標準的である
\cite{SekiEtAl1991,Kanazawa2019Ogden,Yoshinaka2011}；大きさの上界は
そこには述べられておらず，ここで導く．
$\tuple w=(w_1,\ldots,w_d)$ と空でない $I\subseteq[d]$ について，$\pi_I(\tuple w)$ は
$I$ で添字付けられた成分を順に残し，$\sigma\tuple w$ は成分を置換し，
$\psi(\tuple w)$ は成分の連続したかたまりを連接する；したがって
$\|\pi_I(\tuple w)\|_1\le\|\tuple w\|_1=\|\sigma\tuple w\|_1=\|\psi(\tuple w)\|_1$．

\begin{lemma}[段階ごとのタプル言語]
\label{lem:normalization-stages}
$G_0$ を既約な $f$-MCFG で，ランク高々 $r_{\max}$ のものとする．以下の各変換は
線形性を保ち，ファンアウト，ランクおよびどの雛形も増やさず，
非終端記号と規則の個数を，$f$ と $r_{\max}$ のみに依存する因子倍にする．
\begin{enumerate}[label=\textup{(\alph*)},leftmargin=*]
\item \emph{非削除}（
\cite[Lemma~2.2]{SekiEtAl1991} の証明における構成 (f3)）．$A^I$ を空でない $I\subseteq[\fan(A)]$ について導入する．
$R:A\to\rho(B_1,\ldots,B_q)$ について，$J_j$ を $B_j$ の成分で，
$I$ で添字付けられた出力（$\rho$ の）が用いるものとする；これらの出力を残し，$B_j^{J_j}$ を用い（$J_j\ne\emptyset$ ならば），
そうでなければ $B_j$ を省く．規則は非削除であり，$L_{A^I}=\pi_I(L_A(G_0))$．
\item \emph{非置換．}ファンアウト $d$ の非終端記号を
$\sigma\in\mathfrak S_d$（$\{1,\ldots,d\}$ 上の対称群）で飾る．
子の置換は，親の雛形がその変数を用いる順序によって定まる．
飾った非終端記号は $\sigma L_A$ を生成する．
\item \emph{空の成分．}$\mathsf{Pat}(A)$ を
$L_A$ の台の型の集合とし，飽和によって計算する．実現可能な子の型について，
空の成分の変数を $\eps$ で置き換え，空の親の成分を削除し，すべてが空の子を省き，
$A^I$ を各空でない $I\in\mathsf{Pat}(A)$ について導入する．このとき
$L_{A^I}=\pi_I\{\tuple u\in L_A:\operatorname{supp}(\tuple u)=I\}$ であり，結果は
$\eps$-free であり，$S^{\{1\}}$ は言語の空でない部分を生成する．
\item \emph{非併合}（
\cite[Section~2.3, Lemma~1]{Yoshinaka2011} の証明における構成）．非終端記号を，そのファンアウトの合成
$\psi$ で飾る．子の合成は，その
連続する変数の極大な連なりで親の雛形におけるものによって与えられる．飾った非終端記号は
像 $\psi(L_A)$ を生成する．
\end{enumerate}
段階 (b)--(d) は，それ以前の段階で確立した性質を保つ．
\end{lemma}

\begin{proof}
(a) 任意の子のタプルについて
$\pi_I(\rho(\tuple v_1,\ldots,\tuple v_q))=\rho^I((\pi_{J_j}\tuple v_j)_{J_j\ne\emptyset})$；
帰納法は $\subseteq$ を与え，$\supseteq$ については，省いた子を
終端導出で完成させる．それは $G_0$ が既約であるので存在する．各非終端記号と規則の複製は高々 $2^f$
個である．
(b) 線形性と非削除性により子の置換は well-defined である；帰納法は
$\sigma L_A$ を与える．複製は高々 $f!$ 個である．
(c) 残した変数は一回ずつ順に現れる；帰納法は恒等式を与え，
省いたすべてが空の子は，その空の型の導出によって完成させる．
残した子 $B^J$ は $J$ で添字付けられた成分のみをもち，その各々は
親の雛形によってちょうど一回，元の順序で用いられるので，非削除性と
非置換性は持続する．型は非終端記号ごとに
高々 $2^f$ 通りであり，規則ごとの選び方は $2^{fr_{\max}}$ 通りである．
(d) 非置換性により連なりは各子の順序と両立し，非削除性により
それらはその成分を覆う；このとき
$\rho^\psi(\psi_1(\tuple v_1),\ldots,\psi_q(\tuple v_q))=\psi(\rho(\tuple v_1,\ldots,\tuple v_q))$
であり，帰納法は像の恒等式を与える（併合は単射である必要はないので，
逆像についての主張はしない）．空でない成分を連接すると空でない
成分が得られ，子の併合した各ブロックは一回ずつ順に用いられるので，
$\eps$-free 性，非削除性および非置換性は持続する．大きさは固定した $f,r_{\max}$ について線形である
\cite[Section~2.3, after Example~2]{Yoshinaka2011}．
\end{proof}

\begin{proof}[命題~\ref{prop:good-normalization} の証明]
$L(G_{\mathrm{in}})\cap\Sigma^+=\emptyset$ ならば，$G_{\mathrm{good}}$ を
開始記号のみからなり規則をもたないものとする；それは良く，
$\emptyset=L(G_{\mathrm{in}})\setminus\{\eps\}$ を生成する．そうでなければ，(a)--(d) を刈り込んだ
入力に適用し，再び刈り込む．各段階は大きさを
$f,r_{\max}$ のみに依存する因子倍にするので，
$\|G_{\mathrm{good}}\|=O_{f,r_{\max}}(\|G_{\mathrm{in}}\|)$ であり，
(c) により $L(G_{\mathrm{good}})=L(G_{\mathrm{in}})\setminus\{\eps\}$．この言語が
空でなければ，最後の刈り込みは既約な文法を残し，刈り込みは良さを保つ．
\end{proof}


\section{三文字の倍増族}
\label{app:bracketed-family}
\label{app:three-letter-family}

この付録は命題~\ref{prop:three-letter-family} を証明する；~(ii) の証明は
\cite[Appendix~C]{kuriyamaCFG} の決定的パーザの議論を手直ししたものである．
$n\ge1$ を固定し，$x,y\in\Sigma^*$ について $x\equiv y$ と書くのは
$\{(u,v):uxv\in L_n\}=\{(u,v):uyv\in L_n\}$ であるときとする；これはモノイド合同である．

\paragraph{パーザ．}
\emph{スタック}とは，$E_0,\ldots,E_n$ と閉じ標識
$\hat{\mathtt b}_i,\hat{\mathtt c}_i$（$1\le i\le n$）の上の語であり，その頂上を左に置く；$\eps$
は空のスタックである．すべてのスタック $\mathcal Z$ と $1\le i\le n$ について，遷移は次である．
\[
\begin{aligned}
E_0\mathcal Z&\xrightarrow{\mathtt a}\mathcal Z,&
E_i\mathcal Z&\xrightarrow{\mathtt b}E_{i-1}E_{i-1}\hat{\mathtt b}_i\mathcal Z,&
E_i\mathcal Z&\xrightarrow{\mathtt c}E_{i-1}\hat{\mathtt c}_i\mathcal Z,\\
\hat{\mathtt b}_i\mathcal Z&\xrightarrow{\mathtt b}\mathcal Z,&
\hat{\mathtt c}_i\mathcal Z&\xrightarrow{\mathtt c}\mathcal Z,
\end{aligned}
\]
そして他の遷移はない；特に，どの遷移も $\eps$ から出ない．
$\mathcal X\xrightarrow{w}\mathcal Y$ と書くのは，$w$ を $\mathcal X$ から読むと
$\mathcal Y$ で終わるときとする．三つの性質を用いる．
\begin{enumerate}[label=\textup{(D\arabic*)},leftmargin=2.6em]
\item \emph{決定性．}頂上の記号と文字が遷移を定めるので，
すべての語はスタックの上の部分写像として作用する．
\item \emph{重み．}$E_i$ に重み $4^i$ を，各標識に重み $1$ を与え，
$\operatorname{wt}(\mathcal X)$ を重みの総和とする．すべての遷移は
$\operatorname{wt}$ を下げる．$4^i>2\cdot4^{i-1}+1$ だからである；したがって
$\mathcal X\xrightarrow{w}\mathcal X$ は $w=\eps$ についてのみ成り立つ．
\item \emph{局所性．}遷移は頂上の記号のみを読み書き換え，標識は
取り出されるだけである．したがって $\mathcal X\xrightarrow{w}\mathcal Y$ は
$\mathcal X\mathcal Z\xrightarrow{w}\mathcal Y\mathcal Z$ を導く．逆に，
$\mathcal X_1\ne\eps$ かつ $\mathcal X_1\mathcal Z\xrightarrow{w}\mathcal Y$ ならば，次のいずれかである：
\textup{(a)} 実行が $\mathcal Z$ に達する：$w=w_1w_2$ で
$\mathcal X_1\xrightarrow{w_1}\eps$ かつ $\mathcal Z\xrightarrow{w_2}\mathcal Y$，したがって
$\operatorname{wt}(\mathcal Y)\le\operatorname{wt}(\mathcal Z)$；または \textup{(b)} それが
達しない：$\mathcal Y=\mathcal Y_1\mathcal Z$ で $\mathcal X_1\xrightarrow{w}\mathcal Y_1\ne\eps$，
したがって $\operatorname{wt}(\mathcal Y)>\operatorname{wt}(\mathcal Z)$ であり，$\mathcal Y_1$ は，
$\mathcal X_1$ の最後の記号が標識であるときにはいつでも，その記号で終わる．
固定した $\mathcal Z$ について，$\operatorname{wt}(\mathcal Y)$ を
$\operatorname{wt}(\mathcal Z)$ と比べるとどちらの場合かが決まる．
\end{enumerate}
文字 $\mathtt b$ または $\mathtt c$ は開きと閉じの両方を行うので，語は
すべてのスタックの上で単射に作用する必要はない；単射性は\emph{到達可能な}スタック，
すなわち $\mathcal X$ で $E_n\xrightarrow{u}\mathcal X$ がある $u$ について成り立つものの上で成り立つ．

\begin{lemma}\label{lem:app-br-recognition}
\label{lem:app-tl-reachable}
\textup{(i)} $0\le i\le n$ とすべてのスタック $\mathcal Z$ について，
$E_i\mathcal Z\xrightarrow{w}\mathcal Z$ であるのは $w\in L^{(i)}$ であるとき，かつそのときに限る．

\textup{(ii)} $\mathcal F_i:=\{E_{i-1}\hat{\mathtt b}_i,\hat{\mathtt b}_i,\hat{\mathtt c}_i\}$ とおく．
すべての到達可能なスタックは $\eps$ であるか，または $\mathcal YF_{j+1}\cdots F_n$ であって
$F_i\in\mathcal F_i$ かつ $\mathcal Y\in\{E_j,\hat{\mathtt b}_j,\hat{\mathtt c}_j\}$ なるものである
\textup($\mathcal Y=E_j$ で $0\le j\le n$，または標識で $1\le j\le n$\textup)．
すべての語は到達可能なスタックの上で部分単射として作用する．
\end{lemma}

\begin{proof}
(i) 実行 $E_i\mathcal Z\xrightarrow{w}\mathcal Z$ は (D3) の場合~(a) にある．なぜなら
場合~(b) は $\mathcal Y_1\mathcal Z\ne\mathcal Z$ で終わるからである；したがって $w=w_1w_2$ で
$E_i\xrightarrow{w_1}\eps$ かつ $\mathcal Z\xrightarrow{w_2}\mathcal Z$ であり，(D2) により $w_2=\eps$．
(D3) における持ち上げとあわせて，$E_i\xrightarrow{w}\eps$ であるのは
$w\in L^{(i)}$ であるとき，かつそのときに限ることを，$i$ に関する帰納法で示せば十分である．$E_0$ からは $\mathtt a$-遷移のみが適用でき，
$\eps$ はどの遷移も許さない．$i\ge1$ については最初の文字は $\mathtt b$ または $\mathtt c$ である．
$\mathtt c$ の後，スタック $E_{i-1}\hat{\mathtt c}_i$ は空にされなければならない；(D3) の場合~(a) を
$\mathcal X_1=E_{i-1}$，$\mathcal Z=\hat{\mathtt c}_i$ として用いると，残りは $zz'$ と分かれ，
$E_{i-1}\xrightarrow{z}\eps$ かつ $z'=\mathtt c$ である．$\mathtt b$ の後，
$E_{i-1}E_{i-1}\hat{\mathtt b}_i$ のそのような二回の分割は $z_1z_2\mathtt b$ で
$E_{i-1}\xrightarrow{z_j}\eps$ なるものを与える．

(ii) そのようなスタックの集合は $E_n$ を含み，遷移について閉じている：
$E_jF_{j+1}\cdots F_n$ から，$\mathtt b$ は
$E_{j-1}(E_{j-1}\hat{\mathtt b}_j)F_{j+1}\cdots F_n$ を与え，$\mathtt c$ は
$E_{j-1}\hat{\mathtt c}_jF_{j+1}\cdots F_n$ を与える；頂上を取り出すと
$F_{j+1}\cdots F_n$ が残り，これは $\eps$ であるか，再び述べた形である．ここで最初のフレーム
$E_j\hat{\mathtt b}_{j+1}$ は，$\mathcal Y=E_j$ にフレーム
$\hat{\mathtt b}_{j+1}$ が続くものとして読む．

単射性については，文字を考えれば十分である．遷移は到達可能な
スタックを到達可能なスタックに写すからである．\emph{フレーム積}とはスタック
$F_{i+1}\cdots F_n$ で $F_k\in\mathcal F_k$ なるものである（$i=n$ ならば空のスタック）．各
フレーム $F_k$ はちょうど一つの標識を含み，その添字は $k$ であるので，フレーム積は
$i$ を定める：それは $n$ である（$\eps$ については）．そうでなければ最小の
標識の添字より一つ小さい．フレーム積において，記号 $E$ の後には常に標識
$\hat{\mathtt b}$ が続く．到達可能なスタックの上での文字の遷移と，
その結果の形は次のとおりである：
\begin{center}
\footnotesize
\renewcommand{\arraystretch}{1.15}
\begin{tabular}{@{}llll@{}}
\hline
文字 & 遷移（到達可能な元） & 結果の先頭 & 元の復元方法\\
\hline
$\mathtt a$ & $E_0\mathcal Z\to\mathcal Z$ & フレーム積 & $\mathcal Z$\\
$\mathtt b$，開き & $E_i\mathcal Z\to E_{i-1}E_{i-1}\hat{\mathtt b}_i\mathcal Z$ & 二つの記号 $E$ & $i$ と $\mathcal Z$ を読み取る\\
$\mathtt b$，閉じ & $\hat{\mathtt b}_i\mathcal Z\to\mathcal Z$，$\mathcal Z$ はフレーム積 & 標識，$E\hat{\mathtt b}$，または $\eps$ & $i$ を $\mathcal Z$ から\\
$\mathtt c$，開き & $E_i\mathcal Z\to E_{i-1}\hat{\mathtt c}_i\mathcal Z$ & $E$ の後に $\hat{\mathtt c}$ & $i$ と $\mathcal Z$ を読み取る\\
$\mathtt c$，閉じ & $\hat{\mathtt c}_i\mathcal Z\to\mathcal Z$，$\mathcal Z$ はフレーム積 & 標識，$E\hat{\mathtt b}$，または $\eps$ & $i$ を $\mathcal Z$ から\\
\hline
\end{tabular}
\end{center}
閉じの行において $\mathcal Z$ がフレーム積であるのは，到達可能なスタック
で標識 $\hat{\mathtt b}_i$ または $\hat{\mathtt c}_i$ を頂上にもつものは $\mathcal YF_{i+1}\cdots
F_n$ で $\mathcal Y$ がその標識であるものだからである．$\mathtt b$ については，開きの結果は
二つの記号 $E$ で始まり，閉じの結果は決してそうならない．フレーム積においては
記号 $E$ の後に標識が続くからである；$\mathtt c$ については，開きの結果は
$E\hat{\mathtt c}$ で始まり，これはフレーム積には決して現れない．したがって一つの文字の開きと
閉じは互いに素な像をもち，各行はその最後の
列により単射であり，すべての文字，したがってすべての語は，
到達可能なスタックの上で部分単射として作用する．
\end{proof}

\begin{lemma}\label{lem:app-br-exchange}
\label{lem:app-tl-recognition}
$\mathcal X$ が到達可能であり，$w\in L^{(i)}$ かつ $\mathcal X\xrightarrow{w}\mathcal Y$ ならば，
$\mathcal X=E_i\mathcal Y$．したがって $w\equiv w'$ がすべての $w,w'\in L^{(i)}$ について成り立つ．
\end{lemma}

\begin{proof}
$i$ に関する帰納法による；$i=0$ については，$\mathtt a$ は $E_0$ を取り出すだけである．$i\ge1$ とする；出会うすべてのスタックは
到達可能である．$w=\mathtt cv\mathtt c$ で $v\in L^{(i-1)}$ であり，最初の
$\mathtt c$ が $E_j$ を開くならば，
$v$ についての帰納法の仮定（$E_{j-1}\hat{\mathtt c}_j\mathcal Z$ における）は $j=i$ を与え，最後の $\mathtt c$ は
$\hat{\mathtt c}_i$ を取り出す．最初の $\mathtt c$ が標識を取り出すならば，残りは $\eps$ であるか
フレームの積である；$v$ を読むためには，帰納法の
仮定によりそれは $E_{i-1}$ で始まらなければならないので，その最初のフレームは $E_{i-1}\hat{\mathtt b}_i$ であり，$v$ の後の頂上
は $\hat{\mathtt b}_i$ であり，その上で $\mathtt c$ は遷移をもたない．
$w=\mathtt bv_1v_2\mathtt b$ であり，最初の $\mathtt b$ が $E_j$ を開くならば，帰納法の
仮定は $j=i$ を与え，語 $v_1,v_2$ は $E_{i-1}$ の二つの複製を取り出し，
最後の $\mathtt b$ は $\hat{\mathtt b}_i$ を取り出す．最初の $\mathtt b$ が標識を取り出すならば，
前と同様に $v_1$ は頂上 $\hat{\mathtt b}_i$ を残し，帰納法の仮定により，そこから $v_2$ は
読めない．帰結については，$uwv\in L_n$ ならば，実行
$E_n\xrightarrow{u}\mathcal X\xrightarrow{w}\mathcal Y\xrightarrow{v}\eps$ は
$\mathcal X=E_i\mathcal Y$ をもち，
補題~\ref{lem:app-br-recognition}(i) により $E_i\mathcal Y\xrightarrow{w'}\mathcal Y$；したがって $uw'v\in L_n$ であり，逆も同様である．
\end{proof}

\begin{lemma}[分岐の同期]\label{lem:app-br-branch}
$\mathcal X\xrightarrow{x}\mathcal Y$ かつ $\mathcal X\xrightarrow{y}\mathcal Y$ とし，ここで $x,y$
は空でなく，異なる文字で始まるとする．このとき $\mathcal X=E_i\mathcal Z$ が
ある $1\le i\le n$ とスタック $\mathcal Z$ について成り立ち，$x=wx'$，$y=w'y'$ であって，
$w,w'\in L^{(i)}$，$\mathcal Z\xrightarrow{x'}\mathcal Y$ かつ $\mathcal Z\xrightarrow{y'}\mathcal Y$ である．
\end{lemma}

\begin{proof}
頂上 $E_i$（$i\ge1$）だけが二つの異なる文字，すなわち $\mathtt b$ と
$\mathtt c$ を許す．$x=\mathtt cx_0$ かつ $y=\mathtt by_0$ とする．したがって
$E_{i-1}\hat{\mathtt c}_i\mathcal Z\xrightarrow{x_0}\mathcal Y$ かつ
$E_{i-1}E_{i-1}\hat{\mathtt b}_i\mathcal Z\xrightarrow{y_0}\mathcal Y$．(D3) により，両方の実行
は同じ場合にある．場合~(b) は不可能である：それは
$\mathcal Y=\mathcal Y_1\mathcal Z=\mathcal Y_2\mathcal Z$ で，$\mathcal Y_1$ が
$\hat{\mathtt c}_i$ で終わり $\mathcal Y_2$ が $\hat{\mathtt b}_i$ で終わるものを与えるからである．場合~(a) では，
$x_0=x_1x'$ で $E_{i-1}\hat{\mathtt c}_i\xrightarrow{x_1}\eps$ かつ
$\mathcal Z\xrightarrow{x'}\mathcal Y$ であるので，$w:=\mathtt cx_1$ は $E_i\xrightarrow{w}\eps$ を満たし，
補題~\ref{lem:app-br-recognition}(i) により $w\in L^{(i)}$．同じ議論が
$y=w'y'$ を与える．
\end{proof}

\begin{proposition}\label{prop:app-br-congruence}
\label{prop:app-tl-congruence}
$\mathcal X\xrightarrow{x}\mathcal Y$ かつ $\mathcal X\xrightarrow{y}\mathcal Y$ ならば，$x\equiv y$．
したがって $L_n$ は Clark--Eyraud の意味で代入可能である．
\end{proposition}

\begin{proof}
$|x|+|y|$ に関する帰納法による．$x$ または $y$ が空ならば，$\mathcal X=\mathcal Y$ であり，
(D2) により他方も空である．$x=cx'$ かつ $y=cy'$ が文字 $c$ について成り立つならば，(D1) は
$\mathcal X\xrightarrow{c}\mathcal X'$ で $\mathcal X'\xrightarrow{x'}\mathcal Y$ かつ
$\mathcal X'\xrightarrow{y'}\mathcal Y$ なるものを与え，帰納法が適用できる．そうでなければ，
補題~\ref{lem:app-br-branch} は $x=wx'$，$y=w'y'$ を与える；帰納法は
$x'\equiv y'$ を与え，補題~\ref{lem:app-br-exchange} は $w\equiv w'$ を与え，$\equiv$ は
合同である．

帰結については，$x,y\in\Sigma^+$ で $uxv,uyv\in L_n$ とする．(D1) と
補題~\ref{lem:app-br-recognition}(i) により，実行
$E_n\xrightarrow{u}\mathcal X\xrightarrow{x}\mathcal Y\xrightarrow{v}\eps$ と
$E_n\xrightarrow{u}\mathcal X\xrightarrow{y}\mathcal Y'\xrightarrow{v}\eps$ がある；$\mathcal Y$ と
$\mathcal Y'$ は到達可能であるので，補題~\ref{lem:app-br-recognition}(ii) は
$\mathcal Y=\mathcal Y'$ を与え，$x\equiv y$．$\{\mathtt a,\mathtt b,\mathtt c\}$ の外の文字は
空の分布しかもたない．
\end{proof}

\begin{proof}[命題~\ref{prop:three-letter-family} の証明]
(i) すべての雛形は空でなく，各変数を一回ずつ用いる；ファンアウト1では，
非置換性と非併合性は自動的に成り立つ．文法は既約である．なぜなら
$E_i\Rightarrow^*\mathtt c^i\mathtt a\mathtt c^i$ だからである．それは $n+2$ 個の非終端記号と，
大きさ $3$（開始），$2$（基底），$7$ と $5$（各水準の二つの規則）の規則をもつ
（定義~\ref{def:standard-mcfg}）．これより $\|G_n\|=13n+7$．
$E_i$ の最短の産出は長さ $2i+1$ をもち，$E_i$ より上の水準は少なくとも $2(n-i)$ 個の
文字を加え，それは規則 $E_j\to\mathtt cE_{j-1}\mathtt c$ によって達成される．したがって開始，基底
および単項の規則は規則の厚さ $2n+1$ をもち，水準 $i$ の二分の規則は
$2(n-i)+2+2(2i-1)=2n+2i$ をもつ；したがって $\bar\theta_{G_n}=4n$．

(ii) 最初の主張は命題~\ref{prop:app-br-congruence} である．それは，
自明な型付けについての $f=1$ における定義~\ref{def:tuple-substitutable} の条件であり，
有限言語 $L_n$ は $1$-MCFL である；
定義~\ref{def:class-cmcf} の後の項目~(a) は $L_n\in\Cmcf{1}{s}$ を与える．

(iii) 導出が $\mathtt c$ を含まない語を産出するのは，それが規則
$E_i\to\mathtt cE_{i-1}\mathtt c$ を用いないとき，かつそのときに限る．そのような導出（$E_i$ からの）は，
高さ $i$ の完全二分木であり，$W_i$ を産出し，$|W_i|=2|W_{i-1}|+2$，$|W_0|=1$ である．

(iv) (ii) により $L_n\in\Cmcf{1}{s_U}$ であり，(iii) により $1\in \mathsf S_{L_n}$ で
$\mu_1(L_n)=|W_n|$ である．$m=1$ とした命題~\ref{prop:rare-type} は
型 $1$ の語を強いるが，それは $W_n$ でしかあり得ない．
\end{proof}

\section{ファンアウト1での孤立した文字}
\label{app:two-letter-fanout-one}

この付録は定理~\ref{thm:two-letter-fanout-one} と系~\ref{cor:two-letter-fanout-one} を証明する．
以下では $s:\Sigma^+\twoheadrightarrow\mathsf S$ は型付け，$\mathtt a\in\Sigma$ は $s$ について孤立した文字（定義~\ref{def:isolated-letter}），
$L\in\Cmcf{1}{s}$ とし，$\eps\notin L$ と仮定してよい．
$\Sigma^+$ の語は，$\mathtt a^+$ に入るとき\emph{純}，$\mathtt a$ を含まないとき\emph{$\mathtt a$ なし}，
それ以外，すなわち $\mathtt a$ と他の文字を含むとき\emph{$\mathtt a$ 混合}という．
孤立性は次の二つの事実を通してのみ使われる：
\begin{enumerate}[label=\textup{(I\arabic*)},leftmargin=*]
\item すべての純な語は型 $s(\mathtt a)$ をもち，この型の語は純な語である；
\item すべての $\mathtt a$ 混合の語は一つの型をもつ．
\end{enumerate}
$x\in\Sigma^+$ について $\D(x):=\D^{(1)}_L(x)$ と書き，$\D(x)=\D(y)$ のとき $x\sim y$ と書く．
これは $\Sigma^+$ 上の合同関係であり，分布が空の語どうしも $\sim$ で同値である．
したがって $x\equiv_{L,s}y$ であるのは，$x\sim y$ かつ $s(x)=s(y)$ のとき，かつそのときに限る．

\subsection{準備}

\begin{lemma}[消去]
\label{lem:app-h-cancel}
$x,y\in\Sigma^+$ は同じ型をもち，$z\in\Sigma^+$，$\D(zx)\ne\emptyset$，$zx\sim zy$ とする．
このとき $x\sim y$ である．右側についても同様である．
\end{lemma}

\begin{proof}
$(l,r)\in\D(zx)$ について $lzyr\in L$ となるので，$x$ と $y$ は文脈 $(lz,r)$ を共有する．
\end{proof}

この補題は，純な $x,y$ か $\mathtt a$ 混合の $x,y$ にだけ適用する．それらはそれぞれ (I1)，(I2) から同じ型をもつ．
非終端記号 $X$ が導出的な文脈 $(l,r)$ をもてば，$X$ の同じ型の二つの生成語は $(l,r)$ を共有し，$\sim$ で同値である（補題~\ref{lem:typed-homogeneity}）．
$X$ が実現する型 $t$ について，型付き非終端記号 $X_t$ の代表元の長さ，すなわち $X$ の型 $t$ の生成語と $\sim$ で同値な型 $t$ の語の最短の長さを $\tau(X,t)$ と書く．

\paragraph{二分化．}
$G$ を $L$ のファンアウト1の reduced good な表示とし，$\sigma_G$ をその規則に現れる文字の個数とする．したがって $\sigma_G\le\min\{|\Sigma|,\|G\|\}$ である．
$t\ge3$ のすべての規則 $X\to Y_1\cdots Y_t$ を新しい非終端記号を通る右線形の鎖に置き換え，
長さ2以上の右辺に現れる終端記号を前終端記号 $T_{\mathtt x}\to\mathtt x$（使われる文字 $\mathtt x$ についてだけ）に置き換え，trim する．
得られる文法 $G'$ は reduced かつ good であり，規則は $X\to YZ$，$X\to Y$，$X\to\mathtt x$（$\mathtt x\in\Sigma$）の形で，$L$ を生成する．
$G$ の規則 $X\to\gamma$ は高々 $\max\{0,|\gamma|-2\}$ 個の新しい非終端記号と高々 $\max\{1,|\gamma|-1\}$ 個の規則を与え，
どちらも $\|G\|$ でのその規則の大きさ $1+r+|\gamma|$ 以下なので
\[
 k_V:=|V'|\le\|G\|+\sigma_G\le2\|G\|,\qquad |P'|\le\|G\|+\sigma_G
\]
であり，$|\Sigma|=2$ ならば $k_V\le\|G\|+2$，$|P'|\le\|G\|+2$ である．
さらに $\theta:=\bar\theta_{G'}\le\bar\theta_G$ である．新しい規則を使う導出は，それが由来する $G$ の規則を使う導出に対応し，
前終端の規則 $T_{\mathtt x}\to\mathtt x$ を使う導出は，$\mathtt x$ を書く $G$ の規則を使う導出に対応するからである．
$G'$ のすべての非終端記号は，長さ $\le\theta$ の生成語と，$|l|+|r|\le\theta$ の導出的な文脈 $(l,r)$ をもつ：
その非終端記号を左辺に，あるいは右辺にもつ規則を通る最短の文をとればよい．
$G$ の非終端記号は，同じ生成語をもつ $G'$ の非終端記号である．

\subsection{木と基底}

純な生成語をもつ $G'$ の非終端記号 $A$ を固定する．
その最短のものを $\mathtt a^n$，$A$ からの $\mathtt a^n$ の導出木で節点の数が最小のものを $T$，$|l|+|r|\le\theta$ をみたす $A$ の導出的な文脈を $(l,r)$ とする．
$T$ の部分木を同じラベルからの導出で置き換えると $A$ の生成語 $w$ が得られ，したがって $lwr\in L$ である．
生成語が $\theta$ より長い節点を\emph{重い}といい，重い子をもたない重い節点を\emph{種}という．

\begin{lemma}
\label{lem:app-h-tree}
\begin{enumerate}[label=\textup{(\alph*)},leftmargin=*]
\item $T$ の根から葉への道でラベルは繰り返さない；そのような道の非終端記号の節点は高々 $k_V$ 個である．
\item $T$ の同じラベル $Y$ の節点はすべて同じ生成語 $\mathtt a^{n_Y}$ をもつ．
\item 重い節点のラベル $Y$ はすべて，長さ $\le\theta$ の生成語 $\omega_Y\notin\mathtt a^+$ をもつ．
\item 種の生成語の長さは $\le2\theta$ である；種どうしは互いに素である；
$k\ge1$ 個の種を含む部分木の生成語の長さは $\le k(k_V+2)\theta$ である；
隣り合う二つの種の左端の差は高々 $2k_V\theta$ である．
\end{enumerate}
\end{lemma}

\begin{proof}
(a) 祖先と子孫が同じラベルをもてば，祖先の部分木を子孫の部分木に置き換えると，$A$ のより短い純な生成語か，より小さい木が得られる．
(b) そうでなければ長い方の部分木を短い方に置き換える．
(c) ラベル $Y$ の重い節点での規則を $\rho$ とする．
$\rho$ を使う長さ $\le\theta$ の文は，$|y|\le\theta<n_Y$ をみたす $Y$ の生成語 $y$ を含む；
$y\in\mathtt a^+$ ならば，それを $T$ に挿入すると $A$ のより短い純な生成語が得られる．
(d) 種は単位規則の節点ではなく，二つの軽い子をもつ．
種 $v$ が種 $u$ の下にあれば，$v$ の上にある $u$ の子は重い．
$k$ 個の種をもつ部分木の重い節点は種への道の上にあるので高々 $kk_V$ 個であり，種でないものはそれぞれ高々一つの軽い子をもつ．
隣り合う種 $u,v$ の間には，それらの最小共通祖先からの二本の道からぶら下がる部分木の生成語があり，それは高々 $2(k_V-1)$ 個である；
重いものがあれば $u$ と $v$ の間に種を含むので，どれも軽い．
\end{proof}

$X_0$ を $T$ の根とし，$X_{i+1}$ を $X_i$ の子のうち生成語の長い方とする．すると $n_{X_i}\le2n_{X_{i+1}}$ である．
根の部分木が種をもたなければ根は軽く $n\le\theta$ である；種を少なくとも一つ，高々 $k_V$ 個もてば $n\le k_V(k_V+2)\theta$ である．
そうでなければ，部分木が少なくとも $k_V+1$ 個の種を含む最大の添字を $i^*$ とする；
このとき $X_{i^*+1}$ が存在し，補題~\ref{lem:app-h-tree}(a) から $i^*\le k_V-2$ であり，$B_0:=n_{X_{i^*}}\le2k_V(k_V+2)\theta$ である；
$n_{X_{i^*}}<n$ と仮定してよく，$T$ は二項の節点をもつので $k_V\ge2$ である．
$X_{i^*}$ の下の種は $T$ の中で連続している；連続する $k_V+1$ 個のうち二つは同じラベル $Y$ をもち，
位置 $x_1<x_2$ にあって $d:=x_2-x_1\le2k_V^2\theta$ である．
$g_0:=d-n_Y\ge0$，$\omega:=\omega_Y$，$x:=x_1$，$y:=n-x_2-n_Y$，$y_1:=n-x_1-n_Y=d+y$ とおく．
二つの種の両方，最初のもの，二番目のものを $\omega$ で置き換えると，$A$ の生成語
\[
 W_1=\mathtt a^x\omega\mathtt a^{g_0}\omega\mathtt a^y,\qquad
 W_2=\mathtt a^x\omega\mathtt a^{d+y},\qquad
 W_3=\mathtt a^{x+d}\omega\mathtt a^y
\]
が得られる．ここで $x+g_0+y\ge1$ である：そうでなければ二つの種が，少なくとも $k_V+1\ge3$ 個の種を含む根の生成語を覆ってしまう．
$d\ge n_Y\ge1$ かつ $\omega\notin\mathtt a^+$ なので，語 $W_1,W_2,W_3$ は $\mathtt a$ 混合である．

\begin{lemma}[基底]
\label{lem:app-h-base}
$g_0\ge1$ または $\omega$ が $\mathtt a$ 混合ならば $z_0:=\omega\mathtt a^{g_0}$，$c_0:=d$ とおき，
$g_0=0$ かつ $\omega$ が $\mathtt a$ なしならば $z_0:=\omega\mathtt a^d$，$c_0:=2d$ とおく．
このとき $z_0$ は $\mathtt a$ 混合で，$\D(z_0)\ne\emptyset$，$d\mid c_0$ であり，
\[
 \mathtt a^dz_0\sim z_0\mathtt a^d\qquad\text{かつ}\qquad z_0z_0\sim\mathtt a^{c_0}z_0.
\]
\end{lemma}

\begin{proof}
(I2) から $W_1,W_2,W_3$ は一つの型をもつ；それらは $(l,r)$ を共有するので $\sim$ で同値である．
$W_2\sim W_3$ の左の $\mathtt a^x$ と右の $\mathtt a^y$ を消去すると（補題~\ref{lem:app-h-cancel}；残る語は $\mathtt a^d$ と $\omega$ を含むので $\mathtt a$ 混合であり，
消去の前の語は $W_2$ の因子なので分布は空でない），$\omega\mathtt a^d\sim\mathtt a^d\omega$ を得る．
$g_0\ge1$ または $\omega$ が $\mathtt a$ 混合ならば，$W_1\sim W_2$ で同じ消去を行うと（このとき $\omega\mathtt a^{g_0}\omega$ は $\mathtt a$ 混合），
$\omega\mathtt a^{g_0}\omega\sim\omega\mathtt a^d$ となり，$z_0z_0\sim\omega\mathtt a^{d+g_0}=z_0\mathtt a^d$ と $\mathtt a^dz_0\sim z_0\mathtt a^d$ が従う．
$g_0=0$ かつ $\omega$ が $\mathtt a$ なしならば，$n_Y=d$ かつ $x+y\ge1$ である．
$y\ge1$ ならば，$\mathtt a^x$ と $\mathtt a^{y-1}$ を消去して $\omega\omega\mathtt a\sim\omega\mathtt a^{d+1}$ を得るので，
$z_0z_0\sim\omega\omega\mathtt a^{2d}\sim\omega\mathtt a^{3d}=z_0\mathtt a^{2d}$ である；
$y=0$ ならば，$\mathtt a^{x-1}$ を消去して $\mathtt a\omega\omega\sim\mathtt a\omega\mathtt a^d$ を得るので，
$\omega\omega$ の前に文字 $\mathtt a$ を一つ残して
$z_0z_0\sim\mathtt a^{2d}\omega\omega=\mathtt a^{2d-1}(\mathtt a\omega\omega)\sim\mathtt a^{2d-1}(\mathtt a\omega\mathtt a^d)=\mathtt a^{2d}z_0$ である；
$\mathtt a$ なしの語 $\omega\omega$ は何とも比較されない．
ここで比較する語は $\mathtt a$ と $\omega$ を含むので $\mathtt a$ 混合であり，消去の前の語は $W_1$ の因子である．
どちらの場合も $z_0\mathtt a^{2d}\sim\mathtt a^{2d}z_0$ かつ $\mathtt a^dz_0\sim z_0\mathtt a^d$ である．
最後に，$z_0$ は $W_1$ または $W_2$ の因子である．
\end{proof}

$\mathtt a$ が型 $1$ の唯一の文字であるような $U_1$ の上への型付けでは，$\mathtt a$ なしの語は $\mathtt a$ 混合の語と同じ型をもち，
二番目の場合は避けられる；これは使わない．

\emph{注釈付きの符号}とは，$e\ge0$，$d\mid e$，$z\sim\mathtt a^ez_0$ をみたす組 $(z,e)$ であり，その\emph{価値}は $e+c_0$ である．
$d\mid R$ をみたす純な塊 $\mathtt a^R$ の価値は $R$ とする．
注釈付きの符号とそのような塊の指定された積 $u$ について，$\operatorname{val}(u)$ はその因子の価値の和である．
（一般に指数 $e$ は $z$ から定まらないので，価値は注釈に付ける．）

\begin{lemma}
\label{lem:app-h-collapse}
$u$ が少なくとも一つの符号を含めば，$u\sim\mathtt a^{\operatorname{val}(u)-c_0}z_0$ である．
\end{lemma}

\begin{proof}
各符号を $\mathtt a^ez_0$ で置き換え，$\mathtt a^dz_0\sim z_0\mathtt a^d$ ですべての塊を左へ移し，
$z_0z_0\sim\mathtt a^{c_0}z_0$ と $d\mid c_0$ から帰納法で従う $z_0^t\sim\mathtt a^{(t-1)c_0}z_0$ を使う．
\end{proof}

\subsection{重い道に沿った符号}

$n_{X_i}<n$ となる最小の添字を $i_0$ とする；$1\le i_0\le i^*$ である．
$i_0\le i\le i^*$ について，$T$ の中での $X_i$ の文脈を $(\mathtt a^{p_i},\mathtt a^{q_i})$ とすると $p_i+q_i\ge1$ である；
$X_i$ の部分木は二つの種を含む；その中で最初の種の左と右にある文字の数を $L_i,R_i$ とすると $R_i\ge g_0+n_Y\ge1$ であり，
$X_i$ の $\mathtt a$ 混合の生成語 $\Omega_i:=\mathtt a^{L_i}\omega\mathtt a^{R_i}$ をとる．

\begin{lemma}
\label{lem:app-h-borrow}
$\omega_{X_i}$ が $\mathtt a$ 混合ならば $\omega_{X_i}\sim\Omega_i$ である．
$\omega_{X_i}$ が $\mathtt a$ なしで $p_i\ge1$ ならば $\mathtt a\omega_{X_i}\sim\mathtt a\Omega_i$ である；
$\omega_{X_i}$ が $\mathtt a$ なしで $p_i=0$ ならば $\omega_{X_i}\mathtt a\sim\Omega_i\mathtt a$ である．
\end{lemma}

\begin{proof}
$l\mathtt a^{p_i}\omega_{X_i}\mathtt a^{q_i}r$ と $l\mathtt a^{p_i}\Omega_i\mathtt a^{q_i}r$ はともに $L$ に入る．
一番目の場合，$\omega_{X_i}$ と $\Omega_i$ は $\mathtt a$ 混合で文脈 $(l\mathtt a^{p_i},\mathtt a^{q_i}r)$ を共有する．
二番目の場合，$\mathtt a\omega_{X_i}$ と $\mathtt a\Omega_i$ は $\mathtt a$ 混合で文脈 $(l\mathtt a^{p_i-1},\mathtt a^{q_i}r)$ を共有する；
三番目の場合は $q_i\ge1$ として対称である．(I2) で結論する．
\end{proof}

$s'_i,s_i\in[0,d]$ を次のようにとる：$d\mid L_i+s'_i$；
補題~\ref{lem:app-h-base} の第一の場合は $R_i+s_i\equiv g_0\pmod d$，第二の場合は $R_i+s_i\equiv0\pmod d$；
補題~\ref{lem:app-h-borrow} で一文字借りるときは，それぞれ $s'_i\ge1$，$s_i\ge1$．
$\hat z_i:=\mathtt a^{s'_i}\omega_{X_i}\mathtt a^{s_i}$ とおく．
補題~\ref{lem:app-h-borrow} から $\hat z_i\sim\mathtt a^{s'_i+L_i}\omega\mathtt a^{R_i+s_i}$ である．
$R_i\ge g_0+n_Y$ なので，第一の場合は $k\ge0$ として $R_i+s_i=g_0+kd$，第二の場合は $k\ge0$ として $R_i+s_i=d+kd$ である；
したがって $e_i:=s'_i+L_i+kd$ について $\hat z_i\sim\mathtt a^{e_i}z_0$ であり，$(\hat z_i,e_i)$ は注釈付きの符号で，
どちらの場合も（$n_{X_i}=L_i+n_Y+R_i$ と，それぞれ $d-g_0=n_Y$，$n_Y=d$ を使って）その価値は
\[
 v_i=n_{X_i}+s_i+s'_i,\qquad|\hat z_i|\le\theta+2d
\]
である．$i<i_0$ については $v_i:=n$ とおく（符号はない）．
$X_i$ は二つの種を含む生成語をもつ $X_{i^*}$ の祖先なので，$i\le i^*$ について $n_{X_i}\ge d+n_Y$ である．
したがって，すべての $i\le i^*$ について
\begin{equation}
\label{eq:app-h-star}
 v_i\ge n_{X_i}>d,\qquad v_i\le2v_{i+1}+2d\ \ (i<i^*)
\end{equation}
であり，$v_{i^*}\le B_0+2d$ である；さらに $n\le2n_{X_{i_0}}\le2v_{i_0}$ である．

$A$ の生成語 $\mathtt a^{x_1}\omega\mathtt a^{y_1}$ は，$h:=y_1-g_0$ あるいは $h:=y_1-d$ として $\mathtt a^{x_1}z_0\mathtt a^{h}$ に等しい．
$0\le\rho,\rho'<d$，$d\mid e^*$ として $h=qd+\rho$，$x_1+qd=e^*+\rho'$ と書くと，
$l_\circ:=l\mathtt a^{\rho'}$，$r_\circ:=\mathtt a^\rho r$ について $l_\circ\mathtt a^{e^*}z_0r_\circ\in L$ を得，
$V^*:=e^*+c_0=n-\rho-\rho'$ である．

\subsection{周期性と所属}

\begin{lemma}
\label{lem:app-h-period}
$P_0:=\{p\ge0:\mathtt a^pz_0\sim z_0\}$ とおく．
\begin{enumerate}[label=\textup{(\roman*)},leftmargin=*]
\item ある $p_0\ge1$ について $P_0=\{0\}$ または $P_0=p_0\mathbb Z_{\ge0}$ である．
\item ある $f$ について $l'\mathtt a^fz_0r'\in L$ ならば，
$\{f\ge0:l'\mathtt a^fz_0r'\in L\}$ は，$P_0=\{0\}$ のとき一点であり，そうでないとき $p_0$ を法とする剰余類全体である．
\item $P_0\ne\{0\}$ ならば $\mathtt a\sim\mathtt a^{1+p_0}$ であり，したがってすべての $\mathtt a^k$ は $1\le k'\le p_0$ をみたすある $\mathtt a^{k'}$ と $\equiv_{L,s}$ で同値である．
\end{enumerate}
\end{lemma}

\begin{proof}
語 $\mathtt a^fz_0$ はすべて $\mathtt a$ 混合なので，(I2) から一つの型をもつ．
(i) $P_0$ は和について閉じている．
$P_0$ の $q>p>0$ について，差 $q-p$ についても閉じている：$\D(\mathtt a^pz_0)=\D(z_0)\ne\emptyset$ なので，
$\mathtt a^p(\mathtt a^{q-p}z_0)\sim\mathtt a^pz_0$ に補題~\ref{lem:app-h-cancel} を適用すればよい；$p=0$ の場合は自明である．
(ii) 二つの元 $f<f'$ から $\mathtt a^fz_0\sim\mathtt a^{f'}z_0$，したがって補題~\ref{lem:app-h-cancel} から $f'-f\in P_0$ を得る；
逆に，この集合は $P_0$ の元を足すことについて閉じ，$\mathtt a^{p_0}z_0\sim z_0$ から $p_0$ を引くことについても閉じる．
(iii) $(l',r')\in\D(z_0)$ について $l'\mathtt a^{p_0}z_0r'$ と $l'\mathtt a^{2p_0}z_0r'$ はともに $L$ に入るので，
(I1) から一つの型をもつ純な語 $\mathtt a$ と $\mathtt a^{1+p_0}$ は文脈 $(l'\mathtt a^{p_0-1},z_0r')$ を共有する．
\end{proof}

$D:=\operatorname{lcm}(d,p_0)$ とし，$P_0=\{0\}$ ならば $D:=\infty$ とする．
補題~\ref{lem:app-h-collapse} と補題~\ref{lem:app-h-period}(ii) から，符号を含む上のような指定された積 $u$ すべてについて
\begin{equation}
\label{eq:app-h-member}
 l_\circ ur_\circ\in L\iff\operatorname{val}(u)\equiv V^*\pmod D
\end{equation}
である（$D=\infty$ のときは等号）；価値と $V^*$ は $d$ の倍数である．
同じ語の二つの注釈は $D$ の倍数だけ異なるので，所属は注釈の選び方によらない．

\subsection{数え上げ}

\begin{lemma}
\label{lem:app-h-count}
$(z_1,e_1),\ldots,(z_M,e_M)$（$M\ge6$）は注釈付きの符号で，$|z_j|\le\ell_z$，その価値 $w_j$ の部分和は相異なり，
$\Sigma_w:=\sum_jw_j<D$ とし，$Q$ は $\Sigma_w+\operatorname{val}(Q)=V^*$ をみたす符号と塊の固定した指定された積とする．
$m:=|l_\circ|+|r_\circ|+|Q|+\sum_j|z_j|$ とおく．このとき
\[
 (27/4)^{\lfloor M/3\rfloor}\le(M+1)^2k_V(m+1)^2(M\ell_z+1)^2.
\]
\end{lemma}

\begin{proof}
置換 $\pi$ について，語 $w_\pi:=l_\circ z_{\pi(1)}\cdots z_{\pi(M)}Qr_\circ$ は \eqref{eq:app-h-member} から $L$ に入る．
$G'$ での $w_\pi$ の導出木をとり，単位節点を縮め，各節点に「符号 $z_{\pi(1)},\ldots,z_{\pi(M)}$ のうち最初の文字がその生成語の中にあるものの個数」という重みを与える．
葉の重みは $\le1$，根の重みは $M$ である；重い方の子へ下りると，重み $\kappa\in(M/3,2M/3]$ の節点 $u$ に着く．そのラベルを $X$，生成語の区間を $[p,q)$ とする．
$z_j$ が $[p,q)$ の内部にあるような添字 $j$ の集合を $I(\pi)$ とする；$\kappa-1\le|I(\pi)|\le\kappa$ である．
$\pi$ の\emph{記録}は，$X$，$p$，$q$，および $[p,q)$ の端のうち $M$ 個の符号のどれかの真の内部にあるものそれぞれについての，その符号の添字とずれからなる；
記録は高々 $k_V(m+1)^2(M\ell_z+1)^2$ 通りである．
三つの塊 $l_\circ$，$z_{\pi(1)}\cdots z_{\pi(M)}$，$Qr_\circ$ は，すべての $\pi$ で同じ位置を占める．
最初の文字が $[p,q)$ の中にある符号は，$I(\pi)$ の符号と，内部で始まり $q$ の後で終わる高々一つの符号である；$p$ を通る符号は $p$ の前で始まる．
これが $\kappa-1\le|I(\pi)|$ を与える．$[p,q)$ の各端は，$l_\circ$ の中，$Qr_\circ$ の中，二つの符号の境目，あるいは符号の真の内部にある；
最後の場合にだけそこで符号が切れ，そのときその添字とずれは記録に含まれる．

$\pi,\pi'$ が同じ記録をもつとする．
$w_\pi$ の木の $u$ の部分木を $w_{\pi'}$ の対応する部分木で置き換えると，
$w_\pi[0,p)\,w_{\pi'}[p,q)\,w_\pi[q,|w_\pi|)\in L$ の導出が得られる．
符号の内部にある端は両側で同じ添字とずれをもつので，断片は同じ符号の完全な写しにつながり，
この語は $l_\circ u''Qr_\circ$ となる．ここで $u''$ は完全な符号の積で，$z_j$ はその中に
$1-\mathbf 1[j\in I(\pi)]+\mathbf 1[j\in I(\pi')]$ 回現れる；
これは端が $l_\circ$ や $Qr_\circ$ の中にある場合も含み，$\kappa\ge2$ なので両端が $Qr_\circ$ の中にあることはない．
\eqref{eq:app-h-member} から $\sum_{I(\pi')}w_j\equiv\sum_{I(\pi)}w_j\pmod D$ である；
両辺の和は $[0,\Sigma_w)\subseteq[0,D)$ にあるので等しく，$I(\pi')=I(\pi)$ である．

したがって記録と $I=I(\pi)$（$|I|=j$）は，内側の順序（$\le j!$ 通り）とこの塊の周りの残りの符号の並べ方（$\le(M-j+1)!$ 通り）を除いて $\pi$ を定める．
$j\in\{\kappa-1,\kappa\}$ かつ $M/3<\kappa\le2M/3$ なので $j$ は $J:=[\lfloor M/3\rfloor,\lfloor2M/3\rfloor]$ に入り，
ある $j\in J$ について $\binom{M+1}{j}\le(M+1)k_V(m+1)^2(M\ell_z+1)^2$ である．
$k:=\lfloor M/3\rfloor$ とおく．$\binom{M+1}{\cdot}$ は対称で単峰であり $\lfloor M/3\rfloor+\lfloor2M/3\rfloor\le M$ なので，
すべての $j\in J$ は $\binom{M+1}{j}\ge\binom{M+1}{k}\ge\binom{3k}{k}$ をみたす．
最後に，$0\le i\le3k$ について $t_i:=\binom{3k}{i}3^{-i}(2/3)^{3k-i}$ の和は $1$ であり，
$t_{i+1}/t_i=(3k-i)/(2i+2)$ が $1$ 以上となるのは $i\le k-1$ のとき，かつそのときに限る；
したがって $t_k$ はそれらの最大であり $t_k\ge1/(3k+1)\ge1/(M+1)$ である．
$t_k=\binom{3k}{k}4^k/27^k$ なので，これは $\binom{3k}{k}\ge(27/4)^k/(M+1)$ を与え，補題が従う．
\end{proof}

\subsection{選択と埋め草}

$\Lambda:=\min(V^*,D)$ とおく．したがって $\Lambda\le V^*\le n$ である．

\begin{lemma}[選択]
\label{lem:app-h-select}
値 $w_t:=v_{j_t}$ が超増加（$w_t>w_1+\cdots+w_{t-1}$）であり，$w_1+\cdots+w_M<\Lambda$ かつ $\Lambda\le3^M(v_{i^*}+d)$ をみたす
$M\ge0$ 個の添字 $i^*=j_1>j_2>\cdots>j_M\ge i_0$ がある．
\end{lemma}

\begin{proof}
$v_{i^*}\ge\Lambda$ ならば $M:=0$ ととる．そうでなければ $j_1:=i^*$ ととる．
$j_t$ と $S_t:=w_1+\cdots+w_t<\Lambda$ が与えられたとき，$j_t$ より下で $v_i>S_t$ をみたす最大の添字を $i$ とする；
$v_0=n\ge\Lambda$ なのでそれは存在する（$i_0\ge1$ を思い出す）．
$S_t+v_i\ge\Lambda$ ならば $M:=t$ として止める；そうでなければ $j_{t+1}:=i$ とおく．
値は超増加であり，その和は作り方から $\Lambda$ 未満である．
とられる値は $\Lambda\le n$ 未満であるが，$i<i_0$ については $v_i=n$ である；したがってとられる添字はすべて $i_0$ 以上である．
$i$ の選び方から $v_{i+1}\le S_t$ である：$i+1<j_t$ ならば最大性から，$i+1=j_t$ ならば $v_{j_t}=w_t$ だからである．
\eqref{eq:app-h-star} から $v_i\le2S_t+2d$ である．
したがって $i$ がとられるときは $S_{t+1}+d\le3(S_t+d)$ であり，過程が止まるときは $\Lambda\le S_M+v_i<3(S_M+d)$ である．
$S_1=v_{i^*}$ から $\Lambda<3^M(v_{i^*}+d)$ を得る．
\end{proof}

\begin{lemma}[埋め草]
\label{lem:app-h-filler}
$0<X<n$ を $d$ の倍数とする．
$\hat z_{i_0},\ldots,\hat z_{i^*}$ のうちの高々 $3k_V$ 個の符号と，$d\mid X'$，$X'<v_{i^*}$ をみたす高々一つの純な塊 $\mathtt a^{X'}$ からなる指定された積 $Q$ で，
$\operatorname{val}(Q)=X$ かつ $|Q|\le(3k_V+2)\ell_z$ をみたすものがある．ここで $\ell_z:=5k_V^2\theta$ である．
\end{lemma}

\begin{proof}
$i=i_0,\ldots,i^*$ について順に，$\hat z_i$ を $q_i:=\lfloor X/v_i\rfloor$ 個加えて $X$ を $X-q_iv_i$ で置き換える．
$X<n\le2v_{i_0}$ なので $q_{i_0}\le1$ である．
$i>i_0$ については現在の $X$ は $v_{i-1}$ 未満であり，\eqref{eq:app-h-star} から $v_{i-1}\le2v_i+2d<4v_i$ なので $q_i\le3$ である．
最後の余り $X'<v_{i^*}$ は，すべての価値が $d$ の倍数なので $d$ の倍数であり，塊として加える．
$i^*-i_0\le k_V-3$ なので，符号は高々 $1+3(k_V-3)<3k_V$ 個で，それぞれ長さ $\le\theta+2d\le\ell_z$ であり，
$X'<v_{i^*}\le B_0+2d\le(6k_V^2+4k_V)\theta\le2\ell_z$ である．
\end{proof}

\subsection{定理~\ref{thm:two-letter-fanout-one} の証明}

$A$ が純な生成語をもつとする．(I1) から型 $s(\mathtt a)$ の $A$ の生成語はその純な生成語なので $\tau(A,s(\mathtt a))\le n$ であり，
$\tau(A,s(\mathtt a))\le11\cdot3^{M^*}k_V^2\theta$ を示せば十分である．ここで
\[
 M^*:=\max\bigl(\{5\}\cup\{M\ge6:(27/4)^{\lfloor M/3\rfloor}
 \le(M+1)^4k_V(3k_V+M+4)^2\ell_z^{\,4}\}\bigr)
\]
であり，左辺は $M$ について指数的に増加するのでこれは有限である．
根の部分木の種が高々 $k_V$ 個か，$n_{X_{i^*}}=n$ ならば $n\le2k_V(k_V+2)\theta$ である．
そうでなければ補題~\ref{lem:app-h-select} の添字をとる．

\emph{数え上げ．}
$M\ge6$ とする．$X:=V^*-\sum_tw_t$ とおく；$\sum_tw_t<\Lambda\le V^*\le n$ かつ $w_1>0$ なので $0<X<n$ であり，$X$ は $d$ の倍数である．
$Q$ を補題~\ref{lem:app-h-filler} で与えられるものとする．
符号 $\hat z_{j_1},\ldots,\hat z_{j_M}$ は長さ $\le\ell_z$ で価値が超増加なので，その部分和は相異なり，$\sum_tw_t<\Lambda\le D$ である．
$|l_\circ|+|r_\circ|\le\theta+2d\le\ell_z$ から $m\le(3k_V+M+3)\ell_z$ を得，
$(m+1)^2(M\ell_z+1)^2\le(3k_V+M+4)^2(M+1)^2\ell_z^{\,4}$ とともに補題~\ref{lem:app-h-count} は $M\le M^*$ を与える．
したがってどの場合も $M\le M^*$ であり，補題~\ref{lem:app-h-select} から
\[
 \Lambda\le3^{M^*}(v_{i^*}+d)\le3^{M^*}(B_0+3d)\le10\cdot3^{M^*}k_V^2\theta
\]
である．

\emph{結び．}
$\Lambda=V^*$ ならば $n=V^*+\rho+\rho'<\Lambda+2d$ であり $\tau(A,s(\mathtt a))\le n$ である．
$\Lambda=D<V^*$ ならば $P_0\ne\{0\}$ であり，補題~\ref{lem:app-h-period}(iii) から $n'\le p_0\le D$ について $\mathtt a^n\equiv_{L,s}\mathtt a^{n'}$ となるので，
$\tau(A,s(\mathtt a))\le n'$ である．
小さい場合と合わせて，$\tau(A,s(\mathtt a))\le\max\{2k_V(k_V+2)\theta,\,10\cdot3^{M^*}k_V^2\theta+4k_V^2\theta\}\le11\cdot3^{M^*}k_V^2\theta$ である．

\emph{指数．}
$\mu:=(27/4)^{1/3}$，$\gamma:=\ln3/\ln\mu<1.7260$，$\eta:=1/80$，$C_\eta:=\sup_{M\ge0}(M+1)^6\mu^{-\eta M}<e^{34}$ とおく
（$x>0$ についての $6\ln x-(x-1)\eta\ln\mu$ の最大は $x=6/(\eta\ln\mu)$ でとられる）．
$M^*$ を定める集合の $M\ge6$ について，$\mu^{M-2}\le(27/4)^{\lfloor M/3\rfloor}$ と $3k_V+M+4\le5k_V(M+1)$ から
\[
 \mu^M\le15625\,\mu^2(M+1)^6k_V^{11}\theta^4,\qquad\text{したがって}\qquad
 \mu^{(1-\eta)M}\le15625\,\mu^2C_\eta k_V^{11}\theta^4
\]
を得，$3^M=\mu^{\gamma M}\le(15625\,\mu^2C_\eta k_V^{11}\theta^4)^{\gamma/(1-\eta)}$ である；これは $M=5$ についても成り立つ．
$\gamma/(1-\eta)=80\gamma/79<1.7479$ なので
\[
 \tau(A,s(\mathtt a))\le11\cdot3^{M^*}k_V^2\theta\le C_0k_V^{2+19.23}\theta^{1+6.992}
 \le C_0k_V^{22}\theta^8,\qquad C_0:=11(15625\,\mu^2C_\eta)^{80\gamma/79}<10^{36}
\]
を得る．
$k_V\le2\|G\|$，$\theta\le\bar\theta_G$ であり，$G$ の非終端記号は同じ生成語をもつ $G'$ の非終端記号なので，
$C:=2^{22}C_0<10^{43}$ として定理が従う．$\eta\to0$ とすると，指数は $2+11\gamma<20.99$ と $1+4\gamma<7.91$ にいくらでも近くとれる．\qed

\subsection{他の型と系~\ref{cor:two-letter-fanout-one} の証明}

\begin{lemma}
\label{lem:app-h-other}
$G'$ を上のとおりとし，$A$ を $G'$ の非終端記号とする．
\begin{enumerate}[label=\textup{(\roman*)},leftmargin=*]
\item $s:\Sigma^+\twoheadrightarrow U_1$ で $\mathtt a$ が型 $1$ の唯一の文字であり，$A$ が型 $\zeta$ の生成語をもてば，
$\tau(A,\zeta)\le(k_V-1)\theta+1$ である．
\item $\Sigma=\{\mathtt a,\mathtt b\}$，$s=c$ で，$A$ が内容 $\Sigma$ の生成語をもてば，$\tau(A,\Sigma)\le3k_V\theta+2$ である．
\end{enumerate}
\end{lemma}

\begin{proof}
(i) 型 $\zeta$ の $A$ の生成語の導出で節点の数が最小のものをとり，文字が $\mathtt a$ でない葉に印をつける．
根から印をつけた葉への道の上の節点の部分木を，同じ道の上の同じラベルのより低い節点の部分木で置き換えても，印をつけた葉は残り，したがって型も残る；
よってこの道でラベルは繰り返さず，道は高々 $k_V$ 個の非終端記号の節点をもつ．
そこからぶら下がる高々 $k_V-1$ 個の部分木を長さ $\le\theta$ の生成語で置き換えると，型 $\zeta$ で長さ $\le(k_V-1)\theta+1$ の $A$ の生成語が得られる．
(ii) 内容 $\Sigma$ の生成語の導出で節点の数が最小のものをとり，$\mathtt a$ の葉と $\mathtt b$ の葉に印をつける．
根，分岐点，二つの葉の間の三つの道の区間のそれぞれの中で（異なる区間の間ではなく）同じ縮約を行うと，骨格の節点は高々 $3k_V$ 個である；
そこからぶら下がる高々 $3k_V$ 個の部分木を長さ $\le\theta$ の生成語で置き換えると，長さ $\le3k_V\theta+2$ の内容 $\Sigma$ の生成語が得られる．
どちらの場合も新しい語は同じ型の $A$ の生成語であり，したがって同じ類に入る（補題~\ref{lem:typed-homogeneity}）．
\end{proof}

\begin{proof}[系~\ref{cor:two-letter-fanout-one} の証明]
どちらの場合も $G'$ は $L\cap\Sigma^+$ のファンアウト1，ランク $2\le r_{\max}$ の reduced good な表示で $\bar\theta_{G'}\le\bar\theta_G$ であり，
系~\ref{cor:distribution-representatives} を $G'$ に適用すると
$\|K\|_+\le N_{\mathrm{typ}}(G',s)(1+2\bar\theta_G+r_{\max}\tau(G',s))+\mathbf 1_{\{\eps\in L\}}$ を得る．
残るのは $\tau(G',s)$ と $N_{\mathrm{typ}}(G',s)$ の評価である．

(i) $s$ の正の像は半束の中で $s(\mathtt a)$ と $s(\mathtt b)$ によって生成されるので，$s(w)$ は $c(w)$ のみによる；
したがって $s\preceq c$，$\Cmcf{1}{s}\subseteq\Cmcf{1}{c}$ であり，$\equiv_{L,c}$ は $\equiv_{L,s}$ を細分する．
$G'$ の型付き非終端記号 $A_t$ と，その生成語の一つの $c$ の型 $t'$ について，$\tau_{A_t}(G',s)\le\tau(A,t')$ である．
両方の文字が $c$ について孤立しているので，$\mathtt a$ と $\mathtt b$ のそれぞれについて $c$ に適用した定理~\ref{thm:two-letter-fanout-one} と，
補題~\ref{lem:app-h-other}(ii) から $\tau(G',s)\le C(\|G\|+2)^{22}\bar\theta_G^{\,8}$ である．
$G'$ の型付きの細分は高々 $3k_V$ 個の型付き非終端記号と，規則一つあたり高々9個の型付き規則をもつので，
$N_{\mathrm{typ}}(G',s)\le3k_V+9|P'|\le12(\|G\|+2)$ である．

(ii) $G'$ の型付き非終端記号は，$A$ が $\mathtt a^+$ の生成語をもつ $A_1$ と，$A_\zeta$ の形である．
定理~\ref{thm:two-letter-fanout-one} と補題~\ref{lem:app-h-other}(i) から $\tau(G',s)\le C(\|G\|+2)^{22}\bar\theta_G^{\,8}$ である．
型付き非終端記号は高々 $2k_V$ 個，型付き規則は規則一つあたり高々4個なので
$N_{\mathrm{typ}}(G',s)\le2k_V+4|P'|\le6(\|G\|+\sigma_G)\le12\|G\|$ である．

どちらの場合も $1+2\bar\theta_G\le3\bar\theta_G$ から，$C':=12(C+2)<10^{45}$ として評価を得る．
\end{proof}

\section{ファンアウト1での \texorpdfstring{$U_1$}{U1} の上への型付け}
\label{app:u1-typings}

\subsection{有限の代入可能な言語の文字が互いに素な像}

\begin{lemma}\label{lem:app-u1-insertion}
$L\subseteq\Delta^+$ は有限で Clark--Eyraud 代入可能とする．$PQ\in L$ かつ $PTQ\in L$ となる
$P,Q\in\Delta^*$ と $T\in\Delta^+$ は存在しない．
\end{lemma}

\begin{proof}
そうでないとする．$P\ne\eps$ ならば，語 $P$ と $PT$ は文脈 $(\eps,Q)$ を共有するので
$\D_L(P)=\D_L(PT)$ である；すると $PT^jQ\in L$ から $PT^{j+1}Q\in L$ が従い，
$L$ はあらゆる長さの語 $PT^jQ$ を含み，有限性に反する．$P=\eps$ ならば，
文脈 $(\eps,\eps)$ を共有する $Q$ と $TQ$ について同様に論じる．
\end{proof}

\begin{lemma}\label{lem:app-u1-code}
$h:\Delta^*\to\Sigma^*$ を，すべての $h(\delta)$（$\delta\in\Delta$）が空でなく，
すべての $h(\delta)$ の連接の中で同じ文字が二度現れないような準同型とする．
$L\subseteq\Delta^+$ が有限で Clark--Eyraud 代入可能ならば，$h(L)$ もそうである．
\end{lemma}

\begin{proof}
ある $h(\delta)$ に現れる各文字は，$\delta$ と $h(\delta)$ の中でのその位置を定める．
したがって $h$ は単射であり，語 $h(D)$ の各文字の出現は，符号語 $h(\delta)$ の定まった出現に属する；
隣り合う二つの文字が一つの符号語の出現の中にあるか，ある文字が符号語の出現を始めるか終えるかは，文字そのものから読み取れる．

$x,y\in\Sigma^+$，$x\ne y$ で，$h(L)$ の中で $pxq=h(D_1)$，$pyq=h(D_2)$ とする．
$p$ の内部にある符号語の出現と $q$ の内部にあるものは，両方の語で同じである；
したがって $D_i=PM_iQ$ と書ける．ここで $M_i$ はそれぞれ $x$，$y$ と交わる符号語の出現からなる．
$p$ が符号語の出現 $h(\delta)$ の内部で終われば $M_1$ と $M_2$ はともに $\delta$ で始まり，
$q$ が符号語の出現 $h(\delta')$ の内部で始まれば両者は $\delta'$ で終わる．

$x$ が一つの符号語の出現 $h(\delta)$ の内部にあり，その全体ではないとする；例えば $p$ がその内部で終わるとする（他の場合は対称である）．
このとき $M_1=\delta$ であり，$M_2$ は $\delta$ で始まる；そうでなければ $y=x$ となるので $M_2\ne M_1$ であり，
$T\in\Delta^+$ として $M_2=\delta T$ である．$D_1=P\delta Q$，$D_2=P\delta TQ$ は補題~\ref{lem:app-u1-insertion} に反する．
$x$ と $y$ を入れ替えても同じである．したがって，どちらも符号語全体でない限り一つの符号語の出現の内部にはない．
このとき $M_1=[\delta]X[\delta']$，$M_2=[\delta]Y[\delta']$ である．ここで $X,Y\in\Delta^*$ はそれぞれ $x$，$y$ の内部にある符号語であり，
括弧の文字は対応する切れ目が符号語の内部にあるとき，かつそのときに限り現れる．
$X=\eps\ne Y$ または $Y=\eps\ne X$ ならば，再び補題~\ref{lem:app-u1-insertion} に反し，$X=Y=\eps$ ならば $x=y$ となる．
したがって $X,Y\in\Delta^+$ であり，
\[
 x=\alpha\,h(X)\,\beta,\qquad y=\alpha\,h(Y)\,\beta
\]
である．ここで $\alpha$ は $p$ の後の $h(\delta)$ の部分（または $\eps$），$\beta$ は $q$ の前の $h(\delta')$ の部分（または $\eps$）である．
語 $X$ と $Y$ は $L$ の中で文脈を共有するので $\D_L(X)=\D_L(Y)$ である．

次に $p'xq'\in h(L)$ とする．$\alpha$ と $h(X)$ の間の切れ目，$h(X)$ と $\beta$ の間の切れ目は符号語の出現の境界であり，
それは $x$ の文字から読み取れる；したがってある $P',Q'$ について $p'\alpha=h(P')$，$\beta q'=h(Q')$ であり，$P'XQ'\in L$ である．
すると $P'YQ'\in L$ であり，$h(P'YQ')=p'yq'$ である．対称性から $\D_{h(L)}(x)=\D_{h(L)}(y)$ である．
\end{proof}


\subsection{二文字の型付けについての2次のデータ}

命題~\ref{prop:binary-u1-quadratic} を証明する．
$k:=2^m$，$M:=2^n$，$\psi(w):=|w|_{\mathtt b}+|w|_{\mathtt a}/k$ とおく．$\psi$ は加法的である．
語 $w$ について，$\lambda(w)$ と $\varrho(w)$ を $\mathtt a^*$ に入るその最長の接尾辞と最長の接頭辞の長さとする．

\begin{proof}[命題~\ref{prop:binary-u1-quadratic} の証明]
(i) $\Sigma^+$ の語が $L_{m,n}$ に入るのは，その $\mathtt a$ の極大の連の長さがすべて $k$ で割り切れ，$\psi(w)=M$ であるとき，かつそのときに限る：
そのような語は塊 $\mathtt a^k$ と $\mathtt b$ に一意に分解され，$\psi$ は塊を数える．
$x,y\in\Sigma^+$ が同じ型をもち $uxv,uyv\in L_{m,n}$ とする．加法性から $\psi(x)=\psi(y)$ である．
$x,y$ が純ならば，これは $x=y$ を与える．
そうでなければ，$x'$ を $\mathtt b$ で始まり $\mathtt b$ で終わる語として $x=\mathtt a^{\alpha}x'\mathtt a^{\beta}$ と書き，
同様に $y=\mathtt a^{\alpha'}y'\mathtt a^{\beta'}$ と書く．
$x'$ と $y'$ の内部の $\mathtt a$ の極大の連は $uxv$ と $uyv$ の極大の連なので，その長さは $k$ で割り切れる；
$x$ の両端を通る連の長さは $\lambda(u)+\alpha$ と $\beta+\varrho(v)$ なので，$\alpha\equiv\alpha'$，$\beta\equiv\beta'\pmod k$ である．
任意の文脈 $(p,q)$ について，$pxq$ の $\mathtt a$ の極大の連は，$p$ の末尾にないもの，長さ $\lambda(p)+\alpha$ の連，
$x'$ の内部の連，長さ $\beta+\varrho(q)$ の連，$q$ の先頭にないもの，である；
$pyq$ についても $\alpha',\beta',y'$ で同じことが成り立つ．
したがって割り切れる条件は一致し，$\psi(pxq)=\psi(pyq)$ であり，$pxq\in L_{m,n}\iff pyq\in L_{m,n}$ である．
有限言語は文脈自由である．

(ii) 非終端記号 $C_0,\ldots,C_m$，$A_0,\ldots,A_n$，開始記号 $A_n$ と，
規則 $C_0\to\mathtt a$，$C_j\to C_{j-1}C_{j-1}$（$1\le j\le m$），
$A_0\to C_m\mid\mathtt b$，$A_i\to A_{i-1}A_{i-1}$（$1\le i\le n$）をとる．
このとき $L(C_j)=\{\mathtt a^{2^j}\}$，$L(A_0)=\{\mathtt a^k,\mathtt b\}$ であり，$L(A_i)$ は
$L(A_0)$ の語 $2^i$ 個の連接の集合なので $L(A_n)=L_{m,n}$ である；文法は good かつ reduced である．
非終端記号は $m+n+2$ 個で，規則の大きさは $2$，$5$（$m$ 個），$3$，$2$，$5$（$n$ 個）である（定義~\ref{def:standard-mcfg}）．
これは $6m+6n+9$ を与える．
ファンアウト1では，規則 $R:B\to\rho(B_1,\ldots,B_q)$ を使う最短の文の長さは
$\theta_G(R)=o_B+c_R+\sum_jd_{B_j}$ である．ここで $d_C$ は $C$ の最短の生成語の長さ，
$o_C$ は文の中での $C$ の出現の周りの文字の最小の個数，$c_R$ は $R$ が書く終端記号の個数である；
ここでは $d_{A_0}=1$，$d_{A_i}=2^i$，$d_{C_j}=2^j$，$o_{A_i}=M-2^i$，$o_{A_0}=o_{C_m}=M-1$，
$o_{C_j}=M-1+k-2^j$ である．
これは $A_i$ の規則と $A_0\to\mathtt b$ について $\theta_G(R)=M$ を，$A_0\to C_m$ と $C_j$ の規則について
$\theta_G(R)=M+k-1$ を与える；したがって $\bar\theta_{G_{m,n}}=2^m+2^n-1$ である．

(iii) $L_{m,n}$ の唯一の純な語は $\mathtt a^{kM}$ なので，命題~\ref{prop:rare-type} からすべての特徴標本はそれを含み，
$\|F\|_+\ge kM+1>2^{m+n}$ である．$m=n$ のとき，$\bar\theta_{G_{n,n}}<2^{n+1}$ から
$\bar\theta_{G_{n,n}}^{\,2}/4<4^n<\|F\|_+$ である．評価 $p(\|G\|)\bar\theta_G$ はすべての $n$ について
$4^n\le p(12n+9)\,2^{n+1}$ を与えるが，これはどの多項式 $p$ についても成り立たない．
\end{proof}

\section{ファンアウト2以上での二文字の内容型付け}
\label{app:two-letter-content}

この付録は定理~\ref{thm:two-letter-content} を証明する．
以下では $\Sigma=\{\mathtt a,\mathtt b\}$，$c:=c_\Sigma$，$f\ge2$ を固定し，$L\in\Cmcf{f}{c}$ とし，
$G=(V,\Sigma,P,S,\fan)$ を $L(G)=L\cap\Sigma^+$ をみたすファンアウト $\le f$，ランク $\le r:=r_{\max}$ の reduced good な文法とする；
$v:=|V|$，$\theta:=\bar\theta_G$，$D_f:=(6f+1)3^f$ とおく．
考えている型付き非終端記号のアリティの条件を使う補題~\ref{lem:app-two-basic}(i) を除き，
議論は $(1,c)$- と $(2,c)$-タプル代入可能性のみを使う．これらは $f\ge2$ なので成り立つ．
各 $A\in V$ について文脈 $E_A$ は補題~\ref{lem:thickness-equivalence}(iv) のものとする．したがって $|E_A|_{\mathrm{fix}}\le\theta$
かつ $E_A[L_A(G)]\subseteq L(G)$ である；各 $B\in V$ について $\|\omega_B\|_1\le\theta$ をみたす $\omega_B\in L_B(G)$ を固定し（項目~(ii)），
すべての規則が高々 $\theta$ 個の終端記号を書くことを思い出す（項目~(i)）．
導出木の中で，節点 $u$ の\emph{生成語}はそれが導くタプル，$\ell_u$ はその総長，$B_u$ はその非終端記号であり，
$\ell_u>\theta$ のとき $u$ は\emph{重い}という；節点はその雛形の終端記号を\emph{書く}．
文字 $x,y$ は常に $\Sigma$ の二つの文字をある順序で表す．
語の\emph{連}は極大の因子 $x^k$（$k\ge1$）である；最初の連でも最後の連でもない連を\emph{内部の}連という．
したがって内部の $y$-連の前後には $x$ がある．
\emph{$x$-骨格} $\sigma_x(w)$ は $w$ の $x$-連をそれぞれ一つの $x$ で置き換えたものであり，
$|w|_x$ が $w$ の $x$-連の個数を超えるとき，$w$ は\emph{余分な} $x$ をもつという．
$L$ のある語が因子 $x^k$ をもつとき \emph{$x^k$ が現れる}といい，
同じ Parikh ベクトル $\Psi(u)=\Psi(w)$ をもつ $u,w\in\Sigma^+$ について $u\in L$ と $w\in L$ が同値であるとき，
$L$ は \emph{Parikh 飽和}であるという．

\paragraph{証明の地図．}
表は，各段階について，それが使うものと，それが使われる場所を示す．
\begin{center}
\footnotesize
\setlength{\tabcolsep}{3pt}
\renewcommand{\arraystretch}{1.15}
\begin{tabularx}{\linewidth}{@{}>{\raggedright\arraybackslash}p{0.24\linewidth}>{\raggedright\arraybackslash}X>{\raggedright\arraybackslash}p{0.26\linewidth}@{}}
\hline
段階 & 使うもの & 使われる場所\\
\hline
補題~\ref{lem:app-two-basic}(i) 型付きの閉包
 & 補題~\ref{lem:typed-homogeneity}，アリティ $d$ の代入可能性
 & 補題~\ref{lem:app-two-saturated}\\
補題~\ref{lem:app-two-basic}(ii) 恒等の道の切断
 & $G$ が good であること
 & (P2)，補題~\ref{lem:app-two-saturated} の混合の場合\\
補題~\ref{lem:app-two-runs} 連の交換
 & $(1,c)$，$(2,c)$-代入可能性，$x^4$ が現れる
 & 補題~\ref{lem:app-two-alternate}，命題~\ref{prop:app-two-nonsaturated}\\
補題~\ref{lem:app-two-alternate}
 & 補題~\ref{lem:app-two-runs}(i),(iii)
 & 命題~\ref{prop:app-two-nonsaturated}\\
補題~\ref{lem:app-two-insertion} 置換が作る連
 & 補題~\ref{lem:witness-length-template}(i)，節点数最小の最短の $T$
 & 命題~\ref{prop:app-two-nonsaturated}\\
命題~\ref{prop:app-two-nonsaturated}（飽和でない）
 & 上の三つの補題；道の上の型付きのラベル
 & 定理~\ref{thm:two-letter-content}\\
補題~\ref{lem:app-two-saturated}，和集合が $\Sigma$
 & 補題~\ref{lem:app-two-basic}；道に沿った床関数
 & 定理~\ref{thm:two-letter-content}\\
(P1) 行の剛性
 & $(1,c)$-代入可能性，$K$ の最小性
 & (P3)，(P4)\\
(P2) 高さ，混合の節点
 & 補題~\ref{lem:app-two-basic}(ii)，$K$ の最小性
 & (P3)，(P6)\\
(P3) 重い純な節点
 & (P1)，(P2)
 & (P4)，(P5)\\
(P4) 共通の比 $\alpha$
 & (P3)
 & (P5)，(P6)\\
(P5) $\alpha$ の評価
 & (P3)，(P4)
 & (P6)\\
(P6) 数え上げ
 & (P2)，(P4)，(P5)
 & 定理~\ref{thm:two-letter-content}\\
\hline
\end{tabularx}
\end{center}

\begin{lemma}\label{lem:app-two-basic}
\textup{(i)} アリティ $d$ のすべての型付き非終端記号 $X=A_{\mathbf t}$ について，
$\operatorname{Cl}_{L,c}(X)=\{\tuple z\in(\Sigma^+)^d:c^{(d)}(\tuple z)=\mathbf t,\
E_A[\tuple z]\in L\}$ である；したがって $\tau_X$ はそのようなタプルの最小の総長である．

\textup{(ii)} $G$ の導出木の中で，$u$ が同じ非終端記号をもつ $u'$ の真の祖先であり，
それらの生成語の総 Parikh ベクトルが等しいとする．
このとき $u$ から $u'$ への道の上のすべての規則はランク1で，終端記号を書かず，成分ごとの恒等写像である．
したがって $u$ の部分木を $u'$ の部分木で置き換えても導かれるタプルは変わらず，少なくとも一つの節点が除かれる．
\end{lemma}

\begin{proof}
(i) $\tuple u\in L_X$ を固定する．補題~\ref{lem:typed-homogeneity} から，
$\operatorname{Cl}_{L,c}(X)$ は $\tuple u$ の型と分布をもつタプルの集合であり，それらは $E_A\in\D_L(\tuple u)$ に受理される．
$E_A$ に受理される型 $\mathbf t$ のタプル $\tuple z$ はこの文脈を $\tuple u$ と共有するので，代入可能性から $\D_L(\tuple z)=\D_L(\tuple u)$ である．

(ii) 線形性と非削除性から，$u'$ の生成語のすべての文字は $u$ の生成語に現れるので，
Parikh ベクトルが等しいことは，道の上で書かれる終端記号と，道から外れる兄弟を排除する．兄弟は空でない生成語をもつからである（$\eps$ なし）．
したがって道の上のすべての規則は，あるファンアウト $d'$ の子を一つもち，その $d'$ 個の変数はそれぞれ一度（非削除性），
順に（非置換性）現れ，親の各成分は空でない（$\eps$ なし）．
一つの成分の中の二つの変数は連続することになり，それは非併合性が禁じる；
したがってすべての成分は一つの変数であり，ファンアウトは一致し，規則は恒等写像である．
\end{proof}

\begin{lemma}[連の交換]\label{lem:app-two-runs}
$x^4$ が現れるとする．
\begin{enumerate}[label=\textup{(\roman*)},leftmargin=*]
\item $w,w'\in\Sigma^+$ が $\sigma_x(w)=\sigma_x(w')$ かつ $|w|_x=|w'|_x$ をみたせば，$w\in L$ と $w'\in L$ は同値である．
\item $L$ のある語が因子 $xxyx$ または $xyxx$ をもてば，$L$ は Parikh 飽和である．
\item $L$ が Parikh 飽和でなければ，因子 $xx$ をもつ $L$ のすべての語の中で，すべての内部の $y$-連の長さは少なくとも $2$ である．
\item $L$ が Parikh 飽和でなく $y^4$ が現れなければ，因子 $xx$ をもつ $L$ のすべての語は高々一つの内部の $y$-連をもつ．
\end{enumerate}
\end{lemma}

\begin{proof}
(i) $c(x)$ は冪等なので，補題~\ref{lem:power-transfer} が $u=x$，$i_0=p=1$，$i=\delta=j=1$ で適用できる：
$x^4$ が現れるので，$w_{\mathrm M}\ne\eps$ ならば
$w_{\mathrm L}xw_{\mathrm M}xxw_{\mathrm R}\in L$ と $w_{\mathrm L}xxw_{\mathrm M}xw_{\mathrm R}\in L$ は同値である．
$R$ と $R'$ を $|R|\ge2$ である語 $w$ の $x$-連とし，例えば $R'$ が $R$ の左にあるとする（他の場合は対称である）．
最初に示した $x$ を $R'$ の最後の文字，$xx$ を $R$ の最初の二文字として $w=U\,x\,Z\,xx\,V$ と書く；
因子 $Z$ は二つの連を隔てるので $Z\ne\eps$ である．
したがって $U\,xx\,Z\,x\,V\in L$ と $w\in L$ は同値である．
これは一つの $x$ を $R$ から $R'$ へ移し，$x$-骨格と $|w|_x$ を保つ．
同じ $x$-骨格をもつ語は同じ個数 $\rho$ の $x$-連をもち，その連の長さは $|w|_x$ の $\rho$ 個の正の部分への合成である；
そのような任意の二つの合成は，大きさ少なくとも $2$ の部分から別の部分へ一単位を移す操作でつながる．

(ii) 例えば $UxxyxV\in L$ とする（他の場合は対称である）．(i) から $UxyxxV\in L$ である．
文脈 $(Ux,xV)$ はどちらも内容 $\Sigma$ の $xy$ と $yx$ を受理するので $\D_L(xy)=\D_L(yx)$ である：
すべての文脈で $xy$ を $yx$ に置き換えることができ，その逆もできる．
隣接する互換は一つの Parikh ベクトルをもつすべての語をつなぐ．

(iii) $w\in L$ が因子 $xx$ と長さ $1$ の内部の $y$-連をもつとする．
このとき $w$ は余分な $x$ をもち，(i) から余分な文字をその $y$-連の前の $x$-連へ移すことができる；
得られる $L$ の語は因子 $xxyx$ をもつので，(ii) から $L$ は Parikh 飽和である．

(iv) $w\in L$ が因子 $xx$ と二つの内部の $y$-連をもつとする．
このとき $w$ は，$x$-連 $R$ で隔てられた連続する二つの内部の $y$-連 $y^e$ と $y^{e'}$ をもつ；$A_0$ を $y^e$ の前の $x$-連とする．
(iii) から $e,e'\ge2$ である．
$w$ は余分な $x$ をもつので，(i) は $w$ の $x$-骨格と文字 $x$ の個数をもつ $L$ の二つの語を与える：
$R$ の長さが $2$ である語 $W_2=U\,x\,y^e\,xx\,V'$（$y^e$ の前に示した $x$ は $A_0$ の最後の文字であり，$V'$ は $y^{e'}$ で始まる）と，
$R$ の長さが $1$ である語 $W_1=U_1\,x\,y^e\,x\,y^{e'}V_1$ である．
文脈 $(U,xV')$ は $W_2$ を与える $xy^ex$ と，$U\,xx\,y^e\,x\,V'$ を与える $xxy^e$ を受理する．
後者は $W_2$ の $x$-骨格と $x$ の個数をもち，(i) から $L$ に入る．
どちらのタプルも内容 $\Sigma$ をもつので $\D_L(xy^ex)=\D_L(xxy^e)$ である．
文脈 $(U_1,y^{e'}V_1)$ は $W_1$ を与える $xy^ex$ を受理するので $xxy^e$ も受理し，$U_1\,xx\,y^{e+e'}V_1\in L$ である．
$e+e'\ge4$ なので，この語は因子 $y^4$ をもち，矛盾である．
\end{proof}

\begin{lemma}\label{lem:app-two-alternate}
$L$ が Parikh 飽和でなく，$\mathtt a^4$ と $\mathtt b^4$ がともに現れるならば，
少なくとも四つの連をもつ $L$ のすべての語のすべての連の長さは $1$ である．
\end{lemma}

\begin{proof}
$w\in L$ が少なくとも四つの連をもつとする．その二番目と三番目の連は内部の連で異なる文字からなるので，
各文字 $y$ について，$w$ は内部の $y$-連 $\beta$ と少なくとも二つの $y$-連をもつ．
$y$ についての補題~\ref{lem:app-two-runs}(i) により，文字 $y$ を $\beta$ から別の $y$-連へ移すと，
$\beta$ の長さが $1$ で $x$-連が $w$ のものである $L$ の語が得られる．
$w$ の $x$-連の長さが少なくとも $2$ ならば，この語は因子 $xx$ と長さ $1$ の内部の $y$-連をもち，
補題~\ref{lem:app-two-runs}(iii) に反する．
したがって $x$ のどちらの選び方についても，$w$ のすべての $x$-連の長さは $1$ である．
\end{proof}

\begin{lemma}[短い置換が作る連]\label{lem:app-two-insertion}
$X=A_{\mathbf t}$ をアリティ $d\le f$ の型付き非終端記号とし，$T$ を，$X$ からの $\widetilde G_0$ の導出で，
その生成語 $\tuple u$ の総長が最小で，その中で節点の数が最小のものとする．
$T$ が $|Q|\ge6f+2$ をみたす重い節点の反鎖 $Q$ をもてば，ある文字 $x$ について次が成り立つ：
\textup{(a)} $\tuple u$ のある成分は内容 $\{x\}$ と少なくとも $5$ の長さをもち，したがって $x^5$ が現れる；
\textup{(b)} $L$ のある語 $W$ は因子 $xx$ と少なくとも四つの連をもち，$\mathtt a^5$ と $\mathtt b^5$ がともに現れるのでなければ，
$W$ は二つの内部の $y$-連をもつ．
\end{lemma}

\begin{proof}
内容 $\{x\}$ の各成分 $u_i$（\emph{純な}成分）の最初と最後の $\min\{2,|u_i|\}$ 文字に印をつけ，
内容 $\Sigma$ の各成分の一つの $\mathtt a$ と一つの $\mathtt b$ に印をつける；印は高々 $4f$ 個である．
各印は一つの節点が書き，$Q$ の節点は互いに素な部分木をもつので，
部分木が印を書かない $Q$ の節点の集合 $Q_0$ は $|Q_0|\ge|Q|-4f\ge2f+2$ をみたす．

$h\in Q_0$ の非終端記号を $C_h$ とする．型の添字を消し，$h$ の部分木を $\omega_{C_h}$ の導出で置き換える．
新しい根のタプルは $L_A(G)$ に入り，より短い．非削除性から $h$ の生成語のすべての文字は $\tuple u$ に現れ，
$\|\omega_{C_h}\|_1\le\theta<\ell_h$ だからである．
印は内容 $\Sigma$ のすべての成分の内容を保つ．
すべての純な成分が純なままならば，新しいタプルは型 $\mathbf t$ をもち，
$\lambda:=\widehat c$ とした補題~\ref{lem:witness-length-template}(i) から $L_X$ に入るが，これは $T$ の選び方に反する．
したがって $\omega_{C_h}$ は，純な成分の中にある $h$ の塊（$h$ の生成語の成分が占める根の成分の区間）に他方の文字を入れる．
塊は印を含まないので，その純な成分の長さは少なくとも $5$ であり，塊はその最初の二文字の後，最後の二文字の前にある．
$E_A[\tuple u]\in L$ なので，これは (a) を示す．

各 $h\in Q_0$ は高々 $f$ 個の塊をもつ．
一つの根の成分の中でそれらの塊が間に文字なしに隣り合うとき，二つの所有者を結ぶ．
各塊は高々二つの隣り合う塊をもつので，このグラフの最大次数は高々 $2f$ であり，
$|I|\ge|Q_0|/(2f+1)>1$，したがって $|I|\ge2$ をみたす独立集合 $I$ をもつ．
すべての $h\in I$ の部分木を同時に置き換えると，$\tuple u'\in L_A(G)$ と $W:=E_A[\tuple u']\in L$ を得る．
文字 $x$ の純な成分を固定し，他方の文字を $y$ と書く．
$I$ の塊の外にあるその文字は変わらず，したがって $x$ に等しい；
これは $I$ の外の所有者の塊（保たれる）と，最初の二文字と最後の二文字（印がついている）を含む．
したがって新しい成分の各極大の $y$-連の前後には $x$ があり，それは間に文字なしに連続する $I$ の塊の和の中にある．
そのような塊は一つの所有者に属する：
間に $I$ の塊がない $I$ の相異なる所有者 $h\ne h'$ の二つの塊は，変わらない文字で隔てられる．
その隙間は空でない（そうでなければ $h$ と $h'$ は結ばれる）か，$I$ の外の所有者の塊を含むからである．
したがって新しい成分の各 $y$-連は $I$ の中に一意な\emph{作る}所有者をもち，前の段落から，各 $h\in I$ は，
その文字を $x_h$ と呼ぶ純な成分の中に少なくとも一つの連を作る．
$I$ の中の $h\ne h'$ について，それらが作る連は異なる：
同じ純な成分の中にあれば作る所有者が異なり，異なる純な成分の中にあれば $\tuple u'$ の異なる成分の中にある．

$I$ の中の $h_1\ne h_2$ をとる．$x_{h_1}=x_{h_2}=:x$ ならば，$W$ は因子 $xx$（長さ少なくとも $5$ の純な成分の最初の二文字）と
二つの内部の $y$-連をもち，したがって少なくとも五つの連をもつ．
そうでなければ，$h_1$ と $h_2$ に (a) の議論を適用して $\mathtt a^5$ と $\mathtt b^5$ はともに現れ，
$x:=x_{h_1}$，$y:=x_{h_2}$ とすると，二つの純な成分はそれぞれ連 $x,y,x$ と $y,x,y$ を含む；
これらは $W$ の六つの相異なる連であり，$W$ は因子 $xx$ をもつ．
\end{proof}

\begin{proposition}\label{prop:app-two-nonsaturated}
$L$ が Parikh 飽和でなければ $t_{\mathrm{typ}}(G,c)\le D_f(r+1)v\theta$ である．
\end{proposition}

\begin{proof}
$X$ と $T$ を補題~\ref{lem:app-two-insertion} のとおりとする．
まず，$T$ の重い節点のすべての反鎖が高々 $6f+1$ 個の節点をもつことを示す．
そうでなければ補題~\ref{lem:app-two-insertion} が適用できる．
$\mathtt a^4$ も $\mathtt b^4$ も現れなければ，(a) は不可能である．
両方が現れれば，(b) の語 $W$ は少なくとも四つの連と長さ少なくとも $2$ の連をもち，補題~\ref{lem:app-two-alternate} に反する．
ちょうど一方，例えば $x_0^4$ が現れれば，(a) から $x=x_0$ であり $y^5$ は現れないので，
(b) から $W$ は因子 $xx$ と二つの内部の $y$-連をもち，補題~\ref{lem:app-two-runs}(iv) に反する．

型付きのラベルは非終端記号とその成分の内容の組なので，高々 $3^fv$ 個ある．
$T$ の道の上で型付きのラベルは繰り返さない．二つの等しいラベルの間で切ると，
生成語が長くならず節点の少ない $X$ からの導出が得られるからである；したがって道は高々 $3^fv$ 個の節点をもつ．
重い節点は祖先について閉じており，一つの深さにあるものは反鎖をなすので，それらは高々 $(6f+1)3^fv=D_fv$ 個である．
根が重くなければ $\|\tuple u\|_1\le\theta$ である．
そうでなければ，$\tuple u$ のすべての文字は重い節点が書くか，その重くない子の生成語の中にあり，
各重い節点は高々 $\theta$ 文字を書き，高々 $r$ 個の子をもつ．したがって $\|\tuple u\|_1\le D_f(r+1)v\theta$ である．
\end{proof}

\begin{lemma}\label{lem:app-two-saturated}
$L$ が Parikh 飽和ならば $\tau(G,c)\le C_{f,r}v^2\theta^4H_\theta$ である．ここで $C_{f,r}:=(r+1)(4r(3f+1)+D_f)+3$ である．
\end{lemma}

\begin{proof}
アリティ $d$ の $X=A_{\mathbf t}$ を固定し，$(p,q):=\Psi(E_A)$ とおく．したがって $p+q\le\theta$ である．
また $\mathcal P:=\Psi(L\cap\Sigma^+)$ とおく．
補題~\ref{lem:app-two-basic}(i) と Parikh 飽和から，$\tau_X$ は
$(p,q)+\sum_i\Psi(z_i)\in\mathcal P$ をみたす型 $\mathbf t$ のタプル $\tuple z$ についての $\|\tuple z\|_1$ の最小値である．
$(x_B,y_B):=\Psi(\omega_B)$ と書き，導出木の節点 $u$ について，その生成語の Parikh ベクトルを $(a_u,b_u)$ とする．
$u$ が語 $w$ の成功した導出の中にあれば，その部分木を $\omega_{B_u}$ の導出で置き換えると，
Parikh ベクトル $\Psi(w)-(a_u,b_u)+(x_{B_u},y_{B_u})$ をもつ $L$ の語が得られる；
$B_{u'}=B_u$ である真の子孫 $u'$ について，$u$ の部分木を $u'$ の部分木で置き換えると，
$\Psi(w)-(a_u-a_{u'},b_u-b_{u'})$ をもつ語が得られる．
以下の比較はすべてこの種類のものである：それは，固定した導出からそのような置換で得られる完全な成功した導出が生成する $L$ の語の
Parikh ベクトルに関するものであり，Parikh 飽和から，語の $L$ への所属はその Parikh ベクトルのみによる．
内容 $t_1,\dots,t_d$ の和集合は $\Sigma$ であるか，そうでなければそれらはすべて一つの単集合に等しい．

\emph{和集合が $\Sigma$ の場合．}
$t_{\mathtt a}$ と $t_{\mathtt b}$ を，内容がそれぞれ $\mathtt a$，$\mathtt b$ を含む成分の個数とする；どちらも $\{1,\dots,f\}$ に入る．
$n_{\mathtt a}$ 個の文字 $\mathtt a$ と $n_{\mathtt b}$ 個の文字 $\mathtt b$ をもつ型 $\mathbf t$ のタプルが存在するのは
$n_{\mathtt a}\ge t_{\mathtt a}$ かつ $n_{\mathtt b}\ge t_{\mathtt b}$ のとき，かつそのときに限るので，
$\tau_X$ はこれらの下界と $(p+n_{\mathtt a},q+n_{\mathtt b})\in\mathcal P$ をみたす最小の $n_{\mathtt a}+n_{\mathtt b}$ である．
最小にする組を固定し，$c_{\mathtt a}:=n_{\mathtt a}-t_{\mathtt a}+1\ge1$，
$c_{\mathtt b}:=n_{\mathtt b}-t_{\mathtt b}+1\ge1$ とおき，
$W:=\mathtt a^{p+n_{\mathtt a}}\mathtt b^{q+n_{\mathtt b}}$ の成功した導出 $T$ で節点の数が最小のものを固定する．
重い節点 $u$ が $a_u<c_{\mathtt a}$ かつ $b_u<c_{\mathtt b}$ をみたせば，その置換は
$n'_{\mathtt a}\ge t_{\mathtt a}$，$n'_{\mathtt b}\ge t_{\mathtt b}$，
$n'_{\mathtt a}+n'_{\mathtt b}<n_{\mathtt a}+n_{\mathtt b}$ をみたす Parikh ベクトル $(p+n'_{\mathtt a},q+n'_{\mathtt b})\in\mathcal P$ を与える；
したがってすべての重い節点は $a_u\ge c_{\mathtt a}$ または $b_u\ge c_{\mathtt b}$ をみたす．
$u'$ が同じ非終端記号をもつ $u$ の真の子孫ならば，$a_u-a_{u'}\ge c_{\mathtt a}$ または $b_u-b_{u'}\ge c_{\mathtt b}$ である：
差が零ならば節点数の最小性に反し（補題~\ref{lem:app-two-basic}(ii)），他の小さい差は組の選び方に反する．
道に沿って $(\lfloor a_u/c_{\mathtt a}\rfloor,\lfloor b_u/c_{\mathtt b}\rfloor)$ は成分ごとに非増加であり，
それが一定の間は非終端記号は繰り返さない．
$n\ge t\ge1$，$p\ge0$ について $\lfloor(p+n)/(n-t+1)\rfloor\le p+t$ である
（$(p+t)(n-t+1)-(p+n)=(n-t)(p+t-1)\ge0$）ので，それは高々 $1+(p+t_{\mathtt a})+(q+t_{\mathtt b})\le\theta+2f+1$ 個の値をとり，
道は高々 $v(\theta+2f+1)$ 個の節点をもつ．
一つの深さにある重い節点は互いに素な生成語をもつので，それらは高々
$(p+n_{\mathtt a})/c_{\mathtt a}+(q+n_{\mathtt b})/c_{\mathtt b}\le\theta+2f$ 個である．
重い節点は祖先について閉じており，すべての文字は重い節点が書くか，その高々 $r$ 個の重くない子の一つの中にある；したがって
$\tau_X\le|W|\le\max\{\theta,(r+1)\theta(\theta+2f)(\theta+2f+1)v\}\le(2f+1)(2f+2)(r+1)v\theta^3$ である．

\emph{すべての内容が $\{\mathtt a\}$ に等しい場合}（$\{\mathtt b\}$ の場合は対称である）．$h:=d\in\{1,\dots,f\}$ とする．
$K:=\min\{k\ge h:(p+k,q)\in\mathcal P\}$ とおく；このとき $\tau_X=K$ であり，
そうでなければ $\tau_X\le(2f-1)\theta$ なので，$K>\theta+2(f-1)\ge p+2(h-1)$ と仮定してよい．
$W:=\mathtt b^q\mathtt a^{p+K}$ の成功した導出 $T$ で節点の数が最小のものを固定する．

(P1) \emph{行の剛性}：$c_0\ge1$，$1\le\delta\le K-h$ をみたす $(c_0,y),(c_0+\delta,y)\in\mathcal P$ は存在しない．
そうでなければ，文脈 $(\mathtt b^y\mathtt a^{c_0-1},\eps)$ は $\mathtt a$ と $\mathtt a^{1+\delta}$ を受理するので，
$(1,c)$-代入可能性から $\D_L(\mathtt a)=\D_L(\mathtt a^{1+\delta})$ である（どちらも内容 $\{\mathtt a\}$ をもつ）；
文脈 $(\mathtt b^q\mathtt a^{p+K-\delta-1},\eps)$ は $\mathtt a^{1+\delta}$ を受理するので $\mathtt a$ も受理し，
$h\le K-\delta<K$ について $(p+K-\delta,q)\in\mathcal P$ となり，$K$ の選び方に反する．

(P2) \emph{高さ．}
$u'$ が同じ非終端記号をもつ $u$ の真の子孫で $b_u=b_{u'}$ ならば $a_u-a_{u'}\ge K-h+1$ である：
値 $0$ は節点数の最小性に反し（補題~\ref{lem:app-two-basic}(ii)），$[1,K-h]$ の値は $K$ の選び方に反する．
道を下ると $b_u$ は非増加である（非削除性から子の生成語は親の生成語の中に現れる）ので，
$b_u$ の一つの値をもつ道の節点は区間をなす；$b_u$ は高々 $q+1$ 個の値をとり，
一つの非終端記号と $b_u$ の一つの値をもつ三つの節点は $p+K\ge a_u\ge2(K-h+1)$ を与えるが，$K>p+2h-2$ なのでこれは偽である．
したがって道は高々 $2v(q+1)$ 個の節点をもつ．
生成語が $\mathtt b$ を含む節点を\emph{混合}という．
$1\le k\le q$ について，$b_u=k$ をみたし，そのような真の祖先をもたない混合の節点は互いに素な生成語をもつので，
高々 $\lfloor q/k\rfloor$ 個である；そのうちの一つの下では，そのような二つの節点は比較可能であり（そうでなければその生成語は $2k$ 個の文字 $\mathtt b$ を含む），
したがって鎖をなし，それは上の議論から高々 $2v$ 個の節点をもつ．したがって混合の節点は
$M\le2v\sum_{k=1}^q\lfloor q/k\rfloor\le2vqH_q\le2v\theta H_\theta$ 個である．

(P3) \emph{重い純な節点．}
混合でない節点 $u$ で $\ell_u=a_u>\theta+f-1$ をみたすものを\emph{重い純な}節点と呼び，
$g_u:=\ell_u-x_{B_u}\ge f$ とおく．
比較可能な二つの重い純な節点は異なる非終端記号をもつ．そうでなければ (P2) から上の節点は
$\ell_u\ge K-h+1+\theta+f>p+K=|W|_{\mathtt a}$ をもつからである．
特に，一つの非終端記号をもつ重い純な節点は互いに素な生成語をもつ．
重い純な節点の反鎖 $J$ について，そのすべてを置き換えると Parikh ベクトル $(p+K-\sigma(J),q+Y(J))$ をもつ語が得られる．
ここで $\sigma(J):=\sum_{u\in J}g_u$，$Y(J):=\sum_{u\in J}y_{B_u}$ である；
この語は完全な成功した導出が生成するので，その Parikh ベクトルは $\mathcal P$ に入る．
したがって $\sigma(J),\sigma(J')\le K-h$ かつ $Y(J)=Y(J')=:Y$ をみたす二つの反鎖 $J,J'$ は
$(p+K-\sigma(J),q+Y),(p+K-\sigma(J'),q+Y)\in\mathcal P$ を与える．
その第一座標は少なくとも $p+h\ge1$ で，差は $|\sigma(J)-\sigma(J')|\le K-h$ である；(P1) から $\sigma(J)=\sigma(J')$ である．

(P4) \emph{繰り返す非終端記号．}
$B$ が $m_B\ge2$ 個の重い純な節点のラベルならば，(P3) からそれらの生成語は互いに素であり，
これらの節点は互いに比較不能で，そのどの集合も反鎖である．したがって各々は $\ell_u\le p+K-(\theta+f)\le K-f$，
$g_u\le K-f\le K-h$ をみたす；(P3) で二つの単集合を比べると共通の値 $g_B$ を得，したがって共通の $\ell_B$ を得る．
さらに $y_B=0$ ならば，一つの置換は $h\le K-g_B<K$ について $(p+K-g_B,q)\in\mathcal P$ を与え，これは不可能である；したがって $y_B\ge1$ である．
$m_B>\theta$ のとき $B$ を\emph{頻出}という．
頻出の $B,C$ について，$B$ をラベルとする $y_C$ 個の節点と，別に $C$ をラベルとする $y_B$ 個の節点を選ぶ．
$m_B>\theta\ge y_C$，$m_C>\theta\ge y_B$ なのでこれは可能である；
どちらの集合も $Y=y_By_C$ をもつ反鎖であり，$y_Cg_B\le(m_B-1)\ell_B\le p+K-\ell_B\le K-f\le K-h$ であり，$C$ についても同様である．
(P3) から $y_Cg_B=y_Bg_C$ なので，すべての頻出の $B$ は共通の比 $\alpha:=g_B/y_B$ をもち，$\ell_B=x_B+\alpha y_B\le(1+\alpha)\theta$ である．

(P5) \emph{比．}
$\alpha_*:=(r+1)(\theta+f-1)(2\theta+1)v\le3f(r+1)v\theta^2$ とおく；
頻出の非終端記号が存在するとき $\alpha\le\alpha_*$ を示す．$\alpha\le\theta$ ならば明らかである．
そうでなければ，頻出の非終端記号 $B$ をラベルとする重い純な節点 $u$ を固定し，その部分木の重い純な節点の個数を $H_u$ とする．
これらの節点はその部分木の中で祖先について閉じている．
$u$ の生成語のすべての文字は，重い純な節点が書くか，その重くない子の生成語に属する．
各重い純な節点は高々 $\theta$ 文字を書き，そのような子を高々 $r$ 個もち，その各々の総生成語の長さは高々 $\theta+f-1$ である；
したがって $\ell_B\le(r+1)(\theta+f-1)H_u$ である．
これらの節点を，部分木の中での重い純な真の祖先の個数で分類する．
重い純な節点の鎖は (P3) から互いに異なる非終端記号をもつので高々 $v$ 個の節点をもち，
したがってこの個数は $v$ 未満で，分類は高々 $v$ 個である；比較可能な二つの節点は異なる個数をもつので，各分類は反鎖である．
一つの分類の中で，頻出のラベルをもつ節点は長さ少なくとも $\alpha$ で，$u$ の生成語の中に互いに素な生成語をもつので，
高々 $(1+\alpha)\theta/\alpha<\theta+1$ 個である；他のラベルをもつ節点は $T$ 全体で高々 $v\theta$ 個である．
したがって $H_u\le v(2\theta+1)$ であり，$\alpha\le g_B\le\ell_B\le\alpha_*$ である．

(P6) \emph{数え上げ．}
混合でない節点で $\ell_u>T_*:=(\alpha_*+1)\theta\le(3f+1)(r+1)v\theta^3$ をみたすものを\emph{大きい}という．
(P4)--(P5) から頻出の非終端記号は大きい節点のラベルにならないので，大きい節点は高々 $v\theta$ 個であり，
それらは混合でない節点の中で祖先について閉じている．
非削除性から節点の生成語は親の生成語の中に現れるので，混合の節点は祖先について閉じている．
$W$ の文字 $\mathtt a$ が，混合でも大きくもない節点 $w$ によって書かれるとする．
二つの閉包性から，混合でも大きくもない $w$ の祖先（$w$ を含む）は，根から $w$ への道の最後の区間をなす；その最上の節点を $z$ とする．
このとき $z$ は $T$ の根か，混合または大きい節点の子であり，$z$ の部分木は混合の節点も大きい節点も含まないので $\ell_z\le T_*$ である．
したがって $W$ のすべての文字 $\mathtt a$ は，混合の節点が書くか，大きい節点が書くか，そのような $z$ の生成語の中にある；
節点 $z$ は高々 $1+rM+rv\theta$ 個である．したがって
\[
\begin{aligned}
 p+K&\le\theta M+(rM+rv\theta+1)T_*+\theta\cdot v\theta\\
 &\le3v\theta^2H_\theta+4rv\theta H_\theta\cdot(3f+1)(r+1)v\theta^3
 \le\bigl(4r(3f+1)(r+1)+3\bigr)v^2\theta^4H_\theta
\end{aligned}
\]
である．ここで $M\le2v\theta H_\theta$ と $rM+rv\theta+1\le4rv\theta H_\theta$ を使った．
$(2f+1)(2f+2)\le D_f$ なので，どちらの場合も，また $K\le\theta+2(f-1)$ のときも，$\tau_X\le C_{f,r}v^2\theta^4H_\theta$ である．
\end{proof}

\begin{proof}[定理~\ref{thm:two-letter-content} の証明]
$L$ が Parikh 飽和ならば補題~\ref{lem:app-two-saturated} が適用できる．
そうでなければ命題~\ref{prop:app-two-nonsaturated} は $\tau(G,c)\le t_{\mathrm{typ}}(G,c)\le D_f(r+1)v\theta$ を与える．
どちらの場合も $\tau(G,c)\le C_{f,r}v^2\theta^4H_\theta$ である．
$c$ の像は三つの元をもつので $|\mathsf S^\sharp|=4$ であり，$N_{\mathrm{typ}}\le C_c(G)\le4^{fr}\|G\|$ である．
系~\ref{cor:distribution-representatives} から，
$\|K\|_+\le N_{\mathrm{typ}}(1+2\theta+r\tau(G,c))+\mathbf 1_{\{\eps\in L\}}$ をみたす共通の特徴標本が存在する．
$v\le\|G\|$ とすると，これは $1+2\theta\le3\|G\|^2\theta^4H_\theta$ と $\tau\le C_{f,r}\|G\|^2\theta^4H_\theta$ を使って最初の評価を，
$1+2\theta\le3\|G\|\theta$ と $\tau\le D_f(r+1)\|G\|\theta$ を使って二番目の評価を与える．
\end{proof}

\section{すべてのファンアウトにおける線形因子}
\label{app:linear-factor-fanout}

この付録は命題~\ref{prop:linear-factor-all-fanout} を証明する．$\Psi$ は Parikh 写像を表す．

\begin{proof}[命題~\ref{prop:linear-factor-all-fanout} の証明]
$t:=2f$ かつ $L:=L^{(f)}_{q,n}$ とおく．

(i) 言語 $R:=\mathtt a_1^*\mathtt d_1\mathtt a_2^*\mathtt d_2\cdots
\mathtt a_t^*\mathtt d_t\mathtt z^*$ は道の言語（$0$ から $t$ への）であって，グラフは
$\{0,\ldots,t\}$ 上のもので，辺 $\mathtt a_h:h-1\to h-1$，$\mathtt d_h:h-1\to h$ および
$\mathtt z:t\to t$ をもち，$R$ の語はその Parikh ベクトルによって定まる．
$w_i$ において標識 $\mathtt d_h$ は一回，位置 $h(i+1)$ に現れるので，
一つの標識の位置，および二つの標識の距離 $(h'-h)(i+1)$ は $i$ を定める；すべての語は
長さ $W$ をもつ．

$J\subseteq L$ かつ $1\le d\le f$ とし，$\tuple x,\tuple y$ を空でない
$d$-タプルとし，$E=(v_0,\ldots,v_d)$ を $d$-文脈で $E[\tuple x]=w_i\in J$ かつ
$E[\tuple y]=w_j\in J$ なるものとし，$F$ を $d$-文脈で $F[\tuple x]\in J$ なるものとする．
$F[\tuple y]\in J$ を示す；逆は対称である．
補題~\ref{lem:app-lf-path} により $F[\tuple y]\in R$．$i=j$ ならば，
$\sum_h\Psi(x_h)=\sum_h\Psi(y_h)$ であるので，$F[\tuple y]$ は
$F[\tuple x]$ の Parikh ベクトルをもち，それに等しい．$i\ne j$ とする．$v_0$ にある標識は，
$w_i$ と $w_j$ で同じ位置をもち，$v_d$ にある標識は末尾から同じ距離をもち，
一つの $v_h$ にある二つの標識は同じ距離をもつことになる；そのそれぞれが $i=j$ を強いる．したがって
$E$ の固定部分は高々 $d-1$ 個の標識を運び，$\tuple x$ は少なくとも
$t-d+1\ge d+1$ 個の標識を運ぶ．ある成分 $x_h$ は二つを運び，$L$ の語で
因子 $x_h$ をもつものはすべて $w_i$ である．したがって $F[\tuple x]=w_i=E[\tuple x]$ であり，
$F$ と $E$ の固定部分は等しい Parikh ベクトルをもち，$\Psi(F[\tuple y])=\Psi(E[\tuple y])$ であり，
$F[\tuple y]=w_j\in J$．有限言語は文脈自由であり，包含は
定義~\ref{def:class-cmcf} の後の項目~(a) である．

(ii) 各 $i$ について，(i) により $L\setminus\{w_i\}\in\Cmcf{f}{s}$ であるので，すべての
特徴標本は，補題~\ref{lem:distinguishing} により $w_i$ を含む．したがって
$F=L$．逆に，
$K=L$ は有限の $K$ で $L\subseteq K\subseteq L$ なる唯一のものであり，標準の
作用素は，標本無矛盾性と健全性により，その上で $L$ を返す．

(iii) 非終端記号 $P_0,\ldots,P_q$（ファンアウト $f$）と $S_0,\ldots,S_q$，
$C_0,\ldots,C_n$（ファンアウト1），開始記号 $S_q$，および次の規則をとる．
\[
\begin{aligned}
 &P_0\to(\mathtt d_1,\mathtt d_3,\ldots,\mathtt d_{2f-1}),\qquad
 C_0\to(\mathtt z),\qquad
 C_h\to(\xi_1^1\xi_2^1)(C_{h-1},C_{h-1})\ (1\le h\le n),\\
 &P_i\to(\mathtt a_1\xi_1^1\mathtt a_2,\ \mathtt a_3\xi_1^2\mathtt a_4,\ \ldots,\
 \mathtt a_{2f-1}\xi_1^f\mathtt a_{2f})(P_{i-1})\ (1\le i\le q),\\
 &S_j\to(\xi_1^1\mathtt d_2\xi_1^2\mathtt d_4\cdots\xi_1^f\mathtt d_{2f}\xi_2^1)(P_j,C_n)
 \ (0\le j\le q),\qquad
 S_j\to(\xi_1^1\mathtt z^{2f})(S_{j-1})\ (1\le j\le q).
\end{aligned}
\]
帰納法により $L(C_h)=\{\mathtt z^{2^h}\}$，
\[
 L(P_i)=\{(\mathtt a_1^i\mathtt d_1\mathtt a_2^i,\ldots,
 \mathtt a_{2f-1}^i\mathtt d_{2f-1}\mathtt a_{2f}^i)\},\qquad
 L(S_j)=\{\mathtt a_1^i\mathtt d_1\cdots\mathtt a_{2f}^i\mathtt d_{2f}
 \mathtt z^{2f(j-i)+2^n}:0\le i\le j\},
\]
したがって $L(S_q)=L$．すべての変数は一回ずつ
順に現れ，$P_j$ の連続する成分は
$\mathtt d_2,\mathtt d_4,\ldots$ によって区切られ，どの成分も空でない；したがって文法は良く，
既約である．それは $2q+n+3$ 個の非終端記号と，$1,q,q+1,q,1,n$ 個の規則で
大きさ $f+1$，$3f+2$，$2f+4$，$2f+3$，$2$，$5$ のものを，それぞれ $P_0$，$P_i$，$S_j$ の最初と
二番目の規則，$C_0$ および $C_h$ についてもつ（定義~\ref{def:standard-mcfg}）；
これは述べた大きさを与え，$n\le q$ は $(7f+17)(q+1)$ を与える．既約な
表示ではすべての規則はある成功する導出に現れ，すべての語は
長さ $W$ をもつので，$\bar\theta_G=W$．最後に $G$ を任意の既約な良い
ファンアウト $\le f$ の表示（$L$ の）とする．自明な型付けについて，
定理~\ref{thm:exact-reconstruction} は $\Wit(\widetilde G_0)$ を
$\mathcal B_{f,s_{\mathrm{triv}}}$ に対して特徴的にし，それは高々
$|V|+|P|\le\|G\|$ 個の語をもつ（定義~\ref{def:witness-set} と，
定理~\ref{thm:data-from-typed-thickness} の後の注意）；$s_{\mathrm{triv}}$ についての (ii) により，それは
$q+1$ 個の語をもつ．最後の評価は (ii)，(iii) と $|w_i|=W$ から従う．

(iv) 同一視はどの長さもどの雛形の形も変えないので，文法は
同じ大きさと厚さで既約かつ良いままであり，
(ii) と (iii) の議論は，(i) を示せばそのまま適用できる．(i) については，$J\subseteq L$ とし，
$u\,x\,v=w_i\in J$ かつ $u\,y\,v=w_j\in J$ で $x,y\in\Sigma^+$ とする；$i=j$ ならば
$x=y$．$i\ne j$ とする．すべての語は長さ $W$ をもつ．$u$ が $\mathtt d$ を含むならば，
最初の $\mathtt d$ は位置 $i+1$ を $w_i$ において，$j+1$ を $w_j$ においてもつことになる；$v$ が
$\mathtt d$ を含むならば，最後の $\mathtt d$ の後には $2(q-i)+2^n$ 個と
$2(q-j)+2^n$ 個の文字が続くことになる．したがって $x$ は両方の文字 $\mathtt d$（$w_i$ の）を含み，
因子 $\mathtt d\mathtt a^i\mathtt d$ をもち，それは $L$ の他のどの語にも現れず，$w_i$ には
一回だけ現れる．したがって $p\,x\,q'\in J$ は $(p,q')=(u,v)$ を強い，そのとき
$p\,y\,q'=w_j\in J$；逆は対称である．
\end{proof}





\section*{原稿作成過程における生成 AI および AI 支援技術の使用に関する申告}
本研究の準備中，著者は ChatGPT（OpenAI）と Claude
（Anthropic）を，原稿の構成，文章の
表現と翻訳，参考文献と記法の照合，および
\LaTeX{} ソースの編集の補助に用い，さらに証明の案の検討・検証や査読の補助にも用いた．
全ての数学的内容は著者が確認した．これらの道具を用いた後，著者は
必要に応じて内容を確認・編集し，本出版物の内容について全責任を負う．

\begin{thebibliography}{99}

\bibitem{Gold1967}
E.~M.~Gold.
\newblock Language identification in the limit.
\newblock {\em Information and Control}, 10(5):447--474, 1967.

\bibitem{Angluin1980}
D.~Angluin.
\newblock Inductive inference of formal languages from positive data.
\newblock {\em Information and Control}, 45(2):117--135, 1980.

\bibitem{ClarkEyraud2007}
A.~Clark and R.~Eyraud.
\newblock Polynomial identification in the limit of substitutable context-free
languages.
\newblock {\em Journal of Machine Learning Research}, 8:1725--1745, 2007.

\bibitem{deLaHiguera1997}
C.~de~la~Higuera.
\newblock Characteristic sets for polynomial grammatical inference.
\newblock {\em Machine Learning}, 27:125--138, 1997.

\bibitem{CosteTyping2004}
F.~Coste, D.~Fredouille, C.~Kermorvant, and C.~de~la~Higuera.
\newblock Introducing domain and typing bias in automata inference.
\newblock In \emph{Grammatical Inference: Algorithms and Applications -- ICGI 2004},
LNCS 3264, pp.~115--126, Springer, 2004.

\bibitem{Takada1988}
Y.~Takada.
\newblock Grammatical inference for even linear languages based on control sets.
\newblock \emph{Information Processing Letters}, 28(4):193--199, 1988.
\newblock doi:10.1016/0020-0190(88)90208-6.

\bibitem{Takada1995}
Y.~Takada.
\newblock Learning formal languages based on control sets.
\newblock In K.~P. Jantke and S.~Lange (eds.), \emph{Algorithmic Learning
for Knowledge-Based Systems}, LNCS 961, pp.~316--339, Springer, 1995.

\bibitem{Clark2013}
A.~Clark.
\newblock Learning trees from strings: A strong learning algorithm for some context-free grammars.
\newblock \emph{Journal of Machine Learning Research}, 14:3537--3559, 2013.

\bibitem{CosteGaretNicolas2012}
F.~Coste, G.~Garet, and J.~Nicolas.
\newblock Local substitutability for sequence generalization.
\newblock In {\em Proceedings of the 11th International Conference on
Grammatical Inference}, JMLR Workshop and Conference Proceedings 21,
pp.~97--111, 2012.

\bibitem{EyraudHeinzYoshinaka2016}
R.~Eyraud, J.~Heinz, and R.~Yoshinaka.
\newblock Efficiency in the identification in the limit learning paradigm.
\newblock In J.~Heinz and J.~M.~Sempere (eds.), {\em Topics in Grammatical
Inference}, pp.~25--46. Springer, 2016.
\newblock doi:10.1007/978-3-662-48395-4\_2.

\bibitem{CosteNicolas2019}
F.~Coste and J.~Nicolas.
\newblock Learning local substitutable context-free languages from positive
examples in polynomial time and data by reduction.
\newblock In {\em Proceedings of the 14th International Conference on
Grammatical Inference}, Proceedings of Machine Learning Research 93,
pp.~155--168, 2019.

\bibitem{Yoshinaka2008}
R.~Yoshinaka.
\newblock Identification in the limit of $k,l$-substitutable context-free
languages.
\newblock In {\em Grammatical Inference: Algorithms and Applications},
LNCS 5278, pp.~266--279. Springer, 2008.

\bibitem{Yoshinaka2009}
R.~Yoshinaka.
\newblock Learning mildly context-sensitive languages with multidimensional
 substitutability from positive data.
\newblock In {\em Algorithmic Learning Theory}, LNCS 5809, pp.~278--292.
 Springer, 2009.
\newblock doi:10.1007/978-3-642-04414-4\_24.

\bibitem{Yoshinaka2011}
R.~Yoshinaka.
\newblock Efficient learning of multiple context-free languages with
multidimensional substitutability from positive data.
\newblock {\em Theoretical Computer Science}, 412(19):1821--1831, 2011.
\newblock doi:10.1016/j.tcs.2010.12.058.

\bibitem{Yoshinaka2015General}
R.~Yoshinaka.
\newblock General perspective on distributionally learnable classes.
\newblock In \emph{Proceedings of the 14th Meeting on the Mathematics
of Language (MoL 2015)}, pp.~87--98.
Association for Computational Linguistics, 2015.
\newblock doi:10.3115/v1/W15-2308.

\bibitem{Yoshinaka2010}
R.~Yoshinaka.
\newblock Polynomial-time identification of multiple context-free languages
from positive data and membership queries.
\newblock In {\em Grammatical Inference: Theoretical Results and Applications},
LNCS 6339, pp.~230--244. Springer, 2010.
\newblock doi:10.1007/978-3-642-15488-1\_19.

\bibitem{ClarkYoshinaka2016}
A.~Clark and R.~Yoshinaka.
\newblock Distributional learning of context-free and multiple context-free grammars.
\newblock In J.~Heinz and J.~M.~Sempere (eds.), \emph{Topics in Grammatical
Inference}, pp.~143--172. Springer, 2016.
\newblock doi:10.1007/978-3-662-48395-4\_6.

\bibitem{ClarkYoshinaka2014}
A.~Clark and R.~Yoshinaka.
\newblock Distributional learning of parallel multiple context-free grammars.
\newblock {\em Machine Learning}, 96(1--2):5--31, 2014.
\newblock doi:10.1007/s10994-013-5403-2.

\bibitem{OatesEtAl2006}
T.~Oates, T.~Armstrong, L.~Becerra-Bonache, and M.~Atamas.
\newblock Inferring grammars for mildly context-sensitive languages in polynomial-time.
\newblock In {\em Grammatical Inference: Algorithms and Applications -- ICGI 2006},
LNCS 4201, pp.~137--147. Springer, 2006.

\bibitem{KanazawaYoshinaka2023}
M.~Kanazawa and R.~Yoshinaka.
\newblock Extending distributional learning from positive data and membership queries.
\newblock In {\em Proceedings of the 16th International Conference on Grammatical Inference},
Proceedings of Machine Learning Research 217, pp.~8--22, 2023.

\bibitem{DuncanEtAl2026Substitution}
A.~Duncan, M.~Elder, L.~Frenkel, and M.~Lyu.
\newblock A substitution lemma for multiple context-free languages.
\newblock \emph{International Journal of Algebra and Computation},
2026 (online publication).
\newblock doi:10.1142/S0218196726500463.

\bibitem{RambowSatta1999}
O.~Rambow and G.~Satta.
\newblock Independent parallelism in finite copying parallel rewriting systems.
\newblock {\em Theoretical Computer Science}, 223(1--2):87--120, 1999.
\newblock doi:10.1016/S0304-3975(97)00190-4.

\bibitem{kuriyamaCFG}
T.~Kuriyama.
\newblock Distributional learning of context-free languages under fixed
finite-semigroup typing.
\newblock Manuscript under review for \emph{Theoretical Computer Science};
\url{https://arxiv.org/abs/1409.6247v5}.

\bibitem{KuriyamaSCL2026}
T.~Kuriyama.
\newblock Finite-monoid compression in syntactic concept lattices: arity
hierarchies and a pseudovariety trichotomy.
\newblock \url{https://arxiv.org/abs/2608.29639v1}.

\bibitem{Pin1986}
J.-E. Pin.
\newblock {\em Varieties of Formal Languages}.
\newblock North Oxford Academic, London, and Plenum, New York, 1986.

\bibitem{Costa2002LI}
J.~C.~Costa.
\newblock Some pseudovariety joins involving locally trivial semigroups.
\newblock \emph{Semigroup Forum}, 64(1):12--28, 2002.
\newblock doi:10.1007/s002330010060.

\bibitem{Costa2002LIG}
J.~C.~Costa.
\newblock Some pseudovariety joins involving groups and locally trivial semigroups.
\newblock In \emph{Semigroups, Algorithms, Automata and Languages},
pp.~341--348. World Scientific, 2002.

\bibitem{SekiEtAl1991}
H.~Seki, T.~Matsumura, M.~Fujii, and T.~Kasami.
\newblock On multiple context-free grammars.
\newblock {\em Theoretical Computer Science}, 88(2):191--229, 1991.
\newblock doi:10.1016/0304-3975(91)90374-B.

\bibitem{Kallmeyer2010}
L.~Kallmeyer.
\newblock {\em Parsing Beyond Context-Free Grammars}.
\newblock Springer, Berlin, Heidelberg, 2010.
\newblock doi:10.1007/978-3-642-14846-0.

\bibitem{Kanazawa2019Ogden}
M.~Kanazawa.
\newblock Ogden's lemma, multiple context-free grammars, and the control language hierarchy.
\newblock {\em Information and Computation}, 269:104449, 2019.
\newblock doi:10.1016/j.ic.2019.104449.

\bibitem{GomezRodriguezEtAl2009}
C.~G\'omez-Rodr\'iguez, M.~Kuhlmann, G.~Satta, and D.~Weir.
\newblock Optimal reduction of rule length in linear context-free rewriting systems.
\newblock In {\em Proceedings of Human Language Technologies: The 2009 Annual
Conference of the North American Chapter of the Association for Computational
Linguistics}, pp.~539--547. Association for Computational Linguistics, 2009.

\bibitem{GeroldingerHalterKoch2006}
A.~Geroldinger and F.~Halter-Koch.
\newblock {\em Non-Unique Factorizations: Algebraic, Combinatorial and
Analytic Theory}, volume 278 of {\em Pure and Applied Mathematics}.
\newblock Chapman \& Hall/CRC, Boca Raton, 2006.

\bibitem{AlmeidaEtAl2009}
J.~Almeida, S.~Margolis, B.~Steinberg, and M.~Volkov.
\newblock Representation theory of finite semigroups, semigroup radicals and
formal language theory.
\newblock {\em Transactions of the American Mathematical Society},
361(3):1429--1461, 2009.
\newblock doi:10.1090/S0002-9947-08-04712-0.

\bibitem{Straubing1985}
H.~Straubing.
\newblock Finite semigroup varieties of the form $\mathbf V*\mathbf D$.
\newblock {\em Journal of Pure and Applied Algebra}, 36:53--94, 1985.
\newblock doi:10.1016/0022-4049(85)90062-3.

\bibitem{Tilson1987}
B.~Tilson.
\newblock Categories as algebra: an essential ingredient in the theory of monoids.
\newblock {\em Journal of Pure and Applied Algebra}, 48:83--198, 1987.
\newblock doi:10.1016/0022-4049(87)90108-3.

\end{thebibliography}
\normalsize

\end{document}
